%% notes-tex/database/database-expansion-hosts/database-expansion-hosts.tex
%% [[experiments.database.expansion-hosts]]
%% https://github.com/Mjvolk3/torchcell/tree/main/notes-tex/database/database-expansion-hosts/database-expansion-hosts.tex

\documentclass[10pt,a4paper]{article}
\usepackage[draft]{tcdoc}
\usepackage{float}

\usepackage{pdflscape}

\newcolumntype{L}[1]{>{\raggedright\arraybackslash}p{#1}}

\title{\vspace{-1.5em}Adding 200 Datasets Across Hosts}
\author{Michael Volk}
\date{\today}

\begin{document}
\maketitle
\thispagestyle{empty}

\noindent
The served graph at release 4.0, dated 2026-10-17, holds 170 schematized and
L0--L4-verified datasets over three organisms, 70 for \org{Saccharomyces cerevisiae}, 68 for
\org{Escherichia coli} and 32 for \org{Pseudomonas putida}, and 73,146,575 genotype by
environment records. The yeast corpus carries 65,889,109 of those records and the two
bacteria carry 4,561,865 and 2,695,603.

This pass continues the two rankings before it and adds three axes neither of them had. The
yeast pass ranked 154 candidates on one organism and the bacterial pass ranked 92 on two
hosts, and every row of both sat inside the three-member organism vocabulary the served
schema already holds. Here a second yeast enters, \org{Pichia kudriavzevii}. The pseudomonad
axis widens from KT2440 alone to the eleven pseudomonad rows of one downloadable artifact plus
\org{Pseudomonas aeruginosa}. And the pipeline question is taken up, meaning whether an
automated route from paper artifact to provenance-pinned ontology record is the right thing to
build next. What that section reaches is a prerequisite rather than a verdict: the served
schema has no field that can hold the justification such a pipeline would produce
(Sec.~\ref{sec:pipeline}).

The new host's result is a measured negative and it comes first. No published genome-wide
essentiality or fitness screen exists for \org{P. kudriavzevii} under any of its four species
names, and the search surface behind that conclusion is reported with it: a species-name by
screen-keyword query over Europe PMC returned 53 hits with 50 inspected, a sweep of 74
preprints across five preprint servers found no screen, GEO queried by taxid 4909 returned 29
records and four series, and lab-name searches covered the three groups that build such
libraries in other yeasts (Sec.~\ref{sec:pk}). The whole ingestible mass of the host is roughly
25 small rows. Its
largest phenotype resource is a preprint matrix of 161 isolates against 57 stressors whose
accession has not been retrieved, and its best fully public benchmark is 128 drug MIC values
on 32 sequenced isolates.

Measurements and rows come from different organisms in this pass, and that asymmetry sets its
shape. Three \org{S. cerevisiae} regulatory-DNA rows carry more \term{measurements} than
everything else here combined, meaning records times phenotype dimensionality
(Sec.~\ref{sec:measurements}), at 102,371,025, 51,339,035 and 6,739,258 measured promoter
sequences, and all three ingest through a schema that already serves their sibling
(Sec.~\ref{sec:regdna}). The row count comes from the pseudomonads, where one 10.4 GB
\texttt{sqlite} file holds 1,793 genome-wide fitness experiments over 11 pseudomonad
organisms under one schema and one condition vocabulary (Sec.~\ref{sec:pseudo}). The ranked
pool does not reach 200 rows, and the gap is not padded: what is missing is open discovery
plus the named retrieval asks of Table~\ref{tab:asks}, and Sec.~\ref{sec:state-gaps} names the
four axes the next discovery pass would search.

Table~\ref{tab:counts} gives the composition of the pool by class and by host.
Table~\ref{tab:final} is the recommended set in rank order with ten extra rows, so a row that
proves unreachable has a named replacement, and Table~\ref{tab:sources} gives each recommended
row's citation, link, data location, and the reason for its position. Three tables hold the
findings that sit outside the ranking: Table~\ref{tab:namespaces} records the four
gene-identifier namespaces live for \org{P. kudriavzevii} and what would join them,
Table~\ref{tab:mirror} measures what the two provenance mirrors store, and
Table~\ref{tab:schema} lists every change this wave needs in the served vocabulary and the
candidate machinery. Table~\ref{tab:asks} lists the asks and who each one goes to, and
Table~\ref{tab:excluded} records what was considered and dropped. Sec.~\ref{sec:summary}
closes with the summary statistics of Table~\ref{tab:summary}, and Sec.~\ref{sec:full} prints
the whole pool as Table~\ref{tab:candidates}, with a source for every row in
Table~\ref{tab:sources-all} and each row's candidate-gate verdict in
Table~\ref{tab:verdicts}.

\vspace{0.5em}
\tccontents
\clearpage

%% notes-tex/database/database-expansion-hosts/sections/state.tex
%% [[experiments.database.expansion-hosts]]
%% https://github.com/Mjvolk3/torchcell/tree/main/notes-tex/database/database-expansion-hosts/sections/state.tex
%% SOURCE: hand-authored prose. Served dataset and record counts come from release
%% 2026.10.17-4295f7ca and the per-organism tables of the thesis database chapter
%% (experiments/database/results/organism_summary.json, probe variant probes-v1); the
%% candidate figures come from the triage report
%% /scratch/projects/torchcell-scratch/tmp/deep-research/reports/Pk and next dataset wave triage.md
%% and the research notes beside it.

\section{State across hosts\secstatus{todo}}\label{sec:state}

A \term{host} here is an organism the served graph can declare, meaning one that
\texttt{torchcell/data/experiment\_dataset.py} holds as an \texttt{Organism} member and the
genomes tier resolves to pinned bytes. Three hosts are served today and the asymmetry between
them is the starting condition of this pass: \org{S. cerevisiae} carries 70 datasets and
65,889,109 records, \org{E. coli} 68 datasets and 4,561,865 records, and \org{P. putida} 32
datasets and 2,695,603 records. The three-organism total of 73,146,575 is computed against a
yeast figure of 65,889,107, two records below the figure the yeast table itself reports, so
the discrepancy is in the served tables and is carried here rather than resolved.

Where each host is thin is a property of its phenotype mix rather than of its record count.
\org{E. coli}'s largest section is environmental and chemogenomic response at 22 datasets,
followed by fitness, metabolite, and protein abundance at 8 each, expression and viability at
5 each, product titer at 3, and flux, gene interaction, and protein synthesis and turnover at
2 each, with one dataset each for morphology, phosphoproteome, and promoter activity. For
\org{P. putida} protein abundance leads with 14 datasets, then environmental and chemogenomic
response with 7, product titer with 6, metabolite with 2, and expression, fitness, and
viability with one each. So the newer host of the two carries almost no gene-perturbation
fitness, which is the layer the yeast corpus is built on, and the layer a transfer claim would
be judged against.

The plan this serves is stated in the thesis conclusions chapter: ``The first pillar is the
database, which already holds \org{Escherichia coli} and extends next to ten further
gram-negative bacteria and ten yeasts, each with its own virtual cell, with \org{Pseudomonas
putida} and \org{Pichia kudriavzevii} as the first hosts with little data, on which the models
are judged.'' The same passage records the test as a comparison of label counts. A Cell Graph
Transformer trained on
\org{S. cerevisiae} and fine-tuned on \org{P. kudriavzevii}, both budding yeasts, or trained
on \org{E. coli} and fine-tuned on \org{P. putida}, both in the class Gammaproteobacteria,
could match a model trained from scratch on the target with several-fold fewer labels, and the
threshold recorded beside it is five-fold fewer labels at matched Pearson correlation.

That test fixes what a low-data host needs before a model can be judged on it, and the list is
short and ordered. It needs a reference genome pinned in the genomes tier, and
\org{P. kudriavzevii} SD108 already has one, because
\file{torchcell/sequence/genome/registry.py} pins the assembly set
\texttt{pkudriavzevii\_SD108\_v1} and \file{torchcell/candidates/gates.py} already lists that
set among the readable ones. It needs an \texttt{Organism} member, which that host
does not have, so no loader can declare it and \texttt{dataset\_adapter\_map} can hold nothing
for it. It needs a host token in the gate, so that G2 can resolve a row's strain to that
assembly set. And it needs held-out labels of the same kind as the source host's labels, and
the subsections below measure how far each host of this wave is from that.

\subsection{What the served graph holds per host\secstatus{todo}}\label{sec:state-served}

On the gene-perturbation fitness axis the served graph is one host deep. The yeast corpus's
fitness and gene-interaction maps supply 67 percent of all 65,889,109 yeast records, and the
metabolite and titer screens that bear directly on strain engineering are three to four orders
of magnitude smaller. \org{P. putida} contributes one fitness dataset and \org{E. coli} eight,
so a cross-host fitness comparison has a dense source and two sparse targets.

On the sequence-to-readout axis the served graph is thinner still, on every host. One dataset
serves a varied-sequence library at scale, \texttt{vaishnav2022.py} with 71,158 records of
\texttt{VariablePromoterReporterPerturbation} crossed with
\texttt{PooledSortSeqExpressionPhenotype}, and \org{E. coli} carries a single promoter-activity
dataset. Everything else perturbs a gene list. That matters for a low-data host because a
promoter, a terminator, and an integration locus are the parts a new host is engineered with
first, before any deletion collection for it exists.

\subsection{What this wave adds per host\secstatus{todo}}\label{sec:state-rows}

\notesec{\org{S. cerevisiae}, measurements rather than rows} The three regulatory-DNA rows add
102,371,025, 51,339,035 and 6,739,258 measured promoter sequences, all on the same sort-seq
platform and all under classes the served schema already holds (Sec.~\ref{sec:regdna}). Below
them the same axis supplies a 1,000,000-sequence transcription-start-site library, a
489,348-sequence 5' UTR library, a 13,000-sequence 3' end library, and the position-effect
panels, each of which needs one new leaf rather than a new record type.

\notesec{\org{P. putida}, two ready-made matrices} The additions are dense and small in row
count: 332 RB-TnSeq samples as a 4,732 by 332 gene fitness matrix with t statistics and
per-sample metadata, 321 RNA-seq profiles as a 5,564 by 321 matrix over 118 typed conditions,
a 71-condition nitrogen BarSeq screen with 672 genes showing a significant phenotype, and a
24-strain by 205-substrate respiration panel that is the only resource letting this host enter
as a species with intraspecific variation rather than as one reference strain
(Sec.~\ref{sec:pseudo}).

\notesec{Nine further pseudomonads, inside one artifact} The nine non-KT2440 pseudomonad rows
of the Fitness Browser hold 1,479 genome-wide fitness experiments between them, measured, and
roughly $6.7\times10^{6}$ measurements, inferred from the per-gene matrix shape. They arrive
in the same \texttt{feba.db} download, under the same schema, with the same condition
vocabulary and the same RefSeq mapping files as the KT2440 rows already served, so the
marginal work is genome registration rather than loader work.

\notesec{\org{P. aeruginosa}, essentiality and isolate panels} This host enters for the
layer the others lack: 199 curated essentiality calls scoped by strain and medium, a
414-isolate clinical panel with a genome, a transcriptome, and 1,656 agar-dilution MIC records
per the CLSI standard, a 443-isolate longitudinal panel with 8 phenotypes, and a
125-isolate survey with 13 readouts (Sec.~\ref{sec:pseudo}).

\notesec{\org{P. kudriavzevii}, about 25 small rows} The host's additions are a 161-isolate by
57-stressor growth matrix, 9,177 measurements inferred, a 12-strain low-pH expression atlas at
two sampling points, 128 four-drug MIC values on 32 sequenced isolates, a 47-guide by
21-locus landing-pad screen with four readouts, and four small GEO series
(Sec.~\ref{sec:pk}). None of them is a screen.

\subsection{What is still missing after this wave\secstatus{todo}}\label{sec:state-gaps}

Four gaps survive the wave, and each is a measured absence rather than an unexamined axis.

There is no measured essentiality or genome-wide fitness screen in \org{P. kudriavzevii} of
any kind, so a model trained on yeast essentiality has no target to be judged against in that
host. Its only measured gene-perturbation fitness data are the 47-guide landing-pad screen,
one validated paralog-pair deletion, and the targeted knockouts inside each engineering paper.
Essentiality exists only as output of the iIsor850 reconstruction over 850 genes, which the
yeast pool already excludes under the rule \texttt{not-a-dataset}. The one remaining lead is
unresolved rather than negative: NCBI holds 713 BioProjects for the species and 357 of them
match \texttt{resequencing}, and three sampled ones registered by the DOE Joint Genome
Institute on 2026-05-11 are titled as SD108 \texttt{IoSL} resequencing with no publication
attached.

There is no baseline genome-wide essentiality call for KT2440 from the 2018 to 2026 window at
all, because RB-TnSeq by construction cannot score a gene it cannot mutate. The dataset that
would complete it, a 78,932-member CRISPRi library with 10 to 15 guides per gene and 798
non-targeting controls screened on glucose, acetate, and p-coumarate at two concentrations,
exists only as a DOE abstract with no publication, no repository, and no accession, so it is
tracked rather than planned for.

There is no \org{P. putida} or environmental-pseudomonad isolate collection with paired
sequence and growth data beyond the 24-strain respiration panel, and that panel's
data-availability statement could not be read. So the intraspecific genotype axis for this host
rests on one preprint.

And there is no ribosome profiling, CAGE, transcription-start-site mapping, or arrayed
overexpression dataset for \org{P. kudriavzevii}, by a targeted query and the GEO taxid sweep.
Chinese-language indices, ArrayExpress and BioStudies, and the Zenodo, Dryad, and figshare
APIs were not searched, so absence on those routes is unmeasured rather than negative.

%% notes-tex/database/database-expansion-hosts/sections/rule.tex
%% [[experiments.database.expansion-hosts]]
%% https://github.com/Mjvolk3/torchcell/tree/main/notes-tex/database/database-expansion-hosts/sections/rule.tex
%% SOURCE: hand-authored prose. The rule, the rank quantity and the pool guard are read from
%% experiments/database/scripts/build_candidate_datasets_table.py and
%% experiments/database/scripts/build_bacteria_candidate_datasets_table.py; the per-row
%% figures come from the triage report
%% /scratch/projects/torchcell-scratch/tmp/deep-research/reports/Pk and next dataset wave triage.md

\section{The selection rule\secstatus{todo}}\label{sec:rule}

\subsection{Sequence reconstruction is the gate, not a preference\secstatus{todo}}\label{sec:gate}

A phenotype is only trainable against a genotype that can be written down, so every row names a
\term{sequence basis}, meaning the route by which the total genomic content of one strain would
be reconstructed. The bases are named per host, because the reference differs per host and because
a \texttt{PP\_} locus tag, a JGI \texttt{gNNNN} gene model, and a PA14 locus tag are different
namespaces that must stay distinguishable rather than being merged into one string field.

Two of the bases this pass needs do not exist in either pool today. The yeast
\texttt{SeqBasis} literal holds 11 values and every one of them is S288C-relative or
yeast-panel-relative, so \texttt{isolate-WGS} transfers to a \org{P. kudriavzevii} isolate
panel unchanged while a designed edit in the SD108 background has no value to take. The
bacterial literal is KT2440- and K-12-relative in the same way, so a transposon insertion in
\org{Stutzerimonas stutzeri} RCH2 or a guide library in PAO1 needs its own basis rather than a
borrowed one. Table~\ref{tab:schema} lists both additions with the policy each carries.

What fails the gate is a strain whose genome is unknown even in principle. The three
acetic-acid evolved \org{P. kudriavzevii} lines, PkAC-7, PkAC-8 and PkAC-9, land there: they
are characterized phenotypically and transcriptionally across six stresses, and no
resequencing, variant table, or accession surfaced in either search pass, so the genotype is
undefined and the strains are usable only as named strains with stress phenotypes. Rows in
that state are dropped in Table~\ref{tab:excluded} rather than ranked low, because no amount
of phenotypic value substitutes for a missing genotype.

\subsection{The rank quantity is measurements, not records\secstatus{todo}}\label{sec:measurements}

Rows are ordered by \term{measurements}, meaning records times phenotype dimensionality,
descending, which is the quantity \texttt{build\_candidate\_datasets\_table.py} already sorts
on. \term{Records} is the number of genotype $\times$ environment $\times$ replicate rows a
dataset can yield, and \term{dimensionality} counts the values each row carries. A scalar MIC
panel of $n$ isolates against $k$ drugs contributes $nk$ numbers, and a transcriptome panel of
$n$ isolates at 5,564 genes contributes $5.6\times10^{3}n$.

Ranking on records alone gets the order visibly wrong here, and the two clearest cases are
worth naming rather than asserting the rule. The 414-isolate clinical panel of Khaledi et al.
2020 holds 414 records, which would rank it below a 5,459-mutant arrayed screen, and yet each
of those records is a genome plus a transcriptome, so the row is roughly $2.4\times10^{6}$
measurements, inferred from the panel shape. The nine new pseudomonad rows hold 1,479 fitness
experiments, measured, against roughly $6.7\times10^{6}$ measurements, inferred, because each
experiment is a per-gene vector. The product places both correctly, and the supported-datasets
table applies that same normalization through its compressed-signal column.

\subsection{Band comes before tier, and why\secstatus{todo}}\label{sec:band}

Scale is the right default ordering and it is the wrong first question when a specific host is
being brought up. Each row therefore carries a \term{band} applied before tier and before
scale, and the composition by band is Table~\ref{tab:counts}.

The reason is measured in this pass rather than argued. The three regulatory-DNA rows at
$10^{8}$, $5.1\times10^{7}$ and $6.7\times10^{6}$ measured promoter sequences carry more
measurements than every other row combined, and none of the three bears on a low-data host,
because all three are \org{S. cerevisiae} and their perturbation is an 80 bp promoter insert in
a fixed scaffold. Conversely the largest \org{P. kudriavzevii} phenotype resource is 9,177
measurements, inferred, which is four orders of magnitude below the smallest pseudomonad
fitness row, so a pure measurement order puts every row of the new host below any plausible
cut. Both facts are consequences of what the hosts have been measured with rather than of what
they are worth to a transfer experiment, and a band corrects for exactly that difference.

\subsection{The per-host floor, and why a floor rather than a score bonus\secstatus{todo}}\label{sec:floor}

The recommended set carries a floor of rows per host, applied after ranking and reported row by
row, and every promotion it makes is printed in Table~\ref{tab:final} rather than folded into a
score.

A floor is the right instrument because the setup a new host needs is per host and is done
once. An \texttt{Organism} member, a host token in \file{torchcell/candidates/gates.py}, and a
reference release pinned in the genomes tier are each written the first time any row for that
host is built and never again, so the setup for a host with one built row is the same as the
setup for a host with twenty. A floor guarantees that every host in the wave starts, so that
one-time setup yields served records in this pass rather than in the next one.

A score bonus would do the same arithmetic and lose the audit. Folding a host weight into the
rank quantity makes a promoted row indistinguishable from a row that earned its position, and
the whole point of the measurement order is that a surprising position can be checked against
two numbers. The pin mechanism of the bacterial pass is the precedent: a stakeholder priority
is a real criterion, so it is applied and reported rather than absorbed into the score.

\subsection{Where the lines fall\secstatus{todo}}\label{sec:lines}

The long-run cut is 200 additions against the 170 datasets served, so the target is 370. The
pool ranked here does not reach the cut and the shortfall is reported rather than padded.
Sec.~\ref{sec:state-gaps} names the four axes the next discovery pass would search, and
Table~\ref{tab:asks} lists the nine asks that would convert unreachable rows into reachable
ones. Table~\ref{tab:counts} gives the composition of the pool by class and by host, so the
cost of the band order and of the floor is visible per host.

%% GENERATED FILE -- do not hand-edit.
%% SOURCE: experiments/database/scripts/build_hosts_candidate_datasets_table.py
\begin{table}[H]\centering
\small
\caption[]{Candidates by class and by host, split at the two wave lines.
\emph{Genotypes}, \emph{Instances} and \emph{Meas.} sum the per-row axes over the whole
pool; a row with no count contributes nothing, so every total is a lower bound.
\emph{Meas.} is instances times phenotype dimensionality, which is why the per-host
ordering of the measurement column does not follow the row counts: one yeast
regulatory-DNA row carries more measurements than every pseudomonad row together.}
\label{tab:counts}
\begin{tabular}{@{}l r r r r r r@{}}
\toprule
 & Wave 1 & Wave 2 & Reserve & Genotypes & Instances & Meas. \\
\midrule
CRISPR library screen & 1 & 0 & 6 & 96,055 & $1.3\times 10^{5}$ & $1.3\times 10^{5}$ \\
Deep mutational scan & 3 & 0 & 7 & 925,159 & $1.5\times 10^{6}$ & $1.5\times 10^{6}$ \\
Engineering parts & 3 & 1 & 1 & 3,116 & $3.1\times 10^{3}$ & $3.1\times 10^{3}$ \\
Essentiality & 0 & 0 & 7 & 7,385 & $2.3\times 10^{3}$ & $2.3\times 10^{3}$ \\
Fitness / chemical genomics & 0 & 9 & 10 & 2,753,324 & $2.5\times 10^{7}$ & $3.2\times 10^{7}$ \\
Metabolome / flux & 0 & 0 & 8 & 359 & $9.6\times 10^{3}$ & $1.9\times 10^{6}$ \\
Modality / backbone & 0 & 1 & 2 & 14,247 & $1.4\times 10^{4}$ & $1.4\times 10^{4}$ \\
Natural variation & 4 & 0 & 15 & 28,291 & $4.1\times 10^{5}$ & $4.2\times 10^{5}$ \\
Production campaign & 0 & 5 & 26 & 471 & $1.5\times 10^{3}$ & $1.7\times 10^{3}$ \\
Proteome & 1 & 1 & 13 & 4,028 & $5.8\times 10^{3}$ & $7.9\times 10^{5}$ \\
Regulatory DNA & 27 & 1 & 0 & 166,789,446 & $1.7\times 10^{8}$ & $2.1\times 10^{8}$ \\
Tolerance / robustness & 0 & 1 & 11 & 8,267 & $9.5\times 10^{3}$ & $2.0\times 10^{5}$ \\
Transcriptome & 7 & 1 & 18 & 477 & $3.4\times 10^{3}$ & $1.6\times 10^{7}$ \\
Transposon fitness & 4 & 0 & 7 & 105,459 & $3.4\times 10^{6}$ & $9.8\times 10^{6}$ \\
\midrule
\multicolumn{7}{@{}l}{\emph{The same rows, by host}}\\
E. coli & 9 & 0 & 16 & 2,916,904 & $6.0\times 10^{6}$ & $5.5\times 10^{7}$ \\
P. aeruginosa & 0 & 1 & 23 & 109,342 & $1.1\times 10^{5}$ & $2.6\times 10^{6}$ \\
P. fluorescens & 0 & 4 & 1 & 100,000 & 745 & $3.4\times 10^{6}$ \\
P. kudriavzevii & 8 & 10 & 21 & 391 & $9.9\times 10^{3}$ & $3.6\times 10^{5}$ \\
P. putida & 8 & 0 & 41 & 10,341 & $5.5\times 10^{4}$ & $5.7\times 10^{6}$ \\
P. simiae & 0 & 1 & 0 & 0 & 201 & $9.1\times 10^{5}$ \\
P. syringae & 0 & 2 & 0 & 0 & 69 & $3.1\times 10^{5}$ \\
S. cerevisiae & 25 & 1 & 29 & 167,599,106 & $1.9\times 10^{8}$ & $2.0\times 10^{8}$ \\
S. stutzeri & 0 & 1 & 0 & 0 & 350 & $1.6\times 10^{6}$ \\
\midrule
\multicolumn{7}{@{}l}{\emph{The same rows, by band}}\\
regulatory DNA & 34 & 1 & 0 & 166,797,405 & $1.7\times 10^{8}$ & $2.1\times 10^{8}$ \\
cross-host fitness & 0 & 9 & 0 & 0 & $1.4\times 10^{3}$ & $6.2\times 10^{6}$ \\
host bootstrap & 8 & 10 & 21 & 391 & $9.9\times 10^{3}$ & $3.6\times 10^{5}$ \\
scale & 8 & 0 & 110 & 3,938,288 & $3.0\times 10^{7}$ & $5.5\times 10^{7}$ \\
\midrule
Total & 50 & 20 & 131 & 170,736,084 & $2.0\times 10^{8}$ & $2.7\times 10^{8}$ \\
\bottomrule
\end{tabular}
\end{table}


One line is set by the machinery rather than by judgment. The yeast pool holds 169
\texttt{Candidate} literals, of which 154 are unbuilt and ranked, and its \texttt{main()}
raises on a pool over 250 rows, so a 200-row wave cannot go into one pool at all. The guard
forces the split by organism, and the gate needs the same split for its own reasons. The yeast
\texttt{Candidate} model has no \texttt{organism} field at all, the bacterial one admits two
values, and \texttt{verdict.CandidateTable} admits only the two strings
\texttt{bacteria} and \texttt{yeast}, which is the \texttt{table} field of all 52 stored
verdicts. Table~\ref{tab:schema} records each of those edits.

%% notes-tex/database/database-expansion-hosts/sections/pk.tex
%% [[experiments.database.expansion-hosts]]
%% https://github.com/Mjvolk3/torchcell/tree/main/notes-tex/database/database-expansion-hosts/sections/pk.tex
%% SOURCE: hand-authored prose. Every count comes from the triage report
%% /scratch/projects/torchcell-scratch/tmp/deep-research/reports/Pk and next dataset wave triage.md
%% and from the two search passes recorded in
%% .../research_notes/Pk and next dataset wave triage/pichia_kudriavzevii.md, which carries the
%% per-finding verified, snippet and inferred labels this section preserves.

\section{Pichia kudriavzevii\secstatus{todo}}\label{sec:pk}

No published genome-wide essentiality or fitness screen exists for this species. There is no
Tn-seq, SATAY, Hermes, or piggyBac mutagenesis library, no genome-wide CRISPR knockout,
interference, or activation screen, and no arrayed deletion collection, for
\org{Pichia kudriavzevii}, \org{Issatchenkia orientalis}, \org{Candida krusei}, or
\org{Candida glycerinogenes} across 2010 to 2026. Douglass et al. 2018 established that those
four names are one species, with genomes collinear and 99.6 percent identical, so a search run
under all four names is the complete search for the host.

Each search axis is reported with its own counts, because a negative result is only citable
through the surface that produced it. A Europe PMC query crossing both current species names against ``CRISPR screen'',
``genome-wide screen'', ``essential genes'', ``transposon'', ``Tn-seq'', ``sgRNA library'',
``ribosome profiling'' and ``overexpression library'' returned 53 hits, of which 50 were
inspected at title level, with zero screens in this organism. A sweep of 74 preprints across
bioRxiv, medRxiv, Research Square, Preprints.org and Authorea found no screen, library, or
essentiality preprint. GEO DataSets queried by taxid 4909 returned 29 records and four series,
none of them a screen, and taxid 573 must not be used for the clinical name because it returns
1,328 unrelated records for a different genus. Lab-name searches established that Wheeldon's
genome-scale CRISPR platform exists for \org{Yarrowia lipolytica} and \org{Komagataella
phaffii} only, the latter as a six-fold-coverage library of 30,848 active guides covering more
than 99 percent of that genome, and has never been extended to this host. Essentiality in this
species exists only as output of the iIsor850 reconstruction over 850 genes, which the yeast
pool already excludes under the rule \texttt{not-a-dataset} because a reconstruction is not an
experiment.

The consequence for what the host can be used for is a limit on the label rather than on the
host. Gene-level fitness, essentiality, and genetic interaction are
unavailable as held-out labels in this species, so a transfer claim of the form trained on
yeast fitness and judged on \org{P. kudriavzevii} fitness has nothing to be judged against
today. What is available is isolate-level stress tolerance, antifungal susceptibility,
strain-level expression under acid stress, and per-locus integration fitness. Four axes remain
unsearched and are therefore unmeasured rather than negative: Chinese-language indices,
ArrayExpress and BioStudies, the Zenodo, Dryad, and figshare APIs, and the per-lab publication
pages of Alper, Blenner, and Jeffries.

\subsection{What does exist, led by the MAESTRO panel\secstatus{todo}}\label{sec:pk-have}

\notesec{MAESTRO, 170 isolates sequenced and 161 phenotyped} MAESTRO is the machine-learning
association method for small cohorts that the preprint is named for, and the panel it was run
on is whole-genome sequencing of 170
isolates with growth of 161 of them assessed across 57 stressors spanning heat, low pH,
organic acids, lignocellulosic inhibitors, and antifungals across azoles, polyenes, and
echinocandins, which is 9,177 growth measurements, inferred as the product of the two reported
axes. Copy-number variation in under 3.5 percent of genes has pleiotropic effects on stress
tolerance, a copy-number feature model reaches median $R^{2}$ 0.67 and median Pearson
$r$ 0.92 across the 57 conditions, and deletion of the tandem paralog pair
\gene{gene4260}/\gene{gene4261} produced effects across several stresses, which is the one
CRISPR-validated arm in the host. The panel is a preprint of 2025-09-04, bioRxiv
10.1101/2025.08.29.673034, corresponding author Ping-Hung Hsieh at LBNL. Its numbers conflict
across DOE pages, which report 162 strains, 305,435 SNPs, and 58 conditions against the
preprint's 170 strains and a project poster's 305,405 SNPs and 16,177 INDELs, so the figures
carried here are the preprint's and the disagreement is recorded rather than averaged. The
phenotype matrix itself has not been retrieved: bioRxiv returns 403 to every automated client
and Europe PMC's full-text route returns 500 for the preprint identifier.

\notesec{The low-pH expression atlas, 12 strains} Twelve isolates chosen for varying acid
tolerance, six tolerant and six susceptible, profiled at control pH near 5.5 and at pH 1.5,
sampled at 0 h and 1 h, with further arms testing 3-hydroxypropionate and citramalate, giving
roughly $2.4\times10^{5}$ expression measurements, inferred. It is the one omics resource in
the host built as a strain by condition matrix with a released processed count table, and the
table sits at Illinois Data Bank DOI 10.13012/B2IDB-1336177\_V1, which returns 403 to every
automated client.

\notesec{Douglass 2018, the cleanest public benchmark} Thirty-two sequenced isolates, 20
clinical and 12 environmental, with per-isolate MICs in mg/L for fluconazole, flucytosine,
amphotericin B, and micafungin, which is 128 MIC values read in full from the paper's own
table. Fluconazole MICs span 8 to 128 with a mode of 32. The same table gives ploidy, MAT
genotype, Illumina coverage, \gene{ERG11} and \gene{FKS1} variants, and copy status across the
\gene{ABC11}-\gene{ABC1} region per strain. Reads are at SRA SRP139299, the CBS573 annotated
reference at GenBank CP028773 to CP028778, and the second type strain at CP028531 to CP028535,
so this is the only row in the host that is fully public and fully provenanced today.

\notesec{The landing-pad screen, the only measured per-locus fitness} Forty-seven guide RNAs
across 21 intergenic regions, each characterized for cutting efficiency, integration
efficiency, cellular fitness, and GFP expression level, which is 188 readouts. Its artifacts
are a 23 KB package on a lab-hosted Box URL plus a GitHub repository, with no DOI and no
sequence accession, which is the retrieval risk Table~\ref{tab:asks} records.

\notesec{Four GEO series and the SD108 product line} GSE202655 holds six samples on L-malate
engineering and GSE159457 six samples on C4-dicarboxylate transporters in strain CY902 with raw
reads at SRP287255. GSE262836 and GSE282019 are two Delhi neonatal-intensive-care outbreak
series that together pair genome, expression, and fluconazole MIC across 10 RNA-seq samples.
The engineering line on SD108 covers succinate, 3-hydroxypropionate, citramalate,
glycolic acid, and xylose utilization, and its scale was overestimated on the first pass: Tran
et al. 2023 contributes about 10 distinct genotypes countable from the body text, with
Supplementary Table 5 unread and no sequence accession deposited at all, so the 100 to 250
genotype figure for the whole lineage is an upper-range estimate and is labeled one. The parts
panel of Cao et al. 2020 names 16 promoters and then ``and other tested promoters'', so its
total is unstated, and its deposit is one of at least eight Illinois Data Bank records for the
species, every one of which returned 403 on both passes.

\subsection{Four gene-identifier namespaces for one species\secstatus{todo}}\label{sec:pk-ids}

Four annotation namespaces are live for this one species, and the gap between the largest two
is wide enough that they are not inter-mappable by inspection. Table~\ref{tab:namespaces}
records each one with its gene count, the data that uses it, and the route that would join it
to the others.

The JGI MycoCosm assembly Issorie2, SD108 v2.0, is 21.69 Mbp over 130 scaffolds with 4,925
gene models, and 119 of those 130 scaffolds are predicted to be an alternate haplotype, with
separate primary and secondary allele gene-model tracks served. The whole engineering line
writes its loci in that namespace's \texttt{gNNNN} form. The MAESTRO preprint writes
\texttt{geneNNNN} instead, and whether \texttt{g4260} and \texttt{gene4260} are the same locus
is unverified, which is the kind of silent mismatch that corrupts a join without failing a
load. A third set is live in Ensembl Fungi for the same strain, assembly ASM76445v1, with 7,107
coding genes and 328 non-coding genes imported from INSDC. The clinical work uses a fourth,
the CBS573 reference annotated with YGAP by similarity and synteny to \org{S. cerevisiae},
about 10.8 Mb over five nuclear chromosomes.

The 4,925 against 7,107 difference is the reason a crosswalk has to be built from protein
sequence rather than from identifier similarity, and the haplotype split is the reason a naive
gene count in this species reads low against the 6,607 genes of the S288C reference. Cao et al.
2020 itself uses 10.8 Mbp as the genome size in its copy-number calculation, the haploid figure
rather than Issorie2's 21.69 Mbp, so two genome sizes already coexist in one lab's output. Under
the rule that the data is the data, a loader stores the released name verbatim in whichever
namespace the release used and records which namespace that is, and resolution across the four
is a downstream utility rather than a load-time gate.

%% GENERATED FILE -- do not hand-edit.
%% SOURCE: experiments/database/scripts/build_hosts_candidate_datasets_table.py
\begingroup
\footnotesize
\begin{longtable}{@{}L{22mm} L{34mm} L{20mm} L{30mm} L{66mm}@{}}
\caption[]{The gene-identifier namespaces this wave's released data speaks, per host.
\emph{Carries} is the model or gene count the namespace itself holds and \emph{Joins} is
what maps it to another namespace. One species in this wave has four live namespaces
whose model counts differ by more than two thousand genes, so a released gene id is not
a well-formed identifier until the namespace is named. A \emph{Joins} entry that says a
crosswalk has to be built is saying no join exists today: it is not a pointer to one.}
\label{tab:namespaces}\\
\toprule
Host & Namespace & Spelling & Carries & Joins \\
\midrule
\endfirsthead
\toprule
Host & Namespace & Spelling & Carries & Joins \\
\midrule
\endhead
\bottomrule
\endfoot
\org{P. kudriavzevii} & JGI Issorie2 gene models & \texttt{gNNNN} & 4,925 gene models & Nothing published. A crosswalk to the INSDC set has to be built by protein BLAST, and 119 of 130 scaffolds are predicted to be an alternate haplotype, so the gap is not only an identifier gap \\
\addlinespace[5pt]
\org{P. kudriavzevii} & MAESTRO gene spelling & \texttt{geneNNNN} & the same JGI model set under a different spelling & Assumed equal to the JGI spelling and not verified: g4260 must be checked against gene4260 rather than assumed equal, which is exactly the pair the MAESTRO CNV finding rests on \\
\addlinespace[5pt]
\org{P. kudriavzevii} & INSDC / Ensembl Fungi SD108 & \texttt{locus tag} & 7,107 coding genes & The 4,925-versus-7,107 gap is large enough that the two sets are not trivially inter-mappable. This is the set the genomes tier pins as pkudriavzevii\_SD108\_v1 \\
\addlinespace[5pt]
\org{P. kudriavzevii} & CBS573 / YGAP annotation & \texttt{YGAP id} & the clinical-isolate annotation & Used by the clinical work, including the MIC panel and the ERG11 and FKS1 variant calls, so a clinical row and an SD108 engineering row do not share a gene id at all \\
\addlinespace[5pt]
\org{P. putida} & KT2440 locus tags & \texttt{PP\_NNNN} & 5,729 genes in the served reference & Already served: the bacterial layer namespaces its five perturbation types by locus tag, so a PP\_ row writes the record type the schema holds \\
\addlinespace[5pt]
\org{P. putida} & Fitness Browser locusId & \texttt{orgId, locusId} & one row per gene per browser organism & refseq\_mapping.tsv maps org and locusId to a RefSeq accession and refseq\_assemblies maps orgId to an NCBI assembly, so the identifier bridge is supplied by the source rather than built \\
\addlinespace[5pt]
\org{P. aeruginosa} & PAO1 and PA14 reference numbers & \texttt{PA NNNN} & the two reference annotations the Tn-seq literature uses & GOLD\_84 carries PAO1 locus tags and GOLD\_115 is PA14 in LB, so the two gene lists are in different namespaces and must not be merged into one boolean \\
\addlinespace[5pt]
\org{P. aeruginosa} & Bakta de novo annotation & \texttt{per-isolate tag} & one annotation per newly sequenced isolate & Nothing maps a Bakta tag to a PA number, so an isolate survey annotated this way forces the alias layer rather than joining on a reference id \\
\addlinespace[5pt]
\org{S. cerevisiae} & SGD S288C R64-4-1 & \texttt{YALNNNC} & 6,607 genes in the shared vocabulary & The served namespace. A regulatory-DNA row is not gene-indexed at all: its identifier is the measured sequence, which is why those rows carry a reporter-locus or plasmid basis rather than a gene id \\
\addlinespace[5pt]
\end{longtable}
\endgroup


\subsection{The open leads\secstatus{todo}}\label{sec:pk-leads}

Three leads are unresolved rather than negative, and the first is time-sensitive. NCBI holds
713 BioProjects for this species and 357 of them match \texttt{resequencing}. Three sampled
ones, PRJNA1464946, PRJNA1464947 and PRJNA1464948, are titled as \org{Issatchenkia
orientalis} SD108 \texttt{IoSL10}, \texttt{IoSL16} and \texttt{IoSL19} resequencing,
registered on 2026-05-11 by the DOE Joint Genome Institute with no publication attached. A
series sequenced per isolate at that scale is either a mutagenized or an evolved SD108
library, and if it carries phenotypes it is the genome-wide resource this host currently
lacks. That is a hypothesis about an unopened deposit, not a finding.

The per-isolate deposition route for the MAESTRO panel is the second lead and is inferred
rather than verified. Three projects in the \texttt{IO\_NNN} naming exist, PRJNA571188 for
IO\_179, PRJNA571352 for IO\_048, and PRJNA583207 for IO\_044 gene expression profiling, each
registered as one BioProject per isolate rather than under an umbrella, which would explain
why no single accession for the panel appears in any search. Enumerating the complete
isolate-to-accession map needs either the preprint's data statement or a pass over all 357
summaries.

The third lead is a 69-condition expression matrix in this species, 5,140 genes by 69
conditions, described in a senior thesis from the Sinha group at UIUC that carries no
data-availability, code, deposition, or supplementary section of any kind. The provenance of
the expression and transcription-factor concentration data is one sentence that cites nothing.
So it is reachable by asking and by nothing else, and with a real provenance record a
69-condition expression matrix in this host would be worth more to the host's program than
anything else now public.

%% notes-tex/database/database-expansion-hosts/sections/regdna.tex
%% [[experiments.database.expansion-hosts]]
%% https://github.com/Mjvolk3/torchcell/tree/main/notes-tex/database/database-expansion-hosts/sections/regdna.tex
%% SOURCE: hand-authored prose. Counts come from the triage report
%% /scratch/projects/torchcell-scratch/tmp/deep-research/reports/Pk and next dataset wave triage.md
%% and from .../research_notes/Pk and next dataset wave triage/scerevisiae_gene_level_readouts.md
%% and .../yeast_expression_fm_corpus.md; the served Vaishnav figures are read from
%% torchcell/datasets/scerevisiae/vaishnav2022.py.

\section{Regulatory DNA\secstatus{todo}}\label{sec:regdna}

The \term{sequence-to-readout} tier is the set of datasets whose perturbation is a specified
DNA sequence, meaning a promoter, a 5' UTR, a core promoter, a transcription start site, a 3'
end, a terminator, or an integration coordinate, and whose readout is a reporter level, a
translation output, or the fitness of the strain carrying it. The served graph holds one such
dataset at scale, \texttt{vaishnav2022.py} with 71,158 records, against 70 yeast datasets whose
perturbation space is a gene list. The fitness and gene-interaction maps alone supply 67
percent of all 65,889,109 yeast records, so adding to that axis adds to the two-thirds the
corpus already has, while this tier's perturbation space is continuous rather than enumerable,
which is the perturbation space a model of regulatory grammar is fitted on, and no deletion
collection supplies it.

\notesec{The de Boer and Vaishnav lineage} de Boer et al. 2020 measured the expression output
of 102,371,025 promoter sequences across five experiments. The design is an 80 bp fully random
insert at $-170$ to $-90$ relative to the transcription start site inside a fixed scaffold,
driving YFP against a constitutive RFP on an episomal vector, with cells sorted into 18 uniform
bins by the log ratio of the two reporters and the promoters in each bin sequenced, so the
expression estimate is a read-weighted bin average. The media are complex and there are three
of them, glucose, galactose, and glycerol, which makes the medium a genuine environment
variable rather than a single condition. Vaishnav et al. 2022 ran the same platform with two
training libraries, 30,722,376 random sequences in complex medium and 20,616,659 in defined
medium, and the graph already serves the Vaishnav child of those two. The DREAM 2022 release is
the third member at 6,739,258 training sequences plus 71,103 held out, and every published
yeast expression model is scored on that held-out set.

The per-record cost of the largest row is characterized only in aggregate and belongs in the
record as such. Because each library's complexity exceeded the number of sorted cells, 78
percent of de Boer's promoter sequences appear in only one bin, often as one read from one
cell, which the authors assess as roughly 24 percent error in the expression estimate on
held-out data. That is a typed gap or a read-count field rather than a reason to decline the
row, and \texttt{vaishnav2022.py} already carries the same information as \texttt{read\_count}.
Two provenance items are open: the paper names GEO GSE104903 and GSE104878 and never states
which accession holds which library, and it never itemizes the five experiments that sum to
102,371,025, so one GEO page check settles both.

\notesec{The superset admission the existing loader already makes possible} The served
Vaishnav store holds 71,158 records of \texttt{VariablePromoterReporterPerturbation} crossed
with \texttt{PooledSortSeqExpressionPhenotype}, of which 71,054 carry an estimate and 104
\texttt{seqID} groups carry a typed \texttt{ProvenanceGap} for verbatim \texttt{NA} cells. The
two deferred training libraries are the complex-medium and defined-medium training files on
Zenodo, 3.59 GB and 2.47 GB, plus a 0.67 GB table on GEO,
and the loader's own docstring already calls them ``a size decision for the owner,
not a schema or a provenance one.'' Admitting them is the superset case the incremental path
was built for: the served records are a proven subset, so the admission check can prove the
development store still produces every served experiment id plus more, and the increment
matches served nodes by content id and adds only the new records. Admitting de Boer needs no
new class either, since the same two classes serve its child. The DREAM release is the one
member of the three that needs more than a download, because it is a curated slice of the same
platform and therefore needs an \texttt{AggregationRecord} at G1 rather than entering as a
primary measurement.

\notesec{The 3' end and the position-effect rows} Shalem et al. 2015 measured 13,000 designed
3' end sequences, 102 bp variable inside a 150 bp oligo, by FACS-seq into 16 bins of a
two-reporter ratio, and the library sequences and their measured expression values sit together
in one supplementary XLSX with no accession. It extends sequence-level coverage from the
promoter to the terminator and sits among the eight benchmark tasks the species-aware DNA
language model is scored on, so ingesting it serves a dataset a published model is measured
against. Its only uncertainty figure is a library-wide technical relative standard deviation of
13.2 percent rather than a per-construct error. The position-effect rows are the
one-construct-many-loci shape, where the sequence variable is purely the integration
coordinate. There are four of them: a 1,044-locus panel with a greater-than-13-fold expression
range and location dominating promoter, reporter, and carbon source, a panel of about 250 loci
each with a full transcriptome rather than one reporter value, a 482-locus panel from an
earlier study, and a genome-wide scan reporting 428 high-expressed and 444 low-expressed
sites. Not one of the four has a located accession, and the largest states that its table is
available from the corresponding author on request, so the whole sub-tier is supplement-bound
and two of its rows are asks.

\notesec{What the rest of the tier holds} Four further rows carry $10^{4}$ measurements or
more each and each needs exactly one new leaf. A Pol II MASTER library measures initiation
efficiency at 1,000,000 transcription start sites, deposited at GEO GSE185290 and BioProject
PRJNA766624, and its own preprint reported about 80,000 variants against the published
1,000,000, so the release has to be pinned rather than the paper cited. A random 5' UTR library
detected 489,348 of more than 500,000 designed variants upstream of \gene{HIS3}, read out by
competitive growth rather than by sorting, which by the rule that a different technology is its
own leaf makes it a separate phenotype class from a sort-seq bin average. A promoter-swap
library places 130 synthetic promoters spanning a 500-fold range upstream of named native
genes, giving about 10,000 strains with a per-strain pooled fitness, where the released gene
count is reported as 81, 96, and about 100 in three places and has to be reconciled before a
loader is written. And a natural-variant MPRA measures 5,832 variants in the promoters of 2,503
genes with the tested 144 bp windows released, which is the only row found that pins a native
variant to a measured cis effect on a named gene. The terminator panels are the weakest tier
for machine readability: 100 terminators span 0.0613 to 1.8002 relative to the \gene{PGK1t}
standard and a second study characterized over 3,000 genomic terminators, and no repository
deposit was confirmed for any terminator panel, so all of them are supplement-only.

\notesec{The expression foundation models are not the candidates, and the libraries under them are}
Two 2025 yeast models assemble a corpus, and reading their corpora returns mostly data the
graph already serves. The fungal DNA language model's supervised corpus is 5,215 tracks, of
which 1,014 are the 1,000-isolate RNA-seq the graph serves as Caudal 2024, 1,128 ChIP-exo and
20 ChIP-MNase tracks come from one 2021 deposit at GEO GSE147927, and 3,053 are new induction
RNA-seq tracks generated at Calico. The single-cell model trains on expression matrices with no
DNA sequence anywhere, and six of its named sources already have loaders. So the models supply
a triage shortlist rather than a dataset, and their three non-served items are each retrievable
on their own terms. The 3,053 new tracks are the exception that is worth an ask rather than a
parse: they are distributed only as BigWig and TFRecord on a Google Cloud Storage bucket with
no GEO, SRA, Zenodo, or Hugging Face accession anywhere in the paper, so a sha256-pinned mirror
is the only reproducible route and there is no archive fallback if the bucket moves. Two
further rows are dropped from this tier rather than ranked in it, and both are recorded in
Table~\ref{tab:excluded}. The synthetic-chromosome phenotype claims appear as figures, and the
matching deposit lists 231 samples across genome sequencing, proteome, and transcriptome with
no phenotype data type at all. The aggregate deep-mutational-scanning route needs a live query
before it can be counted, because the curated aggregate holds 217 assays over 186 proteins and
its yeast-only subset has to be read off the live service after a second aggregate's search
endpoint returned 404.

%% notes-tex/database/database-expansion-hosts/sections/pseudomonads.tex
%% [[experiments.database.expansion-hosts]]
%% https://github.com/Mjvolk3/torchcell/tree/main/notes-tex/database/database-expansion-hosts/sections/pseudomonads.tex
%% SOURCE: hand-authored prose. Per-organism experiment counts, file sizes and HTTP responses
%% come from the triage report
%% /scratch/projects/torchcell-scratch/tmp/deep-research/reports/Pk and next dataset wave triage.md
%% and from .../research_notes/Pk and next dataset wave triage/pseudomonads.md, which records
%% which figures were read on the live page and which came from a search snippet.

\section{Pseudomonads\secstatus{todo}}\label{sec:pseudo}

The row count of this wave comes from one downloadable file. The public Fitness Browser held 74
organisms when its organism list was read on 2026-11-19, of which 11 rows are pseudomonads
counting two backgrounds of one \org{Pseudomonas syringae} strain, and those 11 rows carry
1,793 genome-wide fitness experiments between them. KT2440 holds 314 of them, over 54 carbon
sources, 52 nitrogen sources, 2 stress compounds, and 27 other conditions. The nine
non-KT2440 rows hold 1,479: \org{Stutzerimonas stutzeri} RCH2 at 350, \org{Aquipseudomonas
helvetica} FW300-N2E3 at 211, FW300-N2E2 at 203, \org{A. simiae} WCS417 at 201, FW300-N2C3 at
184, FW300-N1B4 at 147, \org{A. jessenii} GW456-L13 at 106, and \org{A. fragariae} pv.
\emph{syringae} B728a at 46 in the wild type and 23 in a \gene{mexB} deletion background. The
eleventh row is \org{Pseudomonas aeruginosa} PA14 at 8 experiments and 3 other conditions, with
no carbon, nitrogen, or stress panel, so that host's presence in this resource is a token row
rather than a compendium.

Retrieval is one request and the terms are the open question. The whole-browser
\texttt{sqlite3} database at \texttt{cgi\_data/feba.db} returns HTTP 200 to a plain
\texttt{curl} at 10,433,708,032 bytes with a Last-Modified of 2026-10-27, and the companion
protein-sequence file returns 118,142,489 bytes at the same release date. The CGI layer is the
opposite: the help page and the organism list return 403 with a Cloudflare challenge header to
every non-browser client, as does the archival record page, so no license string has been read
for either the live database or the 2017 release. That is the cheapest ask in
Table~\ref{tab:asks} and it gates nine organisms. Counts for this resource also have to be
dated rather than quoted: search summaries circulate 50 organisms and 7,593 experiments, a
KBase mirror reports 48 organisms, and a separate non-public instance reports 104 organisms
with 989 KT2440 experiments. The live public page supersedes all of them and the non-public
instance's figures must never be mixed in. The resource is a rolling release and the bulk
download page states that ``data for sucrose and D-mannitol from this publication should be
disregarded'' in favor of fresh stock data, so every ingested fitness value carries the release date and
sha256 of the file it came from, and two builds months apart will not agree.

\subsection{The ready-made KT2440 matrix, and what it overlaps\secstatus{todo}}\label{sec:pseudo-kt}

Borchert et al. 2024 assembled 332 publicly available KT2440 RB-TnSeq datasets collected
between 2017 and 2022 into a 4,732 gene by 332 sample fitness matrix with t statistics and
per-sample metadata, released on GitHub with gene identifiers as \texttt{PP\_} locus tags, at
$1.6\times10^{6}$ measurements. The matrix drops the genes with no data across all samples,
4,732 of 5,564 protein-coding genes, so its shape is the released shape rather than a product
of axes. Two things about it decide the build order rather than the rank. 254 of the 332
samples come from the Fitness Browser itself, so de-duplication has to be by experiment
identity and not by count, and the remaining samples are in SRA under PRJNA809672,
PRJNA856070, and PRJNA1011287. And the paper disagrees with itself on the condition count, 179
unique conditions in the abstract and results against 183 in the methods, so neither figure can
be quoted without the other. It also states no genome assembly version for the locus tags it
uses, which is a provenance gap to close before a loader is written.

Two further KT2440 rows pair with it on the same axis. A transcriptome compendium holds 321
high-quality RNA-seq profiles selected from 541 collected samples, 305 published and 16 new, as
a 5,564 gene by 321 sample matrix over 118 unique conditions with 84 independent components, on
GitHub under an MIT license, which is structurally parallel to the fitness matrix so one
condition vocabulary serves both. Its per-sample GEO accession list sits in a supplementary
spreadsheet rather than in the text. A nitrogen BarSeq screen adds 52 nitrogen-containing
compounds plus 19 amino-acid dropout conditions, 71 in all, with significant fitness phenotypes
in 672 genes, of which 529 encode proteins unreviewed in UniProt and 256 had no significant
phenotype in previous BarSeq studies. A respiration panel of 24 strains against 205
substrates, 190 from two Biolog plates plus 15 aromatics off-plate, is the only resource that
would let this host enter as a species with intraspecific variation rather than as one
reference strain, and its data-availability statement is unread because the preprint server
rate-limited four attempts.

\subsection{The essentiality gold standards are in the wrong species\secstatus{todo}}\label{sec:pseudo-ess}

The one curated essentiality benchmark in this triage is in \org{P. aeruginosa}, not in the
host that needs it. A 2026 paper intersects five published Tn-seq studies into two gene sets,
GOLD\_115 for PA14 grown in LB and GOLD\_84 as a cross-strain core carrying PAO1 locus tags,
199 gene calls in all, shipped as a supplementary table with its own new reads at SRA
PRJNA1240204 against reference NC\_008463.1. The published essential-gene counts across that
literature are 335, 352, 434, 634, 321 core, 115 gold, and 84 core-gold, and the same paper
measures gold-standard recovery ranging from 0 to 100 percent depending on the calling tool.
The threshold and the calling tool therefore account for part of that spread alongside biology,
which makes the placement rule binding: every ingested essentiality call carries its calling
method, medium, and strain, and is never merged across studies into one boolean. GOLD\_115 is
already scoped that way by construction, PA14 in LB, so it fits the schema's essentiality layer
as released.

KT2440 is the inverse case and the gap is measured. It has dense quantitative fitness data for
its non-essential genes and no baseline genome-wide essentiality call from 2018 to 2026 at all,
because RB-TnSeq cannot score a gene it cannot mutate. The complement exists and is not
public: a 78,932-member CRISPRi library with 10 to 15 guides per gene and 798 non-targeting
controls as internal standards, screened on glucose, acetate, and p-coumarate at two
concentrations, appears only as a DOE abstract with no essential-gene count, no repository, no
accession, and no journal. The only pseudomonad CRISPRi score table plausibly downloadable
today is a PAO1 library of more than 80,000 guides over coding sequences and intergenic
regions, with more than a thousand genes showing a phenotype on depletion, deposited as three
Zenodo records whose contents are unlabeled, so which of the three holds gene-level scores is
the ask recorded against that row. CRISPRi and RB-TnSeq are not interchangeable perturbation
types for a typed schema in any case: interference is partial, polar within an operon, and
dependent on guide efficiency, so the guide sequence is part of the genotype and guide-level
and gene-level scores are different record types.

\org{P. aeruginosa}'s other contribution is sequenced isolates with measured phenotypes. One
414-isolate clinical panel carries a genome and a transcriptome per isolate with MICs by agar
dilution to the CLSI 2018 standard for four drugs, run in triplicate with up to five replicates
when discordant and an inclusion rule of three concordant results, which is about 1,656 MIC
records, and every layer sits under its own accession with PA14 locus tags on a stated
reference, so it is the one panel here that would need no provenance reconstruction. A
443-isolate longitudinal panel adds 8 phenotypes in 4 categories with the whole dataset in one
spreadsheet, and its isolate count is internally inconsistent on its own page at 474, 373, 389,
and 443, so it has to be reconciled against that spreadsheet before ingestion. A 125-isolate
survey adds 13 readouts and uses Bakta locus tags rather than PAO1 \texttt{PA} numbers, which
forces the alias layer rather than exercising it. And an arrayed screen of 5,459 mutants against
19 antibiotics would be the densest arrayed-library phenotype matrix in the species if the full
matrix was released, which is unconfirmed, making it the highest-value unresolved row of this
section.

\subsection{Admitting a new pseudomonad host\secstatus{todo}}\label{sec:pseudo-admit}

Three of the nine clear the bar of a sequenced genome plus a large, downloadable,
multi-condition phenotype set produced by one uniform pipeline: RCH2 at 350 experiments,
FW300-N2E3 at 211, and WCS417 at 201, with FW300-N2E2 at 203, FW300-N2C3 at 184, and
FW300-N1B4 at 147 behind them on the same evidence. The marginal cost of adding them is genome
and annotation registration rather than loader work, because they arrive in the same
\texttt{feba.db} artifact under the same schema and condition vocabulary, and the identifier
bridge is supplied by the source: \texttt{refseq\_mapping.tsv} maps organism and locus id to a
RefSeq accession and \texttt{refseq\_assemblies} maps the organism id to an NCBI assembly.

Two costs are real and both are policy rather than engineering. Each host needs an
\texttt{Organism} member and a reference genome in the genomes tier, so ten hosts means ten of
each, and Table~\ref{tab:schema} records that as one deliberate change rather than ten
accidental ones. And the names have moved: the browser has adopted GTDB-style renaming, so
KT2440 is listed as \org{Aquipseudomonas alloputida}, RCH2 as \org{Stutzerimonas
stutzeri}\_AE, and B728a as \org{A. fragariae} pv. \emph{syringae}, which means the organism table carries both
the historical and the current binomial as aliases or a row will not resolve to the literature
that produced it.

Orthology stays a modeling-time operation, and the evidence for that is in the resource itself.
It ships conserved functional associations and per-organism paralog tables beside the fitness
matrices, so a loader stores the released locus tag verbatim and transfer is computed
downstream, which the data-is-the-data rule requires in any case. The counter-example
is specific enough to settle the question: 20 PA14 essential genes have no PAO1 ortholog, only
192 genes are shared across all PA14 screens, and only 17 are shared across all PAO1 screens.
A load-time ortholog merge would turn each of those disagreements into a silently wrong record.

%% notes-tex/database/database-expansion-hosts/sections/pipeline.tex
%% [[experiments.database.expansion-hosts]]
%% https://github.com/Mjvolk3/torchcell/tree/main/notes-tex/database/database-expansion-hosts/sections/pipeline.tex
%% SOURCE: hand-authored prose. The mirror figures are measured by
%% experiments/database/scripts/audit_provenance_mirror.py and printed in tables/mirror.tex; the
%% quote, gap and schema-field counts come from
%% /scratch/projects/torchcell-scratch/tmp/deep-research/research_notes/Pk and next dataset wave triage/mirror_storage_audit.md,
%% and the API constraints and cost arithmetic from
%% .../research_notes/Pk and next dataset wave triage/artifact_to_ontology_pipeline.md.

\section{From artifact to ontology\secstatus{todo}}\label{sec:pipeline}

%% GENERATED FILE -- do not hand-edit.
%% SOURCE: experiments/database/scripts/audit_provenance_mirror.py
\begin{table}[H]\centering
\footnotesize
\caption[]{What the two provenance mirrors actually store, measured from the trees on disk and from every \texttt{manifest.json} in them, validated through \texttt{torchcell.literature.manifest.Manifest}. \emph{Library} is \texttt{\$DATA\_ROOT/torchcell-library/} (paper PDFs, OCR markdown, SI) and \emph{Raw} is \texttt{\$DATA\_ROOT/torchcell-raw/} (the raw dataset files a loader consumed). Each cell is the share of the population carrying the property, with the count and the population size. Rows are ordered by consequence, worst first: the first block is what a rebuild cannot reproduce, the second what a drift check cannot see, the third the manifest metadata, the fourth shape drift, the last the size of each tree. \emph{Files in no manifest} is the integrity hole: bytes no sha256 pins and no recorded command re-fetches. Byte totals sum file sizes and exclude directory entries, so they read a little under \texttt{du}. \texttt{provenance\_complete} counts manifests asserting a complete chain, so a false value counts as unfilled.}
\label{tab:mirror}
\begin{tabular}{@{}L{74mm} L{48mm} L{48mm}@{}}
\toprule
Property & Library & Raw \\
\midrule
File records with a retrieval record & 2.4\% (599/24,950) & 75.6\% (472/624) \\
Keys where every file record has one & 0.0\% (0/717) & 78.1\% (75/96) \\
Retrieval source ever re-checked (\texttt{last\_check}) & 0.0\% (0/599) & 3.6\% (17/472) \\
File records with a processing record & 1.1\% (278/24,950) & 5.1\% (32/624) \\
\texttt{paper.md} with reproducible MinerU OCR & 15.7\% (112/715) & 0.0\% (0/1) \\
Files on disk named in a manifest & 97.9\% (24,950/25,483) & 92.2\% (624/677) \\
Manifest \texttt{si\_expected} (SI completeness checklist) & 0.0\% (0/717) & 92.7\% (89/96) \\
Manifest \texttt{si\_data\_sources} & 10.5\% (75/717) & 96.9\% (93/96) \\
Manifest \texttt{zotero\_item\_key} & 99.4\% (713/717) & 7.3\% (7/96) \\
Manifest \texttt{collections} & 98.7\% (708/717) & 0.0\% (0/96) \\
Manifest \texttt{doi} & 90.0\% (645/717) & 99.0\% (95/96) \\
Manifest \texttt{provenance\_complete} true & 99.9\% (716/717) & 97.9\% (94/96) \\
File record \texttt{original\_filename} & 2.4\% (590/24,950) & 28.4\% (177/624) \\
File record \texttt{zotero\_md5} & 3.1\% (766/24,950) & 0.0\% (0/624) \\
File record \texttt{source} & 9.2\% (2,291/24,950) & 100.0\% (624/624) \\
File record \texttt{sha256} & 100.0\% (24,950/24,950) & 100.0\% (624/624) \\
Retrieval \texttt{retriever} resolves in \texttt{RETRIEVERS} & 100.0\% (599/599) & 87.5\% (413/472) \\
Retrieval \texttt{source\_url} & 100.0\% (599/599) & 99.8\% (471/472) \\
\midrule
Files in no manifest & 533 files, 0.79 GB (4.9\% of tree) & 53 files, 16.84 GB (48.2\% of tree) \\
Manifest records whose file is gone & 0 & 0 \\
Unresolvable \texttt{retriever} strings & 0 distinct, 0 records & 5 distinct, 59 records \\
File records in the pre-nesting retrieval shape & 2 records in 1 key & 0 records in 0 keys \\
Manifests with fields outside the model & 1 & 1 \\
File roles outside the documented vocabulary & 1 distinct, 2 records & 1 distinct, 1 record \\
Retrieval methods in use & 4 & 9 \\
\midrule
Bytes on disk & 16.10 GB & 34.92 GB \\
Citation-key directories & 717 & 99 \\
Keys carrying a manifest & 100.0\% (717/717) & 97.0\% (96/99) \\
File records across all manifests & 24,950 & 624 \\
Files on disk in key directories & 25,483 & 677 \\
\bottomrule
\end{tabular}
\end{table}


Table~\ref{tab:mirror} measures what the two provenance mirrors store, and the measurement
decides the order of work. In the library mirror 2.4 percent of 24,950 file records carry a
retrieval record, 0 of 717 keys carry one on \texttt{paper.pdf}, whose entire provenance is the
string \texttt{zotero:attachment:<key>}, and 15.7 percent of 715 \texttt{paper.md} files carry
the parser version, arguments, and DPI that produced them. So for 84 percent of the corpus a
re-parse would move every quote anchor with no recorded recipe to reproduce the text that was
quoted. In the raw mirror 48.2 percent of the bytes, 16.84 GB across 19 keys, sit in no
manifest at all, including the 10.2 GB fitness database of Sec.~\ref{sec:pseudo} and all 6.6 GB
of one 2025 single-cell key, and three keys carry no manifest. Nothing re-verifies either tree:
\texttt{last\_check} is set on none of the 599 library retrieval records and on 3.6 percent of
the 472 raw ones, 17 of 1,071 in all, and while \texttt{verify\_artifact} and
\texttt{check\_source} exist and are tested, no
cron, slurm, or CI job calls either, the cron list being five jobs of which none hashes or
re-fetches anything. \texttt{si\_expected}, the per-key completeness checklist, is 0 of 717 in
the library, so no automated check can report that a supplementary item is missing rather than
absent.

What survived into the built stores is the other half of the measurement. Across all 180 built
stores the entire quote corpus is 890 distinct triples of citation key, source sha256, and
verbatim quote, over 95 citation keys and 139 pinned source hashes, against 334,083 typed
\texttt{ProvenanceGap} instances. Absence is recorded at roughly 375 times the rate of
evidence, which is correct policy under the typed-absence rule and is also the measure of how
much of the ontology is unsourced. The reverse edge is thinner still: 1 of 133
\texttt{Publication} constructions sets \texttt{identifier} with a sha256 and the other 132
identify their paper by a DOI or PMID string, \texttt{ArtifactRef} appears in 5 of 123
registered loader modules and reaches 6 of 181 built stores, and \texttt{build\_manifest.json}
carries no mirror reference at all. A served record therefore cannot be resolved to the exact
bytes it was derived from, and a mirror artifact cannot be queried for the records it produced.

\subsection{The nine fields that may carry a sourced value\secstatus{todo}}\label{sec:pipeline-slot}

The binding constraint is narrower and sharper than any of those figures. The schema check
\texttt{sourced\_value\_fields()}, run live, reports exactly nine schema fields that can carry
a \texttt{SourcedValue}, and each of the nine is the \texttt{provenance} field of one class:
\texttt{Media}, \texttt{MediaComponent}, \texttt{PreCulture}, \texttt{CultureFormat},
\texttt{StrainBackground}, \texttt{EukaryoticStrainBackground}, \texttt{BackgroundAllele},
\texttt{IntegratedCassette}, and \texttt{EssentialityStatement}. All nine are media, strain
background, culture format, or essentiality. Nothing on the phenotype value, the perturbation,
or the experiment carries a quote.

That is exactly the output an extraction pipeline produces. A sample count, an uncertainty type,
a unit, and a threshold are the four values the dataset-addition rule requires to be sourced
from the paper or supplement with a verbatim quote, a sha256, and a section or line, and all
four remain loader literals whose justification never leaves the source file. The Kuzmin
combined-mutant sample count is the worked case and it is still a bare literal, with the
back-solve that produced it documented only in a Dendron note, while \texttt{StatDerivation},
the model built for exactly that resolution, is used by 2 of 123 loaders. A pipeline whose
whole value is producing those sourced values has nowhere typed to write them, so the first
thing to build is the slot. A verbatim quote also reaches the LMDB experiment record in only 27
of 181 stores, and for the other 154 it lives in the sibling
\file{preprocess/experiment_reference_index.json}, which is not the served artifact.

\subsection{The staged design\secstatus{todo}}\label{sec:pipeline-stages}

The design that follows from that audit does the unglamorous work first and puts the model only
where its output can be checked. Eight stages, two of them with no model in them at all:

\begin{itemize}[leftmargin=1.2em,itemsep=1pt,topsep=2pt]
  \item \textbf{0. Integrity floor.} No model. Bring the 16.84 GB into the manifest, replace the
    hand-rolled sidecars that no program reads, fix the seven retrieval records that cannot be
    dispatched because their recorded retriever is a shell string, a prose sentence, a
    loader-private function, or a class method, and install the missing re-verify job. A quote
    written against unhashed bytes cannot be checked against anything.
  \item \textbf{1. Parse and pin.} Record the parser name, version, arguments, DPI, and output
    sha256 per artifact. The convention already exists and 112 of 715 keys use it. The parser is
    a replaceable pinned stage rather than a choice, because the benchmark suites disagree with
    each other and the headline figures are vendor-run.
  \item \textbf{2. Grounded read.} A citations-enabled read of the mirrored markdown as a
    \texttt{text/plain} document, emitting claims with character spans into the mirrored bytes.
  \item \textbf{3. Schema-enforced structuring.} A second call that carries the schema, with the
    pydantic surface behind a cache breakpoint.
  \item \textbf{4. Ontology mapping.} A cheap lexical mapper first for recall, then a model
    adjudicating its candidates against a retrieved ontology slice rather than a dumped one.
  \item \textbf{5. Deterministic gate.} No model. Re-hash each cited span against the mirrored
    bytes, and detect truncation, refusal, cache misconfiguration, and ungroundedness from the
    response fields rather than from a judgment. This is \texttt{audit\_sourced\_value}, which
    already exists, is called from 31 of 123 loader modules, and is skipped in CI by a sentinel
    data root.
  \item \textbf{6. Adjudication.} The expensive model, on the disagreements only, so its spend
    stays proportional to the failure rate.
  \item \textbf{7. Human review.} Unavoidable, and quantified below.
\end{itemize}

Three API constraints determine where those stage boundaries fall. Citations and structured
outputs are mutually exclusive in one request and return a 400 together, so grounding and
schema enforcement are two calls and not one. Native PDF input is vision plus per-page
extracted text rather than an OCR service, and it cannot cite images in a PDF at all, so a
figure-borne number cannot be grounded that way. Plain text chunks into sentences and yields
character spans into bytes the mirror pins, where a PDF yields page spans into bytes the
publisher controls. And the schema handed to the API is not the pydantic model verbatim,
because recursive schemas, \texttt{oneOf}, \texttt{not}, numerical constraints, and
string-length constraints each return a 400, so the SDK helpers strip the unsupported
constraints and re-validate locally. A supplementary spreadsheet is rejected outright in a
document block, so it becomes CSV text, which keeps the citation span, or an upload into code
execution, which loses it.

\subsection{The cost model, and its assumptions\secstatus{todo}}\label{sec:pipeline-cost}

The figures below are arithmetic over the documented prices, and the assumptions are stated
because the result moves with them:

\begin{itemize}[leftmargin=1.2em,itemsep=1pt,topsep=2pt]
  \item 300 papers, each a 12-page main PDF plus a 20-page supplement plus about 30,000 tokens
    of converted supplementary tables.
  \item 3,500 tokens per page, the midpoint of the documented 1,500 to 3,000 text tokens plus
    up to 1,568 visual tokens, which is about 150,000 input tokens of artifact per paper.
  \item A 20,000-token schema prefix, cached.
  \item Outputs of 8,000, 4,000, 2,000, and 15,000 tokens for the four model stages.
  \item The adjudication stage firing on 20 percent of papers, and a 30 percent retry
    overhead on everything.
\end{itemize}

On those assumptions 300 papers cost about \$120 to \$175 with a mid-tier worker and a
top-tier adjudicator run in batch, and the planning figure is \$190 to \$250. Running every
stage on the top-tier model synchronously at high effort costs about \$810 to \$1,000 for the
same work, roughly four to five times as much. Input tokens on the artifact are 70 to 80
percent of the bill in every variant, so one design detail is load-bearing. Caching the
20,000-token schema prefix saves about \$0.02 per paper, while caching the 150,000-token
document between stages saves about \$0.14, about a third of the total. The documentation states
that an ephemeral cache entry written during batch processing would likely expire before the
follow-up request runs, so a batch-based multi-pass design cannot count on cross-batch hits and
either puts all passes for one paper in a single request or budgets the higher figure.
Batch-first is a throughput choice before it is a cost one, since the batch API is documented
as exempt from the ramp limits that return a 429 on a burst even while under the token limits,
and cache reads do not count toward the input-token limit either.

The accuracy regimes are the number worth planning against. Published pipelines of this shape
land at 0.94 to 0.98 F1, with one entity and relation extraction at 0.96 and 0.94 and a
protein-degradation workflow at record-level F1 0.98 with 92 percent and 82.5 percent of
recovered records validated on expert review. A bare model on strict whole-record extraction
lands at 55 to 60 percent, measured at 57 percent and 56 percent for two current general
models on a nanocrystal-synthesis benchmark. At 300 papers and about 20 fields each, the good
regime still implies 120 to 360 wrong fields, so the pipeline's value is in making a wrong
field findable rather than in removing the review queue. The gap between the two regimes is the
engineering and not the model.

Two targeting decisions follow from the audit rather than from the API. The pipeline points at
the raw tree and this wave's keys rather than at the library mirror, because 124 of 717 library
keys are named anywhere in \file{torchcell/} source and only 95 are quoted by a built store,
so roughly 554 keys and about 12 GB are reading material served through the literature API
rather than substrate the ontology depends on, while 96 of 99 raw keys are named in code. And
the cheapest genuinely useful first output is not new records: \texttt{ProvenanceGap.resolve\_with},
the field that names the artifact which would close a deferred gap, is populated on 0.8 percent
of gaps against 4,933 gaps typed \texttt{deferred\_pending\_source\_review}. Filling it needs
no schema change, has a deterministic check, and turns an existing 4,933-item worklist into an
actionable one.

\subsection{What this wave needs changed\secstatus{todo}}\label{sec:pipeline-schema}

Table~\ref{tab:schema} lists every change in the served vocabulary and the candidate machinery
that the rows of this document require, with the policy each one carries. Two of them gate
every row rather than one: a host token the gate can resolve, without which
\texttt{hosts\_named()} raises on any \org{P. kudriavzevii} row so that G2 cannot even be
measured, and an \texttt{Organism} member, without which no loader can declare the host and the
adapter map can hold nothing for it. The most expensive entry is the slot of
Sec.~\ref{sec:pipeline-slot}: extending \texttt{SourcedValue} to the phenotype value and the
perturbation changes classes that many served closures hold, so its rebuild scope reads
\texttt{full\_kg\_rebuild} and its record carries the \texttt{schema-breaking} label. Every
other entry on the table is additive, meaning a new leaf or a new member that changes no record
any served dataset already holds.

%% GENERATED FILE -- do not hand-edit.
%% SOURCE: experiments/database/scripts/build_hosts_candidate_datasets_table.py
\begingroup
\footnotesize
\begin{longtable}{@{}L{50mm} L{48mm} L{74mm}@{}}
\caption[]{Every schema or candidate-machinery change this wave requires, with the
symbol it lands on and the policy it triggers. \emph{Symbol} is the exact module path
and name, so the change is addressable rather than described. A change marked additive
adds a class or a vocabulary member that only its own closure holds, so it enters the
served graph by incremental admission; a change to a class other closures hold is a full
rebuild of every dataset whose closure holds it. The first two are the critical path:
without them no loader can declare a host and the candidate gate cannot even be
measured.}
\label{tab:schema}\\
\toprule
Symbol & What changes & Policy it triggers \\
\midrule
\endfirsthead
\toprule
Symbol & What changes & Policy it triggers \\
\midrule
\endhead
\bottomrule
\endfoot
\file{torchcell/datamodels/schema.py::Genotype (perturbation identity)} & Admit a perturbation whose identity is NOT a gene. Eleven of the twenty largest yeast rows in this fill, and most of the E. coli regulatory-DNA rows, carry no gene identifier anywhere in their bulk data: the record key is a lineage barcode, a raw measured sequence, or a residue substitution string. A barcoded lineage is the hardest case, because the perturbation is UNKNOWN by construction and is identified only by its tag, which is a typed absence on the genotype rather than a missing field & Served-vocabulary change. The honest route is a typed ProvenanceGap on the gene-identity field rather than a forced mapping to a gene, which would invent a locus the release never names. Schema-change record and impact report; additive if it lands as a new perturbation sibling, breaking if it loosens a required field on the served Genotype Not additive. \\
\addlinespace[5pt]
\file{torchcell/datamodels/schema.py::SourceType} & Add a patent member. The enum holds journal\_article, dissertation, preliminary\_report and in\_house, and the identity rule for every non-journal member is title plus the deposited document's mirror-relative path and sha256. A patent fits that rule exactly and has no DOI, so the three ranked 3-HP patents can be deposited and hashed but not yet declared. Without the member the only honest alternatives are to drop eleven strains with titer and yield, or to file a patent as a journal article, which would be false & Served-vocabulary change, but additive: no served record carries the new member, so the closure of every built dataset is unchanged. Schema-change record and impact report; no schema-breaking label and no rebuild, since nothing served reserializes Additive. \\
\addlinespace[5pt]
\file{torchcell/data/experiment_dataset.py::Organism} & Add a pkudriavzevii member. The enum holds exactly scerevisiae, ecoli and pputida today, and ALL\_ORGANISMS is frozenset(Organism), so no loader can declare the host and dataset\_adapter\_map can hold nothing for it & Served-vocabulary change: schema-change record, impact report, the served datasets named in the PR, and the schema-breaking label if a served closure changes. The genome side is already done, which is what makes this the whole cost: registry.py pins PKUDRIAVZEVII\_SD108 = pkudriavzevii\_SD108\_v1 and gates.py already lists it in READABLE\_ASSEMBLY\_SETS with the comment that no host token names it yet. Additive. \\
\addlinespace[5pt]
\file{torchcell/candidates/gates.py::HOSTS, _ORGANISM_OF_TABLE_ROW} & Add HostTokens for SD108 and CBS573, and admit 'P. kudriavzevii' to the set of organisms a table row may name, which holds E. coli, P. putida and S. cerevisiae today & Gate-code change, not a schema change. Until it lands hosts\_named() raises ValueError on any P. kudriavzevii row, so G2 cannot even be measured and no verdict can be written for one. The eight pseudomonad hosts need the same treatment, one token per reference strain. Additive. \\
\addlinespace[5pt]
\file{torchcell/data/experiment_dataset.py::Organism (pseudomonads)} & Nine or ten further members for the Fitness Browser organisms and P. aeruginosa, plus the matching reference genomes in the genomes tier & The same served-vocabulary policy, ten times over. Worth staging as one deliberate change rather than ten accidental ones, because the Fitness Browser rows arrive in one artifact and the marginal work per organism after the first is a genome registration. Additive. \\
\addlinespace[5pt]
\file{torchcell/candidates/verdict.py::CandidateTable} & A third literal for this pool. The type is Literal['bacteria', 'yeast'] and it is also the table field of every one of the 52 stored verdicts & Record-model change on a frozen, extra-forbidding model, so it is a policy decision rather than a widening: until it lands, every row in Table 11 reads 'no verdict' and candidate-gate can read this pool's rows but cannot write a verdict for one. gates.py::CandidateRow.table and gates.py::TEXT\_FIELDS carry the same literal and the same keys. Not additive. \\
\addlinespace[5pt]
\file{torchcell/candidates/cli.py::TABLE_SCRIPTS} & Register this script as a third pool, so candidate-gate can load its rows by name the way it loads the bacteria and yeast pools & CLI change only. Done. The gate subcommand's --table choices still admit two values, so passing this pool's name needs the CandidateTable decision above first. Additive. \\
\addlinespace[5pt]
\file{experiments/database/scripts/build_candidate_datasets_table.py::CANDIDATES} & A separate pool rather than more rows in the yeast one. The yeast Candidate has no organism field at all, the bacteria one admits only E. coli and P. putida, and the yeast main() raises on a pool over 250 rows against 169 literals held & Splitting by wave is forced by the guard, not merely tidier. This script is that split, and its own guard is 400 because the pool spans eight hosts. Additive. \\
\addlinespace[5pt]
\file{experiments/database/scripts/build_hosts_candidate_datasets_table.py::SeqBasis} & Host-specific sequence bases. The yeast SeqBasis is eleven values, every one S288C- or yeast-panel-relative, so isolate-WGS transfers to a Pichia or pseudomonad isolate panel and an SD108 designed edit does not & Script-level vocabulary, enforced as a validation error on any unlisted value. Declared explicitly per host so a bacterial row cannot be written against a yeast-shaped basis by accident. Additive. \\
\addlinespace[5pt]
\file{torchcell/datamodels/schema.py (phenotype leaves)} & New leaves for TSS initiation efficiency, antifungal and antibiotic MIC, Biolog respiration per well, plate-reader reporter ratio, and 13C flux & By the placement rule a different technology is its own leaf with an explicit conversion, so a plate-reader GFP ratio and a sort-seq bin average are different leaves rather than one class with a note. Each is additive and admissible on its own and none weakens an existing required field. Additive. \\
\addlinespace[5pt]
\file{torchcell/datamodels/schema.py (perturbation leaves)} & New leaves for a 5' UTR variant, a 3' end variant, a promoter swap upstream of a native gene, a tRNA point substitution, and an amino-acid substitution for the ProteinGym subset & Additive leaves, so each should clear the admission check on its own closure. In the served graph the perturbation subclasses collapse into one Perturbation node whose kind is a property, so a query naming every deletion names the kinds one by one. Additive. \\
\addlinespace[5pt]
\file{torchcell/datamodels/schema.py (GenePerturbation namespaces)} & Gene-identifier namespaces for a JGI gNNNN model id, an INSDC SD108 locus tag, a PA or PA14 number, a Bakta locus tag and a Fitness Browser locusId & The same additive change the bacterial layer already made for b-numbers and PP\_ tags. The genome never gates a record, so a released name the reference does not hold is a report-class verification finding and never a dropped row. Additive. \\
\addlinespace[5pt]
\file{torchcell/verification/sourced.py::SourcedValue} & Extend it to the phenotype value and the perturbation. ontology\_checks.sourced\_value\_fields() reports exactly nine fields that can carry one today, all of them media, strain background, culture format or essentiality, so a quoted n\_samples, uncertainty type, unit or threshold has nowhere typed to be written & The most expensive change here: it moves classes many served closures hold, so the record's rebuild\_scope reads full\_kg\_rebuild and the PR carries the schema-breaking label. It is also the blocking change for an artifact-to-ontology pipeline, whose entire output is that class of value. Not additive. \\
\addlinespace[5pt]
\file{notes/paper.north-star.dataset-triage.md (guide-level exclusion)} & An explicit decision on gRNA-library rows. The standing rule excludes a dataset whose primary label is per-guide activity or cleavage efficiency rather than a per-strain phenotype & The landing-pad screen's cutting and integration efficiency fall under the rule while its fitness and GFP readouts do not, and the PAO1 and KT2440 CRISPRi libraries are guide-level by construction, where the guide sequence is part of the genotype. This needs a decision, not an inherited rule. Additive. \\
\addlinespace[5pt]
\file{experiments/database/scripts/build_candidate_datasets_table.py::EXCLUDED} & Revisit the standing off-species exclusion of the K. phaffii screen and the Y1000+ panel, whose note reads 'Revisit if the generalization axis widens to other hosts' & A cross-host wave is that widening, so the note needs a decision. The iIsor850 exclusion stands on its own merits under not-a-dataset: a reconstruction is not an experiment. Additive. \\
\addlinespace[5pt]
\file{$DATA_ROOT/torchcell-raw/priceMutantPhenotypesThousands2018} & Bring the key into the manifest before extending it. Six of its nine files are unhashed over 10.21 GB, and a parallel non-pydantic convention sits beside it: DEPOSIT.md, feba.db.sha256 and SHA256SUMS.txt, which no program reads & A precondition of the cross-host fitness band rather than a schema change: those rows extend this key, and an ingested fitness value has to carry the release date and sha256 of the db file it came from, because the bulk page itself says the 2018 release's sucrose and D-mannitol data should be disregarded in favor of fresh stock data. Additive. \\
\addlinespace[5pt]
\end{longtable}
\endgroup


%% notes-tex/database/database-expansion-hosts/sections/final.tex
%% [[experiments.database.expansion-hosts]]
%% https://github.com/Mjvolk3/torchcell/tree/main/notes-tex/database/database-expansion-hosts/sections/final.tex
%% SOURCE: generated tables only; each table file names its own generating script in its header.

\section{The first fifty, and ten extra\secstatus{todo}}\label{sec:list}

%% GENERATED FILE -- do not hand-edit.
%% SOURCE: experiments/database/scripts/build_hosts_candidate_datasets_table.py
\begin{landscape}
\begingroup
\footnotesize
\setlength{\tabcolsep}{3pt}
\renewcommand{\arraystretch}{1.15}
\begin{longtable}{@{}r@{\hspace{3pt}} L{38mm} L{11mm} L{14mm} L{18mm} L{20mm} r@{\hspace{4pt}} r@{\hspace{4pt}} c@{\hspace{4pt}} L{20mm} L{66mm}@{}}
\caption[]{The recommended first fifty of the cross-host wave, then ten extra in rank
order in case a row above proves unreachable. \emph{Band, tier}: \emph{R} is regulatory
DNA, the one axis that adds measurements rather than rows; \emph{F} is cross-host
fitness, the Fitness Browser rows that supply the row count; \emph{H} is host bootstrap,
which exists because \org{P. kudriavzevii} cannot win a measurement contest against a
pseudomonad fitness compendium under any ordering; \emph{S} is scale. The band is applied
before tier, and inside a band the tier rule orders, then measurements. \emph{Genotypes} and \emph{Env} are the perturbation and condition axes.
\emph{Inst.} is genotype$\times$environment records, $\dagger$ where it is the product of
the two axes rather than a reported count and $\ddagger$ where it is an
order-of-magnitude estimate. \emph{Meas.} is instances times phenotype dimensionality,
the quantity rows are ranked on. \emph{Sequence basis} is the route to each strain's
total genomic content; a row with no route is excluded rather than ranked
(Table~\ref{tab:excluded}). Superscript \textbf{M} marks a row whose values need a human with a browser and \textbf{R} a row that needs a request to a named party (Table~\ref{tab:asks}); \textbf{B} marks a row whose per-record values are not released and \textbf{A} a corpus that re-serves other papers. \emph{Ver.} carries the triage report's own verification label for the counts: \textbf{V} read on the source in session, \textbf{S} a search snippet only, \textbf{I} inferred, which includes every count that is the product of two verified numbers. Citations, links and data locations are in
Table~\ref{tab:sources}; the gene-identifier namespace each host's ids speak is in
Table~\ref{tab:namespaces}; what the wave changes in the schema is in
Table~\ref{tab:schema}.}
\label{tab:final}\\
\toprule
\textbf{\#} & \textbf{Dataset (class; phenotype)} & \textbf{Band, tier} & \textbf{Host} & \textbf{Genotypes} & \textbf{Env} & \textbf{Inst.} & \textbf{Meas.} & \textbf{Ver.} & \textbf{Sequence basis} & \textbf{Why} \\
\midrule
\endfirsthead
\multicolumn{11}{@{}l}{\footnotesize\emph{Table~\ref{tab:final}, continued}}\\
\toprule
\textbf{\#} & \textbf{Dataset (class; phenotype)} & \textbf{Band, tier} & \textbf{Host} & \textbf{Genotypes} & \textbf{Env} & \textbf{Inst.} & \textbf{Meas.} & \textbf{Ver.} & \textbf{Sequence basis} & \textbf{Why} \\
\midrule
\endhead
\bottomrule
\endfoot
1 & \textbf{de Boer 2020 (GPRA sort-seq, 100M random promoters)}\newline {\scriptsize Regulatory DNA; expression driven, 18-bin sort-seq YFP/RFP} & R, 1 & \org{S. cerevisiae} & 102,371,025 promoter sequences & YPD, YPGal, YPGly & $1.0\times 10^{8}$ & $1.0\times 10^{8}$ & V & plasmid-\allowbreak reporter-\allowbreak library & The largest sequence-to-phenotype resource in yeast by two orders of magnitude, and every genotype is an exactly known 80mer at a known position in a known scaffold. The perturbation space is continuous rather than a gene list, which is what a model of regulatory grammar needs and no deletion collection can supply. \\
\addlinespace[5pt]
2 & \textbf{Vaishnav 2022 deferred training libraries}\newline {\scriptsize Regulatory DNA; expression driven, 18-bin sort-seq YFP} & R, 1 & \org{S. cerevisiae} & 30,722,376 complex + 20,616,659 defined & YPD complex, SD-\allowbreak Ura defined & $5.1\times 10^{7}$ & $5.1\times 10^{7}$ & V & plasmid-\allowbreak reporter-\allowbreak library & The cleanest possible superset admission: the loader, the schema, the medium objects, the strain reasoning and the provenance record all exist, and the 71,158 served records are a proven subset, so the admission check can prove the superset by content id. The two files are 3.59 GB and 2.47 GB on Zenodo, which makes it a download decision rather than a schema one. \\
\addlinespace[5pt]
3 & \textbf{DREAM 2022 random promoter challenge (Rafi 2024)}\newline {\scriptsize Regulatory DNA; expression driven, FACS, about 100 cells per sequence} & R, 1 & \org{S. cerevisiae} & 6,739,258 train + 71,103 test & 1 medium & $6.8\times 10^{6}$ & $6.8\times 10^{6}$ & V & plasmid-\allowbreak reporter-\allowbreak library & The held-out benchmark every published yeast expression model is scored on, and the only count in this wave with two independent confirmations. It is a curated slice of the same GPRA platform rather than an independent experiment, so it needs an AggregationRecord at G1 and partial nesting with de Boer and Vaishnav is expected. \\
\addlinespace[5pt]
4 & \textbf{Kotopka and Smolke 2020 (artificial promoter libraries)}\newline {\scriptsize Regulatory DNA; promoter activity} & R, 1 & \org{S. cerevisiae} & over 675,000 constitutive + over 327,000 inducible & constitutive and inducible arms & $1.0\times 10^{6}$ & $1.0\times 10^{6}$ & V & plasmid-\allowbreak reporter-\allowbreak library & A million-scale synthetic promoter library from a platform and a group independent of the de Boer lineage, so it is not nested with the three rows above, and it carries an inducible arm that the GPRA libraries lack. Both counts are lower bounds as the paper words them. \\
\addlinespace[5pt]
5 & \textbf{Zhu 2024 (Pol II MASTER, TSS initiation efficiency)}\newline {\scriptsize Regulatory DNA; transcription initiation efficiency per TSS} & R, 1 & \org{S. cerevisiae} & 1,000,000 TSS sequences & defined promoter context, NTP levels varied & $1.0\times 10^{6}$ & $1.0\times 10^{6}$ & V & plasmid-\allowbreak reporter-\allowbreak library & A million exactly specified TSS-window sequences read out mechanistically rather than through a reporter, deposited as both processed values and raw reads. The 2021 preprint of the same work describes about 80,000 variants, so the published figure is the one to pin and the discrepancy travels with the row. \\
\addlinespace[5pt]
6 & \textbf{Cuperus 2017 (random 5' UTR library)}\newline {\scriptsize Regulatory DNA; competitive growth enrichment, HIS3} & R, 1 & \org{S. cerevisiae} & 489,348 detected random 50-\allowbreak nt 5' UTRs & medium lacking histidine & $4.9\times 10^{5}$ & $4.9\times 10^{5}$ & V & S288C+\allowbreak designed-\allowbreak edit & Half a million exactly specified 5' UTR sequences, and it changes both axes at once: the sequence variable moves from the promoter to the 5' UTR and the readout moves from a sorted reporter to competitive growth. A separate leaf by the measurement rule, since growth and a bin average are different technologies. \\
\addlinespace[5pt]
7 & \textbf{Gui 2021 (position effects with transcriptomes)}\,\textsuperscript{\textbf{M}}\newline {\scriptsize Regulatory DNA; full transcriptome per integration locus} & R, 2 & \org{S. cerevisiae} & about 250 GFP integration loci & 1 & 250 & $1.5\times 10^{6}$ & I & S288C+\allowbreak integration-\allowbreak locus & The only position-effect panel with a gene-level readout rather than one reporter value, which is what makes the integration coordinate a usable sequence variable. Changes in neighboring genes rather than in the reporter itself altered growth rate, and no accession was located. \\
\addlinespace[5pt]
8 & \textbf{Keren 2016 (promoter panel x named gene, pooled fitness)}\newline {\scriptsize Regulatory DNA; pooled growth fitness per strain} & R, 2 & \org{S. cerevisiae} & 130 promoters over a 500-\allowbreak fold range x about 81 genes & glucose, galactose & $2.0\times 10^{4}$$\dagger$ & $2.0\times 10^{4}$ & I & S288C+\allowbreak designed-\allowbreak edit & A sequence-level promoter swap at a named native locus against a per-strain fitness, which is the shape the existing typed record fits most closely of anything in this band. The gene count is reported inconsistently as 81, 96 and about 100 across the paper, the commentary and the two deposit records, so it is unresolved and carried as such. \\
\addlinespace[5pt]
9 & \textbf{Shalem 2015 (3' end FACS-seq library)}\,\textsuperscript{\textbf{M}}\newline {\scriptsize Regulatory DNA; YFP/mCherry ratio, 16-bin FACS-seq} & R, 2 & \org{S. cerevisiae} & 13,000 designed 3' ends, 102 bp variable & SC-\allowbreak URA + galactose, 30 C & $1.3\times 10^{4}$ & $1.3\times 10^{4}$ & V & plasmid-\allowbreak reporter-\allowbreak library & Sequences and measured values sit in one XLSX, it extends coverage from the promoter to the terminator, and it is one of the eight benchmark tasks published models are scored on. Its only uncertainty figure is a library-wide 13.2 percent technical RSD rather than a per-construct error, and no raw reads were deposited anywhere. \\
\addlinespace[5pt]
10 & \textbf{Renganaath 2020 (native promoter variant MPRA)}\newline {\scriptsize Regulatory DNA; MPRA mRNA/DNA allelic cis effect} & R, 2 & \org{S. cerevisiae} & 5,832 variants in 2,503 gene promoters & 1 & $5.8\times 10^{3}$ & $5.8\times 10^{3}$ & V & plasmid-\allowbreak reporter-\allowbreak library & The only dataset found that pins a native variant to a measured cis effect on a named gene with the tested 144 bp window released, which is what joins the synthetic libraries to the natural-isolate panels. 451 variants were called causal, 359 were tested twice in opposite orientation and 527 in both library designs, so the measurement count exceeds the variant count. \\
\addlinespace[5pt]
11 & \textbf{Arg-CCU tRNA deep mutational scan}\,\textsuperscript{\textbf{M}}\newline {\scriptsize Deep mutational scan; fitness under selection} & R, 2 & \org{S. cerevisiae} & 4,176 of 5,184 genotypes after filtering & high temperature, 1 M NaCl & $4.2\times 10^{3}$ & $4.2\times 10^{3}$ & V & S288C+\allowbreak designed-\allowbreak edit & A regulatory-element deep mutational scan, complete over a 14-substitution combinatorial space with a direct fitness readout, and outside the protein-substitution scope of the aggregates. It needs its own loader for exactly that reason. \\
\addlinespace[5pt]
12 & \textbf{Keren 2013 (native promoter panel across conditions)}\,\textsuperscript{\textbf{M}}\newline {\scriptsize Regulatory DNA; promoter activity, YFP with a fixed ADH1 terminator} & R, 3 & \org{S. cerevisiae} & about 900 native promoters & 10 conditions & $9.0\times 10^{3}$$\dagger$ & $9.0\times 10^{3}$ & S & plasmid-\allowbreak reporter-\allowbreak library & Native promoters measured across ten conditions is the only genuine environment axis on native regulatory sequence in this band, and it is one of the eight published benchmark tasks. Both counts reach this table through two secondary sources rather than the paper, and one validation used 625 of the promoters. \\
\addlinespace[5pt]
13 & \textbf{Lubliner 2015 (core promoter variant panel)}\,\textsuperscript{\textbf{M}}\newline {\scriptsize Regulatory DNA; YFP/mCherry ratio, 16-bin FACS} & R, 3 & \org{S. cerevisiae} & 7,536 core promoter variants, several hundred native & 1 & $7.5\times 10^{3}$ & $7.5\times 10^{3}$ & V & plasmid-\allowbreak reporter-\allowbreak library & Isolates the core promoter as the sole variable under a constant strong upstream region, which no other row in the band does. The per-variant values are reported to be in Supplemental Table 1, which was not opened, so the release is very likely rather than confirmed. \\
\addlinespace[5pt]
14 & \textbf{Sharon 2012 (designed promoter library)}\,\textsuperscript{\textbf{M}}\newline {\scriptsize Regulatory DNA; expression, sorting plus barcode sequencing} & R, 3 & \org{S. cerevisiae} & 6,500 designed promoters, 55 percent Gal4/Gcn4 sites & GAL1-\allowbreak 10 and HIS3 reporter backgrounds & $6.5\times 10^{3}$ & $6.5\times 10^{3}$ & V & plasmid-\allowbreak reporter-\allowbreak library & The origin of the GPRA vector the three largest rows use, and the nearest thing to a synthetic transcription-factor library in yeast, which the triage records as a genuine gap rather than a missed search. Designed rather than random, so site number, affinity, spacing and orientation vary one at a time. \\
\addlinespace[5pt]
15 & \textbf{Yamanishi 2013 (genomic terminatome)}\,\textsuperscript{\textbf{M}}\newline {\scriptsize Engineering parts; GFP output under a fixed promoter} & R, 3 & \org{S. cerevisiae} & over 3,000 genomic terminators & fixed TDH3 promoter, GFP & $3.0\times 10^{3}$$\ddagger$ & $3.0\times 10^{3}$ & S & plasmid-\allowbreak reporter-\allowbreak library & The largest terminator panel found and one of the eight published benchmark tasks. The count is reported as over 3,000 by one route and over 5,000 by another, with a validation set of 4,005, so the figure is unresolved and the primary table was never located. \\
\addlinespace[5pt]
16 & \textbf{Dvir 2013 (5' UTR positions -10 to -1)}\,\textsuperscript{\textbf{M}}\newline {\scriptsize Regulatory DNA; reporter fluorescence} & R, 3 & \org{S. cerevisiae} & 2,041 5' UTR variants at positions -\allowbreak 10 to -\allowbreak 1 & 1 & $2.0\times 10^{3}$ & $2.0\times 10^{3}$ & V & plasmid-\allowbreak reporter-\allowbreak library & The cleanest minimal-window 5' UTR panel available: ten bases immediately before the start codon, with the -3 to -1 context dominant, then secondary structure, then out-of-frame start codons, explaining two-thirds of expression variation. A second independently constructed library was held out as a test set. \\
\addlinespace[5pt]
17 & \textbf{Wu 2017 (genome-wide position effects, 1,044 loci)}\,\textsuperscript{\textbf{M}}\newline {\scriptsize Regulatory DNA; RFP intensity, plate reader} & R, 3 & \org{S. cerevisiae} & 1,044 evenly spaced integration loci & promoter, reporter and carbon-\allowbreak source control arms & $1.0\times 10^{3}$ & $1.0\times 10^{3}$ & V & S288C+\allowbreak integration-\allowbreak locus & The largest position-effect panel, with a greater than 13-fold range and location dominating promoter, reporter and carbon source. Each strain also carries a knockout-collection deletion, so the genotype is a deletion plus an insertion rather than a regulatory element, and Table S2 is stated to be available from the corresponding author on request. \\
\addlinespace[5pt]
18 & \textbf{ProteinGym yeast subset}\newline {\scriptsize Deep mutational scan; per-assay deep mutational scanning score} & R, 3 & \org{S. cerevisiae} & yeast subset of 217 assays over 186 proteins; yeast count needs a live query & one selection per assay & --$\ddagger$ & -- & V & S288C+\allowbreak designed-\allowbreak edit & The cheapest route to coverage of the yeast mutational-scan tier: filtering to the YEAST identifier token isolates the subset with uniform formatting and a fixed reference per assay and no per-paper retrieval. The yeast-only count was not obtained, so it is reported as needing a live query rather than as a number. \\
\addlinespace[5pt]
19 & \textbf{Zhang 2016 (genome-wide GFP position scan)}\,\textsuperscript{\textbf{M}}\newline {\scriptsize Regulatory DNA; GFP by FACS} & R, 4 & \org{S. cerevisiae} & 428 high-\allowbreak  and 444 low-\allowbreak expressed sites & 1 & 872 & 872 & V & S288C+\allowbreak integration-\allowbreak locus & A genome-wide integration-site scan whose reported sites split cleanly into high and low classes, adding a FACS readout to a tier that is otherwise plate-reader. The total number of sites screened is not stated and the journal is Chinese-language, so a browser fetch comes first. \\
\addlinespace[5pt]
20 & \textbf{Kita 2017 (context-stratified yeast eQTLs)}\,\textsuperscript{\textbf{M}}\newline {\scriptsize Natural variation; eQTL effect on expression} & R, 4 & \org{S. cerevisiae} & 683 selected of 1,640 eQTLs & not reported & 683 & 683 & V & isolate-\allowbreak WGS & The only eQTL set found that is stratified by regulatory context, across promoter, 5' UTR, 3' UTR and open reading frame, which makes it usable as a variant-effect evaluation set rather than a mapping table. It is a summary table rather than per-record measurements, so it is a borderline row and is ranked as one. \\
\addlinespace[5pt]
21 & \textbf{Chen (earlier position-effect study, 482 loci)}\,\textsuperscript{\textbf{M}}\newline {\scriptsize Regulatory DNA; reporter expression per locus} & R, 4 & \org{S. cerevisiae} & 482 loci, about a third on chromosomes I and VI & not reported & 482 & 482 & S & S288C+\allowbreak integration-\allowbreak locus & A third independent position-effect panel, which matters because the tier's whole claim is that chromosomal location dominates the designed sequence. The count reaches this table only through Wu 2017's citation of it, so the primary paper has to be identified before the row can be filled. \\
\addlinespace[5pt]
22 & \textbf{Wei 2017 (synthetic terminator panel)}\,\textsuperscript{\textbf{M}}\newline {\scriptsize Engineering parts; eGFP strength relative to the PGK1t standard} & R, 4 & \org{S. cerevisiae} & 100 terminators, 45 strong, 31 moderate, 24 weak & 1 & 100 & 100 & V & plasmid-\allowbreak reporter-\allowbreak library & A parts panel whose values are fold-strengths against a named standard part, spanning 0.0613 to 1.8002 relative to PGK1t, with the linker sequence as the designed variable. No repository deposit was found for any terminator panel, so every one of them is supplement-only. \\
\addlinespace[5pt]
23 & \textbf{Lubliner 2013 (native core promoter maximal strengths)}\,\textsuperscript{\textbf{M}}\newline {\scriptsize Regulatory DNA; native core promoter maximal strength} & R, 4 & \org{S. cerevisiae} & native core promoters, count not reported & not reported & --$\ddagger$ & -- & S & plasmid-\allowbreak reporter-\allowbreak library & The primary source behind one of the published core-promoter models, measured over a 90 bp window from -70 to +20 of the transcription start site. No count, accession or deposit was located, so the row exists to keep the source identified rather than to be built next. \\
\addlinespace[5pt]
24 & \textbf{Morse (short synthetic terminators)}\,\textsuperscript{\textbf{M}}\newline {\scriptsize Engineering parts; fluorescent protein output} & R, 4 & \org{S. cerevisiae} & short synthetic terminators, 35 to 70 bp, count not reported & 1 & --$\ddagger$ & -- & S & plasmid-\allowbreak reporter-\allowbreak library & The best short-terminator panel named anywhere in the triage, with the best member reported at 3.7-fold over CYC1. The source is the weakest in the whole wave, an author aggregator page, so none of its numbers may be promoted and the primary paper has to be identified first. \\
\addlinespace[5pt]
25 & \textbf{MaveDB yeast score sets}\,\textsuperscript{\textbf{M}}\newline {\scriptsize Deep mutational scan; variant effect score} & R, 4 & \org{S. cerevisiae} & yeast score sets; count needs a live query & one selection per score set & --$\ddagger$ & -- & S & S288C+\allowbreak designed-\allowbreak edit & The second curated aggregate into the same tier, with an organism filter already built in and over a million variant effect scores across the whole database. The documented search endpoint returns 404, so the yeast count has to come from the live site and the target count is itself reported inconsistently as 39 and 45. \\
\addlinespace[5pt]
26 & \textbf{Fitness Browser RCH2 (Stutzerimonas stutzeri\_AE)}\newline {\scriptsize Fitness / chemical genomics; BarSeq competition fitness per gene, with a t statistic} & F, 1 & \org{S. stutzeri} & pooled barcoded transposon library, gene count not reported per organism & 30 carbon, 19 nitrogen, 50 stress, 38 other & 350 & $1.6\times 10^{6}$ & I & fitbrowser+\allowbreak transposon & The largest single-organism experiment count in the browser, 350 experiments with the broadest stress panel of any pseudomonad in it, and the strain the historical name P. stutzeri points at. It arrives in the same feba.db as every other row in this band, so its marginal cost is one Organism member and one reference genome. \\
\addlinespace[5pt]
27 & \textbf{Fitness Browser FW300-N2E3 (Aquipseudomonas helvetica)}\newline {\scriptsize Fitness / chemical genomics; BarSeq competition fitness per gene, with a t statistic} & F, 1 & \org{P. fluorescens} & pooled barcoded transposon library, gene count not reported per organism & 46 carbon, 30 nitrogen, 42 stress, 12 other & 211 & $9.6\times 10^{5}$ & I & fitbrowser+\allowbreak transposon & The second-largest pseudomonad row in the browser and the most balanced of the four FW300 isolates across all four condition classes. It also carries 115 updated annotations, which is the browser's own measure of how much a row taught about unannotated genes. \\
\addlinespace[5pt]
28 & \textbf{Fitness Browser FW300-N2E2 (Aquipseudomonas fluorescens\_Q)}\newline {\scriptsize Fitness / chemical genomics; BarSeq competition fitness per gene, with a t statistic} & F, 1 & \org{P. fluorescens} & pooled barcoded transposon library, gene count not reported per organism & 41 carbon, 25 nitrogen, 34 stress, 7 other & 203 & $9.2\times 10^{5}$ & I & fitbrowser+\allowbreak transposon & A sibling FW300 isolate at 203 experiments, which is what makes the block a comparative resource rather than nine unrelated screens: four isolates of one species over overlapping condition panels is a within-species transfer test the wave has nowhere else. \\
\addlinespace[5pt]
29 & \textbf{Fitness Browser WCS417 (Aquipseudomonas simiae)}\newline {\scriptsize Fitness / chemical genomics; BarSeq competition fitness per gene, with a t statistic} & F, 1 & \org{P. simiae} & pooled barcoded transposon library, gene count not reported per organism & 51 carbon, 11 nitrogen, 30 stress, 8 other & 201 & $9.1\times 10^{5}$ & I & fitbrowser+\allowbreak transposon & The widest carbon panel of any pseudomonad in the browser at 51 sources, and a plant-associated strain, so it is the one row in the block whose native context is a root rather than groundwater. 118 updated annotations. \\
\addlinespace[5pt]
30 & \textbf{Fitness Browser FW300-N2C3 (Aquipseudomonas fluorescens\_AA)}\newline {\scriptsize Fitness / chemical genomics; BarSeq competition fitness per gene, with a t statistic} & F, 1 & \org{P. fluorescens} & pooled barcoded transposon library, gene count not reported per organism & 45 carbon, 19 nitrogen, 36 stress, 16 other & 184 & $8.3\times 10^{5}$ & I & fitbrowser+\allowbreak transposon & A third FW300 isolate at 184 experiments with 122 updated annotations, the most of any row in the block. The four FW300 isolates together are the only place in this wave where one species is screened four times over the same condition vocabulary. \\
\addlinespace[5pt]
31 & \textbf{Fitness Browser FW300-N1B4 (Aquipseudomonas sp.)}\newline {\scriptsize Fitness / chemical genomics; BarSeq competition fitness per gene, with a t statistic} & F, 2 & \org{P. fluorescens} & pooled barcoded transposon library, gene count not reported per organism & 39 carbon, 21 nitrogen, 34 stress, 10 other & 147 & $6.7\times 10^{5}$ & I & fitbrowser+\allowbreak transposon & The fourth FW300 isolate, and the one whose current binomial the browser leaves at genus level, so it is the row that forces an Organism member for a strain with no species name. 147 experiments over all four condition classes. \\
\addlinespace[5pt]
32 & \textbf{Fitness Browser B728a (Aquipseudomonas fragariae pv. syringae)}\newline {\scriptsize Fitness / chemical genomics; BarSeq competition fitness per gene, with a t statistic} & F, 2 & \org{P. syringae} & pooled barcoded transposon library, gene count not reported per organism & 17 stress, 8 other; no carbon or nitrogen panel & 46 & $2.1\times 10^{5}$ & I & fitbrowser+\allowbreak transposon & A plant pathogen rather than an environmental isolate, screened on stress only, and with a deletion background screened beside it. The B728a pair is the browser's own worked example that it already models genotype by condition across two genotypes of one strain over the same stress set. \\
\addlinespace[5pt]
33 & \textbf{Fitness Browser B728a mexB deletion background}\newline {\scriptsize Fitness / chemical genomics; BarSeq competition fitness per gene, with a t statistic} & F, 3 & \org{P. syringae} & pooled barcoded transposon library in a mexB deletion background & 17 stress, 1 other & 23 & $1.0\times 10^{5}$ & I & fitbrowser+\allowbreak transposon & The only transposon library in the browser screened in a defined deletion background against its own parent over the same stress panel, which makes the pair a background-by-perturbation interaction rather than two independent screens. 23 experiments. \\
\addlinespace[5pt]
34 & \textbf{Fitness Browser PA14 (Pseudomonas aeruginosa)}\newline {\scriptsize Fitness / chemical genomics; BarSeq competition fitness per gene, with a t statistic} & F, 4 & \org{P. aeruginosa} & pooled barcoded transposon library, gene count not reported per organism & 3 other; no carbon, nitrogen or stress panel & 8 & $3.6\times 10^{4}$ & I & fitbrowser+\allowbreak transposon & A token presence rather than a compendium, at 8 experiments and no updated annotations, and it is the evidence that the browser is not a usable source of fitness for this species. It is ranked because the species enters the wave through its essentiality benchmark instead. \\
\addlinespace[5pt]
35 & \textbf{Dubinkina 2024 (12-strain low-pH transcriptomic atlas)}\,\textsuperscript{\textbf{M}}\newline {\scriptsize Transcriptome; gene expression count matrix} & H, 2 & \org{P. kudriavzevii} & 12 strains, 6 acid-\allowbreak tolerant and 6 susceptible & pH 5.5 control and pH 1.5 at 0 h and 1 h, plus two organic-\allowbreak acid arms & 48$\dagger$ & $2.4\times 10^{5}$ & I & Pk-\allowbreak isolate-\allowbreak WGS & The one omics resource in this host built as a strain-by-condition matrix with a released processed count table rather than raw reads only, which makes it the most directly ingestible thing the host has. Stb5, Mac1 and Rtg1/Rtg3 are the implicated regulators, and the sample and replicate counts are not stated on the deposit page. \\
\addlinespace[5pt]
36 & \textbf{MAESTRO 2025 (170-isolate WGS, 57-stressor growth panel)}\,\textsuperscript{\textbf{R}}\newline {\scriptsize Natural variation; growth under each stressor} & H, 2 & \org{P. kudriavzevii} & 170 WGS isolates, 161 phenotyped & 57 stressors: heat, low pH, organic acids, inhibitors, antifungals & $9.2\times 10^{3}$$\dagger$ & $9.2\times 10^{3}$ & I & Pk-\allowbreak isolate-\allowbreak WGS & The single highest-measurement phenotype resource in this host, with whole-genome sequence per genotype and a copy-number model reaching a median Pearson r of 0.92 across the 57 conditions. Copy-number variation of tandem paralogs is the reported mechanism, validated by a gene4260 and gene4261 deletion, and no accession is in hand. \\
\addlinespace[5pt]
37 & \textbf{Fatma 2023 (landing-pad locus screen)}\,\textsuperscript{\textbf{R}}\newline {\scriptsize CRISPR library screen; cutting efficiency, integration efficiency, fitness, GFP level} & H, 3 & \org{P. kudriavzevii} & 47 gRNAs across 21 intergenic regions & 4 readouts per locus & 188$\dagger$ & 188 & I & SD108+\allowbreak guide & The only measured per-locus fitness data in this species, and genotype-precise: the locus and the cassette are both named. Two of its four readouts, cutting and integration efficiency, fall under the standing guide-level exclusion while its fitness and GFP readouts do not, so the row forces that decision rather than inheriting it. \\
\addlinespace[5pt]
38 & \textbf{Douglass 2018 (32 sequenced isolates, 4-antifungal MIC)}\newline {\scriptsize Natural variation; minimum inhibitory concentration, mg/L} & H, 3 & \org{P. kudriavzevii} & 32 sequenced isolates, 20 clinical and 12 environmental & fluconazole, flucytosine, amphotericin B, micafungin & 128$\dagger$ & 128 & I & Pk-\allowbreak isolate-\allowbreak WGS & The cleanest fully public, fully provenanced genotype-to-phenotype benchmark this host has, with ploidy, mating type, Illumina coverage and ERG11, FKS1 and ABC11 variants per strain beside the MIC table. Fluconazole MICs span 8 to 128 mg/L with a mode of 32, and the whole table was read in session. \\
\addlinespace[5pt]
39 & \textbf{GSE282019 (Delhi NICU five-year outbreak)}\newline {\scriptsize Transcriptome; transcript abundance} & H, 4 & \org{P. kudriavzevii} & 3 isolates with differing fluconazole MICs & 1 & 6 & $3.0\times 10^{4}$ & V & Pk-\allowbreak isolate-\allowbreak WGS & The second of the two drug-stratified expression series in this host, with ERG5 and FCR1 differences reported and a published identifier behind it. Whole-genome sequencing covers bloodstream outbreaks over five years in two genetically distinct clusters. \\
\addlinespace[5pt]
40 & \textbf{GSE202655 (L-malic acid, EcSthA transhydrogenase)}\newline {\scriptsize Transcriptome; transcript abundance, with titer and conversion reported separately} & H, 4 & \org{P. kudriavzevii} & MA009-\allowbreak 5 original and MA009-\allowbreak 10 carrying E. coli SthA & low pH & 6 & $3.0\times 10^{4}$ & V & engineered-\allowbreak chassis & A fifth product on this host, L-malic acid, in a background outside the SD108 lineage entirely, and the only engineered-strain expression series not from that lab. The MA009 background has no reference assembly anywhere in the triage, so its strain record carries a provenance gap rather than an assembly pin. \\
\addlinespace[5pt]
41 & \textbf{GSE159457 (CY902 C4 dicarboxylic acid transporters)}\newline {\scriptsize Transcriptome; transcript abundance} & H, 4 & \org{P. kudriavzevii} & CY902 ura3 & pH 5.0 on glucose, succinic acid, or both & 6 & $3.0\times 10^{4}$ & V & engineered-\allowbreak chassis & A clean one-strain carbon-source by expression matrix in a third background, and the only series in this host with an explicit raw-read study identifier. CY902 has no reference assembly in the triage. \\
\addlinespace[5pt]
42 & \textbf{GSE262836 (Delhi NICU, fluconazole-resistance expression)}\newline {\scriptsize Transcriptome; transcript abundance} & H, 4 & \org{P. kudriavzevii} & one fluconazole-\allowbreak susceptible and one resistant isolate & 1 & 4 & $2.0\times 10^{4}$ & V & Pk-\allowbreak isolate-\allowbreak WGS & With its sibling series it is one of only two deposits in this host that pair a drug phenotype with expression, and both surfaced only through a taxid query rather than a name search. 178 genes were called differentially expressed, which is a derived contrast rather than per-record data, so the counts matrix is what is ingestible. \\
\addlinespace[5pt]
43 & \textbf{Thompson 2019 (KT2440 lysine metabolism)}\,\textsuperscript{\textbf{M}}\newline {\scriptsize Transposon fitness; RB-TnSeq gene fitness} & S, 3 & \org{P. putida} & 39 genes with fitness below -\allowbreak 2 on a lysine route and above -\allowbreak 0.5 on glucose & 5-\allowbreak aminovalerate, D-\allowbreak lysine, L-\allowbreak lysine, glucose & 4 & $2.2\times 10^{4}$ & S & KT2440+\allowbreak transposon & Small, but it is the only one of the condition-panel papers in this host whose quantitative table format is named: Table S1 is a released CSV of the RB-TnSeq results. That makes it the cheapest row in the band to build against a stated file. \\
\addlinespace[5pt]
44 & \textbf{Isoprenol CRISPRi DBTL proteome ladder (7 PRIDE deposits)}\newline {\scriptsize Proteome; relative protein abundance} & S, 3 & \org{P. putida} & one engineered genotype set per design round; counts not reported per round & 7 design-\allowbreak build-\allowbreak test-\allowbreak learn rounds & 7 & $3.9\times 10^{4}$ & I & engineered-\allowbreak chassis+\allowbreak guide & The most database-shaped proteomics series in this host: seven deposits indexed one per design round, so the engineered genotype and the molecular readout are both public and already separated by round. It is the best candidate in the wave for a genotype by environment by molecular-phenotype record, and the campaign is CRISPRi-driven, so the guide is part of the genotype. \\
\addlinespace[5pt]
45 & \textbf{Thompson 2020 (KT2440 fatty acids and alcohols)}\newline {\scriptsize Transposon fitness; RB-TnSeq gene fitness} & S, 3 & \org{P. putida} & strong phenotypes for hundreds of genes; no count stated & 13 fatty acids and 10 alcohols & 23 & $1.3\times 10^{5}$ & S & KT2440+\allowbreak transposon & A named carbon-class panel that most likely supplies part of this host's browser carbon-source experiments, and the row with the clearest de-duplication question in the band. A published correction exists, so the loader has to pin which version it read. \\
\addlinespace[5pt]
46 & \textbf{GSE198395 (putidaPRECISE321 carbon sources)}\newline {\scriptsize Transcriptome; RNA-seq expression} & S, 3 & \org{P. putida} & KT2440 and engineered strains; genotype count not reported & different carbon sources, count not reported & 81 & $4.5\times 10^{5}$ & V & engineered-\allowbreak chassis & The largest single deposit inside the host's expression compendium, at 81 samples, and one of only two of its deposits named in the text rather than in a spreadsheet. It is nested inside the compendium row, so admitting both would double-count unless the compendium is split by deposit. \\
\addlinespace[5pt]
47 & \textbf{Mueller 2025 (24-strain P. putida Biolog panel)}\,\textsuperscript{\textbf{R}}\newline {\scriptsize Natural variation; Biolog respiration per well} & S, 2 & \org{P. putida} & 24 strains including KT2440 & Biolog PM01 and PM02A, 190 carbon compounds, plus 15 aromatics & $4.9\times 10^{3}$$\dagger$ & $4.9\times 10^{3}$ & V & pseudomonad-\allowbreak isolate-\allowbreak WGS & The only resource that lets this host enter as a species with intraspecific variation rather than one reference strain, and the only multi-strain Biolog matrix for any pseudomonad: 123 of the 205 compounds supported activity in at least one strain. bioRxiv rate-limited four attempts, so the deposit is unconfirmed. \\
\addlinespace[5pt]
48 & \textbf{Schmidt 2022 (KT2440 nitrogen BarSeq)}\newline {\scriptsize Transposon fitness; BarSeq competition fitness per gene} & S, 2 & \org{P. putida} & 672 genes with a significant phenotype, 529 of them UniProt-\allowbreak unreviewed & 52 nitrogen compounds plus 19 amino-\allowbreak acid dropouts & 71 & $4.0\times 10^{5}$ & I & KT2440+\allowbreak transposon & The largest single-axis nitrogen panel for any pseudomonad, and it nearly doubled the public BarSeq data for this host, with 256 of its hits having had no significant phenotype in any previous study. It ships an interactive view that links into the Fitness Browser, so the values are very likely inside feba.db, but no accession is stated. \\
\addlinespace[5pt]
49 & \textbf{Borchert 2024 (KT2440 fitness compendium)}\newline {\scriptsize Transposon fitness; RB-TnSeq gene fitness with t statistics} & S, 1 & \org{P. putida} & 4,732 genes with data of 5,564 protein-\allowbreak coding & 332 samples; unique conditions reported as 179 in the abstract and 183 in the methods & 332 & $1.6\times 10^{6}$ & V & KT2440+\allowbreak transposon & A ready-made 4,732 by 332 matrix with per-sample metadata in one repository, which makes it the most ingestible fitness artifact for this host. 254 of the 332 samples come from the Fitness Browser, so de-duplication must be by experiment identity and never by count, and the paper states no genome assembly version. \\
\addlinespace[5pt]
50 & \textbf{putidaPRECISE321 (Lim 2022)}\newline {\scriptsize Transcriptome; RNA-seq expression} & S, 1 & \org{P. putida} & KT2440 and engineered strains; strain-\allowbreak genotype count not reported & 118 unique conditions across 21 projects & 321 & $1.8\times 10^{6}$ & I & KT2440-\allowbreak reference & A 5,564 by 321 matrix in one MIT-licensed repository with 118 typed conditions, structurally parallel to the fitness compendium, so one condition vocabulary can serve both. The per-sample accession list lives in a supplementary spreadsheet rather than the text, which is the one provenance item still open. \\
\addlinespace[5pt]
\midrule \multicolumn{11}{@{}l}{\textbf{Ten extra, in rank order: rows 51--60 replace a row above that proves unreachable.}}\\ \midrule
51 & \textbf{Cao 2020 (pooled four-condition RNA-seq)}\,\textsuperscript{\textbf{M}}\newline {\scriptsize Transcriptome; transcript abundance} & H, 4 & \org{P. kudriavzevii} & SD108 wild type & glucose and glucose plus inhibitors, each aerobic and anaerobic & 1 & $4.9\times 10^{3}$ & I & SD108-\allowbreak reference & Supplies the reference expression layer the promoter panel was selected from, under four named inhibitor and oxygen conditions with recipes given to the gram. It is a single library pooled from all four conditions, so there is no per-condition matrix and the honest record is one row whose environment is a pool. \\
\addlinespace[5pt]
52 & \textbf{PXD069498 (extracellular vesicle proteome)}\newline {\scriptsize Proteome; protein identification in purified vesicles} & H, 4 & \org{P. kudriavzevii} & wild type, strain not reported & vesicle purification & 189 & 189 & S & Pk-\allowbreak reference & 189 proteins identified from 33 raw files, which is thin on purpose: an identification list with no genotype and no environment axis is a protein-presence set rather than a genotype-by-environment record. Held so the deposit is accounted for, and deferred rather than built. \\
\addlinespace[5pt]
53 & \textbf{Douglass 2018 (per-strain SNP density and loss of heterozygosity)}\newline {\scriptsize Modality / backbone; SNP density and loss-of-heterozygosity extent per strain} & H, 4 & \org{P. kudriavzevii} & 32 isolates & none; a genome-\allowbreak architecture statistic & 64$\dagger$ & 64 & I & Pk-\allowbreak isolate-\allowbreak WGS & The per-isolate genomic-background descriptors that make the MIC rows interpretable, alongside the ploidy and mating-type calls in the paper's main table. Not a growth phenotype, which is why it is a backbone row rather than a screen. \\
\addlinespace[5pt]
54 & \textbf{Tran 2023 (succinate end-to-end pipeline)}\,\textsuperscript{\textbf{M}}\newline {\scriptsize Production campaign; titer, rate and yield} & H, 4 & \org{P. kudriavzevii} & at least 10 genotypes countable from the body text & sugar media at pH 3 fed-\allowbreak batch, sugarcane juice, pilot batch & 30$\ddagger$ & 30 & I & SD108+\allowbreak designed-\allowbreak edit & The deepest genotype-precise engineering record in this host, with named deletions and chromosomal integrations that are reconstructable, reaching 63.1 g/L at pilot scale. Supplementary Table 5 lists all strains and was not read, so ten is a floor and the often-quoted 100 to 250 genotype figure for the whole lineage is an upper-range estimate. \\
\addlinespace[5pt]
55 & \textbf{Cao 2020 (constitutive promoter-GFP panel)}\,\textsuperscript{\textbf{M}}\newline {\scriptsize Engineering parts; GFP level per promoter construct} & H, 4 & \org{P. kudriavzevii} & 16 named promoters plus an unstated number more & different culture conditions, count not reported & 16$\ddagger$ & 16 & I & SD108-\allowbreak reference & The only promoter-strength panel in this host, and the row type is exactly the one the schema already serves: a named promoter driving a reporter under stated media. It also anchors the partial map from JGI gNNNN ids to gene names, since g247 is TDH3, g853 is GPM1 and g5025 is SED1. \\
\addlinespace[5pt]
56 & \textbf{Xiao 2014 (SD108 reductive-TCA succinate)}\newline {\scriptsize Production campaign; titer, yield, productivity, and acid-tolerance growth calls} & H, 4 & \org{P. kudriavzevii} & SD108 wild type and a uracil-\allowbreak auxotroph rTCA integrant & 4 acid-\allowbreak tolerance conditions plus shake-\allowbreak flask batch & 8$\ddagger$ & 8 & S & SD108-\allowbreak reference & The provenance origin of the whole SD108 lineage: this is where the strain was isolated and typed, and it supplies the acid-tolerance thresholds the later work builds on, at 30 g/L succinic acid, 30 g/L itaconic, 23 g/L adipic and 20 g/L acetic. Reported at 11.63 g/L titer in shake-flask batch. \\
\addlinespace[5pt]
57 & \textbf{Xylose utilization across four natural backgrounds}\,\textsuperscript{\textbf{M}}\newline {\scriptsize Production campaign; xylose utilization} & H, 4 & \org{P. kudriavzevii} & SD108, IO21, IO45, IO46, each with the same three-\allowbreak gene cassette & xylose & 4 & 4 & S & SD108+\allowbreak designed-\allowbreak edit & The only row in the whole wave that puts one identical construct into four different natural backgrounds, which is the background-effect axis a cross-host substrate exists to measure. The IO21, IO45 and IO46 naming matches the IO\_NNN resequencing projects, so this is also the bridge between the engineering line and the natural panel. \\
\addlinespace[5pt]
58 & \textbf{Illinois Data Bank IDB-9063411 (high-yield 3-HP)}\,\textsuperscript{\textbf{M}}\newline {\scriptsize Production campaign; titer, yield and productivity} & H, 4 & \org{P. kudriavzevii} & not reported on the deposit page & fed-\allowbreak batch at pH 4 on corn steep liquor & 3$\ddagger$ & 3 & S & SD108+\allowbreak designed-\allowbreak edit & The highest-yield 3-HP record on this host, at 92 g/L, 0.7 g/g and 0.55 g/L/h, on a stated industrial feedstock. Deposit contents are unverified because the repository returns 403 to every automated client. \\
\addlinespace[5pt]
59 & \textbf{Douglass 2018 (CBS573 RNA-seq)}\newline {\scriptsize Transcriptome; transcript abundance} & H, 4 & \org{P. kudriavzevii} & CBS573 type strain & not reported & --$\ddagger$ & -- & V & CBS573-\allowbreak reference & The only expression data pinned to the haploid-resolved CBS573 reference, which makes it the natural bridge for building the clinical side of a gene-identifier crosswalk. It is raw reads, so admitting it is a build rather than an ingest of released values. \\
\addlinespace[5pt]
60 & \textbf{MAESTRO 2025 (gene4260/gene4261 paralog-pair deletion)}\,\textsuperscript{\textbf{R}}\newline {\scriptsize Tolerance / robustness; stress-tolerance effect of the deletion} & H, 4 & \org{P. kudriavzevii} & one tandem paralog pair deleted & several stresses, count not reported & --$\ddagger$ & -- & S & SD108+\allowbreak designed-\allowbreak edit & One of only three sources of measured gene-perturbation phenotype in this host, and the designed-edit arm that validates the copy-number mechanism the isolate panel infers. Whether the deletion is one double mutant or two singles is not stated. \\
\addlinespace[5pt]
\end{longtable}
\endgroup
\end{landscape}


\clearpage
%% GENERATED FILE -- do not hand-edit.
%% SOURCE: experiments/database/scripts/build_hosts_candidate_datasets_table.py
\begin{landscape}
\begingroup
\footnotesize
\setlength{\tabcolsep}{4pt}
\renewcommand{\arraystretch}{1.15}
\begin{longtable}{@{}r@{\hspace{4pt}} L{78mm} L{72mm} L{76mm}@{}}
\caption[]{Sources for the recommended fifty and the ten extra of Table~\ref{tab:final}, in the same order.
\emph{Why it is ordered here} is the band reason for a host-bootstrap, regulatory-DNA or
cross-host-fitness row, and the tier and measurement rule for a scale row. \emph{Where
the per-record data lives} is the deposit, with its retrieval class: a plain entry is
fetched programmatically, \emph{browser} needs a human because the host returns 403 or
renders in JavaScript, and \emph{ask} means no public copy exists and the request is in
Table~\ref{tab:asks}. Every link is clickable.}
\label{tab:sources}\\
\toprule
\textbf{\#} & \textbf{Dataset, citation and link} & \textbf{Why it is ordered here} & \textbf{Where the per-record data lives} \\
\midrule
\endfirsthead
\multicolumn{4}{@{}l}{\footnotesize\emph{Table~\ref{tab:sources}, continued}}\\
\toprule
\textbf{\#} & \textbf{Dataset, citation and link} & \textbf{Why it is ordered here} & \textbf{Where the per-record data lives} \\
\midrule
\endhead
\bottomrule
\endfoot
1 & \textbf{Douglass 2018 (32 sequenced isolates, 4-antifungal MIC)} \org{P. kudriavzevii}\newline Douglass AP, Offei B, Braun-Galleani S, Coughlan AY, Martos AAR, Ortiz-Merino RA, Byrne KP, Wolfe KH. PLoS Pathog 2018;14:e1007138.\newline \href{https://www.ncbi.nlm.nih.gov/pmc/articles/PMC6053246/}{PMC60532\allowbreak 46} & Owner pin, 2026-11-24, to rank 1 from the host bootstrap band: it is the one P. kudriavzevii growth-type panel that can be fetched today with no ask, and fluconazole is the one compound measured on natural isolates in both yeasts, against Peter 2018 at 20 ug/ml on 971 S. cerevisiae isolates and Bloom 2019 at 100 uM on 13,950 segregants, so it is the first cross-host transfer test that has labels at both ends. The other three drugs are measured in S. cerevisiae only as deletion or CRISPRi chemogenomics, in Parsons 2006, Hillenmeyer 2008, and Piotrowski 2017. & SRA SRP1\allowbreak 39299 (3\allowbreak 2 strain\allowbreak s); GenB\allowbreak ank CP02\allowbreak 8773-\allowbreak CP02877\allowbreak 8 (CBS57\allowbreak 3 refere\allowbreak nce) and\allowbreak  CP02853\allowbreak 1-\allowbreak CP02853\allowbreak 5 (CBS51\allowbreak 47); MIC\allowbreak  table i\allowbreak n the op\allowbreak en-\allowbreak access \allowbreak article \allowbreak body \\
\addlinespace[5pt]
2 & \textbf{de Boer 2020 (GPRA sort-seq, 100M random promoters)} \org{S. cerevisiae}\newline de Boer CG, Vaishnav ED, Sadeh R, Abeyta EL, Friedman N, Regev A. Nat Biotechnol 2020;38:56-65.\newline \href{https://doi.org/10.1038/s41587-019-0315-8}{doi:\allowbreak 10.\allowbreak 1038/\allowbreak s41587-\allowbreak 019-\allowbreak 0315-\allowbreak 8} & Regulatory DNA band: one row carrying more measurements than every other row in this wave combined, on a schema that already serves its child, so the admission is a download decision rather than a design one. & GEO GSE1\allowbreak 04878 an\allowbreak d GSE104\allowbreak 903; Add\allowbreak gene 127\allowbreak 546; whi\allowbreak ch acces\allowbreak sion hol\allowbreak ds which\allowbreak  library\allowbreak  is unre\allowbreak solved a\allowbreak nd one p\allowbreak age chec\allowbreak k settle\allowbreak s it \\
\addlinespace[5pt]
3 & \textbf{Vaishnav 2022 deferred training libraries} \org{S. cerevisiae}\newline Vaishnav ED, de Boer CG, Molinet J, Yassour M, Fan L, Adiconis X, Thompson DA, Levin JZ, Cubillos FA, Regev A. Nature 2022;603:455-463.\newline \href{https://doi.org/10.1038/s41586-022-04506-6}{doi:\allowbreak 10.\allowbreak 1038/\allowbreak s41586-\allowbreak 022-\allowbreak 04506-\allowbreak 6} & Regulatory DNA band, and the single lowest-effort large win in the wave: every piece of machinery it needs is already written and serving its subset. & Zenodo 4\allowbreak 436477 (\allowbreak complex\_\allowbreak media\_\allowbreak trainin\allowbreak g\_\allowbreak data\_\allowbreak Glu.\allowbreak txt 3.\allowbreak 59 GB, \allowbreak defined\_\allowbreak media\_\allowbreak trainin\allowbreak g\_\allowbreak data\_\allowbreak SC\_\allowbreak Ura.\allowbreak txt 2.\allowbreak 47 GB);\allowbreak  GEO GSE\allowbreak 163045 a\allowbreak nd GSE16\allowbreak 3866 \\
\addlinespace[5pt]
4 & \textbf{Hollerer 2020 (ribosome binding site library, uASPIre)} \org{E. coli}\newline Hollerer S, Papaxanthos L, Gumpinger AC, Fischer K, Beisel C, Borgwardt K, Benenson Y, Jeschek M. Large-scale DNA-based phenotypic recording and deep learning enable highly accurate sequence-function mapping. Nat Commun 2020;11:3551.\newline \href{https://doi.org/10.1038/s41467-020-17222-4}{doi:\allowbreak 10.\allowbreak 1038/\allowbreak s41467-\allowbreak 020-\allowbreak 17222-\allowbreak 4} & Regulatory DNA: a 300,000-member binding-site library whose readout is a DNA-recorded time integral, a measurement principle no other row in the band carries. & the Jesc\allowbreak hekLab/\allowbreak uASPIre\allowbreak  reposit\allowbreak ory, rea\allowbreak chable \\
\addlinespace[5pt]
5 & \textbf{Hollerer 2023 (5' UTR and CDS library, uASPIre)} \org{E. coli}\newline Hollerer S, Jeschek M. Nucleic Acids Res 2023;51:2377-2396.\newline \href{https://doi.org/10.1093/nar/gkad040}{doi:\allowbreak 10.\allowbreak 1093/\allowbreak nar/\allowbreak gkad040} & Regulatory DNA: the band's whole purpose, an axis that adds measurements rather than rows, and this row adds more of them than any other non-yeast candidate. & Zenodo 1\allowbreak 0.\allowbreak 5281/\allowbreak zenodo.\allowbreak 7536586\allowbreak  and the\allowbreak  Jeschek\allowbreak Lab/\allowbreak uASPIre\_\allowbreak UTR\_\allowbreak CDS rep\allowbreak ository,\allowbreak  both re\allowbreak achable \\
\addlinespace[5pt]
6 & \textbf{DREAM 2022 random promoter challenge (Rafi 2024)} \org{S. cerevisiae}\newline Rafi AM, et al. A community effort to optimize sequence-based deep learning models of gene regulation. Nat Biotechnol 2024.\newline \href{https://www.ncbi.nlm.nih.gov/geo/query/acc.cgi?acc=GSE254493}{ncbi.\allowbreak nlm.\allowbreak nih.\allowbreak gov/\allowbreak geo/\allowbreak query/\allowbreak acc.\allowbreak cgi?\allowbreak acc=\allowbreak GSE2544\allowbreak 93} & Regulatory DNA band: the benchmark slice, which is what makes a model comparison possible at all, and it carries its own eight-category test design. & GEO GSE2\allowbreak 54493; Z\allowbreak enodo 10\allowbreak 633252 (\allowbreak v2, abou\allowbreak t 4.\allowbreak 0 GB), \allowbreak supersed\allowbreak ing Zeno\allowbreak do 73953\allowbreak 97 \\
\addlinespace[5pt]
7 & \textbf{Kuo 2020 (Shine-Dalgarno 9-mer library in four contexts)} \org{E. coli}\newline Kuo ST, Jahn RL, Cheng YJ, Chen YH, Lee YJ, Hollfelder F, Wen JD, Chou HHD. Global fitness landscapes of the Shine-Dalgarno sequence. Genome Res 2020;30:711-723.\newline \href{https://doi.org/10.1101/gr.260182.119}{doi:\allowbreak 10.\allowbreak 1101/\allowbreak gr.\allowbreak 260182.\allowbreak 119} & Regulatory DNA: a complete nine-mer space in four contexts, which is the band's axis measured exhaustively rather than sampled. & Europe P\allowbreak MC suppl\allowbreak ementary\allowbreak  files u\allowbreak nder PMC\allowbreak 7263185,\allowbreak  48,362,\allowbreak 217 byte\allowbreak s; reads\allowbreak  at PRJN\allowbreak A516114;\allowbreak  the pub\allowbreak lisher's\allowbreak  own sit\allowbreak e return\allowbreak s 403 \\
\addlinespace[5pt]
8 & \textbf{Kotopka and Smolke 2020 (artificial promoter libraries)} \org{S. cerevisiae}\newline Kotopka BJ, Smolke CD. Model-driven generation of artificial yeast promoters. Nat Commun 2020;11:2113.\newline \href{https://doi.org/10.1038/s41467-020-15977-4}{doi:\allowbreak 10.\allowbreak 1038/\allowbreak s41467-\allowbreak 020-\allowbreak 15977-\allowbreak 4} & Regulatory DNA band: the only million-scale promoter library outside the GPRA lineage, which is what makes it a check on it rather than more of it. & Zenodo 3\allowbreak 735426 (\allowbreak Promoter\allowbreak  Design \allowbreak Data\_\allowbreak 2020-\allowbreak 03.\allowbreak zip, ab\allowbreak out 1.\allowbreak 1 GB) \\
\addlinespace[5pt]
9 & \textbf{Zhu 2024 (Pol II MASTER, TSS initiation efficiency)} \org{S. cerevisiae}\newline Zhu C, Vvedenskaya IO, Sze SH, Nickels BE, Kaplan CD. Nat Struct Mol Biol 2024;31:190-202.\newline \href{https://doi.org/10.1038/s41594-023-01171-9}{doi:\allowbreak 10.\allowbreak 1038/\allowbreak s41594-\allowbreak 023-\allowbreak 01171-\allowbreak 9} & Regulatory DNA band: a different readout from reporter level, so it extends the axis rather than deepening it, and at a million records it is the third-largest row in the wave. & GEO GSE1\allowbreak 85290 (p\allowbreak rocessed\allowbreak ); BioPr\allowbreak oject PR\allowbreak JNA76662\allowbreak 4 (raw) \\
\addlinespace[5pt]
10 & \textbf{Savinov 2021 (random 50-mer 3' UTR library)} \org{S. cerevisiae}\newline Savinov A, Brandsen BM, Angell BE, Cuperus JT, Fields S. Genome Biol 2021;22:293.\newline \href{https://doi.org/10.1186/s13059-021-02509-6}{doi:\allowbreak 10.\allowbreak 1186/\allowbreak s13059-\allowbreak 021-\allowbreak 02509-\allowbreak 6} & Regulatory DNA: the 3' UTR end of the transcript, the one regulatory axis no other row in the band measures. & a figsha\allowbreak re recor\allowbreak d servin\allowbreak g a 36,9\allowbreak 20,254-\allowbreak byte ar\allowbreak chive th\allowbreak rough it\allowbreak s API; r\allowbreak eads at \allowbreak SRA PRJN\allowbreak A750726 \\
\addlinespace[5pt]
11 & \textbf{Urtecho 2020 (genome-wide promoter MPRA)} \org{E. coli}\newline Urtecho G, et al. Genome-wide functional characterization of Escherichia coli promoters and regulatory elements responsible for their function. eLife Reviewed Preprint 2023;12:92558.\newline \href{https://doi.org/10.7554/eLife.92558.1}{doi:\allowbreak 10.\allowbreak 7554/\allowbreak eLife.\allowbreak 92558.\allowbreak 1} & Regulatory DNA: the only row that tiles a whole bacterial genome for promoter activity, so its sequences are native rather than designed. & GEO GSE1\allowbreak 44621, r\allowbreak eachable\allowbreak  over FT\allowbreak P; the a\allowbreak ccession\allowbreak  display\allowbreak  page is\allowbreak  behind \allowbreak a challe\allowbreak nge and \allowbreak the auth\allowbreak ors' own\allowbreak  databas\allowbreak e fails \allowbreak TLS \\
\addlinespace[5pt]
12 & \textbf{Cuperus 2017 (random 5' UTR library)} \org{S. cerevisiae}\newline Cuperus JT, Groves B, Kuchina A, Rosenberg AB, Jojic N, Fields S, Seelig G. Genome Res 2017;27:2015-2028.\newline \href{https://genome.cshlp.org/content/27/12/2015}{genome.\allowbreak cshlp.\allowbreak org/\allowbreak content\allowbreak /\allowbreak 27/\allowbreak 12/\allowbreak 2015} & Regulatory DNA band: the only large 5' UTR library, and the only one in the band whose readout is growth rather than fluorescence. & GEO GSE1\allowbreak 04252 \\
\addlinespace[5pt]
13 & \textbf{Johns 2018 (29,249 expression constructs across conditions)} \org{E. coli}\newline Johns NI, et al. Metagenomic mining of regulatory elements enables programmable species-selective gene expression. Nat Methods 2018;15:323-329.\newline \href{https://doi.org/10.1038/nmeth.4633}{doi:\allowbreak 10.\allowbreak 1038/\allowbreak nmeth.\allowbreak 4633} & Regulatory DNA: the only library in the band whose sequences were mined from metagenomes rather than designed or tiled from one genome. & two publ\allowbreak isher su\allowbreak pplement\allowbreak ary file\allowbreak s, both \allowbreak reachabl\allowbreak e; no GE\allowbreak O or SRA\allowbreak  accessi\allowbreak on found\allowbreak , and th\allowbreak e full-\allowbreak text se\allowbreak rvice re\allowbreak turns an\allowbreak  error o\allowbreak n the av\allowbreak ailabili\allowbreak ty state\allowbreak ment \\
\addlinespace[5pt]
14 & \textbf{Yu 2021 (promoter library across induction states)} \org{E. coli}\newline Yu H, et al. Dynamic and quantitative characterization of promoters in Escherichia coli. Nat Commun 2021;12:325.\newline \href{https://doi.org/10.1038/s41467-020-20094-3}{doi:\allowbreak 10.\allowbreak 1038/\allowbreak s41467-\allowbreak 020-\allowbreak 20094-\allowbreak 3} & Regulatory DNA: a chromosomal promoter library carrying an induction response per sequence rather than a single level. & GEO GSE1\allowbreak 45630, r\allowbreak eachable\allowbreak  over FT\allowbreak P \\
\addlinespace[5pt]
15 & \textbf{Cambray 2018 (244,000 designed sequences, multi-phenotype)} \org{E. coli}\newline Cambray G, Guimaraes JC, Arkin AP. Evaluation of 244,000 synthetic sequences reveals design principles to optimize translation in Escherichia coli. Nat Biotechnol 2018;36:1005-1015.\newline \href{https://doi.org/10.1038/nbt.4238}{doi:\allowbreak 10.\allowbreak 1038/\allowbreak nbt.\allowbreak 4238} & Regulatory DNA: one designed library of 244,000 sequences read out five different ways, the widest phenotype fan-out in the wave. & an OSF d\allowbreak eposit a\allowbreak t 130,12\allowbreak 2,456 by\allowbreak tes that\allowbreak  answers\allowbreak  a scrip\allowbreak t, where\allowbreak  the pub\allowbreak lisher's\allowbreak  own lin\allowbreak k redire\allowbreak cts; rea\allowbreak ds at SR\allowbreak P086076 \\
\addlinespace[5pt]
16 & \textbf{Fuqua 2024 (promoter-island mutational landscape)} \org{E. coli}\newline Fuqua T, et al. Dense and pleiotropic regulatory information in a developmental enhancer-like bacterial promoter island. eLife 2024;13:98654.\newline \href{https://doi.org/10.7554/eLife.98654}{doi:\allowbreak 10.\allowbreak 7554/\allowbreak eLife.\allowbreak 98654} & Regulatory DNA: a quarter-million mutated native promoter sequences released as per-bin distributions rather than means. & the tfuq\allowbreak ua95/\allowbreak promote\allowbreak r\_\allowbreak islands\allowbreak  reposit\allowbreak ory, rea\allowbreak chable; \allowbreak reads at\allowbreak  SRA PRJ\allowbreak NA107157\allowbreak 2 \\
\addlinespace[5pt]
17 & \textbf{Gui 2021 (position effects with transcriptomes)} \org{S. cerevisiae}\newline Gui, et al. Transcriptome analysis in yeast reveals the externality of position effects. Mol Biol Evol 2021.\newline \href{https://www.biorxiv.org/content/10.1101/2020.04.02.021162}{biorxiv.\allowbreak org/\allowbreak content\allowbreak /\allowbreak 10.\allowbreak 1101/\allowbreak 2020.\allowbreak 04.\allowbreak 02.\allowbreak 021162} & Regulatory DNA band: the sequence variable is purely the integration coordinate, and this is the one row in that tier whose phenotype is a vector. & no acces\allowbreak sion fou\allowbreak nd; need\allowbreak s a brow\allowbreak ser chec\allowbreak k of the\allowbreak  publish\allowbreak ed suppl\allowbreak ement\newline {\scriptsize (browser)} \\
\addlinespace[5pt]
18 & \textbf{Zhou 2025 (massively parallel bacterial promoter library)} \org{E. coli}\newline Zhou Y, et al. Nucleic Acids Res 2025;53:gkaf863.\newline \href{https://doi.org/10.1093/nar/gkaf863}{doi:\allowbreak 10.\allowbreak 1093/\allowbreak nar/\allowbreak gkaf863} & Regulatory DNA: 700,407 promoter measurements, and the clearest case in the wave of a release exceeding its own paper's claim. & Europe P\allowbreak MC suppl\allowbreak ementary\allowbreak  files u\allowbreak nder PMC\allowbreak 12407095\allowbreak ; reads \allowbreak at PRJNA\allowbreak 1189185 \\
\addlinespace[5pt]
19 & \textbf{Keren 2016 (promoter panel x named gene, pooled fitness)} \org{S. cerevisiae}\newline Keren L, et al. Massively parallel interrogation of the effects of gene expression levels on fitness. Cell 2016;166:1282-1294.\newline \href{https://www.ncbi.nlm.nih.gov/geo/query/acc.cgi?acc=GSE83936}{ncbi.\allowbreak nlm.\allowbreak nih.\allowbreak gov/\allowbreak geo/\allowbreak query/\allowbreak acc.\allowbreak cgi?\allowbreak acc=\allowbreak GSE8393\allowbreak 6} & Regulatory DNA band: the only large tabulated panel that links expression level to fitness gene by gene, and the only one whose perturbation is a swap at a native locus. & GEO GSE8\allowbreak 3936; Ar\allowbreak rayExpre\allowbreak ss E-\allowbreak GEOD-\allowbreak 83936; \allowbreak BioProje\allowbreak ct PRJNA\allowbreak 329129 \\
\addlinespace[5pt]
20 & \textbf{Shalem 2015 (3' end FACS-seq library)} \org{S. cerevisiae}\newline Shalem O, et al. 3' end determinants. PLoS Genet 2015;11:e1005147.\newline \href{https://doi.org/10.1371/journal.pgen.1005147}{doi:\allowbreak 10.\allowbreak 1371/\allowbreak journal\allowbreak .\allowbreak pgen.\allowbreak 1005147} & Regulatory DNA band: the terminator end of the same axis, with the sequence and the value in one released file. & PLoS Gen\allowbreak et S3 Ta\allowbreak ble (XLS\allowbreak X), no a\allowbreak ccession\allowbreak ; needs \allowbreak a supple\allowbreak ment fet\allowbreak ch and a\allowbreak  sha256 \allowbreak pin\newline {\scriptsize (browser)} \\
\addlinespace[5pt]
21 & \textbf{Renganaath 2020 (native promoter variant MPRA)} \org{S. cerevisiae}\newline Renganaath K, Cheung R, Day L, Kosuri S, Kruglyak L, Albert FW. eLife 2020;9:e62669.\newline \href{https://doi.org/10.7554/eLife.62669}{doi:\allowbreak 10.\allowbreak 7554/\allowbreak eLife.\allowbreak 62669} & Regulatory DNA band: real alleles rather than sampled sequence, measured on the same readout, which is the bridge between the synthetic and the natural axes. & GEO GSE1\allowbreak 55944 (r\allowbreak aw plus \allowbreak barcode \allowbreak assignme\allowbreak nts); th\allowbreak e MPRA d\allowbreak esign an\allowbreak d proces\allowbreak sed tabl\allowbreak es are e\allowbreak Life Sup\allowbreak plementa\allowbreak ry Files \\
\addlinespace[5pt]
22 & \textbf{Arg-CCU tRNA deep mutational scan} \org{S. cerevisiae}\newline Deep mutational scan of a single-copy Arg-CCU tRNA over 14 substitutions seen across post-whole-genome-duplication Saccharomyces. PMC6193533.\newline \href{https://pmc.ncbi.nlm.nih.gov/articles/PMC6193533}{PMC61935\allowbreak 33} & Regulatory DNA band: point substitutions in a non-coding gene, which is the one DMS shape the protein aggregates cannot hold. & journal \allowbreak suppleme\allowbreak nt, no a\allowbreak ccession\allowbreak  located\allowbreak ; needs \allowbreak a supple\allowbreak ment fet\allowbreak ch\newline {\scriptsize (browser)} \\
\addlinespace[5pt]
23 & \textbf{Keren 2013 (native promoter panel across conditions)} \org{S. cerevisiae}\newline Keren L, et al. Mol Syst Biol 2013.\newline \href{https://spell.yeastgenome.org/}{spell.\allowbreak yeastge\allowbreak nome.\allowbreak org} & Regulatory DNA band: the only row in it with a real condition axis on native sequence. & not loca\allowbreak ted; rea\allowbreak ched onl\allowbreak y throug\allowbreak h the Zr\allowbreak imec and\allowbreak  Karollu\allowbreak s valida\allowbreak tion set\allowbreak s\newline {\scriptsize (browser)} \\
\addlinespace[5pt]
24 & \textbf{Lubliner 2015 (core promoter variant panel)} \org{S. cerevisiae}\newline Lubliner S, Keren L, Segal E, et al. Core promoter sequence in yeast is a major determinant of expression level. Genome Res 2015;25:1008-1017.\newline \href{https://genome.cshlp.org/content/25/7/1008}{genome.\allowbreak cshlp.\allowbreak org/\allowbreak content\allowbreak /\allowbreak 25/\allowbreak 7/\allowbreak 1008} & Regulatory DNA band: the core promoter is a distinct regulatory element from the upstream activating region every other row varies. & Genome R\allowbreak es Suppl\allowbreak emental \allowbreak Table 1;\allowbreak  needs a\allowbreak  supplem\allowbreak ent fetc\allowbreak h\newline {\scriptsize (browser)} \\
\addlinespace[5pt]
25 & \textbf{Sharon 2012 (designed promoter library)} \org{S. cerevisiae}\newline Sharon E, Kalma Y, Sharp A, et al., Segal E. Nat Biotechnol 2012;30:521-530.\newline \href{https://doi.org/10.1038/nbt.2205}{doi:\allowbreak 10.\allowbreak 1038/\allowbreak nbt.\allowbreak 2205} & Regulatory DNA band: the controlled counterpart the random libraries lack, and the platform's own origin. & Nat Biot\allowbreak echnol j\allowbreak ournal s\allowbreak upplemen\allowbreak t, no ac\allowbreak cession;\allowbreak  needs a\allowbreak  supplem\allowbreak ent fetc\allowbreak h\newline {\scriptsize (browser)} \\
\addlinespace[5pt]
26 & \textbf{Yamanishi 2013 (genomic terminatome)} \org{S. cerevisiae}\newline Yamanishi, et al. Characterization of heterologous protein expression using over 3,000 yeast genomic terminators. SGD reference S000206762.\newline \href{https://www.yeastgenome.org/reference/S000206762}{yeastgen\allowbreak ome.\allowbreak org/\allowbreak referen\allowbreak ce/\allowbreak S000206\allowbreak 762} & Regulatory DNA band: the largest native-terminator panel, which is the one element of the axis no library in the band randomizes. & primary \allowbreak table no\allowbreak t locate\allowbreak d; needs\allowbreak  the pap\allowbreak er ident\allowbreak ified, t\allowbreak hen a su\allowbreak pplement\allowbreak  fetch\newline {\scriptsize (browser)} \\
\addlinespace[5pt]
27 & \textbf{Dvir 2013 (5' UTR positions -10 to -1)} \org{S. cerevisiae}\newline Dvir S, Velten L, Sharon E, Zeevi D, Carey LB, Weinberger A, Segal E. Proc Natl Acad Sci USA 2013;110:E2792-E2801.\newline \href{https://doi.org/10.1073/pnas.1222534110}{doi:\allowbreak 10.\allowbreak 1073/\allowbreak pnas.\allowbreak 1222534\allowbreak 110} & Regulatory DNA band: a controlled ten-base window, which is the designed counterpart to the half-million random 5' UTRs above it. & PNAS jou\allowbreak rnal sup\allowbreak plement,\allowbreak  no acce\allowbreak ssion; n\allowbreak eeds a s\allowbreak upplemen\allowbreak t fetch\newline {\scriptsize (browser)} \\
\addlinespace[5pt]
28 & \textbf{Wu 2017 (genome-wide position effects, 1,044 loci)} \org{S. cerevisiae}\newline Wu XL, et al. Genome-wide landscape of position effects on heterogeneous gene expression in Saccharomyces cerevisiae. Biotechnol Biofuels 2017.\newline \href{https://www.ncbi.nlm.nih.gov/pmc/articles/PMC5516366/}{PMC55163\allowbreak 66} & Regulatory DNA band by its sequence variable, the integration coordinate, even though the readout is one reporter value. & Table S2\allowbreak , stated\allowbreak  availab\allowbreak le from \allowbreak the corr\allowbreak espondin\allowbreak g author\allowbreak  on reas\allowbreak onable r\allowbreak equest\newline {\scriptsize (browser)} \\
\addlinespace[5pt]
29 & \textbf{ProteinGym yeast subset} \org{S. cerevisiae}\newline ProteinGym: deep mutational scanning scores from 217 assays over 186 proteins, distributed through ProteinGymR and ProteinGymPy.\newline \href{https://proteingym.org}{proteing\allowbreak ym.\allowbreak org} & Regulatory DNA band as the aggregate route into the mutational-scan tier, which is the sequence-to-function axis in coding space rather than regulatory space. & ProteinG\allowbreak ymR and \allowbreak ProteinG\allowbreak ymPy pac\allowbreak kage loa\allowbreak ders; ar\allowbreak chived o\allowbreak n Zenodo\allowbreak , record\allowbreak  number \allowbreak not capt\allowbreak ured \\
\addlinespace[5pt]
30 & \textbf{Zhang 2016 (genome-wide GFP position scan)} \org{S. cerevisiae}\newline Zhang, et al. Genome-wide GFP position-effect scan using the yeast knockout collection. Chin J Biotech 2016.\newline \href{https://wsws.ijournals.cn/cjbcn/article/pdf/gc16070901}{wsws.\allowbreak ijourna\allowbreak ls.\allowbreak cn/\allowbreak cjbcn/\allowbreak article\allowbreak /\allowbreak pdf/\allowbreak gc16070\allowbreak 901} & Regulatory DNA band by its sequence variable, and the only FACS-read position-effect scan. & no acces\allowbreak sion loc\allowbreak ated; Ch\allowbreak inese-\allowbreak languag\allowbreak e journa\allowbreak l PDF, n\allowbreak eeds a b\allowbreak rowser f\allowbreak etch\newline {\scriptsize (browser)} \\
\addlinespace[5pt]
31 & \textbf{Kita 2017 (context-stratified yeast eQTLs)} \org{S. cerevisiae}\newline Kita R, et al. Proc Natl Acad Sci USA 2017;114.\newline \href{https://www.pnas.org/doi/10.1073/pnas.1717421114}{pnas.\allowbreak org/\allowbreak doi/\allowbreak 10.\allowbreak 1073/\allowbreak pnas.\allowbreak 1717421\allowbreak 114} & Regulatory DNA band: variants stratified by the regulatory element they sit in, which is the only released set with that stratification. & PNAS sup\allowbreak plementa\allowbreak ry data \allowbreak file pna\allowbreak s.\allowbreak 1717421\allowbreak 114.\allowbreak sd01; n\allowbreak eeds a s\allowbreak upplemen\allowbreak t fetch\newline {\scriptsize (browser)} \\
\addlinespace[5pt]
32 & \textbf{Chen (earlier position-effect study, 482 loci)} \org{S. cerevisiae}\newline An earlier Chen study characterizing position effects at 482 loci, reported inside Wu 2017; primary citation not located.\newline \href{https://www.ncbi.nlm.nih.gov/pmc/articles/PMC5516366/}{PMC55163\allowbreak 66} & Regulatory DNA band by its sequence variable, the integration coordinate. & not loca\allowbreak ted; the\allowbreak  primary\allowbreak  paper m\allowbreak ust be i\allowbreak dentifie\allowbreak d first\newline {\scriptsize (browser)} \\
\addlinespace[5pt]
33 & \textbf{Wei 2017 (synthetic terminator panel)} \org{S. cerevisiae}\newline Wei, et al. Yeast synthetic terminators: fine regulation of strength through linker sequences. ChemBioChem 2017.\newline \href{https://doi.org/10.1002/cbic.201700516}{doi:\allowbreak 10.\allowbreak 1002/\allowbreak cbic.\allowbreak 2017005\allowbreak 16} & Regulatory DNA band: the terminator parts tier, whose values are relative to a standard part rather than absolute. & ChemBioC\allowbreak hem jour\allowbreak nal supp\allowbreak lement, \allowbreak no repos\allowbreak itory de\allowbreak posit fo\allowbreak und\newline {\scriptsize (browser)} \\
\addlinespace[5pt]
34 & \textbf{Thompson 2019 (KT2440 lysine metabolism)} \org{P. putida}\newline Thompson MG, et al. Massively parallel fitness profiling reveals multiple novel enzymes in Pseudomonas putida lysine metabolism. mBio 2019.\newline \href{https://pmc.ncbi.nlm.nih.gov/articles/PMC6509195}{PMC65091\allowbreak 95} & Scale band: tier 3, ordered by measurements (22,256). & Table S1\allowbreak , a CSV \allowbreak suppleme\allowbreak ntary fi\allowbreak le; no r\allowbreak epositor\allowbreak y\newline {\scriptsize (browser)} \\
\addlinespace[5pt]
35 & \textbf{Morse (short synthetic terminators)} \org{S. cerevisiae}\newline Morse NJ, et al. Short synthetic terminators, 35 to 70 bp; primary citation not located.\newline \href{https://scite.ai/authors/nicholas-j-morse-1g5Qrn}{scite.\allowbreak ai/\allowbreak authors\allowbreak /\allowbreak nichola\allowbreak s-\allowbreak j-\allowbreak morse-\allowbreak 1g5Qrn} & Regulatory DNA band: the short-terminator corner of the parts tier, held so the lead is not lost. & no repos\allowbreak itory de\allowbreak posit fo\allowbreak und; pri\allowbreak mary pap\allowbreak er not i\allowbreak dentifie\allowbreak d\newline {\scriptsize (browser)} \\
\addlinespace[5pt]
36 & \textbf{MaveDB yeast score sets} \org{S. cerevisiae}\newline Esposito D, et al. MaveDB: an open-source platform to distribute and interpret data from multiplexed assays of variant effect. Genome Biol 2019;20:223.\newline \href{https://doi.org/10.1186/s13059-019-1845-6}{doi:\allowbreak 10.\allowbreak 1186/\allowbreak s13059-\allowbreak 019-\allowbreak 1845-\allowbreak 6} & Regulatory DNA band as the second aggregate route into the mutational-scan tier, held because the two aggregates do not hold the same assays. & mavedb.\allowbreak org; th\allowbreak e docume\allowbreak nted sea\allowbreak rch API \allowbreak endpoint\allowbreak  returns\allowbreak  HTTP 40\allowbreak 4, so th\allowbreak e organi\allowbreak sm filte\allowbreak r needs \allowbreak a live q\allowbreak uery\newline {\scriptsize (browser)} \\
\addlinespace[5pt]
37 & \textbf{Isoprenol CRISPRi DBTL proteome ladder (7 PRIDE deposits)} \org{P. putida}\newline Automation of a CRISPRi platform for enhanced isoprenol production in Pseudomonas putida; PRIDE projects PXD063733, PXD063737, PXD063738, PXD063740, PXD063743, PXD063744, PXD063746.\newline \href{https://www.ebi.ac.uk/pride/archive/projects/PXD063733}{ebi.\allowbreak ac.\allowbreak uk/\allowbreak pride/\allowbreak archive\allowbreak /\allowbreak project\allowbreak s/\allowbreak PXD0637\allowbreak 33} & Scale band: tier 3, ordered by measurements (38,948). & PRIDE PX\allowbreak D063733,\allowbreak  PXD0637\allowbreak 37, PXD0\allowbreak 63738, P\allowbreak XD063740\allowbreak , PXD063\allowbreak 743, PXD\allowbreak 063744 a\allowbreak nd PXD06\allowbreak 3746, al\allowbreak l submit\allowbreak ted 2025\allowbreak -\allowbreak 05-\allowbreak 07 and \allowbreak 08 on an\allowbreak  Orbitra\allowbreak p Explor\allowbreak is 480 \\
\addlinespace[5pt]
38 & \textbf{Thompson 2020 (KT2440 fatty acids and alcohols)} \org{P. putida}\newline Thompson MG, Incha MR, Pearson AN, et al. Appl Environ Microbiol 2020;86:e01665-20; correction at 10.1128/AEM.00177-21.\newline \href{https://pubmed.ncbi.nlm.nih.gov/32826213}{PMID 328\allowbreak 26213} & Scale band: tier 3, ordered by measurements (127,972). & not stat\allowbreak ed; the \allowbreak values m\allowbreak ost like\allowbreak ly reach\allowbreak  feba.\allowbreak db \\
\addlinespace[5pt]
39 & \textbf{GSE198395 (putidaPRECISE321 carbon sources)} \org{P. putida}\newline GEO series GSE198395, KT2440 and engineered strains on different carbon sources, part of putidaPRECISE321, released 2022-04-28.\newline \href{https://www.ncbi.nlm.nih.gov/geo/query/acc.cgi?acc=GSE198395}{ncbi.\allowbreak nlm.\allowbreak nih.\allowbreak gov/\allowbreak geo/\allowbreak query/\allowbreak acc.\allowbreak cgi?\allowbreak acc=\allowbreak GSE1983\allowbreak 95} & Scale band: tier 3, ordered by measurements (450,684). & GEO GSE1\allowbreak 98395; p\allowbreak rogramma\allowbreak tically \allowbreak retrieva\allowbreak ble \\
\addlinespace[5pt]
40 & \textbf{Mueller 2025 (24-strain P. putida Biolog panel)} \org{P. putida}\newline Mueller J, Krishnan J, Wei Q, Hefner Y, Monk JM, Verkler H, Tibocha-Bonilla JD, et al. bioRxiv 2025;2025.11.17.688870.\newline \href{https://doi.org/10.1101/2025.11.17.688870}{doi:\allowbreak 10.\allowbreak 1101/\allowbreak 2025.\allowbreak 11.\allowbreak 17.\allowbreak 688870} & Scale band: tier 2, ordered by measurements (4,920). & unconfir\allowbreak med:\allowbreak  bioRxi\allowbreak v return\allowbreak ed HTTP \allowbreak 429 on f\allowbreak our atte\allowbreak mpts, so\allowbreak  no BioP\allowbreak roject, \allowbreak Zenodo r\allowbreak ecord or\allowbreak  machine\allowbreak -\allowbreak readabl\allowbreak e Table \allowbreak S2 is co\allowbreak nfirmed\newline {\scriptsize (ask)} \\
\addlinespace[5pt]
41 & \textbf{Schmidt 2022 (KT2440 nitrogen BarSeq)} \org{P. putida}\newline Schmidt M, Pearson AN, Incha MR, et al. Appl Environ Microbiol 2022;88:e02430-21.\newline \href{https://pmc.ncbi.nlm.nih.gov/articles/PMC9004399}{PMC90043\allowbreak 99} & Scale band: tier 2, ordered by measurements (395,044). & no acces\allowbreak sion sta\allowbreak ted; the\allowbreak  study l\allowbreak inks int\allowbreak o the Fi\allowbreak tness Br\allowbreak owser, s\allowbreak o its ex\allowbreak periment\allowbreak s are mo\allowbreak st likel\allowbreak y inside\allowbreak  feba.\allowbreak db \\
\addlinespace[5pt]
42 & \textbf{Borchert 2024 (KT2440 fitness compendium)} \org{P. putida}\newline Borchert AJ, Bleem AC, Lim HG, Rychel K, Dooley KD, Kellermyer ZA, Hodges TL, Palsson BO, Beckham GT. mSystems 2024;9:e00942-23.\newline \href{https://doi.org/10.1128/msystems.00942-23}{doi:\allowbreak 10.\allowbreak 1128/\allowbreak msystem\allowbreak s.\allowbreak 00942-\allowbreak 23} & Scale band: tier 1, ordered by measurements (1,571,024). & github.\allowbreak com/\allowbreak beckham\allowbreak -\allowbreak lab/\allowbreak fModule\allowbreak  (fitnes\allowbreak s values\allowbreak , statis\allowbreak tics, me\allowbreak tadata);\allowbreak  the non\allowbreak -\allowbreak browser\allowbreak  samples\allowbreak  are at \allowbreak SRA PRJN\allowbreak A809672,\allowbreak  PRJNA85\allowbreak 6070, PR\allowbreak JNA10112\allowbreak 87 \\
\addlinespace[5pt]
43 & \textbf{putidaPRECISE321 (Lim 2022)} \org{P. putida}\newline Lim HG, Rychel K, Sastry AV, Mueller J, Niu W, Feist AM, Palsson BO. Metab Eng 2022;72.\newline \href{https://doi.org/10.1016/j.ymben.2022.04.004}{doi:\allowbreak 10.\allowbreak 1016/\allowbreak j.\allowbreak ymben.\allowbreak 2022.\allowbreak 04.\allowbreak 004} & Scale band: tier 1, ordered by measurements (1,786,044). & github.\allowbreak com/\allowbreak SBRG/\allowbreak modulom\allowbreak e\_\allowbreak ppu (MI\allowbreak T); name\allowbreak d GEO de\allowbreak posits G\allowbreak SE198395\allowbreak  (81 sam\allowbreak ples) an\allowbreak d GSE193\allowbreak 493 (22)\allowbreak ; the fu\allowbreak ll per-\allowbreak sample \allowbreak GEO list\allowbreak  is in t\allowbreak he SI sp\allowbreak readshee\allowbreak t \\
\addlinespace[5pt]
44 & \textbf{GSE262836 (Delhi NICU, fluconazole-resistance expression)} \org{P. kudriavzevii}\newline Genomic characterization and transcriptomic profile of Candida krusei in a neonatal intensive care unit in Delhi, India. GEO series GSE262836, 2025.\newline \href{https://www.ncbi.nlm.nih.gov/geo/query/acc.cgi?acc=GSE262836}{ncbi.\allowbreak nlm.\allowbreak nih.\allowbreak gov/\allowbreak geo/\allowbreak query/\allowbreak acc.\allowbreak cgi?\allowbreak acc=\allowbreak GSE2628\allowbreak 36} & Host bootstrap: one of two rows in this host pairing a drug phenotype with expression, and openly retrievable. & GEO GSE2\allowbreak 62836, p\allowbreak latform \allowbreak GPL34351\allowbreak ; progra\allowbreak mmatical\allowbreak ly retri\allowbreak evable \\
\addlinespace[5pt]
45 & \textbf{GSE159457 (CY902 C4 dicarboxylic acid transporters)} \org{P. kudriavzevii}\newline Identification and characterization of C4 dicarboxylic acid transporters in Pichia kudriavzevii CY902. GEO series GSE159457, 2021; PMID 33629807.\newline \href{https://www.ncbi.nlm.nih.gov/geo/query/acc.cgi?acc=GSE159457}{ncbi.\allowbreak nlm.\allowbreak nih.\allowbreak gov/\allowbreak geo/\allowbreak query/\allowbreak acc.\allowbreak cgi?\allowbreak acc=\allowbreak GSE1594\allowbreak 57} & Host bootstrap: a carbon-source axis on expression in a third background, fully deposited. & GEO GSE1\allowbreak 59457, p\allowbreak latform \allowbreak GPL29244\allowbreak ; raw re\allowbreak ads at S\allowbreak RA SRP28\allowbreak 7255; pr\allowbreak ogrammat\allowbreak ically r\allowbreak etrievab\allowbreak le \\
\addlinespace[5pt]
46 & \textbf{Dubinkina 2024 (12-strain low-pH transcriptomic atlas)} \org{P. kudriavzevii}\newline Dubinkina V, Bhogale S, Hsieh PH, Dibaeinia P, Nambiar S, Maslov S, Yoshikuni Y, Sinha S. Microbiol Spectr 2024.\newline \href{https://doi.org/10.1128/spectrum.02536-23}{doi:\allowbreak 10.\allowbreak 1128/\allowbreak spectru\allowbreak m.\allowbreak 02536-\allowbreak 23} & Host bootstrap: the only multi-strain expression matrix this host has, and the only deposit that already carries processed values rather than reads. & Illinois\allowbreak  Data Ba\allowbreak nk DOI 1\allowbreak 0.\allowbreak 13012/\allowbreak B2IDB-\allowbreak 1336177\_\allowbreak V1, hol\allowbreak ding str\allowbreak ain regu\allowbreak lators, \allowbreak strains \allowbreak and gene\allowbreak  sets, a\allowbreak nd gene \allowbreak expressi\allowbreak on count\allowbreak s; the p\allowbreak age retu\allowbreak rns 403 \allowbreak to every\allowbreak  automat\allowbreak ed clien\allowbreak t\newline {\scriptsize (browser)} \\
\addlinespace[5pt]
47 & \textbf{MAESTRO 2025 (170-isolate WGS, 57-stressor growth panel)} \org{P. kudriavzevii}\newline Hsieh PH, Sasaki Y, Yang CI, Wu ZY, Steindorff AS, Haridas S, Ke J, Fatma Z, Zhao Z, Opulente DA, Deng S, Hittinger CT, Grigoriev IV, Dien B, Zhao H, Li YP, Yoshikuni Y. bioRxiv 2025;2025.08.29.673034.\newline \href{https://doi.org/10.1101/2025.08.29.673034}{doi:\allowbreak 10.\allowbreak 1101/\allowbreak 2025.\allowbreak 08.\allowbreak 29.\allowbreak 673034} & Host bootstrap: this host has no genome-wide screen of any kind, so its largest resource is a 161-isolate panel that the measurement rule ranks below every pseudomonad screen. The band is what keeps it in the first tranche. & no acces\allowbreak sion ret\allowbreak rieved:\allowbreak  bioRxi\allowbreak v return\allowbreak s 403 an\allowbreak d Europe\allowbreak  PMC ful\allowbreak lTextXML\allowbreak  for PPR\allowbreak 1078985 \allowbreak returns \allowbreak 500; per\allowbreak -\allowbreak isolate\allowbreak  IO\_\allowbreak NNN Bio\allowbreak Projects\allowbreak  are inf\allowbreak erred, n\allowbreak ot confi\allowbreak rmed\newline {\scriptsize (ask)} \\
\addlinespace[5pt]
48 & \textbf{Fatma 2023 (landing-pad locus screen)} \org{P. kudriavzevii}\newline Fatma Z, Tan SI, Boob A, Zhao H. A landing pad system for multicopy gene integration in Issatchenkia orientalis. bioRxiv 2023;2023.05.21.541627.\newline \href{https://doi.org/10.1101/2023.05.21.541627}{doi:\allowbreak 10.\allowbreak 1101/\allowbreak 2023.\allowbreak 05.\allowbreak 21.\allowbreak 541627} & Host bootstrap: the host's only measured per-locus fitness, which no scale ranking would ever place in a first tranche at 188 records. & a 23 KB \allowbreak package \allowbreak at uofi.\allowbreak box.\allowbreak com (st\allowbreak rains, p\allowbreak lasmids,\allowbreak  guide R\allowbreak NAs, pri\allowbreak mers) pl\allowbreak us githu\allowbreak b.\allowbreak com/\allowbreak HuiminZ\allowbreak hao/\allowbreak Landing\allowbreak -\allowbreak pad-\allowbreak model; \allowbreak no DOI a\allowbreak nd no ar\allowbreak chival d\allowbreak eposit\newline {\scriptsize (ask)} \\
\addlinespace[5pt]
49 & \textbf{GSE282019 (Delhi NICU five-year outbreak)} \org{P. kudriavzevii}\newline Genomic analyses reveal high diversity and rapid evolution of Pichia kudriavzevii within a neonatal intensive care unit in Delhi, India. GEO series GSE282019, 2024; PMID 39853119.\newline \href{https://www.ncbi.nlm.nih.gov/geo/query/acc.cgi?acc=GSE282019}{ncbi.\allowbreak nlm.\allowbreak nih.\allowbreak gov/\allowbreak geo/\allowbreak query/\allowbreak acc.\allowbreak cgi?\allowbreak acc=\allowbreak GSE2820\allowbreak 19} & Host bootstrap: the host's only published drug-stratified expression series, and openly retrievable. & GEO GSE2\allowbreak 82019, p\allowbreak latform \allowbreak GPL34351\allowbreak ; progra\allowbreak mmatical\allowbreak ly retri\allowbreak evable \\
\addlinespace[5pt]
50 & \textbf{GSE202655 (L-malic acid, EcSthA transhydrogenase)} \org{P. kudriavzevii}\newline Metabolic engineering of the acid-tolerant yeast Pichia kudriavzevii for efficient L-malic acid production at low pH. GEO series GSE202655, 2023; PMID 36566973.\newline \href{https://www.ncbi.nlm.nih.gov/geo/query/acc.cgi?acc=GSE202655}{ncbi.\allowbreak nlm.\allowbreak nih.\allowbreak gov/\allowbreak geo/\allowbreak query/\allowbreak acc.\allowbreak cgi?\allowbreak acc=\allowbreak GSE2026\allowbreak 55} & Host bootstrap: an engineered-strain expression series in a non-SD108 background, which the host otherwise lacks. & GEO GSE2\allowbreak 02655, p\allowbreak latform \allowbreak GPL32243\allowbreak ; progra\allowbreak mmatical\allowbreak ly retri\allowbreak evable \\
\addlinespace[5pt]
51 & \textbf{Lubliner 2013 (native core promoter maximal strengths)} \org{S. cerevisiae}\newline Lubliner S, Keren L, Segal E. Nucleic Acids Res 2013.\newline \href{https://academic.oup.com/nar}{academic\allowbreak .\allowbreak oup.\allowbreak com/\allowbreak nar} & Regulatory DNA band: the native counterpart of the core-promoter tier, and the only source for its reference values. & not loca\allowbreak ted\newline {\scriptsize (browser)} \\
\addlinespace[5pt]
52 & \textbf{Fitness Browser RCH2 (Stutzerimonas stutzeri\_AE)} \org{S. stutzeri}\newline Price MN, Wetmore KM, Waters RJ, Callaghan M, Ray J, Liu H, Kuehl JV, Melnyk RA, Lamson JS, Suh Y, et al., Deutschbauer AM, Arkin AP. Nature 2018;557:503-509.\newline \href{https://doi.org/10.1038/s41586-018-0124-0}{doi:\allowbreak 10.\allowbreak 1038/\allowbreak s41586-\allowbreak 018-\allowbreak 0124-\allowbreak 0} & Cross-host fitness: nine new hosts arrive inside one artifact under one schema, one condition vocabulary and an identifier bridge the source supplies, so the band exists to keep them together rather than scattered down a measurement ranking. & feba.\allowbreak db at 1\allowbreak 0,433,70\allowbreak 8,032 by\allowbreak tes, Las\allowbreak t-\allowbreak Modifie\allowbreak d 2026-\allowbreak 10-\allowbreak 27, HTT\allowbreak P 200 to\allowbreak  a plain\allowbreak  curl; o\allowbreak rgId psR\allowbreak CH2; ref\allowbreak seq\_\allowbreak mapping\allowbreak .\allowbreak tsv bri\allowbreak dges loc\allowbreak usId to \allowbreak RefSeq \\
\addlinespace[5pt]
53 & \textbf{Fitness Browser FW300-N2E3 (Aquipseudomonas helvetica)} \org{P. fluorescens}\newline Price MN, Wetmore KM, Waters RJ, et al., Deutschbauer AM, Arkin AP. Nature 2018;557:503-509.\newline \href{https://doi.org/10.1038/s41586-018-0124-0}{doi:\allowbreak 10.\allowbreak 1038/\allowbreak s41586-\allowbreak 018-\allowbreak 0124-\allowbreak 0} & Cross-host fitness: part of the single-artifact block of nine hosts, whose shared schema is the reason to admit them together. & feba.\allowbreak db, org\allowbreak Id pseud\allowbreak o3\_\allowbreak N2E3; t\allowbreak he histo\allowbreak rical bi\allowbreak nomial i\allowbreak s P.\allowbreak  fluore\allowbreak scens an\allowbreak d the cu\allowbreak rrent on\allowbreak e is A.\allowbreak  helvet\allowbreak ica \\
\addlinespace[5pt]
54 & \textbf{Fitness Browser FW300-N2E2 (Aquipseudomonas fluorescens\_Q)} \org{P. fluorescens}\newline Price MN, Wetmore KM, Waters RJ, et al., Deutschbauer AM, Arkin AP. Nature 2018;557:503-509.\newline \href{https://doi.org/10.1038/s41586-018-0124-0}{doi:\allowbreak 10.\allowbreak 1038/\allowbreak s41586-\allowbreak 018-\allowbreak 0124-\allowbreak 0} & Cross-host fitness: one of four isolates of one species in the block, which is the within-species comparison the band makes possible. & feba.\allowbreak db, org\allowbreak Id pseud\allowbreak o6\_\allowbreak N2E2; 1\allowbreak 16 updat\allowbreak ed annot\allowbreak ations \\
\addlinespace[5pt]
55 & \textbf{Fitness Browser WCS417 (Aquipseudomonas simiae)} \org{P. simiae}\newline Price MN, Wetmore KM, Waters RJ, et al., Deutschbauer AM, Arkin AP. Nature 2018;557:503-509.\newline \href{https://doi.org/10.1038/s41586-018-0124-0}{doi:\allowbreak 10.\allowbreak 1038/\allowbreak s41586-\allowbreak 018-\allowbreak 0124-\allowbreak 0} & Cross-host fitness: the plant-associated member of the single-artifact block, with its widest carbon panel. & feba.\allowbreak db, org\allowbreak Id WCS41\allowbreak 7; histo\allowbreak rical bi\allowbreak nomial P\allowbreak .\allowbreak  simiae\allowbreak , curren\allowbreak t A.\allowbreak  simiae \\
\addlinespace[5pt]
56 & \textbf{Fitness Browser FW300-N2C3 (Aquipseudomonas fluorescens\_AA)} \org{P. fluorescens}\newline Price MN, Wetmore KM, Waters RJ, et al., Deutschbauer AM, Arkin AP. Nature 2018;557:503-509.\newline \href{https://doi.org/10.1038/s41586-018-0124-0}{doi:\allowbreak 10.\allowbreak 1038/\allowbreak s41586-\allowbreak 018-\allowbreak 0124-\allowbreak 0} & Cross-host fitness: part of the four-isolate within-species block, and the row that taught the most unannotated genes. & feba.\allowbreak db, org\allowbreak Id pseud\allowbreak o5\_\allowbreak N2C3\_\allowbreak 1 \\
\addlinespace[5pt]
57 & \textbf{Fitness Browser FW300-N1B4 (Aquipseudomonas sp.)} \org{P. fluorescens}\newline Price MN, Wetmore KM, Waters RJ, et al., Deutschbauer AM, Arkin AP. Nature 2018;557:503-509.\newline \href{https://doi.org/10.1038/s41586-018-0124-0}{doi:\allowbreak 10.\allowbreak 1038/\allowbreak s41586-\allowbreak 018-\allowbreak 0124-\allowbreak 0} & Cross-host fitness: completes the four-isolate within-species block, and names the case of a host with no species assignment. & feba.\allowbreak db, org\allowbreak Id pseud\allowbreak o1\_\allowbreak N1B4 \\
\addlinespace[5pt]
58 & \textbf{Fitness Browser B728a (Aquipseudomonas fragariae pv. syringae)} \org{P. syringae}\newline Price MN, Wetmore KM, Waters RJ, et al., Deutschbauer AM, Arkin AP. Nature 2018;557:503-509.\newline \href{https://doi.org/10.1038/s41586-018-0124-0}{doi:\allowbreak 10.\allowbreak 1038/\allowbreak s41586-\allowbreak 018-\allowbreak 0124-\allowbreak 0} & Cross-host fitness: the pathogen member of the block, and half of the one strain in it screened in two genetic backgrounds. & feba.\allowbreak db, org\allowbreak Id Syrin\allowbreak gaeB728a\allowbreak ; the GT\allowbreak DB renam\allowbreak e to A.\allowbreak  fragar\allowbreak iae pv.\allowbreak  syring\allowbreak ae means\allowbreak  both bi\allowbreak nomials \allowbreak are need\allowbreak ed as al\allowbreak iases \\
\addlinespace[5pt]
59 & \textbf{Fitness Browser B728a mexB deletion background} \org{P. syringae}\newline Price MN, Wetmore KM, Waters RJ, et al., Deutschbauer AM, Arkin AP. Nature 2018;557:503-509.\newline \href{https://doi.org/10.1038/s41586-018-0124-0}{doi:\allowbreak 10.\allowbreak 1038/\allowbreak s41586-\allowbreak 018-\allowbreak 0124-\allowbreak 0} & Cross-host fitness: the background-crossed half of the B728a pair, which is the only interaction design in the block. & feba.\allowbreak db, org\allowbreak Id Syrin\allowbreak gaeB728a\_\allowbreak mexBdel\allowbreak ta \\
\addlinespace[5pt]
60 & \textbf{Fitness Browser PA14 (Pseudomonas aeruginosa)} \org{P. aeruginosa}\newline Price MN, Wetmore KM, Waters RJ, et al., Deutschbauer AM, Arkin AP. Nature 2018;557:503-509.\newline \href{https://doi.org/10.1038/s41586-018-0124-0}{doi:\allowbreak 10.\allowbreak 1038/\allowbreak s41586-\allowbreak 018-\allowbreak 0124-\allowbreak 0} & Cross-host fitness: the same artifact and schema, kept visible because its size is the finding about this species rather than an omission. & feba.\allowbreak db, org\allowbreak Id PA14;\allowbreak  the bro\allowbreak wser doe\allowbreak s not re\allowbreak name thi\allowbreak s organi\allowbreak sm \\
\addlinespace[5pt]
\end{longtable}
\endgroup
\end{landscape}


%% notes-tex/database/database-expansion-hosts/sections/caveats.tex
%% [[experiments.database.expansion-hosts]]
%% https://github.com/Mjvolk3/torchcell/tree/main/notes-tex/database/database-expansion-hosts/sections/caveats.tex
%% SOURCE: hand-authored prose; the excluded and asks tables are generated and name their own
%% scripts. The verified, snippet and inferred labels, the retrieval classes and the HTTP
%% responses quoted here come from the triage report
%% /scratch/projects/torchcell-scratch/tmp/deep-research/reports/Pk and next dataset wave triage.md
%% and its research notes; the gate and verdict-store figures from
%% torchcell/candidates/ and database/candidates/ as recorded in
%% .../research_notes/Pk and next dataset wave triage/local_triage_scaffold.md

\section{What was dropped\secstatus{todo}}\label{sec:excluded}

%% GENERATED FILE -- do not hand-edit.
%% SOURCE: experiments/database/scripts/build_hosts_candidate_datasets_table.py
\begingroup
\footnotesize
\begin{longtable}{@{}L{58mm} L{26mm} L{88mm}@{}}
\caption[]{Considered and dropped. \emph{no-sequence} is the hard gate: without a route
to the strain's genomic content there is no genotype to map a phenotype from.
\emph{off-host} is a real dataset in an organism outside this wave's hosts.
\emph{not-a-dataset} is a model, a design program, a review or a prior over the genome.
\emph{no-per-record-data} is a real experiment whose released form is a figure or a
summary rather than per-record values. \emph{already-built} is served, or already in the
raw mirror, so it is an overlap rather than a candidate. \emph{derivative} re-serves or
re-processes another dataset's values, so its records are not net new.}
\label{tab:excluded}\\
\toprule
Dataset or group & Rule & Reason \\
\midrule
\endfirsthead
\toprule
Dataset or group & Rule & Reason \\
\midrule
\endhead
\bottomrule
\endfoot
Beltran 2025 aPCA domainome & off-host & 563,534 variants across 522 protein domains, measured in BY4741 yeast cells with a yeast growth readout, which is why it surfaced as a yeast candidate at all. 503 of those domains are HUMAN proteins. The cells are yeast and the genotype is not: the perturbed sequence is a human domain carried on a plasmid, keyed by UniProt accession and Pfam position, so ingesting it under this organism would file half a million human variants as S. cerevisiae genotypes. It is a strong dataset for a human-domain tier and a category error for this one \\
\addlinespace[5pt]
Erijman 2020 (random activation-domain peptides) & no-per-record-data & About a million random peptides scored for transcription activation, which would have ranked tenth by size. Two things block it. The release count is never stated, so the million is the paper's own figure rather than a counted release, and the processed tables are not reachable: the path returns a 1,817-byte challenge page rather than data, and the code repository holds model weights only. The reads are deposited, so the phenotype would have to be recomputed from them, which makes it a reanalysis and not an ingest \\
\addlinespace[5pt]
Li 2019 (Pareto fronts in trait space) & derivative & 259,357 barcode count rows and 3,717 clones scored across four serial-transfer conditions, fully retrievable and well deposited. It is dropped on OVERLAP rather than on quality: 3,048 of its 3,717 clones, which is 82 percent, come from a dataset already on this database's exclusion list, so most of what it would add is a second scoring of strains already accounted for. The four-condition fitness vector on the remaining 669 clones is the only net-new part, which is too thin to carry a row \\
\addlinespace[5pt]
Yeast fill-pass measured negatives (nine sources checked and unusable) & no-per-record-data & Recorded because each was checked and the result is a finding rather than an absence. Dudley 2005 releases a three-level ordinal phenotype in twelve PDF supplements, and its only data pointer is a dead host. Roy 2024 names no processed table and its BioProject returns zero experiments. Bakerlee 2022 has the right shape but every route to its fitness tables is refused and its deposit holds genotyping alone. Kryazhimskiy 2014, Jerison 2017 and Johnson 2021 are not barcoded at all, which was checked by searching their body text, and use flow-cytometry competition instead, so they are good fitness releases for a different wave. Collins 2007, Fiedler 2009 and Bodenmiller 2010 have no locatable data file or a dead host. Further negatives worth keeping: no yeast messenger-RNA-decay library exists, no yeast degron atlas exceeds 23,600 peptides because the large ones are human, and no genome-wide chemical or nutrient panel above roughly 300,000 records appeared for 2015 to 2026 beyond the 102-condition row ranked above \\
\addlinespace[5pt]
ALEdb (Phaneuf 2019), the adaptive-laboratory-evolution database & not-a-dataset & Checked live on 2026-11-24 because it was the one aggregator that could plausibly have closed this wave's row gap on its own, and it releases ZERO phenotype numbers. The site states 3,721,969 instances of 1,035,104 unique mutations from 13,425 isolates across 271 projects, roughly 94 times the 11,792 mutations its own 2019 paper reported, so it has grown enormously. But a grep for growth rate, fitness, inhibitory concentration and titer across every reachable page and export returned nothing: the exports carry mutations, media, temperature, strain and sequencing quality only. Retrieval is also degraded, with no API, an export route that errors even with a fresh token, and only 107 of the 271 projects browsable anonymously. It is a GENOTYPE source for evolved clones, which is how a later pass should use it, not a dataset. One naming trap: one of its project codes is a Plasmodium study, not E. coli \\
\addlinespace[5pt]
COLOMBOS v3, Bastkowski 2026, Yasir 2026, Reg-Seq, Good 2017 and Tenaillon 2016 & no-per-record-data & Six E. coli objects that each looked larger than anything ranked and released no per-record phenotype. COLOMBOS v3 would have been the biggest expression object in the wave at 4,077 condition contrasts, but its domain times out over both protocols, so it has no resolvable URL at all. Bastkowski 2026 deposits only raw reads for its 13 antibiotics at 4 concentrations, with no per-gene fitness table. Yasir 2026 is a 754-sample transposon compendium, the highest-count such object found, and gives no accession of any kind. Reg-Seq releases sequencing reads and information-footprint matrices rather than per-variant values. Good 2017 has 41,962 rows but no phenotype, and its genotype is a population allele frequency rather than a clone, which fails the definition this wave ranks on. Tenaillon 2016 releases no per-clone fitness, only a 791-byte summary file \\
\addlinespace[5pt]
Kosuri 2013, Goodman 2013 and Moradigaravand 2018 & derivative & Three sets reachable only by a route this database does not take. The first two are blocked at the publisher, and the copies that do answer are a third-party compilation and a re-inferred mirror, so ingesting either would store another analysis's numbers rather than the release. Moradigaravand 2018's only surfaced deposit route is Mendeley Data, which is ruled out as a live source. In all three cases the primary release, not the row, is what is missing \\
\addlinespace[5pt]
Kildegaard 2016 and Chen 2026 (figure-bound 3-HP panels) & no-per-record-data & Two real open-access S. cerevisiae 3-HP experiments with no per-strain numbers to read. Kildegaard 2016 (10.1186/s12934-016-0451-5) measured 14 strains and tabulated only the one fed-batch strain in Table 1. Chen 2026 (10.1038/s41467-026-76008-2) reports 48 colonies, and it DOES ship a Source Data file, which is why a first read called it ingestible, but the values in the fig4g sheet are negative, so they are a relative scale rather than titers and cannot be stored as a concentration. Digitizing a bar, or storing a relative value as though it were a measurement, would both invent precision the source never stated. Two passes over this literature agreed on the structural finding behind these rows: of roughly twenty yeast 3-HP papers whose full text was read, only four ship machine-readable per-strain numbers \\
\addlinespace[5pt]
3-HP in adjacent yeast hosts (O. polymorpha, K. phaffii, S. pombe, Y. lipolytica, R. toruloides, K. marxianus, D. hansenii) & off-host & Several reach titers no S. cerevisiae result matches, led by Ogataea polymorpha at 79.6 g/L and Debaryomyces hansenii at 62.4 g/L, the latter from a natural isolate plus ion mutagenesis rather than engineering. That last figure does NOT survive arithmetic: Li 2020 reports 62.42 g/L from 40 g/L of total substrate, which fails mass balance, so it is recorded as a flagged number and not as a target. They are recorded rather than dropped because the 3-HP product line is the reason this wave touched the chemistry at all, and these are the rows a later wave would add. Admitting any of them means seven new Organism members and seven reference releases pinned in the genomes tier, which is a scope decision for the owner and not a fill \\
\addlinespace[5pt]
Illinois Data Bank IDB-9063411 (high-yield 3-HP) & not-a-dataset & Ranked in this wave until the fill pass read it, and it is not a phenotype dataset at all: the deposit holds the technoeconomic model and the genome-scale model code, as two source repositories, not the strain table. The measurements the row was ranked for, 92 g/L at 0.7 g/g and 0.55 g/L/h, are in Tan 2026, which is now ranked on its own identity with the Source Data file behind it. Keeping both would have counted one experiment twice \\
\addlinespace[5pt]
LASER strain-engineering database (Winkler 2015, 2016) & derivative & Counted from the bytes of the only live release, on Bitbucket: 456 curated papers, 638 designs, 441 E. coli and 197 S. cerevisiae, frozen since 2016-04-27, and the web interface the paper names no longer resolves. Every number in it is re-extracted from a primary paper, so ingesting it would re-serve other datasets' values. It is also thinner than it looks: titer is populated on 181 of 197 yeast designs, yield on 58, there is no rate field at all, and no replicates, no n, no uncertainty, no sequence-level edit and no strain accession. Its own total field sums to 3,472,162 strains built against 638 curated, because it keeps champion strains and discards the screen populations. Its value is as an index onto primary papers, which is how the fill pass used it, and three of those papers are ranked above \\
\addlinespace[5pt]
Charusanti 2010 (pgi-deletion adaptive evolution) & no-per-record-data & The largest pure evolution entry the LASER index points at, with ten endpoint clones and nine resequenced genomes, so the GENOTYPES are released. The phenotype is not: the supplement holds a peroxidase assay and primers, and the per-clone growth rates exist only inside Figure 2A. A figure is not a released per-record value, and digitizing one would invent precision the source never stated \\
\addlinespace[5pt]
scYeast (Fan 2027) figshare deposit & derivative & A model and its training corpus rather than a new measurement. Its attention prior is built from the documented-targets subset of YEASTRACT+, about 170,000 pairs, which is a prior over the genome and not a genotype-by-environment record, and the expression it trains on is the built yeast compendia. Nothing in the deposit is a released per-record value this store does not already hold. \\
\addlinespace[5pt]
Shorkie induction RNA-seq tracks (gs://shorkie-paper/) & no-per-record-data & The 3,053 induction RNA-seq tracks generated at Calico are the one genuinely new measurement in the paper and they are distributed only as BigWig and TFRecord on a Google Cloud Storage bucket, with no GEO, SRA, Zenodo or Hugging Face accession anywhere in the paper. A sha256-pinned mirror is the only reproducible route and there is no archive fallback if the bucket moves, so the right action is the ask in Table 8, not a parse. \\
\addlinespace[5pt]
Zrimec 2020 (expression prediction from regulatory sequence) & derivative & A modeling paper over expression values already held: its records are the public compendia re-encoded as model inputs, so admitting it would re-serve measurements the store has rather than add any. \\
\addlinespace[5pt]
Gu 2025 & derivative & A model trained on the regulatory-DNA libraries this wave ranks, not a new measurement. Its contribution is an architecture and a benchmark score, so there is no per-record value to store that is not already in de Boer, Vaishnav or the DREAM release. \\
\addlinespace[5pt]
Lei 2025 & derivative & Same class as Gu 2025: a sequence-to-expression model scored on the existing libraries. The report ranks the libraries and not the models, because a model's output is a prediction and the substrate stores measurements. \\
\addlinespace[5pt]
Cheng 2025 & derivative & A model over the same corpus again. Named here rather than omitted so the boundary is explicit: a benchmark result on a held-out split of a dataset already in the pool is not a second dataset. \\
\addlinespace[5pt]
iIsor850 genome-scale model (Issatchenkia orientalis) & not-a-dataset & A reconstruction, not an experiment. Its 850 gene calls are model output, which is the only form essentiality exists in for this species, and that is exactly why it cannot stand in for a screen. The yeast pool already excludes it under the same rule and the exclusion stands on its own merits even though this wave widens the host axis. \\
\addlinespace[5pt]
Wetmore 2015 RB-TnSeq methodology release & already-built & 3,903 genome-wide fitness experiments over 96,024 non-essential genes in 25 bacteria, and already in the database. It is the method paper the cross-host fitness band's artifact comes from, so admitting it again would re-serve values the store holds. \\
\addlinespace[5pt]
Fitness Browser GW456-L13 (Aquipseudomonas jessenii) & off-host & A real row in the same artifact as the nine hosts this wave admits, at 106 experiments with 80 updated annotations, and its species has no member in this pool's declared Organism vocabulary. Folding it into a neighboring species would be a silent sequence-level error in a substrate whose claim is sequence-level genotype fidelity, so it waits on the Organism decision rather than being mislabeled. Admit it in the same change that adds the other nine. \\
\addlinespace[5pt]
Gong 2018 (48-strain microsatellite panel, 9 antifungals) & no-sequence & 432 potential susceptibility values against nine drugs, which is more than twice the drug breadth of the sequenced panel, and the genotype is a microsatellite marker fingerprint rather than a sequence. There is no route to the strain's genomic content, so it cannot be joined to a genome even in principle. \\
\addlinespace[5pt]
Wiehlmann 2015 (2,921-isolate 58-marker survey) & no-sequence & 2,921 isolates from 1,448 habitats in 323 clone types, and the 58-marker binary profile is a genotype fingerprint rather than a sequence or a measured phenotype. It must not be ingested as either. \\
\addlinespace[5pt]
PkAC-7, PkAC-8 and PkAC-9 acetic-acid evolved strains & no-sequence & Three strains from acetic acid stepped 7 to 9 g/L, with tolerance measured against six stresses, and no resequencing, no variant table and no accession in either search pass. The genotype is undefined, so there is nothing to map a phenotype from; the 70-gene expression panel on PkAC-9 inherits the same gap. \\
\addlinespace[5pt]
RWT glucose osmotic-stress evolved mutant & no-sequence & Two strains, tolerant to 35 percent glucose with best growth at 20 percent, and the mutations were never characterized. Same gap as the other evolved panels, and the osmotic axis is already covered by the 57-stressor isolate panel. \\
\addlinespace[5pt]
Candida glycerinogenes as an ingestible host & no-sequence & A proprietary, independently owned industrial strain with no public assembly, and the species name was folded into P. kudriavzevii by the population-genomics study that showed one species under four names. Its published values cannot be tied to a public genome. \\
\addlinespace[5pt]
Blazeck 2012 (modular UAS and core hybrid promoters) & no-per-record-data & A real S. cerevisiae regulatory-parts panel in two backgrounds with biological triplicate and standard deviations, and every value exists only as a bar height: the supplementary table is primers. Unusable as released, and worth re-checking only if the hybrid-promoter design becomes the question. \\
\addlinespace[5pt]
Nikel 2015 (EDEMP cycle flux) & no-per-record-data & Reports that about 90 percent of consumed sugar is converted to gluconate and about 10 percent of triose phosphates are recycled to hexose phosphates, and the numbers sit in paper figures. Flux is the layer where the provenance discipline bites hardest, so the honest outcome is that the values are not retrievable rather than a thin extraction. \\
\addlinespace[5pt]
Puchalka 2008 (legacy KT2440 model validation) & no-per-record-data & Validated a genome-scale model against continuous culture and substrate utilization assays, and the utilization matrix itself is not confirmed to be a downloadable table. The validation is a model result; the matrix is what would be ingestible and it was not located. \\
\addlinespace[5pt]
PA14 pooled essential-gene knockdown library (dissertation only) & no-per-record-data & A pooled library with one essential gene knocked down per strain, tested in a murine pneumonia model, described only in a dissertation. No deposited score table was confirmed, so it is a lead rather than a citable dataset. \\
\addlinespace[5pt]
Mobile-CRISPRi and P. aeruginosa CRISPRi optimization methods & not-a-dataset & Vector systems and tuning studies: conjugation-transferable integrating CRISPRi vectors, dCas9 expression tuned to avoid toxicity, and guide designs for common laboratory strains. Methods rather than screens, with no per-record score tables. \\
\addlinespace[5pt]
Ozer 2019 (739-genome population structure) & not-a-dataset & 739 isolate genomes with source metadata, reporting that one lineage is 99.4 percent exoS-positive and the other 95.2 percent exoU-positive. Genotype only, with no measured phenotype, so it belongs beside the reference genomes rather than in a phenotype table. \\
\addlinespace[5pt]
43-strain international reference panel (BCCM/LMG) & not-a-dataset & 43 strains spanning cystic fibrosis, other clinical and environmental sources, 33 of them sequenced alongside 7 reference genomes. It is an order-the-strains resource rather than a phenotype table, and the genome accessions were not confirmed. \\
\addlinespace[5pt]
Murepavadin 414-isolate panel & not-a-dataset & Reported as 414 cystic-fibrosis isolates against 14 antibiotics with 412 genomes, and the citation itself is unconfirmed. The isolate count coincides with the best-documented clinical panel's, which is reason to resolve the citation before treating it as a second dataset. \\
\addlinespace[5pt]
SBW25 barcoded populations & not-a-dataset & Chromosomal barcodes for lineage tracking, which the authors explicitly distinguish from transposon sequencing: nothing is mutagenized. A different assay, whose value would be evolution and lineage phenotypes rather than gene fitness. \\
\addlinespace[5pt]
Pseudomonas Genome DB and the Pseudomonas Ortholog Database & not-a-dataset & Community annotation and orthology over 14,229 genomes, with ortholog groups and batch downloads, and no evidence it serves a phenotype table at all. It is the mapping substrate for the alias layer, which is what the Bakta and PA-number rows need, not a phenotype source. \\
\addlinespace[5pt]
GSE139245 (data-integrity flag) & not-a-dataset & A 15-sample series whose index record carries a summary describing mammalian histone-modification profiling, inconsistent with both its own title and its taxon. Treated as suspect rather than ingested. \\
\addlinespace[5pt]
Illinois Data Bank IDB-2926789 (constraint-based organic-acid potential) & not-a-dataset & Predicted production potential for 22 organic acids from a genome-scale model, which is the same class as the model-derived essential-gene list: model output, not measurement. The growth-coupled knockout designs built on the same model fall under this too. \\
\addlinespace[5pt]
Mendeley Data methylotrophic-protein deposit & not-a-dataset & A P. kudriavzevii deposit for open production of methylotrophic protein. Project policy disallows this repository as a source, so the deposit is recorded and not recommended. \\
\addlinespace[5pt]
PXD044970 and PXD048012 & not-a-dataset & The first lists a different species entirely, Saccharomyces kudriavzevii, and the second lists P. kudriavzevii while its displayed abstract belongs to an unrelated glycolysis study, so it may be mislabeled. Both are suspect and neither may be used without a check. \\
\addlinespace[5pt]
Komagataella phaffii methanol cross-tolerance study & off-host & A different species. Named because it surfaced repeatedly in searches for this wave's low-data host and is not it. \\
\addlinespace[5pt]
681-member environmental yeast isolate collection & off-host & Multi-species rather than specific to this wave's host, so its records cannot be assigned to one organism without splitting the collection first. \\
\addlinespace[5pt]
FUN-PROSE expression-prediction preprint & derivative & A modeling paper in the same group's orbit as the host's expression work, not a dataset: its records are the matrices it trains on. \\
\addlinespace[5pt]
Karollus 2024 benchmark deposit & derivative & Model weights plus an unitemized supporting-data bundle. Every measured dataset behind it belongs to the terminator, promoter and 3' end rows this pool already ranks, so the right action is to take those. \\
\addlinespace[5pt]
Martinez-Carranza 2019 essentiality comparison & derivative & A comparison of published essentiality screens, reporting 336 essential in one reference strain and 434 in the other with 20 having no ortholog, and only 192 genes shared by all screens in one strain against 17 in the other. Those figures are the reason a call must carry its method, medium and strain; the comparison itself re-serves the five studies' calls. \\
\addlinespace[5pt]
Rossi 2021 ChIP-exo and ChIP-MNase track set & not-a-dataset & 1,229 datasets over 791 targets, 400 of them reproducible, and not a genotype-by-environment record at all: it is a molecular-layer track set. It belongs in the artifact tier as a feature source if anywhere, which is a features decision rather than a dataset admission. \\
\addlinespace[5pt]
mLife 2026 genome-editing system & not-a-dataset & A genome-editing system in a nonconventional yeast whose organism is not confirmed to be this wave's host, so the row cannot even be assigned a host. One record check resolves it. \\
\addlinespace[5pt]
Lactic-acid stress comparative transcriptomics & no-sequence & A fermented-grain isolate against a type strain under lactic-acid stress, with differential expression in pyruvate metabolism, transporters and amino acid biosynthesis. Comparative rather than evolved, and neither strain's assembly was located, so there is no route to either genotype. \\
\addlinespace[5pt]
Sc2.0 synthetic-genome program; CNGB CNP0002230 & no-per-record-data & The 2023 Cell papers' phenotype claims appear as figures, and CNP0002230 lists 231 samples across genome sequencing, proteome and transcriptome with no phenotype data type at all. A genome-engineering program is not gene-indexed perturbation-phenotype data, and there is no per-record value to ingest today. \\
\addlinespace[5pt]
Genome-wide KT2440 CRISPRi library (78,932 members) & no-per-record-data & The highest-value missing dataset in the survey and effectively request-only: 10 to 15 gRNAs per gene with 798 non-targeting controls, screened on glucose, acetate and p-coumarate at two concentrations, existing only as a Department of Energy abstract with no publication, no repository and no accession. It cannot be ranked against rows whose numbers were read off a file, so it is an ask rather than a row. \\
\addlinespace[5pt]
Werner 2023 and the engineered-strain product literature & no-per-record-data & Supplementary-table-only across the board. In Werner 2023, the one paper whose data-availability statement was read verbatim, every number is in sciadv.adj0053\_data\_s1.zip and the SI PDF, there is no repository of any kind, and the strains are request-only under a material transfer agreement. The tier is worth ingesting for its named genotypes, which is why those rows are in the pool, and not for record count. \\
\addlinespace[5pt]
Fitness Browser fitprivate instance (104 organisms, 989 KT2440 experiments) & not-a-dataset & A non-public instance whose counts disagree with the live public page's 74 organisms and 314 KT2440 experiments. The live public page supersedes every other count, and fitprivate numbers must never be mixed in, so the instance is named here to keep its figures out of any total. \\
\addlinespace[5pt]
Fitness Browser KBase mirror (48 organisms) and the June 2017 release & derivative & Stale mirrors of the same artifact. The bulk page itself says the 2018 release's sucrose and D-mannitol data should be disregarded in favor of fresh stock data, so an older copy is not an independent dataset and ingesting one would serve values the source has retracted. \\
\addlinespace[5pt]
Wheeldon genome-scale CRISPR libraries (Y. lipolytica, K. phaffii) & off-host & Real genome-scale resources, including a 6-fold-coverage K. phaffii library of 30,848 active guides covering more than 99 percent of the genome, and in neither case the host this wave bootstraps. Searched specifically because a Pichia screen would be the single most valuable row available; the platform has never been extended to P. kudriavzevii. \\
\addlinespace[5pt]
ProteomeCentral P. putida project pages & not-a-dataset & JavaScript-rendered, so a crawl returns nothing. This is a retrieval route and not a dataset: the same projects are reachable through the PRIDE v3 REST API, which is how the proteomics rows in the pool are counted. \\
\addlinespace[5pt]
\end{longtable}
\endgroup


Table~\ref{tab:excluded} holds the rows that were considered and never ranked, each with the
rule that dropped it. Four rules do the work. The three evolved acid-tolerant lines fail the
sequence gate of Sec.~\ref{sec:gate}, because no resequencing of them was deposited and neither
search pass found a variant table. The iIsor850 essential-gene list is a reconstruction rather
than an experiment, and that exclusion stands on its own merits whatever happens to the host's
other rows. The synthetic-chromosome phenotype claims have no per-record values to ingest,
because their released form is a figure. And the coverage-track and record-format releases of
the two expression models are a model's own input format rather than a measurement, with the
one genuinely new measurement under them recorded as an ask instead.

Two standing exclusions in the yeast pool need a decision rather than a rule, and this wave
forces it. That pool excludes a \org{Komagataella phaffii} CRISPR fitness screen and a
multi-species yeast panel as off-species, under the standing note ``Revisit if the
generalization axis widens to other hosts.'' A second yeast host is that widening, so the
off-species note is now a judgment to make rather than an inherited default.
Table~\ref{tab:schema} records it as a decision item.

\section{Before a loader is written\secstatus{todo}}\label{sec:caveats}

Every figure in Table~\ref{tab:final} is an axis count rather than a verified record count, and
the label on each figure marks which it is. A figure read on the source page in session is
marked as read. A figure taken from a search snippet without the page being opened is marked as
such and is provisional. A figure that is the product of two reported axes, or an estimate of a
panel depth, is marked inferred. Every measurement count in this document for a vector-valued
row is inferred in that sense, because it multiplies a reported record count by a reported or
estimated dimensionality. A row whose position looks surprising needs its dimensionality
checked before its record count, because the product set the position.

\notesec{Six internal conflicts have to be reconciled, not averaged} Each of these is a
disagreement inside one source or between two pages of one project, and each would silently
change a stored value. The isolate panel of the new host reports 170 strains sequenced and 161
phenotyped in the preprint against 162 strains on a DOE page, 57 conditions against 58, and
305,405 SNPs on a poster against 305,435 in a DOE summary. The KT2440 fitness compendium says
179 unique conditions in its abstract and results and 183 in its methods. The promoter-swap
fitness library's gene count is reported as 81, 96, and about 100 in three places. The
transcription-start-site library's preprint says about 80,000 variants against 1,000,000 in the
published version, so the release has to be pinned. The longitudinal isolate panel gives 474,
373, 389, and 443 isolates on one page. And the served yeast record total itself is 65,889,109
in the prose and the per-dataset table against 65,889,107 in the organism table whose
three-organism total this document quotes.

\subsection{What is retrievable, what needs a browser, and what needs an ask\secstatus{todo}}\label{sec:retrieval}

Sorting the pool by how a row is obtained rather than by what it contains changes what the
campaign does first, and the three classes are unequal in size and in effort.

\notesec{Retrievable now, programmatically} The 10.4 GB fitness database and its
protein-sequence companion both return HTTP 200 to a plain \texttt{curl} with no challenge, as
do the 2018 release's bulk page and its archival copy. Two GitHub repositories hold the KT2440
fitness and transcriptome matrices. The GEO, SRA, and BioProject accessions of the ranked rows
resolve, as do the Zenodo records behind the promoter libraries, the clinical classifier
inputs, and the PAO1 interference screen. Two APIs complete the class: the proteomics archive's
REST interface, which is the only workable route to roughly 35 \org{P. putida} proteomics
projects because the aggregator's pages are rendered in JavaScript, and the
deep-mutational-scanning aggregate through its R and Python packages. One rule applies to the
largest of these. Scrape nothing from the fitness browser's CGI layer, pull the database file
once, and read it locally.

\notesec{Needs a human with a browser} Every Illinois Data Bank deposit for the new host, six
of them, returned 403 to automated fetch on both passes, which matters because that repository
is the de facto home of this organism's phenotype data and is invisible to GEO and SRA
searches. The preprint full text holding the 57-condition phenotype matrix is unreachable from
two directions, 403 at the preprint server and 500 at the literature API's full-text route, and
the multi-strain respiration panel was rate-limited on four attempts. Four pages behind a
challenge or a redirect are the fitness browser's help page and archival record, which carry
the license, the community genome database's downloads page, a phenotype-microarray database,
and the 1,877-strain susceptibility set. The rest of the class is the supplement-only rows: the
terminator panels, the core-promoter and designed-promoter libraries, the 3' end library's
sequence-and-value table, the position-effect tables, the two aneuploid panels, the two isolate
spreadsheets, the essentiality gene lists, and one supplementary strain table.

\notesec{Needs an ask to a named person or group} This is the shortest class and the one with
the largest effect per request, and Table~\ref{tab:asks} prints each ask with what it would
make retrievable. The nine cover the resequencing series with no publication attached, the 161
by 57 phenotype matrix and its isolate-to-accession map, the unprovenanced 69-condition
expression matrix, a deposit for the 3,053 induction tracks, the gene-level scores of the
78,932-member interference library, an archival deposit for the landing-pad package now behind
a lab-hosted URL that can rot, the license for the fitness browser data, which of three
deposits holds gene-level scores, and the meaning of three uncertainty columns in the served
promoter dataset, which is already an open owner question in that loader.

%% GENERATED FILE -- do not hand-edit.
%% SOURCE: experiments/database/scripts/build_hosts_candidate_datasets_table.py
\begin{landscape}
\begingroup
\footnotesize
\setlength{\tabcolsep}{4pt}
\renewcommand{\arraystretch}{1.15}
\begin{longtable}{@{}L{38mm} L{16mm} L{40mm} L{78mm} L{70mm}@{}}
\caption[]{The rows gated on a request to a named person or group, which is the shortest
list in this document and the highest leverage: each of these is one message, and three
of them are the only route to the dataset at all. A dash in \emph{Host} marks an ask
that spans hosts.}
\label{tab:asks}\\
\toprule
\textbf{Dataset} & \textbf{Host} & \textbf{Who to ask} & \textbf{What exactly to ask for} & \textbf{What it unlocks} \\
\midrule
\endfirsthead
\multicolumn{5}{@{}l}{\footnotesize\emph{Table~\ref{tab:asks}, continued}}\\
\toprule
\textbf{Dataset} & \textbf{Host} & \textbf{Who to ask} & \textbf{What exactly to ask for} & \textbf{What it unlocks} \\
\midrule
\endhead
\bottomrule
\endfoot
\textbf{IoSL resequencing series (PRJNA1464946/47/48)} & \org{P. kudriavzevii} & DOE Joint Genome Institute, submitter of PRJNA1464946, PRJNA1464947 and PRJNA1464948, registered 2026-05-11 & How many strains the IoSL series covers, what perturbation produced them, and whether any phenotype was measured alongside the sequencing & The only remaining candidate for a genome-wide P. kudriavzevii strain library. Three sampled BioProjects are titled Issatchenkia orientalis SD108 IoSL10, IoSL16 and IoSL19 resequencing, with no publication attached; a series sequenced at that scale is either a mutagenized or an evolved SD108 library, and if it carries phenotypes it is the resource this wave cannot otherwise find \\
\addlinespace[5pt]
\textbf{MAESTRO 170-isolate panel} & \org{P. kudriavzevii} & Ping-Hung Hsieh (LBNL), corresponding; also Yasuo Yoshikuni (LBNL/JGI), Huimin Zhao (UIUC/CABBI), Chris Todd Hittinger and Dana Opulente for the isolates, Igor Grigoriev (JGI) for the genomes, Andrew Leakey (CABBI) & The 161-isolate by 57-stressor phenotype matrix as a file, the complete isolate-to-accession map, and the per-isolate VCF and CNV tables & This host's largest phenotype resource by measurements, and the per-isolate variant calls behind the tandem-paralog CNV finding, none of which has a public accession. bioRxiv returns 403 for the full text and Europe PMC's fullTextXML for PPR1078985 returns 500, so the supplementary table is unreachable without a person \\
\addlinespace[5pt]
\textbf{I. orientalis 5,140-gene by 69-condition expression matrix} & \org{P. kudriavzevii} & Saurabh Sinha's group at UIUC, with Payam Dibaeinia & Whose data the matrix is, whether a release exists, and under what terms it could be deposited with a provenance record & A 69-condition expression matrix for this host. The senior thesis carrying it has no data-availability, code, deposition or supplementary section of any kind, so it is reachable by asking and by nothing else; with a real provenance record it would be worth more to this host's program than anything else in that corpus \\
\addlinespace[5pt]
\textbf{Shorkie IDEA induction RNA-seq tracks} & \org{S. cerevisiae} & Calico Life Sciences (Chao, Magzoub, Stoops, Hackett, Linder, Kelley) & A GEO or SRA deposit for the 3,053 induction RNA-seq tracks, or a stable archival copy with a checksum manifest & The one genuinely new measurement in that paper, currently distributed only as BigWig and TFRecord on a Google Cloud Storage bucket with no accession anywhere in the paper and no archive fallback if the bucket moves \\
\addlinespace[5pt]
\textbf{Genome-wide KT2440 CRISPRi library (78,932 members)} & \org{P. putida} & CU Boulder, ORNL and the Center for Bioenergy Innovation: Fenster, Hren, Spangler, Alexander, Fox, Cameron, Eckert, Tuskan & Gene-level or guide-level scores from the library, and whether a deposit exists anywhere, for the glucose, acetate and p-coumarate selections & The complement to RB-TnSeq, covering the essential genes a transposon library cannot score by construction. Abstract-only today: no publication, no repository, no accession, so no file exists to hash \\
\addlinespace[5pt]
\textbf{Landing-pad locus screen (Fatma 2023)} & \org{P. kudriavzevii} & Huimin Zhao lab (UIUC) & An archival deposit for the landing-pad package, currently a lab-hosted uofi.box.com URL plus a GitHub repository & The only measured per-locus fitness data in this species, and genotype-precise, behind a URL that can rot. An institutional Box link is not an archive and nothing mirrors it \\
\addlinespace[5pt]
\textbf{Fitness Browser licensing} & -- & LBL Fitness Browser maintainers & The license and terms of use for the live feba.db release and for the June 2017 archival release & The whole cross-host fitness band, nine organisms, gated by a Cloudflare challenge rather than by anything substantive: help.cgi and the organism list return 403 with Cf-Mitigated: challenge to every non-browser client and the Figshare archival record page returns 403 as well, so nobody has read a license string for either release \\
\addlinespace[5pt]
\textbf{PAO1 genome-wide CRISPRi (Kaczmarczyk 2025)} & \org{P. aeruginosa} & Jenal lab (Kaczmarczyk, Klotz, Manfredi, Jenal) & Which of the three Zenodo records holds the gene-level score table, if opening them does not settle it & The only pseudomonad CRISPRi score table plausibly downloadable today, covering the essential genes RB-TnSeq structurally cannot \\
\addlinespace[5pt]
\textbf{Vaishnav 2022 deferred libraries} & \org{S. cerevisiae} & Vaishnav 2022 authors & The meaning of sdEL, relativeSD and sdScale & Already an open owner question in vaishnav2022.py and listed under issue 1172. de Boer's Supplementary Tables may answer it without an ask, since Vaishnav's Methods defer to de Boer, so one page check comes first \\
\addlinespace[5pt]
\textbf{Dotsch 2009 arrayed PA14 screen} & \org{P. aeruginosa} & The study's authors, or whoever now holds the arrayed PA14 library data & Whether the full 5,459-mutant by 19-antibiotic susceptibility matrix was released, and where & If the full matrix exists it is the densest arrayed-library phenotype matrix in the species. This is the highest-value unresolved pseudomonad item and the value location is unconfirmed, so the row cannot be costed until someone answers \\
\addlinespace[5pt]
\end{longtable}
\endgroup
\end{landscape}


\subsection{The candidate gate, and how much of this pool it has measured\secstatus{todo}}\label{sec:ungated}

The gate runs five gates in order with a stop rule. G1 classifies the source as a primary
measurement, an aggregation, a transcription, or a prediction. G2 resolves every host through
the genomes tier. G3 inventories the deposit off the mirrored bytes. G4 tests overlap against
built loaders by identifier and by released value, and G5 tests schema fit. After a failure or
a block every later gate is unmeasured, a typed gap does not stop the run, and the verdict's
outcome is derived from its gates rather than stored, so a stored verdict cannot disagree with
its own gates.

Measured against this pool, the gate has run on none of it. The verdict store holds 52
verdicts, every one with a table field of bacteria or yeast, and no verdict in it mentions this
document's new host under any of its four names or its strain. That is not an oversight: the
gate's organism set is three values, so \texttt{hosts\_named()} raises a \texttt{ValueError} on
any row of the new host, and G2 cannot be measured rather than returning a block. Every
non-KT2440 pseudomonad row is in the same position for the same reason. So the ungated fraction
of this pool is all of it, and the first two rows of Table~\ref{tab:schema} are the two edits
that change it, so the table lists them as gating every row rather than one.

Two further properties of the store bear on the build order. Enforcement requires every
registered dataset class outside a frozen grandfathered tuple of 134 classes to declare a
module-level citation key and to hold a passing verdict, so a row that is built without a
verdict fails that check rather than entering quietly. And a verdict is evidence rather than a
working document: every record model is frozen and forbids extra fields, one file is written
per citation key, and git keeps the history, so a re-gated row produces a new commit rather
than an edited file.

%% notes-tex/database/database-expansion-hosts/sections/summary.tex
%% [[experiments.database.expansion-hosts]]
%% https://github.com/Mjvolk3/torchcell/tree/main/notes-tex/database/database-expansion-hosts/sections/summary.tex
%% SOURCE: hand-authored prose around a generated table; tables/summary.tex names its own
%% generating script. The time-axis figures come from the triage report
%% /scratch/projects/torchcell-scratch/tmp/deep-research/reports/Pk and next dataset wave triage.md
%% and its research notes; the schema fields from torchcell/datamodels/schema.py.

\clearpage
\section{Summary statistics\secstatus{todo}}\label{sec:summary}

%% GENERATED FILE -- do not hand-edit.
%% SOURCE: experiments/database/scripts/build_hosts_candidate_datasets_table.py
\begingroup
\small
\begin{longtable}{@{}l r r r r r@{}}
\caption[]{Summary of the recommended set, split at the fifty line. \emph{Top 50} and
\emph{Extra} count rows; \emph{Genotypes}, \emph{Instances} and \emph{Meas.} sum the row
axes over all 60 and are lower bounds, since a row with no count contributes
nothing. The set names 22 distinct schema classes the substrate does not hold
(Table~\ref{tab:schema}).}
\label{tab:summary}\\
\toprule
 & Top 50 & Extra & Genotypes & Instances & Meas. \\
\endfirsthead
\multicolumn{6}{@{}l}{\footnotesize\emph{Table~\ref{tab:summary}, continued}}\\
\toprule
 & Top 50 & Extra & Genotypes & Instances & Meas. \\
\endhead
\bottomrule
\endfoot
\midrule
\multicolumn{6}{@{}l}{\emph{By band}}\\
regulatory DNA & 34 & 1 & 166,797,405 & $1.7\times 10^{8}$ & $2.1\times 10^{8}$ \\
cross-host fitness & 0 & 9 & 0 & $1.4\times 10^{3}$ & $6.2\times 10^{6}$ \\
host bootstrap & 8 & 0 & 260 & $9.6\times 10^{3}$ & $3.5\times 10^{5}$ \\
scale & 8 & 0 & 5,467 & $5.8\times 10^{3}$ & $4.4\times 10^{6}$ \\
\midrule
\multicolumn{6}{@{}l}{\emph{By tier}}\\
tier 1 & 15 & 5 & 166,022,613 & $1.7\times 10^{8}$ & $1.9\times 10^{8}$ \\
tier 2 & 12 & 2 & 757,235 & $1.2\times 10^{6}$ & $3.1\times 10^{7}$ \\
tier 3 & 13 & 1 & 21,139 & $3.0\times 10^{4}$ & $7.7\times 10^{5}$ \\
tier 4 & 10 & 2 & 2,145 & $2.2\times 10^{3}$ & $1.5\times 10^{5}$ \\
\midrule
\multicolumn{6}{@{}l}{\emph{By host}}\\
E. coli & 9 & 0 & 2,861,834 & $5.2\times 10^{6}$ & $4.5\times 10^{7}$ \\
P. aeruginosa & 0 & 1 & 0 & 8 & $3.6\times 10^{4}$ \\
P. fluorescens & 0 & 4 & 0 & 745 & $3.4\times 10^{6}$ \\
P. kudriavzevii & 8 & 0 & 260 & $9.6\times 10^{3}$ & $3.5\times 10^{5}$ \\
P. putida & 8 & 0 & 5,467 & $5.8\times 10^{3}$ & $4.4\times 10^{6}$ \\
P. simiae & 0 & 1 & 0 & 201 & $9.1\times 10^{5}$ \\
P. syringae & 0 & 2 & 0 & 69 & $3.1\times 10^{5}$ \\
S. cerevisiae & 25 & 1 & 163,935,571 & $1.6\times 10^{8}$ & $1.7\times 10^{8}$ \\
S. stutzeri & 0 & 1 & 0 & 350 & $1.6\times 10^{6}$ \\
\midrule
\multicolumn{6}{@{}l}{\emph{By class}}\\
CRISPR library screen & 1 & 0 & 47 & 188 & 188 \\
Deep mutational scan & 3 & 0 & 4,176 & $4.2\times 10^{3}$ & $4.2\times 10^{3}$ \\
Engineering parts & 3 & 0 & 3,100 & $3.1\times 10^{3}$ & $3.1\times 10^{3}$ \\
Fitness / chemical genomics & 0 & 9 & 0 & $1.4\times 10^{3}$ & $6.2\times 10^{6}$ \\
Natural variation & 4 & 0 & 900 & $1.5\times 10^{4}$ & $1.5\times 10^{4}$ \\
Proteome & 1 & 0 & 0 & 7 & $3.9\times 10^{4}$ \\
Regulatory DNA & 27 & 1 & 166,789,446 & $1.7\times 10^{8}$ & $2.1\times 10^{8}$ \\
Transcriptome & 7 & 0 & 20 & 472 & $2.6\times 10^{6}$ \\
Transposon fitness & 4 & 0 & 5,443 & 430 & $2.1\times 10^{6}$ \\
\midrule
\multicolumn{6}{@{}l}{\emph{By sequence basis}}\\
KT2440+transposon & 4 & 0 & 5,443 & 430 & $2.1\times 10^{6}$ \\
KT2440-reference & 1 & 0 & 0 & 321 & $1.8\times 10^{6}$ \\
MG1655+reporter-locus & 2 & 0 & 329,392 & $8.3\times 10^{5}$ & $8.6\times 10^{5}$ \\
Pk-isolate-WGS & 5 & 0 & 210 & $9.4\times 10^{3}$ & $2.9\times 10^{5}$ \\
S288C+designed-edit & 5 & 0 & 503,524 & $5.1\times 10^{5}$ & $5.1\times 10^{5}$ \\
S288C+integration-locus & 4 & 0 & 2,648 & $2.6\times 10^{3}$ & $1.5\times 10^{6}$ \\
SD108+guide & 1 & 0 & 47 & 188 & 188 \\
engineered-chassis & 3 & 0 & 3 & 93 & $5.1\times 10^{5}$ \\
engineered-chassis+guide & 1 & 0 & 0 & 7 & $3.9\times 10^{4}$ \\
fitbrowser+transposon & 0 & 9 & 0 & $1.4\times 10^{3}$ & $6.2\times 10^{6}$ \\
isolate-WGS & 1 & 0 & 683 & 683 & 683 \\
plasmid-reporter-library & 22 & 1 & 165,961,158 & $1.7\times 10^{8}$ & $2.1\times 10^{8}$ \\
pseudomonad-isolate-WGS & 1 & 0 & 24 & $4.9\times 10^{3}$ & $4.9\times 10^{3}$ \\
\midrule
\multicolumn{6}{@{}l}{\emph{By retrieval class}}\\
retrievable programmatically & 30 & 9 & 166,762,265 & $1.7\times 10^{8}$ & $2.2\times 10^{8}$ \\
needs a human with a browser & 17 & 1 & 40,635 & $4.9\times 10^{4}$ & $1.8\times 10^{6}$ \\
needs an ask to a named party & 3 & 0 & 232 & $1.4\times 10^{4}$ & $1.4\times 10^{4}$ \\
\midrule
\multicolumn{6}{@{}l}{\emph{By verification label}}\\
counts read on the source & 34 & 0 & 166,787,537 & $1.7\times 10^{8}$ & $2.1\times 10^{8}$ \\
counts from a search snippet & 7 & 1 & 4,421 & $1.3\times 10^{4}$ & $1.6\times 10^{5}$ \\
counts inferred, including a product & 9 & 9 & 11,174 & $3.2\times 10^{4}$ & $1.0\times 10^{7}$ \\
\midrule
\multicolumn{6}{@{}l}{\emph{Other attributes}}\\
needs a schema class the substrate lacks & 24 & 10 & 2,536,941 & $2.6\times 10^{6}$ & $1.1\times 10^{7}$ \\
carries a time axis & 1 & 0 & 1,214,438 & $1.2\times 10^{6}$ & $7.3\times 10^{6}$ \\
figures sourced this pass & 43 & 9 & 166,798,711 & $1.7\times 10^{8}$ & $2.2\times 10^{8}$ \\
figures from recall, to confirm & 7 & 1 & 4,421 & $1.3\times 10^{4}$ & $1.6\times 10^{5}$ \\
instances a reported count & 38 & 9 & 166,759,707 & $1.7\times 10^{8}$ & $2.2\times 10^{8}$ \\
instances a product of the axes & 8 & 0 & 40,425 & $1.6\times 10^{5}$ & $4.0\times 10^{5}$ \\
instances an estimate & 4 & 1 & 3,000 & $3.0\times 10^{3}$ & $3.0\times 10^{3}$ \\
\midrule
Total & 50 & 10 & 166,803,132 & $1.7\times 10^{8}$ & $2.2\times 10^{8}$ \\
\end{longtable}
\endgroup


Table~\ref{tab:summary} gives the pool by host, by band, and by tier, with the record,
genotype, and measurement sums for each. Three properties of those sums are worth stating
before they are read as a ranking.

A handful of rows account for the measurement sums, so the hosts are not comparable on them.
Three \org{S. cerevisiae} rows hold more measurements than every other row in the pool
combined, which makes the yeast total a statement about one platform rather than about the
host, and one artifact accounts for the pseudomonad total, its rows sharing a pipeline and a
condition vocabulary. The new host's total is four orders of magnitude below either, which is
the measured condition of the host rather than a property of the ranking.

Record counts and measurement counts disagree in direction, which is the point of the rank
quantity in Sec.~\ref{sec:measurements}. The isolate panels hold hundreds of records each and
millions of measurements, because one record is a genome plus a transcriptome or a
genome plus a proteome, while the arrayed and pooled screens hold the opposite shape. Summing
either column alone orders the hosts the wrong way round.

And the sums are upper bounds in two specific ways that Sec.~\ref{sec:caveats} records. A
count that is the product of two reported axes assumes a dense matrix, and a real screen drops
strains per condition, so every such figure is an upper bound. A measurement count for a
vector-valued row multiplies by a reported or estimated depth, and a label-free proteome or a
sparse count matrix quantifies a different subset in each sample, so the dense assumption fails
there too.

\subsection{Time\secstatus{todo}}\label{sec:time}

Five rows in this pool carry a time dimension, and they differ in kind as well as in
length. The low-pH expression atlas samples at 0 h and 1 h after the acid shift, so its series
is two points per strain. The 3,053 induction tracks of the fungal language model's corpus
typically carry eight time points per induction, which is the same design as the served
induction series. The longitudinal isolate panel is the longest axis in the pool at 0.2 to 10.2
years of colonization across 52 clone types, where the ordering is part of the phenotype rather
than a label on it. The isoprenol design-build-test-learn ladder is indexed per round across
seven deposits, which is a campaign axis and not a time axis, and it appears here only so that
the two are not conflated. And one allele-specific expression study on the diauxic shift carries 19
time points, the densest series found.

What the schema holds for all of them is \texttt{Environment.duration\_hours} and
\texttt{Environment.duration\_generations}, treatment durations that make two exposure lengths
two distinct environments, and the served induction series was ingested that way as 1,752
records with one environment per time point. That is a flattening: a model sees time as a
categorical condition, and nothing on the record states that eight environments are one series
in order. A first-class time field on the record and a series identity that groups its points
are the two additions a time-resolved model needs, and the five rows above would exercise both.
They are recorded here as an open schema item rather than folded into the
ranking, because the rule of Sec.~\ref{sec:rule} does not score a time axis.

%% notes-tex/database/database-expansion-hosts/sections/full.tex
%% [[experiments.database.expansion-hosts]]
%% https://github.com/Mjvolk3/torchcell/tree/main/notes-tex/database/database-expansion-hosts/sections/full.tex
%% SOURCE: generated tables only; each table file names its own generating script in its header.

\clearpage
\section{The full pool\secstatus{todo}}\label{sec:full}

Table~\ref{tab:candidates} is the whole pool the recommended set was cut from, in rank order,
with each row's host, class, both axes, its record and measurement counts, its phenotype, its
sequence basis, and a clickable link. Table~\ref{tab:sources-all} gives every row's citation, link, the
reason for its position, and where the per-record values live, in the same order, with each
unconfirmed location marked as such. Table~\ref{tab:verdicts} gives the candidate-gate verdict
of every ranked row from the verdict store, and prints no verdict for a row the gate has not
run on, which Sec.~\ref{sec:ungated} counts.

%% GENERATED FILE -- do not hand-edit.
%% SOURCE: experiments/database/scripts/build_hosts_candidate_datasets_table.py
\begin{landscape}
\begingroup
\footnotesize
\setlength{\tabcolsep}{4pt}
\renewcommand{\arraystretch}{1.15}
\begin{longtable}{@{}r@{\hspace{4pt}} L{40mm} L{13mm} L{21mm} L{21mm} L{21mm} r@{\hspace{5pt}} r@{\hspace{5pt}} c@{\hspace{5pt}} L{24mm} L{21mm} L{28mm}@{}}
\caption[]{Every candidate in the cross-host pool, in rank order: band first, then tier,
then measurements descending. \emph{Genotypes} and \emph{Env} are the perturbation and
condition axes. \emph{Instances} is the number of genotype$\times$environment records,
marked $\dagger$ where it is the product of the two axes rather than a figure the source
reports, and $\ddagger$ where it is an order-of-magnitude estimate. \emph{Meas.} is
instances times phenotype dimensionality, which is what a vector-valued panel actually
contributes and what rows are ranked on. \emph{Sequence basis} is how the total genomic
content of one strain would be reconstructed; a row with no such route is excluded rather
than ranked (Table~\ref{tab:excluded}). Superscript \textbf{M} marks a row whose values need a human with a browser and \textbf{R} a row that needs a request to a named party (Table~\ref{tab:asks}); \textbf{B} marks a row whose per-record values are not released and \textbf{A} a corpus that re-serves other papers. \emph{Ver.} carries the triage report's own verification label for the counts: \textbf{V} read on the source in session, \textbf{S} a search snippet only, \textbf{I} inferred, which includes every count that is the product of two verified numbers. \emph{Link} resolves to the source and is clickable; the full citation, the
reason for the position and the data location are in Table~\ref{tab:sources-all}. The
RECOMMENDED set is the fifty of Table~\ref{tab:final} and its ten spares, not this table:
rows 1--200 here are the rows above the cut, meaning the additions this wave targets.}
\label{tab:candidates}\\
\toprule
\textbf{\#} & \textbf{Dataset} & \textbf{Host} & \textbf{Class} & \textbf{Genotypes} & \textbf{Env} & \textbf{Instances} & \textbf{Meas.} & \textbf{Ver.} & \textbf{Phenotype} & \textbf{Sequence basis} & \textbf{Link} \\
\midrule
\endfirsthead
\multicolumn{12}{@{}l}{\footnotesize\emph{Table~\ref{tab:candidates}, continued}}\\
\toprule
\textbf{\#} & \textbf{Dataset} & \textbf{Host} & \textbf{Class} & \textbf{Genotypes} & \textbf{Env} & \textbf{Instances} & \textbf{Meas.} & \textbf{Ver.} & \textbf{Phenotype} & \textbf{Sequence basis} & \textbf{Link} \\
\midrule
\endhead
\bottomrule
\endfoot
1 & \textbf{Douglass 2018 (32 sequenced isolates, 4-antifungal MIC)} & \org{P. kudriavzevii} & Natural variation & 32 sequenced isolates, 20 clinical and 12 environmental & fluconazole, flucytosine, amphotericin B, micafungin & 128$\dagger$ & 128 & I & minimum inhibitory concentration, mg/L & Pk-\allowbreak isolate-\allowbreak WGS & \href{https://www.ncbi.nlm.nih.gov/pmc/articles/PMC6053246/}{PMC60532\allowbreak 46} \\
\addlinespace[5pt]
2 & \textbf{de Boer 2020 (GPRA sort-seq, 100M random promoters)} & \org{S. cerevisiae} & Regulatory DNA & 102,371,025 promoter sequences & YPD, YPGal, YPGly & $1.0\times 10^{8}$ & $1.0\times 10^{8}$ & V & expression driven, 18-\allowbreak bin sort-\allowbreak seq YFP/RFP & plasmid-\allowbreak reporter-\allowbreak library & \href{https://doi.org/10.1038/s41587-019-0315-8}{doi:\allowbreak 10.\allowbreak 1038/\allowbreak s41587-\allowbreak 019-\allowbreak 0315-\allowbreak 8} \\
\addlinespace[5pt]
3 & \textbf{Vaishnav 2022 deferred training libraries} & \org{S. cerevisiae} & Regulatory DNA & 30,722,376 complex + 20,616,659 defined & YPD complex, SD-\allowbreak Ura defined & $5.1\times 10^{7}$ & $5.1\times 10^{7}$ & V & expression driven, 18-\allowbreak bin sort-\allowbreak seq YFP & plasmid-\allowbreak reporter-\allowbreak library & \href{https://doi.org/10.1038/s41586-022-04506-6}{doi:\allowbreak 10.\allowbreak 1038/\allowbreak s41586-\allowbreak 022-\allowbreak 04506-\allowbreak 6} \\
\addlinespace[5pt]
4 & \textbf{Hollerer 2020 (ribosome binding site library, uASPIre)} & \org{E. coli} & Regulatory DNA & 303,503 exactly specified ribosome-\allowbreak binding-\allowbreak site variants on a plasmid in TOP10 delta-\allowbreak rhaA & one condition, three independent replicates & $8.4\times 10^{5}$ & $7.6\times 10^{6}$ & V & recombinase-\allowbreak recorded reporter output & plasmid-\allowbreak reporter-\allowbreak library & \href{https://doi.org/10.1038/s41467-020-17222-4}{doi:\allowbreak 10.\allowbreak 1038/\allowbreak s41467-\allowbreak 020-\allowbreak 17222-\allowbreak 4} \\
\addlinespace[5pt]
5 & \textbf{Hollerer 2023 (5' UTR and CDS library, uASPIre)} & \org{E. coli} & Regulatory DNA & 1,214,438 exactly specified 5' UTR and coding-\allowbreak sequence variants on a plasmid in TOP10 delta-\allowbreak rhaA & one condition, read out at 6 timepoints & $1.2\times 10^{6}$ & $7.3\times 10^{6}$ & V & relative translation rate & plasmid-\allowbreak reporter-\allowbreak library & \href{https://doi.org/10.1093/nar/gkad040}{doi:\allowbreak 10.\allowbreak 1093/\allowbreak nar/\allowbreak gkad040} \\
\addlinespace[5pt]
6 & \textbf{DREAM 2022 random promoter challenge (Rafi 2024)} & \org{S. cerevisiae} & Regulatory DNA & 6,739,258 train + 71,103 test & 1 medium & $6.8\times 10^{6}$ & $6.8\times 10^{6}$ & V & expression driven, FACS, about 100 cells per sequence & plasmid-\allowbreak reporter-\allowbreak library & \href{https://www.ncbi.nlm.nih.gov/geo/query/acc.cgi?acc=GSE254493}{ncbi.\allowbreak nlm.\allowbreak nih.\allowbreak gov/\allowbreak geo/\allowbreak query/\allowbreak acc.\allowbreak cgi?\allowbreak acc=\allowbreak GSE2544\allowbreak 93} \\
\addlinespace[5pt]
7 & \textbf{Kuo 2020 (Shine-Dalgarno 9-mer library in four contexts)} & \org{E. coli} & Regulatory DNA & the complete 262,144-\allowbreak member space of 9-\allowbreak nucleotide Shine-\allowbreak Dalgarno sequences, each measured in 4 ribosome-\allowbreak binding-\allowbreak site contexts & 4 ribosome-\allowbreak binding-\allowbreak site contexts & $1.0\times 10^{6}$ & $1.0\times 10^{6}$ & V & log reporter output & plasmid-\allowbreak reporter-\allowbreak library & \href{https://doi.org/10.1101/gr.260182.119}{doi:\allowbreak 10.\allowbreak 1101/\allowbreak gr.\allowbreak 260182.\allowbreak 119} \\
\addlinespace[5pt]
8 & \textbf{Kotopka and Smolke 2020 (artificial promoter libraries)} & \org{S. cerevisiae} & Regulatory DNA & over 675,000 constitutive + over 327,000 inducible & constitutive and inducible arms & $1.0\times 10^{6}$ & $1.0\times 10^{6}$ & V & promoter activity & plasmid-\allowbreak reporter-\allowbreak library & \href{https://doi.org/10.1038/s41467-020-15977-4}{doi:\allowbreak 10.\allowbreak 1038/\allowbreak s41467-\allowbreak 020-\allowbreak 15977-\allowbreak 4} \\
\addlinespace[5pt]
9 & \textbf{Zhu 2024 (Pol II MASTER, TSS initiation efficiency)} & \org{S. cerevisiae} & Regulatory DNA & 1,000,000 TSS sequences & defined promoter context, NTP levels varied & $1.0\times 10^{6}$ & $1.0\times 10^{6}$ & V & transcription initiation efficiency per TSS & plasmid-\allowbreak reporter-\allowbreak library & \href{https://doi.org/10.1038/s41594-023-01171-9}{doi:\allowbreak 10.\allowbreak 1038/\allowbreak s41594-\allowbreak 023-\allowbreak 01171-\allowbreak 9} \\
\addlinespace[5pt]
10 & \textbf{Savinov 2021 (random 50-mer 3' UTR library)} & \org{S. cerevisiae} & Regulatory DNA & 867,386 scored random 50-\allowbreak nucleotide 3' UTR elements across two libraries, from 4.6 million constructed & two library designs under one selection each & $8.7\times 10^{5}$ & $8.7\times 10^{5}$ & V & log2 enrichment, read out as growth through a reporter & plasmid-\allowbreak reporter-\allowbreak library & \href{https://doi.org/10.1186/s13059-021-02509-6}{doi:\allowbreak 10.\allowbreak 1186/\allowbreak s13059-\allowbreak 021-\allowbreak 02509-\allowbreak 6} \\
\addlinespace[5pt]
11 & \textbf{Urtecho 2020 (genome-wide promoter MPRA)} & \org{E. coli} & Regulatory DNA & 321,123 genomic fragments tiling the MG1655 genome, plus 17,635 transcription-\allowbreak start-\allowbreak site constructs across 3 landing pads & 2 media & $7.9\times 10^{5}$ & $7.9\times 10^{5}$ & V & RNA over DNA promoter activity & MG1655+\allowbreak reporter-\allowbreak locus & \href{https://doi.org/10.7554/eLife.92558.1}{doi:\allowbreak 10.\allowbreak 7554/\allowbreak eLife.\allowbreak 92558.\allowbreak 1} \\
\addlinespace[5pt]
12 & \textbf{Cuperus 2017 (random 5' UTR library)} & \org{S. cerevisiae} & Regulatory DNA & 489,348 detected random 50-\allowbreak nt 5' UTRs & medium lacking histidine & $4.9\times 10^{5}$ & $4.9\times 10^{5}$ & V & competitive growth enrichment, HIS3 & S288C+\allowbreak designed-\allowbreak edit & \href{https://genome.cshlp.org/content/27/12/2015}{genome.\allowbreak cshlp.\allowbreak org/\allowbreak content\allowbreak /\allowbreak 27/\allowbreak 12/\allowbreak 2015} \\
\addlinespace[5pt]
13 & \textbf{Johns 2018 (29,249 expression constructs across conditions)} & \org{E. coli} & Regulatory DNA & 29,249 regulatory constructs mined from metagenomes & 4 conditions & $1.2\times 10^{5}$$\dagger$ & $1.2\times 10^{5}$ & V & expression, with a transcription-\allowbreak start-\allowbreak site readout beside it & plasmid-\allowbreak reporter-\allowbreak library & \href{https://doi.org/10.1038/nmeth.4633}{doi:\allowbreak 10.\allowbreak 1038/\allowbreak nmeth.\allowbreak 4633} \\
\addlinespace[5pt]
14 & \textbf{Yu 2021 (promoter library across induction states)} & \org{E. coli} & Regulatory DNA & 8,269 promoters at a chromosomal landing pad between essQ and cspB in an MG1655 derivative & with and without induction & $3.4\times 10^{4}$ & $6.8\times 10^{4}$ & V & promoter activity and its fold change on induction & MG1655+\allowbreak reporter-\allowbreak locus & \href{https://doi.org/10.1038/s41467-020-20094-3}{doi:\allowbreak 10.\allowbreak 1038/\allowbreak s41467-\allowbreak 020-\allowbreak 20094-\allowbreak 3} \\
\addlinespace[5pt]
15 & \textbf{Cambray 2018 (244,000 designed sequences, multi-phenotype)} & \org{E. coli} & Regulatory DNA & exactly 244,000 designed sequences, matching the paper's own headline & 4 growth conditions for the growth readout & $2.4\times 10^{5}$ & $2.5\times 10^{7}$ & V & protein level, polysome profile, RNA over DNA, RNA half-\allowbreak life and growth & plasmid-\allowbreak reporter-\allowbreak library & \href{https://doi.org/10.1038/nbt.4238}{doi:\allowbreak 10.\allowbreak 1038/\allowbreak nbt.\allowbreak 4238} \\
\addlinespace[5pt]
16 & \textbf{Fuqua 2024 (promoter-island mutational landscape)} & \org{E. coli} & Regulatory DNA & 245,639 mutated daughter sequences derived from 25 promoter-\allowbreak island templates & one condition, two reporter channels & $2.5\times 10^{5}$ & $2.0\times 10^{6}$ & V & sort-\allowbreak seq expression in two channels & plasmid-\allowbreak reporter-\allowbreak library & \href{https://doi.org/10.7554/eLife.98654}{doi:\allowbreak 10.\allowbreak 7554/\allowbreak eLife.\allowbreak 98654} \\
\addlinespace[5pt]
17 & \textbf{Gui 2021 (position effects with transcriptomes)}\,\textsuperscript{\textbf{M}} & \org{S. cerevisiae} & Regulatory DNA & about 250 GFP integration loci & 1 & 250 & $1.5\times 10^{6}$ & I & full transcriptome per integration locus & S288C+\allowbreak integration-\allowbreak locus & \href{https://www.biorxiv.org/content/10.1101/2020.04.02.021162}{biorxiv.\allowbreak org/\allowbreak content\allowbreak /\allowbreak 10.\allowbreak 1101/\allowbreak 2020.\allowbreak 04.\allowbreak 02.\allowbreak 021162} \\
\addlinespace[5pt]
18 & \textbf{Zhou 2025 (massively parallel bacterial promoter library)} & \org{E. coli} & Regulatory DNA & 233,469 promoters in the released file, against 112,955 headlined in the text & one condition, three replicates & $7.0\times 10^{5}$ & $7.0\times 10^{5}$ & V & promoter activity & plasmid-\allowbreak reporter-\allowbreak library & \href{https://doi.org/10.1093/nar/gkaf863}{doi:\allowbreak 10.\allowbreak 1093/\allowbreak nar/\allowbreak gkaf863} \\
\addlinespace[5pt]
19 & \textbf{Keren 2016 (promoter panel x named gene, pooled fitness)} & \org{S. cerevisiae} & Regulatory DNA & 130 promoters over a 500-\allowbreak fold range x about 81 genes & glucose, galactose & $2.0\times 10^{4}$$\dagger$ & $2.0\times 10^{4}$ & I & pooled growth fitness per strain & S288C+\allowbreak designed-\allowbreak edit & \href{https://www.ncbi.nlm.nih.gov/geo/query/acc.cgi?acc=GSE83936}{ncbi.\allowbreak nlm.\allowbreak nih.\allowbreak gov/\allowbreak geo/\allowbreak query/\allowbreak acc.\allowbreak cgi?\allowbreak acc=\allowbreak GSE8393\allowbreak 6} \\
\addlinespace[5pt]
20 & \textbf{Shalem 2015 (3' end FACS-seq library)}\,\textsuperscript{\textbf{M}} & \org{S. cerevisiae} & Regulatory DNA & 13,000 designed 3' ends, 102 bp variable & SC-\allowbreak URA + galactose, 30 C & $1.3\times 10^{4}$ & $1.3\times 10^{4}$ & V & YFP/mCherry ratio, 16-\allowbreak bin FACS-\allowbreak seq & plasmid-\allowbreak reporter-\allowbreak library & \href{https://doi.org/10.1371/journal.pgen.1005147}{doi:\allowbreak 10.\allowbreak 1371/\allowbreak journal\allowbreak .\allowbreak pgen.\allowbreak 1005147} \\
\addlinespace[5pt]
21 & \textbf{Renganaath 2020 (native promoter variant MPRA)} & \org{S. cerevisiae} & Regulatory DNA & 5,832 variants in 2,503 gene promoters & 1 & $5.8\times 10^{3}$ & $5.8\times 10^{3}$ & V & MPRA mRNA/DNA allelic cis effect & plasmid-\allowbreak reporter-\allowbreak library & \href{https://doi.org/10.7554/eLife.62669}{doi:\allowbreak 10.\allowbreak 7554/\allowbreak eLife.\allowbreak 62669} \\
\addlinespace[5pt]
22 & \textbf{Arg-CCU tRNA deep mutational scan}\,\textsuperscript{\textbf{M}} & \org{S. cerevisiae} & Deep mutational scan & 4,176 of 5,184 genotypes after filtering & high temperature, 1 M NaCl & $4.2\times 10^{3}$ & $4.2\times 10^{3}$ & V & fitness under selection & S288C+\allowbreak designed-\allowbreak edit & \href{https://pmc.ncbi.nlm.nih.gov/articles/PMC6193533}{PMC61935\allowbreak 33} \\
\addlinespace[5pt]
23 & \textbf{Keren 2013 (native promoter panel across conditions)}\,\textsuperscript{\textbf{M}} & \org{S. cerevisiae} & Regulatory DNA & about 900 native promoters & 10 conditions & $9.0\times 10^{3}$$\dagger$ & $9.0\times 10^{3}$ & S & promoter activity, YFP with a fixed ADH1 terminator & plasmid-\allowbreak reporter-\allowbreak library & \href{https://spell.yeastgenome.org/}{spell.\allowbreak yeastge\allowbreak nome.\allowbreak org} \\
\addlinespace[5pt]
24 & \textbf{Lubliner 2015 (core promoter variant panel)}\,\textsuperscript{\textbf{M}} & \org{S. cerevisiae} & Regulatory DNA & 7,536 core promoter variants, several hundred native & 1 & $7.5\times 10^{3}$ & $7.5\times 10^{3}$ & V & YFP/mCherry ratio, 16-\allowbreak bin FACS & plasmid-\allowbreak reporter-\allowbreak library & \href{https://genome.cshlp.org/content/25/7/1008}{genome.\allowbreak cshlp.\allowbreak org/\allowbreak content\allowbreak /\allowbreak 25/\allowbreak 7/\allowbreak 1008} \\
\addlinespace[5pt]
25 & \textbf{Sharon 2012 (designed promoter library)}\,\textsuperscript{\textbf{M}} & \org{S. cerevisiae} & Regulatory DNA & 6,500 designed promoters, 55 percent Gal4/Gcn4 sites & GAL1-\allowbreak 10 and HIS3 reporter backgrounds & $6.5\times 10^{3}$ & $6.5\times 10^{3}$ & V & expression, sorting plus barcode sequencing & plasmid-\allowbreak reporter-\allowbreak library & \href{https://doi.org/10.1038/nbt.2205}{doi:\allowbreak 10.\allowbreak 1038/\allowbreak nbt.\allowbreak 2205} \\
\addlinespace[5pt]
26 & \textbf{Yamanishi 2013 (genomic terminatome)}\,\textsuperscript{\textbf{M}} & \org{S. cerevisiae} & Engineering parts & over 3,000 genomic terminators & fixed TDH3 promoter, GFP & $3.0\times 10^{3}$$\ddagger$ & $3.0\times 10^{3}$ & S & GFP output under a fixed promoter & plasmid-\allowbreak reporter-\allowbreak library & \href{https://www.yeastgenome.org/reference/S000206762}{yeastgen\allowbreak ome.\allowbreak org/\allowbreak referen\allowbreak ce/\allowbreak S000206\allowbreak 762} \\
\addlinespace[5pt]
27 & \textbf{Dvir 2013 (5' UTR positions -10 to -1)}\,\textsuperscript{\textbf{M}} & \org{S. cerevisiae} & Regulatory DNA & 2,041 5' UTR variants at positions -\allowbreak 10 to -\allowbreak 1 & 1 & $2.0\times 10^{3}$ & $2.0\times 10^{3}$ & V & reporter fluorescence & plasmid-\allowbreak reporter-\allowbreak library & \href{https://doi.org/10.1073/pnas.1222534110}{doi:\allowbreak 10.\allowbreak 1073/\allowbreak pnas.\allowbreak 1222534\allowbreak 110} \\
\addlinespace[5pt]
28 & \textbf{Wu 2017 (genome-wide position effects, 1,044 loci)}\,\textsuperscript{\textbf{M}} & \org{S. cerevisiae} & Regulatory DNA & 1,044 evenly spaced integration loci & promoter, reporter and carbon-\allowbreak source control arms & $1.0\times 10^{3}$ & $1.0\times 10^{3}$ & V & RFP intensity, plate reader & S288C+\allowbreak integration-\allowbreak locus & \href{https://www.ncbi.nlm.nih.gov/pmc/articles/PMC5516366/}{PMC55163\allowbreak 66} \\
\addlinespace[5pt]
29 & \textbf{ProteinGym yeast subset} & \org{S. cerevisiae} & Deep mutational scan & yeast subset of 217 assays over 186 proteins; yeast count needs a live query & one selection per assay & --$\ddagger$ & -- & V & per-\allowbreak assay deep mutational scanning score & S288C+\allowbreak designed-\allowbreak edit & \href{https://proteingym.org}{proteing\allowbreak ym.\allowbreak org} \\
\addlinespace[5pt]
30 & \textbf{Zhang 2016 (genome-wide GFP position scan)}\,\textsuperscript{\textbf{M}} & \org{S. cerevisiae} & Regulatory DNA & 428 high-\allowbreak  and 444 low-\allowbreak expressed sites & 1 & 872 & 872 & V & GFP by FACS & S288C+\allowbreak integration-\allowbreak locus & \href{https://wsws.ijournals.cn/cjbcn/article/pdf/gc16070901}{wsws.\allowbreak ijourna\allowbreak ls.\allowbreak cn/\allowbreak cjbcn/\allowbreak article\allowbreak /\allowbreak pdf/\allowbreak gc16070\allowbreak 901} \\
\addlinespace[5pt]
31 & \textbf{Kita 2017 (context-stratified yeast eQTLs)}\,\textsuperscript{\textbf{M}} & \org{S. cerevisiae} & Natural variation & 683 selected of 1,640 eQTLs & not reported & 683 & 683 & V & eQTL effect on expression & isolate-\allowbreak WGS & \href{https://www.pnas.org/doi/10.1073/pnas.1717421114}{pnas.\allowbreak org/\allowbreak doi/\allowbreak 10.\allowbreak 1073/\allowbreak pnas.\allowbreak 1717421\allowbreak 114} \\
\addlinespace[5pt]
32 & \textbf{Chen (earlier position-effect study, 482 loci)}\,\textsuperscript{\textbf{M}} & \org{S. cerevisiae} & Regulatory DNA & 482 loci, about a third on chromosomes I and VI & not reported & 482 & 482 & S & reporter expression per locus & S288C+\allowbreak integration-\allowbreak locus & \href{https://www.ncbi.nlm.nih.gov/pmc/articles/PMC5516366/}{PMC55163\allowbreak 66} \\
\addlinespace[5pt]
33 & \textbf{Wei 2017 (synthetic terminator panel)}\,\textsuperscript{\textbf{M}} & \org{S. cerevisiae} & Engineering parts & 100 terminators, 45 strong, 31 moderate, 24 weak & 1 & 100 & 100 & V & eGFP strength relative to the PGK1t standard & plasmid-\allowbreak reporter-\allowbreak library & \href{https://doi.org/10.1002/cbic.201700516}{doi:\allowbreak 10.\allowbreak 1002/\allowbreak cbic.\allowbreak 2017005\allowbreak 16} \\
\addlinespace[5pt]
34 & \textbf{Thompson 2019 (KT2440 lysine metabolism)}\,\textsuperscript{\textbf{M}} & \org{P. putida} & Transposon fitness & 39 genes with fitness below -\allowbreak 2 on a lysine route and above -\allowbreak 0.5 on glucose & 5-\allowbreak aminovalerate, D-\allowbreak lysine, L-\allowbreak lysine, glucose & 4 & $2.2\times 10^{4}$ & S & RB-\allowbreak TnSeq gene fitness & KT2440+\allowbreak transposon & \href{https://pmc.ncbi.nlm.nih.gov/articles/PMC6509195}{PMC65091\allowbreak 95} \\
\addlinespace[5pt]
35 & \textbf{Morse (short synthetic terminators)}\,\textsuperscript{\textbf{M}} & \org{S. cerevisiae} & Engineering parts & short synthetic terminators, 35 to 70 bp, count not reported & 1 & --$\ddagger$ & -- & S & fluorescent protein output & plasmid-\allowbreak reporter-\allowbreak library & \href{https://scite.ai/authors/nicholas-j-morse-1g5Qrn}{scite.\allowbreak ai/\allowbreak authors\allowbreak /\allowbreak nichola\allowbreak s-\allowbreak j-\allowbreak morse-\allowbreak 1g5Qrn} \\
\addlinespace[5pt]
36 & \textbf{MaveDB yeast score sets}\,\textsuperscript{\textbf{M}} & \org{S. cerevisiae} & Deep mutational scan & yeast score sets; count needs a live query & one selection per score set & --$\ddagger$ & -- & S & variant effect score & S288C+\allowbreak designed-\allowbreak edit & \href{https://doi.org/10.1186/s13059-019-1845-6}{doi:\allowbreak 10.\allowbreak 1186/\allowbreak s13059-\allowbreak 019-\allowbreak 1845-\allowbreak 6} \\
\addlinespace[5pt]
37 & \textbf{Isoprenol CRISPRi DBTL proteome ladder (7 PRIDE deposits)} & \org{P. putida} & Proteome & one engineered genotype set per design round; counts not reported per round & 7 design-\allowbreak build-\allowbreak test-\allowbreak learn rounds & 7 & $3.9\times 10^{4}$ & I & relative protein abundance & engineered-\allowbreak chassis+\allowbreak guide & \href{https://www.ebi.ac.uk/pride/archive/projects/PXD063733}{ebi.\allowbreak ac.\allowbreak uk/\allowbreak pride/\allowbreak archive\allowbreak /\allowbreak project\allowbreak s/\allowbreak PXD0637\allowbreak 33} \\
\addlinespace[5pt]
38 & \textbf{Thompson 2020 (KT2440 fatty acids and alcohols)} & \org{P. putida} & Transposon fitness & strong phenotypes for hundreds of genes; no count stated & 13 fatty acids and 10 alcohols & 23 & $1.3\times 10^{5}$ & S & RB-\allowbreak TnSeq gene fitness & KT2440+\allowbreak transposon & \href{https://pubmed.ncbi.nlm.nih.gov/32826213}{PMID 328\allowbreak 26213} \\
\addlinespace[5pt]
39 & \textbf{GSE198395 (putidaPRECISE321 carbon sources)} & \org{P. putida} & Transcriptome & KT2440 and engineered strains; genotype count not reported & different carbon sources, count not reported & 81 & $4.5\times 10^{5}$ & V & RNA-\allowbreak seq expression & engineered-\allowbreak chassis & \href{https://www.ncbi.nlm.nih.gov/geo/query/acc.cgi?acc=GSE198395}{ncbi.\allowbreak nlm.\allowbreak nih.\allowbreak gov/\allowbreak geo/\allowbreak query/\allowbreak acc.\allowbreak cgi?\allowbreak acc=\allowbreak GSE1983\allowbreak 95} \\
\addlinespace[5pt]
40 & \textbf{Mueller 2025 (24-strain P. putida Biolog panel)}\,\textsuperscript{\textbf{R}} & \org{P. putida} & Natural variation & 24 strains including KT2440 & Biolog PM01 and PM02A, 190 carbon compounds, plus 15 aromatics & $4.9\times 10^{3}$$\dagger$ & $4.9\times 10^{3}$ & V & Biolog respiration per well & pseudomonad-\allowbreak isolate-\allowbreak WGS & \href{https://doi.org/10.1101/2025.11.17.688870}{doi:\allowbreak 10.\allowbreak 1101/\allowbreak 2025.\allowbreak 11.\allowbreak 17.\allowbreak 688870} \\
\addlinespace[5pt]
41 & \textbf{Schmidt 2022 (KT2440 nitrogen BarSeq)} & \org{P. putida} & Transposon fitness & 672 genes with a significant phenotype, 529 of them UniProt-\allowbreak unreviewed & 52 nitrogen compounds plus 19 amino-\allowbreak acid dropouts & 71 & $4.0\times 10^{5}$ & I & BarSeq competition fitness per gene & KT2440+\allowbreak transposon & \href{https://pmc.ncbi.nlm.nih.gov/articles/PMC9004399}{PMC90043\allowbreak 99} \\
\addlinespace[5pt]
42 & \textbf{Borchert 2024 (KT2440 fitness compendium)} & \org{P. putida} & Transposon fitness & 4,732 genes with data of 5,564 protein-\allowbreak coding & 332 samples; unique conditions reported as 179 in the abstract and 183 in the methods & 332 & $1.6\times 10^{6}$ & V & RB-\allowbreak TnSeq gene fitness with t statistics & KT2440+\allowbreak transposon & \href{https://doi.org/10.1128/msystems.00942-23}{doi:\allowbreak 10.\allowbreak 1128/\allowbreak msystem\allowbreak s.\allowbreak 00942-\allowbreak 23} \\
\addlinespace[5pt]
43 & \textbf{putidaPRECISE321 (Lim 2022)} & \org{P. putida} & Transcriptome & KT2440 and engineered strains; strain-\allowbreak genotype count not reported & 118 unique conditions across 21 projects & 321 & $1.8\times 10^{6}$ & I & RNA-\allowbreak seq expression & KT2440-\allowbreak reference & \href{https://doi.org/10.1016/j.ymben.2022.04.004}{doi:\allowbreak 10.\allowbreak 1016/\allowbreak j.\allowbreak ymben.\allowbreak 2022.\allowbreak 04.\allowbreak 004} \\
\addlinespace[5pt]
44 & \textbf{GSE262836 (Delhi NICU, fluconazole-resistance expression)} & \org{P. kudriavzevii} & Transcriptome & one fluconazole-\allowbreak susceptible and one resistant isolate & 1 & 4 & $2.0\times 10^{4}$ & V & transcript abundance & Pk-\allowbreak isolate-\allowbreak WGS & \href{https://www.ncbi.nlm.nih.gov/geo/query/acc.cgi?acc=GSE262836}{ncbi.\allowbreak nlm.\allowbreak nih.\allowbreak gov/\allowbreak geo/\allowbreak query/\allowbreak acc.\allowbreak cgi?\allowbreak acc=\allowbreak GSE2628\allowbreak 36} \\
\addlinespace[5pt]
45 & \textbf{GSE159457 (CY902 C4 dicarboxylic acid transporters)} & \org{P. kudriavzevii} & Transcriptome & CY902 ura3 & pH 5.0 on glucose, succinic acid, or both & 6 & $3.0\times 10^{4}$ & V & transcript abundance & engineered-\allowbreak chassis & \href{https://www.ncbi.nlm.nih.gov/geo/query/acc.cgi?acc=GSE159457}{ncbi.\allowbreak nlm.\allowbreak nih.\allowbreak gov/\allowbreak geo/\allowbreak query/\allowbreak acc.\allowbreak cgi?\allowbreak acc=\allowbreak GSE1594\allowbreak 57} \\
\addlinespace[5pt]
46 & \textbf{Dubinkina 2024 (12-strain low-pH transcriptomic atlas)}\,\textsuperscript{\textbf{M}} & \org{P. kudriavzevii} & Transcriptome & 12 strains, 6 acid-\allowbreak tolerant and 6 susceptible & pH 5.5 control and pH 1.5 at 0 h and 1 h, plus two organic-\allowbreak acid arms & 48$\dagger$ & $2.4\times 10^{5}$ & I & gene expression count matrix & Pk-\allowbreak isolate-\allowbreak WGS & \href{https://doi.org/10.1128/spectrum.02536-23}{doi:\allowbreak 10.\allowbreak 1128/\allowbreak spectru\allowbreak m.\allowbreak 02536-\allowbreak 23} \\
\addlinespace[5pt]
47 & \textbf{MAESTRO 2025 (170-isolate WGS, 57-stressor growth panel)}\,\textsuperscript{\textbf{R}} & \org{P. kudriavzevii} & Natural variation & 170 WGS isolates, 161 phenotyped & 57 stressors: heat, low pH, organic acids, inhibitors, antifungals & $9.2\times 10^{3}$$\dagger$ & $9.2\times 10^{3}$ & I & growth under each stressor & Pk-\allowbreak isolate-\allowbreak WGS & \href{https://doi.org/10.1101/2025.08.29.673034}{doi:\allowbreak 10.\allowbreak 1101/\allowbreak 2025.\allowbreak 08.\allowbreak 29.\allowbreak 673034} \\
\addlinespace[5pt]
48 & \textbf{Fatma 2023 (landing-pad locus screen)}\,\textsuperscript{\textbf{R}} & \org{P. kudriavzevii} & CRISPR library screen & 47 gRNAs across 21 intergenic regions & 4 readouts per locus & 188$\dagger$ & 188 & I & cutting efficiency, integration efficiency, fitness, GFP level & SD108+\allowbreak guide & \href{https://doi.org/10.1101/2023.05.21.541627}{doi:\allowbreak 10.\allowbreak 1101/\allowbreak 2023.\allowbreak 05.\allowbreak 21.\allowbreak 541627} \\
\addlinespace[5pt]
49 & \textbf{GSE282019 (Delhi NICU five-year outbreak)} & \org{P. kudriavzevii} & Transcriptome & 3 isolates with differing fluconazole MICs & 1 & 6 & $3.0\times 10^{4}$ & V & transcript abundance & Pk-\allowbreak isolate-\allowbreak WGS & \href{https://www.ncbi.nlm.nih.gov/geo/query/acc.cgi?acc=GSE282019}{ncbi.\allowbreak nlm.\allowbreak nih.\allowbreak gov/\allowbreak geo/\allowbreak query/\allowbreak acc.\allowbreak cgi?\allowbreak acc=\allowbreak GSE2820\allowbreak 19} \\
\addlinespace[5pt]
50 & \textbf{GSE202655 (L-malic acid, EcSthA transhydrogenase)} & \org{P. kudriavzevii} & Transcriptome & MA009-\allowbreak 5 original and MA009-\allowbreak 10 carrying E. coli SthA & low pH & 6 & $3.0\times 10^{4}$ & V & transcript abundance, with titer and conversion reported separately & engineered-\allowbreak chassis & \href{https://www.ncbi.nlm.nih.gov/geo/query/acc.cgi?acc=GSE202655}{ncbi.\allowbreak nlm.\allowbreak nih.\allowbreak gov/\allowbreak geo/\allowbreak query/\allowbreak acc.\allowbreak cgi?\allowbreak acc=\allowbreak GSE2026\allowbreak 55} \\
\addlinespace[5pt]
\midrule \multicolumn{12}{@{}l}{\textbf{End of wave 1. Rows 1--50 are the recommended next builds.}}\\ \midrule
51 & \textbf{Lubliner 2013 (native core promoter maximal strengths)}\,\textsuperscript{\textbf{M}} & \org{S. cerevisiae} & Regulatory DNA & native core promoters, count not reported & not reported & --$\ddagger$ & -- & S & native core promoter maximal strength & plasmid-\allowbreak reporter-\allowbreak library & \href{https://academic.oup.com/nar}{academic\allowbreak .\allowbreak oup.\allowbreak com/\allowbreak nar} \\
\addlinespace[5pt]
52 & \textbf{Fitness Browser RCH2 (Stutzerimonas stutzeri\_AE)} & \org{S. stutzeri} & Fitness / chemical genomics & pooled barcoded transposon library, gene count not reported per organism & 30 carbon, 19 nitrogen, 50 stress, 38 other & 350 & $1.6\times 10^{6}$ & I & BarSeq competition fitness per gene, with a t statistic & fitbrowser+\allowbreak transposon & \href{https://doi.org/10.1038/s41586-018-0124-0}{doi:\allowbreak 10.\allowbreak 1038/\allowbreak s41586-\allowbreak 018-\allowbreak 0124-\allowbreak 0} \\
\addlinespace[5pt]
53 & \textbf{Fitness Browser FW300-N2E3 (Aquipseudomonas helvetica)} & \org{P. fluorescens} & Fitness / chemical genomics & pooled barcoded transposon library, gene count not reported per organism & 46 carbon, 30 nitrogen, 42 stress, 12 other & 211 & $9.6\times 10^{5}$ & I & BarSeq competition fitness per gene, with a t statistic & fitbrowser+\allowbreak transposon & \href{https://doi.org/10.1038/s41586-018-0124-0}{doi:\allowbreak 10.\allowbreak 1038/\allowbreak s41586-\allowbreak 018-\allowbreak 0124-\allowbreak 0} \\
\addlinespace[5pt]
54 & \textbf{Fitness Browser FW300-N2E2 (Aquipseudomonas fluorescens\_Q)} & \org{P. fluorescens} & Fitness / chemical genomics & pooled barcoded transposon library, gene count not reported per organism & 41 carbon, 25 nitrogen, 34 stress, 7 other & 203 & $9.2\times 10^{5}$ & I & BarSeq competition fitness per gene, with a t statistic & fitbrowser+\allowbreak transposon & \href{https://doi.org/10.1038/s41586-018-0124-0}{doi:\allowbreak 10.\allowbreak 1038/\allowbreak s41586-\allowbreak 018-\allowbreak 0124-\allowbreak 0} \\
\addlinespace[5pt]
55 & \textbf{Fitness Browser WCS417 (Aquipseudomonas simiae)} & \org{P. simiae} & Fitness / chemical genomics & pooled barcoded transposon library, gene count not reported per organism & 51 carbon, 11 nitrogen, 30 stress, 8 other & 201 & $9.1\times 10^{5}$ & I & BarSeq competition fitness per gene, with a t statistic & fitbrowser+\allowbreak transposon & \href{https://doi.org/10.1038/s41586-018-0124-0}{doi:\allowbreak 10.\allowbreak 1038/\allowbreak s41586-\allowbreak 018-\allowbreak 0124-\allowbreak 0} \\
\addlinespace[5pt]
56 & \textbf{Fitness Browser FW300-N2C3 (Aquipseudomonas fluorescens\_AA)} & \org{P. fluorescens} & Fitness / chemical genomics & pooled barcoded transposon library, gene count not reported per organism & 45 carbon, 19 nitrogen, 36 stress, 16 other & 184 & $8.3\times 10^{5}$ & I & BarSeq competition fitness per gene, with a t statistic & fitbrowser+\allowbreak transposon & \href{https://doi.org/10.1038/s41586-018-0124-0}{doi:\allowbreak 10.\allowbreak 1038/\allowbreak s41586-\allowbreak 018-\allowbreak 0124-\allowbreak 0} \\
\addlinespace[5pt]
57 & \textbf{Fitness Browser FW300-N1B4 (Aquipseudomonas sp.)} & \org{P. fluorescens} & Fitness / chemical genomics & pooled barcoded transposon library, gene count not reported per organism & 39 carbon, 21 nitrogen, 34 stress, 10 other & 147 & $6.7\times 10^{5}$ & I & BarSeq competition fitness per gene, with a t statistic & fitbrowser+\allowbreak transposon & \href{https://doi.org/10.1038/s41586-018-0124-0}{doi:\allowbreak 10.\allowbreak 1038/\allowbreak s41586-\allowbreak 018-\allowbreak 0124-\allowbreak 0} \\
\addlinespace[5pt]
58 & \textbf{Fitness Browser B728a (Aquipseudomonas fragariae pv. syringae)} & \org{P. syringae} & Fitness / chemical genomics & pooled barcoded transposon library, gene count not reported per organism & 17 stress, 8 other; no carbon or nitrogen panel & 46 & $2.1\times 10^{5}$ & I & BarSeq competition fitness per gene, with a t statistic & fitbrowser+\allowbreak transposon & \href{https://doi.org/10.1038/s41586-018-0124-0}{doi:\allowbreak 10.\allowbreak 1038/\allowbreak s41586-\allowbreak 018-\allowbreak 0124-\allowbreak 0} \\
\addlinespace[5pt]
59 & \textbf{Fitness Browser B728a mexB deletion background} & \org{P. syringae} & Fitness / chemical genomics & pooled barcoded transposon library in a mexB deletion background & 17 stress, 1 other & 23 & $1.0\times 10^{5}$ & I & BarSeq competition fitness per gene, with a t statistic & fitbrowser+\allowbreak transposon & \href{https://doi.org/10.1038/s41586-018-0124-0}{doi:\allowbreak 10.\allowbreak 1038/\allowbreak s41586-\allowbreak 018-\allowbreak 0124-\allowbreak 0} \\
\addlinespace[5pt]
60 & \textbf{Fitness Browser PA14 (Pseudomonas aeruginosa)} & \org{P. aeruginosa} & Fitness / chemical genomics & pooled barcoded transposon library, gene count not reported per organism & 3 other; no carbon, nitrogen or stress panel & 8 & $3.6\times 10^{4}$ & I & BarSeq competition fitness per gene, with a t statistic & fitbrowser+\allowbreak transposon & \href{https://doi.org/10.1038/s41586-018-0124-0}{doi:\allowbreak 10.\allowbreak 1038/\allowbreak s41586-\allowbreak 018-\allowbreak 0124-\allowbreak 0} \\
\addlinespace[5pt]
61 & \textbf{Cao 2020 (pooled four-condition RNA-seq)}\,\textsuperscript{\textbf{M}} & \org{P. kudriavzevii} & Transcriptome & SD108 wild type & glucose and glucose plus inhibitors, each aerobic and anaerobic & 1 & $4.9\times 10^{3}$ & I & transcript abundance & SD108-\allowbreak reference & \href{https://doi.org/10.1016/j.ymben.2020.01.005}{doi:\allowbreak 10.\allowbreak 1016/\allowbreak j.\allowbreak ymben.\allowbreak 2020.\allowbreak 01.\allowbreak 005} \\
\addlinespace[5pt]
62 & \textbf{PXD069498 (extracellular vesicle proteome)} & \org{P. kudriavzevii} & Proteome & wild type, strain not reported & vesicle purification & 189 & 189 & S & protein identification in purified vesicles & Pk-\allowbreak reference & \href{https://www.ebi.ac.uk/pride/archive/projects/PXD069498}{ebi.\allowbreak ac.\allowbreak uk/\allowbreak pride/\allowbreak archive\allowbreak /\allowbreak project\allowbreak s/\allowbreak PXD0694\allowbreak 98} \\
\addlinespace[5pt]
63 & \textbf{Tan 2026 (I. orientalis 3-HP strain panel, Nat Commun)} & \org{P. kudriavzevii} & Production campaign & 21 strain configurations off SD108 ura3 and the LPB landing pad, each specified by gene, copy number and integration locus; the total is read off the figure's columns and is never stated & glucose in CSM and in corn steep liquor, across shake flask, 300 mL DASbox and 3 L bioreactor & 25$\ddagger$ & 75 & I & titer, yield and productivity & SD108+\allowbreak designed-\allowbreak edit & \href{https://doi.org/10.1038/s41467-025-67621-8}{doi:\allowbreak 10.\allowbreak 1038/\allowbreak s41467-\allowbreak 025-\allowbreak 67621-\allowbreak 8} \\
\addlinespace[5pt]
64 & \textbf{Douglass 2018 (per-strain SNP density and loss of heterozygosity)} & \org{P. kudriavzevii} & Modality / backbone & 32 isolates & none; a genome-\allowbreak architecture statistic & 64$\dagger$ & 64 & I & SNP density and loss-\allowbreak of-\allowbreak heterozygosity extent per strain & Pk-\allowbreak isolate-\allowbreak WGS & \href{https://www.ncbi.nlm.nih.gov/pmc/articles/PMC6053246/}{PMC60532\allowbreak 46} \\
\addlinespace[5pt]
65 & \textbf{Tran 2023 (succinate end-to-end pipeline)}\,\textsuperscript{\textbf{M}} & \org{P. kudriavzevii} & Production campaign & at least 10 genotypes countable from the body text & sugar media at pH 3 fed-\allowbreak batch, sugarcane juice, pilot batch & 30$\ddagger$ & 30 & I & titer, rate and yield & SD108+\allowbreak designed-\allowbreak edit & \href{https://pmc.ncbi.nlm.nih.gov/articles/PMC10547785/}{PMC10547\allowbreak 785} \\
\addlinespace[5pt]
66 & \textbf{Xylose 3-HP from cellulosic hydrolysates (2026)}\,\textsuperscript{\textbf{M}} & \org{P. kudriavzevii} & Production campaign & 1AP and three xylose derivatives, 1APX, 1APX-\allowbreak X2 and 1APX-\allowbreak X23, off the SD108 lineage of the Tan 2026 parent & xylose and cellulosic hydrolysates without nutrient supplementation & 8$\dagger$ & 24 & V & titer, yield and productivity & SD108+\allowbreak designed-\allowbreak edit & \href{https://doi.org/10.1016/j.biortech.2026.134783}{doi:\allowbreak 10.\allowbreak 1016/\allowbreak j.\allowbreak biortec\allowbreak h.\allowbreak 2026.\allowbreak 134783} \\
\addlinespace[5pt]
67 & \textbf{US9845484B2 (Novozymes and Cargill 3-HP strain tables)} & \org{P. kudriavzevii} & Production campaign & 11 named strains in the CNB1 background built from CD1822, each with locus-\allowbreak explicit edits carrying SEQ ID numbers & one production condition per table & 11 & 22 & V & titer in g/L and yield & Pk-\allowbreak named+\allowbreak designed-\allowbreak edit & \href{https://patents.google.com/patent/US9845484B2/en}{patents.\allowbreak google.\allowbreak com/\allowbreak patent/\allowbreak US98454\allowbreak 84B2/\allowbreak en} \\
\addlinespace[5pt]
68 & \textbf{Cao 2020 (constitutive promoter-GFP panel)}\,\textsuperscript{\textbf{M}} & \org{P. kudriavzevii} & Engineering parts & 16 named promoters plus an unstated number more & different culture conditions, count not reported & 16$\ddagger$ & 16 & I & GFP level per promoter construct & SD108-\allowbreak reference & \href{https://doi.org/10.1016/j.ymben.2020.01.005}{doi:\allowbreak 10.\allowbreak 1016/\allowbreak j.\allowbreak ymben.\allowbreak 2020.\allowbreak 01.\allowbreak 005} \\
\addlinespace[5pt]
69 & \textbf{US10633664B2 (Cargill 3-HP tolerance dose response)} & \org{P. kudriavzevii} & Tolerance / robustness & 2 named CNB1-\allowbreak background strains plus a null control, locus-\allowbreak explicit with SEQ ID numbers & 3-\allowbreak HP in the medium at 0, 35 and 75 g/L & 9$\dagger$ & 9 & V & growth rate under 3-\allowbreak HP, with enzyme activity reported beside it & Pk-\allowbreak named+\allowbreak designed-\allowbreak edit & \href{https://patents.google.com/patent/US10633664B2/en}{patents.\allowbreak google.\allowbreak com/\allowbreak patent/\allowbreak US10633\allowbreak 664B2/\allowbreak en} \\
\addlinespace[5pt]
70 & \textbf{Xiao 2014 (SD108 reductive-TCA succinate)} & \org{P. kudriavzevii} & Production campaign & SD108 wild type and a uracil-\allowbreak auxotroph rTCA integrant & 4 acid-\allowbreak tolerance conditions plus shake-\allowbreak flask batch & 8$\ddagger$ & 8 & S & titer, yield, productivity, and acid-\allowbreak tolerance growth calls & SD108-\allowbreak reference & \href{https://pmc.ncbi.nlm.nih.gov/articles/PMC4244060}{PMC42440\allowbreak 60} \\
\addlinespace[5pt]
\midrule \multicolumn{12}{@{}l}{\textbf{End of wave 2. Rows 51--70 are the bench, promoted as soon as a wave-1 row proves unreachable.}}\\ \midrule
71 & \textbf{Zhang 2026 (green pretreatment, 3-HP and ethanol)}\,\textsuperscript{\textbf{M}} & \org{P. kudriavzevii} & Production campaign & one engineered strain, presumed the NRRL Y-\allowbreak 27441 lineage and not stated in the reachable text & corn stover and brewer's spent grain, each at 5, 10 and 20 wt\% solids, with and without carbonation & 6$\dagger$ & 6 & S & titer & Pk-\allowbreak named+\allowbreak designed-\allowbreak edit & \href{https://doi.org/10.1016/j.biortech.2025.133576}{doi:\allowbreak 10.\allowbreak 1016/\allowbreak j.\allowbreak biortec\allowbreak h.\allowbreak 2025.\allowbreak 133576} \\
\addlinespace[5pt]
72 & \textbf{Xylose utilization across four natural backgrounds}\,\textsuperscript{\textbf{M}} & \org{P. kudriavzevii} & Production campaign & SD108, IO21, IO45, IO46, each with the same three-\allowbreak gene cassette & xylose & 4 & 4 & S & xylose utilization & SD108+\allowbreak designed-\allowbreak edit & \href{https://core.ac.uk/works/286100561}{core.\allowbreak ac.\allowbreak uk/\allowbreak works/\allowbreak 2861005\allowbreak 61} \\
\addlinespace[5pt]
73 & \textbf{US9365875B2 (Novozymes 3-HP per-biomass panel)} & \org{P. kudriavzevii} & Production campaign & 4 named CNB1-\allowbreak background strains with locus-\allowbreak explicit edits & one production condition & 4 & 4 & V & 3-\allowbreak HP per dry cell weight & Pk-\allowbreak named+\allowbreak designed-\allowbreak edit & \href{https://patents.google.com/patent/US9365875B2/en}{patents.\allowbreak google.\allowbreak com/\allowbreak patent/\allowbreak US93658\allowbreak 75B2/\allowbreak en} \\
\addlinespace[5pt]
74 & \textbf{Jeong 2025 (acetate metabolism during xylose fermentation)}\,\textsuperscript{\textbf{M}} & \org{P. kudriavzevii} & Production campaign & IoDY01H, a named production background with no reference assembly in the genomes tier & hemp stalk hydrolysate by separate hydrolysis and fermentation, at pH 3.5 and pH 5.5 & 2$\dagger$ & 2 & V & titer & Pk-\allowbreak named+\allowbreak designed-\allowbreak edit & \href{https://doi.org/10.1016/j.biortech.2025.133113}{doi:\allowbreak 10.\allowbreak 1016/\allowbreak j.\allowbreak biortec\allowbreak h.\allowbreak 2025.\allowbreak 133113} \\
\addlinespace[5pt]
75 & \textbf{Douglass 2018 (CBS573 RNA-seq)} & \org{P. kudriavzevii} & Transcriptome & CBS573 type strain & not reported & --$\ddagger$ & -- & V & transcript abundance & CBS573-\allowbreak reference & \href{https://www.ncbi.nlm.nih.gov/sra/?term=SRP139056}{ncbi.\allowbreak nlm.\allowbreak nih.\allowbreak gov/\allowbreak sra/\allowbreak ?\allowbreak term=\allowbreak SRP1390\allowbreak 56} \\
\addlinespace[5pt]
76 & \textbf{MAESTRO 2025 (gene4260/gene4261 paralog-pair deletion)}\,\textsuperscript{\textbf{R}} & \org{P. kudriavzevii} & Tolerance / robustness & one tandem paralog pair deleted & several stresses, count not reported & --$\ddagger$ & -- & S & stress-\allowbreak tolerance effect of the deletion & SD108+\allowbreak designed-\allowbreak edit & \href{https://doi.org/10.1101/2025.08.29.673034}{doi:\allowbreak 10.\allowbreak 1101/\allowbreak 2025.\allowbreak 08.\allowbreak 29.\allowbreak 673034} \\
\addlinespace[5pt]
77 & \textbf{Illinois Data Bank IDB-0538107 (succinate pipeline deposit)}\,\textsuperscript{\textbf{M}} & \org{P. kudriavzevii} & Production campaign & not reported on the deposit page & fed-\allowbreak batch minimal medium, fed-\allowbreak batch sugarcane juice, pilot batch, all at pH 3 & --$\ddagger$ & -- & S & titer & SD108+\allowbreak designed-\allowbreak edit & \href{https://databank.illinois.edu/datasets/IDB-0538107}{databank\allowbreak .\allowbreak illinoi\allowbreak s.\allowbreak edu/\allowbreak dataset\allowbreak s/\allowbreak IDB-\allowbreak 0538107} \\
\addlinespace[5pt]
78 & \textbf{CABBI inorganic carbon sources (low-pH succinate kinetics)}\,\textsuperscript{\textbf{M}} & \org{P. kudriavzevii} & Production campaign & not reported on the dataset page & inorganic carbon sources at low pH, count not reported & --$\ddagger$ & -- & S & succinate process and kinetic data & SD108+\allowbreak designed-\allowbreak edit & \href{https://cabbi.bio/datasets/}{cabbi.\allowbreak bio/\allowbreak dataset\allowbreak s} \\
\addlinespace[5pt]
79 & \textbf{Glycolic acid overproduction (2026)} & \org{P. kudriavzevii} & Production campaign & not reported & xylose-\allowbreak assisted culture & --$\ddagger$ & -- & S & glycolic acid titer & SD108+\allowbreak designed-\allowbreak edit & \href{https://doi.org/10.3390/jof12100751}{doi:\allowbreak 10.\allowbreak 3390/\allowbreak jof1210\allowbreak 0751} \\
\addlinespace[5pt]
80 & \textbf{Wu 2023 (piggyBac citramalate copy-number and locus variants)}\,\textsuperscript{\textbf{M}} & \org{P. kudriavzevii} & Production campaign & SD108, SD108 ura3, SD108 ura3 leu2, plus integrants varying cimA copy number and locus & glucose, 48 h batch & --$\ddagger$ & -- & S & citramalate titer and molar yield & SD108+\allowbreak designed-\allowbreak edit & \href{https://pmc.ncbi.nlm.nih.gov/articles/PMC9969067}{PMC99690\allowbreak 67} \\
\addlinespace[5pt]
81 & \textbf{Mitochondrial decompartmentalization (2025, with GNPS metabolomics)}\,\textsuperscript{\textbf{M}} & \org{P. kudriavzevii} & Metabolome / flux & not reported & not reported & --$\ddagger$ & -- & S & metabolite intensities & SD108+\allowbreak designed-\allowbreak edit & \href{https://doi.org/10.1038/s41467-025-62304-w}{doi:\allowbreak 10.\allowbreak 1038/\allowbreak s41467-\allowbreak 025-\allowbreak 62304-\allowbreak w} \\
\addlinespace[5pt]
82 & \textbf{GSE39558 (ethanol-stress expression)} & \org{P. kudriavzevii} & Transcriptome & not reported & ethanol stress & --$\ddagger$ & -- & S & transcript abundance & Pk-\allowbreak isolate-\allowbreak WGS & \href{https://www.ncbi.nlm.nih.gov/geo/query/acc.cgi?acc=GSE39558}{ncbi.\allowbreak nlm.\allowbreak nih.\allowbreak gov/\allowbreak geo/\allowbreak query/\allowbreak acc.\allowbreak cgi?\allowbreak acc=\allowbreak GSE3955\allowbreak 8} \\
\addlinespace[5pt]
83 & \textbf{PXD041801 (flocculin expression and floating biofilms)} & \org{P. kudriavzevii} & Proteome & not reported & not reported & --$\ddagger$ & -- & S & protein abundance associated with biofilm formation & Pk-\allowbreak reference & \href{https://www.ebi.ac.uk/pride/archive/projects/PXD041801}{ebi.\allowbreak ac.\allowbreak uk/\allowbreak pride/\allowbreak archive\allowbreak /\allowbreak project\allowbreak s/\allowbreak PXD0418\allowbreak 01} \\
\addlinespace[5pt]
84 & \textbf{PXD039973 (succinylation and amphotericin B resistance)} & \org{P. kudriavzevii} & Proteome & not reported & not reported & --$\ddagger$ & -- & S & protein succinylation site intensity & Pk-\allowbreak reference & \href{https://www.omicsdi.org/dataset/iprox/PXD039973}{omicsdi.\allowbreak org/\allowbreak dataset\allowbreak /\allowbreak iprox/\allowbreak PXD0399\allowbreak 73} \\
\addlinespace[5pt]
85 & \textbf{PXD075355 (Pichia against Botrytis cinerea)} & \org{P. kudriavzevii} & Proteome & not reported & antagonism against Botrytis cinerea & --$\ddagger$ & -- & S & protein abundance under antagonism & Pk-\allowbreak reference & \href{https://www.ebi.ac.uk/pride/archive/projects/PXD075355}{ebi.\allowbreak ac.\allowbreak uk/\allowbreak pride/\allowbreak archive\allowbreak /\allowbreak project\allowbreak s/\allowbreak PXD0753\allowbreak 55} \\
\addlinespace[5pt]
86 & \textbf{Pk E1 toolkit expansion (mycoprotein, 2026)} & \org{P. kudriavzevii} & Engineering parts & not reported & not reported & --$\ddagger$ & -- & S & parts characterization and mycoprotein output & engineered-\allowbreak chassis & \href{https://doi.org/10.1111/1751-7915.70447}{doi:\allowbreak 10.\allowbreak 1111/\allowbreak 1751-\allowbreak 7915.\allowbreak 70447} \\
\addlinespace[5pt]
87 & \textbf{China fermented-vegetable and clinical panel (2024)}\,\textsuperscript{\textbf{B}}\,\textsuperscript{\textbf{M}} & \org{P. kudriavzevii} & Natural variation & isolate count not captured; about a quarter of clinical and 87.5 percent of fermented-\allowbreak vegetable isolates are triploid & fluconazole, voriconazole, itraconazole, amphotericin B, caspofungin, posaconazole & --$\ddagger$ & -- & S & antifungal susceptibility & Pk-\allowbreak isolate-\allowbreak WGS & \href{https://www.ncbi.nlm.nih.gov/pmc/articles/PMC11487970/}{PMC11487\allowbreak 970} \\
\addlinespace[5pt]
88 & \textbf{Indonesian clinical genomes (PRJNA819536)}\,\textsuperscript{\textbf{M}} & \org{P. kudriavzevii} & Natural variation & 16 clinical isolates & resistance phenotypes intended, drug count not reported & --$\ddagger$ & -- & S & resistance phenotype, if a public table exists & Pk-\allowbreak isolate-\allowbreak WGS & \href{https://www.ncbi.nlm.nih.gov/pmc/articles/PMC10320322/}{PMC10320\allowbreak 322} \\
\addlinespace[5pt]
89 & \textbf{Unidentified Illinois Data Bank CRISPR/Cas9 gene-deletion deposit}\,\textsuperscript{\textbf{M}} & \org{P. kudriavzevii} & CRISPR library screen & not reported & not reported & --$\ddagger$ & -- & S & not reported & SD108+\allowbreak designed-\allowbreak edit & \href{https://www.osti.gov/biblio/3013685}{osti.\allowbreak gov/\allowbreak biblio/\allowbreak 3013685} \\
\addlinespace[5pt]
90 & \textbf{Jia 2025 (3-HP from brewer's spent grain)}\,\textsuperscript{\textbf{M}} & \org{P. kudriavzevii} & Production campaign & the IoDY01H background, off the NRRL Y-\allowbreak 27441 lineage rather than SD108; the strain count is not stated & four comparison axes: the pH-\allowbreak adjusting acid, nitrogen supplementation, bicarbonate, and deproteinization of the hydrolysate & --$\ddagger$ & -- & I & titer and yield & Pk-\allowbreak named+\allowbreak designed-\allowbreak edit & \href{https://doi.org/10.1016/j.jbiotec.2025.08.009}{doi:\allowbreak 10.\allowbreak 1016/\allowbreak j.\allowbreak jbiotec\allowbreak .\allowbreak 2025.\allowbreak 08.\allowbreak 009} \\
\addlinespace[5pt]
91 & \textbf{Tan 2026 (corn kernels, 3-HP and ethanol coproduction)}\,\textsuperscript{\textbf{M}} & \org{P. kudriavzevii} & Production campaign & not established; the article could not be reached & corn kernels as feedstock; condition count not established & --$\ddagger$ & -- & S & 3-\allowbreak HP and ethanol titer & Pk-\allowbreak named+\allowbreak designed-\allowbreak edit & \href{https://doi.org/10.1016/j.bcab.2025.103906}{doi:\allowbreak 10.\allowbreak 1016/\allowbreak j.\allowbreak bcab.\allowbreak 2025.\allowbreak 103906} \\
\addlinespace[5pt]
92 & \textbf{Liu 2020 (PPiSeq accessory interactome)} & \org{S. cerevisiae} & Fitness / chemical genomics & 1,660,856 protein pairs as BY4741 by BY4742 diploids, each carrying two cassettes at the HO locus & 9 environments, including peroxide, hydroxyurea, forskolin, doxorubicin, raffinose, FK506, salt and cold & $1.5\times 10^{7}$ & $1.5\times 10^{7}$ & V & relative fitness under selection, with a p-\allowbreak value and a binary call & S288C+\allowbreak integration-\allowbreak locus & \href{https://doi.org/10.7554/eLife.62365}{doi:\allowbreak 10.\allowbreak 7554/\allowbreak eLife.\allowbreak 62365} \\
\addlinespace[5pt]
93 & \textbf{Michel 2017 (SATAY transposon saturation)} & \org{S. cerevisiae} & Transposon fitness & 8 libraries across 7 genetic backgrounds plus a rapamycin arm, at 169,000 to 495,000 mapped insertions each & rich medium and a rapamycin arm & $2.6\times 10^{6}$ & $5.2\times 10^{6}$ & V & insertion density and read count as a fitness proxy & S288C+\allowbreak transposon & \href{https://doi.org/10.7554/eLife.23570}{doi:\allowbreak 10.\allowbreak 7554/\allowbreak eLife.\allowbreak 23570} \\
\addlinespace[5pt]
94 & \textbf{Aggeli 2021 (second-step adaptation, barcoded lineages)} & \org{S. cerevisiae} & Fitness / chemical genomics & 6 barcoded populations from 3 first-\allowbreak step mutant ancestors in replicate; lineage barcodes carry no gene identifier & one environment over 160 generations & $4.4\times 10^{6}$ & $4.4\times 10^{6}$ & V & barcode read-\allowbreak count trajectory, with fitness and establishment time per adaptive lineage & S288C+\allowbreak integration-\allowbreak locus & \href{https://doi.org/10.1038/s41467-021-25440-7}{doi:\allowbreak 10.\allowbreak 1038/\allowbreak s41467-\allowbreak 021-\allowbreak 25440-\allowbreak 7} \\
\addlinespace[5pt]
95 & \textbf{Khaledi 2020 (414-isolate clinical transcriptome panel)} & \org{P. aeruginosa} & Transcriptome & 414 clinical isolates, each with a genome and a transcriptome & 1 culture condition & 414 & $2.4\times 10^{6}$ & I & transcript abundance & pseudomonad-\allowbreak isolate-\allowbreak WGS & \href{https://doi.org/10.15252/emmm.201910264}{doi:\allowbreak 10.\allowbreak 15252/\allowbreak emmm.\allowbreak 2019102\allowbreak 64} \\
\addlinespace[5pt]
96 & \textbf{Blundell 2019 (adaptive diversity in two limitations)} & \org{S. cerevisiae} & Fitness / chemical genomics & 494,066 released barcoded lineages, against about 500,000 claimed & carbon-\allowbreak limited and nitrogen-\allowbreak limited & $2.1\times 10^{6}$ & $2.1\times 10^{6}$ & V & barcode read-\allowbreak count trajectory, with a separate fitness table & S288C+\allowbreak integration-\allowbreak locus & \href{https://doi.org/10.1038/s41559-018-0758-1}{doi:\allowbreak 10.\allowbreak 1038/\allowbreak s41559-\allowbreak 018-\allowbreak 0758-\allowbreak 1} \\
\addlinespace[5pt]
97 & \textbf{Nguyen Ba 2019 (travelling wave of adaptation)} & \org{S. cerevisiae} & Fitness / chemical genomics & about 50,000 barcodes per population, RE-\allowbreak BARCODED every 100 generations so lineages nest & rich medium and rich medium with acetic acid & $1.9\times 10^{6}$ & $1.9\times 10^{6}$ & V & barcode frequency trajectory, with clone fitness and confidence intervals & S288C+\allowbreak integration-\allowbreak locus & \href{https://doi.org/10.1038/s41586-019-1749-3}{doi:\allowbreak 10.\allowbreak 1038/\allowbreak s41586-\allowbreak 019-\allowbreak 1749-\allowbreak 3} \\
\addlinespace[5pt]
98 & \textbf{Karadzas 2026 (SATAY chemogenomics across 20 antifungals)} & \org{S. cerevisiae} & Transposon fitness & 5 drug-\allowbreak sensitive BY-\allowbreak derived backgrounds, transposon-\allowbreak mutagenized, across 27 library-\allowbreak condition screens & 21 named compounds plus controls, which the abstract rounds to 20 & $6.7\times 10^{5}$ & $1.3\times 10^{6}$ & V & transposon count and read count per feature and condition & S288C+\allowbreak transposon & \href{https://doi.org/10.7554/eLife.105333}{doi:\allowbreak 10.\allowbreak 7554/\allowbreak eLife.\allowbreak 105333} \\
\addlinespace[5pt]
99 & \textbf{Chen 2022 (transposon mutagenesis of 16 natural isolates)} & \org{S. cerevisiae} & Transposon fitness & 16 haploid natural isolates including S288C and RM11, each transposon-\allowbreak mutagenized across 5,813 genes & one condition & $9.3\times 10^{4}$$\dagger$ & $1.1\times 10^{6}$ & V & essentiality verdict with a 12-\allowbreak value insertion-\allowbreak statistics vector & isolate-\allowbreak WGS & \href{https://doi.org/10.1038/s41467-022-29228-1}{doi:\allowbreak 10.\allowbreak 1038/\allowbreak s41467-\allowbreak 022-\allowbreak 29228-\allowbreak 1} \\
\addlinespace[5pt]
100 & \textbf{Levy 2015 (high-resolution lineage tracking)} & \org{S. cerevisiae} & Fitness / chemical genomics & 494,066 released barcoded lineages in a BY4709 landing-\allowbreak pad background & one environment, two replicate experiments & $9.9\times 10^{5}$ & $9.9\times 10^{5}$ & V & barcode read-\allowbreak count trajectory & S288C+\allowbreak integration-\allowbreak locus & \href{https://doi.org/10.1038/nature14279}{doi:\allowbreak 10.\allowbreak 1038/\allowbreak nature1\allowbreak 4279} \\
\addlinespace[5pt]
101 & \textbf{Pokusaeva 2019 (HIS3 combinatorial fitness landscape)} & \org{S. cerevisiae} & Deep mutational scan & 875,151 unique amino-\allowbreak acid sequences of one protein, from 4,018,105 nucleotide genotypes & histidine-\allowbreak free medium over 168 hours & $9.6\times 10^{5}$ & $9.6\times 10^{5}$ & V & fitness, with a frequency trajectory & plasmid-\allowbreak coding-\allowbreak variant & \href{https://doi.org/10.1371/journal.pgen.1008079}{doi:\allowbreak 10.\allowbreak 1371/\allowbreak journal\allowbreak .\allowbreak pgen.\allowbreak 1008079} \\
\addlinespace[5pt]
102 & \textbf{Vieitez 2022 (5,362 strains across 102 conditions)} & \org{S. cerevisiae} & Fitness / chemical genomics & 5,362 strains, being 4,888 deletions plus 474 phosphosite point mutants, in two different backgrounds & 102 conditions & $5.4\times 10^{5}$ & $5.4\times 10^{5}$ & V & quantile-\allowbreak normalized colony-\allowbreak size score with a q-\allowbreak value & S288C+\allowbreak designed-\allowbreak edit & \href{https://doi.org/10.1038/s41587-021-01051-x}{doi:\allowbreak 10.\allowbreak 1038/\allowbreak s41587-\allowbreak 021-\allowbreak 01051-\allowbreak x} \\
\addlinespace[5pt]
103 & \textbf{Litsios 2024 (cell-cycle proteome movements, GFP collection)} & \org{S. cerevisiae} & Proteome & 3,902 GFP-\allowbreak tagged open reading frames in BY4741, 3,806 of them carrying a concentration & one growth condition resolved into 6 cell-\allowbreak cycle stages, in triplicate & $3.9\times 10^{3}$ & $5.4\times 10^{5}$ & V & localization class fractions and protein concentration per cell-\allowbreak cycle stage & S288C+\allowbreak designed-\allowbreak edit & \href{https://doi.org/10.1016/j.cell.2024.02.014}{doi:\allowbreak 10.\allowbreak 1016/\allowbreak j.\allowbreak cell.\allowbreak 2024.\allowbreak 02.\allowbreak 014} \\
\addlinespace[5pt]
104 & \textbf{Guenole 2013 (multiconditional DNA-damage E-MAP)} & \org{S. cerevisiae} & Fitness / chemical genomics & 97,578 gene pairs in the released grid, which the prose gives as a slightly different array & 4 conditions, three DNA-\allowbreak damaging agents and a vehicle control & $9.8\times 10^{4}$ & $3.9\times 10^{5}$ & V & colony-\allowbreak size interaction score & S288C-\allowbreak KO & \href{https://doi.org/10.1016/j.molcel.2012.11.023}{doi:\allowbreak 10.\allowbreak 1016/\allowbreak j.\allowbreak molcel.\allowbreak 2012.\allowbreak 11.\allowbreak 023} \\
\addlinespace[5pt]
105 & \textbf{Melnikov 2014 (aminoglycoside resistance deep mutational scan)} & \org{E. coli} & Deep mutational scan & 4,993 of 4,997 designed substitutions detected in an aminoglycoside kinase on a plasmid & 6 aminoglycosides at 4 doses each & $3.0\times 10^{5}$$\dagger$ & $3.0\times 10^{5}$ & V & change in amino-\allowbreak acid frequency under selection & plasmid-\allowbreak coding-\allowbreak variant & \href{https://doi.org/10.1093/nar/gku511}{doi:\allowbreak 10.\allowbreak 1093/\allowbreak nar/\allowbreak gku511} \\
\addlinespace[5pt]
106 & \textbf{BV-BRC E. coli antimicrobial susceptibility panel} & \org{E. coli} & Natural variation & 13,742 distinct sequenced E. coli genomes with a laboratory-\allowbreak measured susceptibility record & 104 antibiotic labels, before vocabulary normalization & $1.8\times 10^{5}$ & $1.8\times 10^{5}$ & V & minimum inhibitory concentration, mg/L & ecoli-\allowbreak isolate-\allowbreak WGS & \href{https://doi.org/10.1093/nar/gkac1003}{doi:\allowbreak 10.\allowbreak 1093/\allowbreak nar/\allowbreak gkac100\allowbreak 3} \\
\addlinespace[5pt]
107 & \textbf{Duan 2025 (Pol II trigger loop across nine conditions)} & \org{S. cerevisiae} & Deep mutational scan & 15,174 substitution variants of one RNA polymerase domain across ten libraries, 6,054 of them viable & 9 conditions including two drugs, manganese, two carbon sources, a lysine dropout and formamide & $1.4\times 10^{5}$$\dagger$ & $1.4\times 10^{5}$ & V & fitness per variant and condition & plasmid-\allowbreak coding-\allowbreak variant & \href{https://doi.org/10.1038/s41467-025-63304-6}{doi:\allowbreak 10.\allowbreak 1038/\allowbreak s41467-\allowbreak 025-\allowbreak 63304-\allowbreak 6} \\
\addlinespace[5pt]
108 & \textbf{Schramm 2023 and Lubrano 2025 (15,120 essential-gene substitutions)} & \org{E. coli} & Deep mutational scan & 15,120 strains, each carrying one distinct single amino-\allowbreak acid substitution across 346 essential proteins, from 16,038 designed & 30 and 42 degrees, plus carbenicillin and gentamicin in the companion study & $6.0\times 10^{4}$$\dagger$ & $6.0\times 10^{4}$ & V & pooled barcoded fitness & ecoli+\allowbreak designed-\allowbreak edit & \href{https://doi.org/10.15252/msb.202311596}{doi:\allowbreak 10.\allowbreak 15252/\allowbreak msb.\allowbreak 2023115\allowbreak 96} \\
\addlinespace[5pt]
109 & \textbf{Chen 2020 (VIM-2 beta-lactamase deep mutational scan)} & \org{E. coli} & Deep mutational scan & 5,535 of 5,607 designed variants covered, on a plasmid in an E. coli cloning host & 3 antibiotics across temperatures, 9 selection conditions & $5.0\times 10^{4}$ & $5.0\times 10^{4}$ & V & log2 enrichment fitness with its standard deviation & plasmid-\allowbreak coding-\allowbreak variant & \href{https://doi.org/10.7554/eLife.56707}{doi:\allowbreak 10.\allowbreak 7554/\allowbreak eLife.\allowbreak 56707} \\
\addlinespace[5pt]
110 & \textbf{Sharon 2018 (CRISPEY natural-variant precision editing)} & \org{S. cerevisiae} & CRISPR library screen & 16,006 natural variants installed one at a time into BY4742, each an RM11-\allowbreak 1a allele, through 32,000 guide and donor pairs & one selective medium over about 25 generations & $4.8\times 10^{4}$ & $4.8\times 10^{4}$ & V & log2 fold change per generation, with standard error and two significance measures & S288C+\allowbreak designed-\allowbreak edit & \href{https://doi.org/10.1016/j.cell.2018.08.057}{doi:\allowbreak 10.\allowbreak 1016/\allowbreak j.\allowbreak cell.\allowbreak 2018.\allowbreak 08.\allowbreak 057} \\
\addlinespace[5pt]
111 & \textbf{Schmidlin 2024 (drug-resistance tradeoffs across 12 environments)} & \org{S. cerevisiae} & Fitness / chemical genomics & 774 adaptive strains evolved in 12 antifungal conditions, on a landing-\allowbreak pad background carrying 4 reported fixed mutations & 12 environments spanning two antifungals & $9.3\times 10^{3}$ & $9.3\times 10^{3}$ & V & fitness per lineage, condition and replicate, with barcode counts beneath it & S288C+\allowbreak integration-\allowbreak locus & \href{https://doi.org/10.7554/eLife.94144}{doi:\allowbreak 10.\allowbreak 7554/\allowbreak eLife.\allowbreak 94144} \\
\addlinespace[5pt]
112 & \textbf{Reichling 2023 (dynamic metabolome across 164 deletions)} & \org{S. cerevisiae} & Metabolome / flux & 164 cataloged deletion mutants from the BY4741 collection, restored to prototrophy by a plasmid & 5 timepoints after rapamycin, at 0, 5, 30, 60 and 90 minutes & $8.9\times 10^{3}$ & $1.7\times 10^{6}$ & V & ion intensities with mass, formula, adduct and pathway identifiers & S288C-\allowbreak KO & \href{https://doi.org/10.7554/eLife.84295}{doi:\allowbreak 10.\allowbreak 7554/\allowbreak eLife.\allowbreak 84295} \\
\addlinespace[5pt]
113 & \textbf{Mateus 2020 (121 mutants, proteome abundance and thermal stability)} & \org{E. coli} & Proteome & 121 mutants, mostly cataloged single-\allowbreak gene deletions in BW25113 & 10 temperatures & 121 & $2.1\times 10^{5}$ & S & protein abundance and thermal stability & ecoli-\allowbreak KO & \href{https://doi.org/10.1038/s41586-020-3002-5}{doi:\allowbreak 10.\allowbreak 1038/\allowbreak s41586-\allowbreak 020-\allowbreak 3002-\allowbreak 5} \\
\addlinespace[5pt]
114 & \textbf{Suzuki 2014 (40 evolved resistant strains, expression and mutations)} & \org{E. coli} & Tolerance / robustness & 40 evolved resistant strains from 10 drugs across 4 independent lines, resequenced & one profiling condition & 41 & $1.8\times 10^{5}$ & V & transcript abundance, with 447 released mutation calls & evolved-\allowbreak WGS & \href{https://doi.org/10.1038/ncomms6792}{doi:\allowbreak 10.\allowbreak 1038/\allowbreak ncomms6\allowbreak 792} \\
\addlinespace[5pt]
115 & \textbf{Zampieri 2017 (190 resequenced evolved populations, metabolome)} & \org{E. coli} & Metabolome / flux & 190 resequenced evolved populations, which the methods also give as 193, across 4 lineages & 3 antibiotics in 2 carbon sources & 190 & $1.5\times 10^{5}$ & S & intracellular and extracellular ion intensities & evolved-\allowbreak WGS & \href{https://doi.org/10.15252/msb.20167028}{doi:\allowbreak 10.\allowbreak 15252/\allowbreak msb.\allowbreak 2016702\allowbreak 8} \\
\addlinespace[5pt]
116 & \textbf{Galardini 2017 ECOREF (894 strains across 214 conditions)} & \org{E. coli} & Natural variation & 869 strains in the released matrix, of 894 phenotyped and 696 sequenced & 214 conditions & $1.1\times 10^{5}$ & $1.1\times 10^{5}$ & V & colony-\allowbreak size S-\allowbreak score & ecoli-\allowbreak isolate-\allowbreak WGS & \href{https://doi.org/10.7554/eLife.31035}{doi:\allowbreak 10.\allowbreak 7554/\allowbreak eLife.\allowbreak 31035} \\
\addlinespace[5pt]
117 & \textbf{Kaczmarczyk 2025 (PAO1 genome-wide CRISPRi)}\,\textsuperscript{\textbf{R}} & \org{P. aeruginosa} & CRISPR library screen & over 80,000 sgRNAs over virtually all coding sequences and intergenic regions & glycolytic and gluconeogenic growth & $8.0\times 10^{4}$$\ddagger$ & $8.0\times 10^{4}$ & I & depletion and enrichment fitness per guide & PAO1+\allowbreak guide & \href{https://doi.org/10.1101/2025.08.12.669819}{doi:\allowbreak 10.\allowbreak 1101/\allowbreak 2025.\allowbreak 08.\allowbreak 12.\allowbreak 669819} \\
\addlinespace[5pt]
118 & \textbf{Lipworth 2024 (8,555 clinical isolates with linked susceptibility)} & \org{E. coli} & Natural variation & 8,555 assembled bloodstream isolates with linked susceptibility, of 9,001 assembled & 7 antibiotics & $6.0\times 10^{4}$$\dagger$ & $6.0\times 10^{4}$ & S & resistant or susceptible call on a clinical standard & ecoli-\allowbreak isolate-\allowbreak WGS & \href{https://doi.org/10.1016/s2666-5247(24)00152-6}{doi:\allowbreak 10.\allowbreak 1016/\allowbreak s2666-\allowbreak 5247(24\allowbreak )00152-\allowbreak 6} \\
\addlinespace[5pt]
119 & \textbf{Stiffler 2015 (TEM-1 beta-lactamase across seven selections)} & \org{E. coli} & Deep mutational scan & 4,997 substitution variants of TEM-\allowbreak 1 on a plasmid & 7 selection strengths & $3.5\times 10^{4}$$\dagger$ & $3.5\times 10^{4}$ & V & log10 enrichment fitness & plasmid-\allowbreak coding-\allowbreak variant & \href{https://doi.org/10.1016/j.cell.2015.01.035}{doi:\allowbreak 10.\allowbreak 1016/\allowbreak j.\allowbreak cell.\allowbreak 2015.\allowbreak 01.\allowbreak 035} \\
\addlinespace[5pt]
120 & \textbf{Furusawa 2020 (95 stressors, 192 resequenced clones, 47-chemical cross-resistance)} & \org{E. coli} & Tolerance / robustness & 192 resequenced clones from 95 stressors across 6 independent lines, plus 64 reconstructed single mutants & cross-\allowbreak resistance measured against 47 chemicals & 192 & $9.0\times 10^{3}$ & V & relative inhibitory concentration & evolved-\allowbreak WGS & \href{https://doi.org/10.1038/s41467-020-19713-w}{doi:\allowbreak 10.\allowbreak 1038/\allowbreak s41467-\allowbreak 020-\allowbreak 19713-\allowbreak w} \\
\addlinespace[5pt]
121 & \textbf{Dotsch 2009 (arrayed PA14 library against 19 antibiotics)}\,\textsuperscript{\textbf{R}} & \org{P. aeruginosa} & Tolerance / robustness & arrayed PA14 non-\allowbreak redundant mutants; the number screened is not stated & 19 clinically used antibiotics & $5.5\times 10^{3}$$\ddagger$ & $5.5\times 10^{3}$ & I & semiautomated antimicrobial susceptibility & PA14+\allowbreak transposon & \href{https://doi.org/10.1128/AAC.00035-09}{doi:\allowbreak 10.\allowbreak 1128/\allowbreak AAC.\allowbreak 00035-\allowbreak 09} \\
\addlinespace[5pt]
122 & \textbf{Bartell 2019 (longitudinal cystic fibrosis panel)}\,\textsuperscript{\textbf{M}} & \org{P. aeruginosa} & Natural variation & 443 longitudinal isolates, 52 clone types, 39 patients & 2 growth media plus the assay conditions for 6 further readouts & 443 & $3.5\times 10^{3}$ & I & growth rate in 2 media, 2 E-\allowbreak test MICs, protease, mucoidity, adhesion, aggregation & pseudomonad-\allowbreak isolate-\allowbreak WGS & \href{https://doi.org/10.1038/s41467-019-08504-7}{doi:\allowbreak 10.\allowbreak 1038/\allowbreak s41467-\allowbreak 019-\allowbreak 08504-\allowbreak 7} \\
\addlinespace[5pt]
123 & \textbf{Genome Medicine 2024 (1,877-strain genome and MIC set)}\,\textsuperscript{\textbf{B}}\,\textsuperscript{\textbf{M}} & \org{P. aeruginosa} & Tolerance / robustness & 1,877 strains plus a 102-\allowbreak strain validation set & antibiotic panel not read & $2.0\times 10^{3}$ & $2.0\times 10^{3}$ & S & susceptibility calls, reclassified per a named standard & pseudomonad-\allowbreak isolate-\allowbreak WGS & \href{https://doi.org/10.1186/s13073-024-01346-z}{doi:\allowbreak 10.\allowbreak 1186/\allowbreak s13073-\allowbreak 024-\allowbreak 01346-\allowbreak z} \\
\addlinespace[5pt]
124 & \textbf{Penaranda 2025 (125-isolate diversity survey)}\,\textsuperscript{\textbf{M}} & \org{P. aeruginosa} & Natural variation & 125 newly sequenced isolates, 97 clinical and 28 environmental, plus PAO1, PA14 and PAK & 5 phenotype assays plus disc-\allowbreak diffusion against 8 antibiotics & 128 & $1.7\times 10^{3}$ & I & pyocyanin, motility, biofilm, cytotoxicity at 4 h and 24 h, 8-\allowbreak drug susceptibility & pseudomonad-\allowbreak isolate-\allowbreak WGS & \href{https://doi.org/10.1101/2025.09.22.677895}{doi:\allowbreak 10.\allowbreak 1101/\allowbreak 2025.\allowbreak 09.\allowbreak 22.\allowbreak 677895} \\
\addlinespace[5pt]
125 & \textbf{Khaledi 2020 (1,656-record agar-dilution MIC panel)} & \org{P. aeruginosa} & Tolerance / robustness & 414 clinical isolates & tobramycin, ceftazidime, ciprofloxacin, meropenem & $1.7\times 10^{3}$$\dagger$ & $1.7\times 10^{3}$ & I & minimum inhibitory concentration, agar dilution per CLSI 2018 & pseudomonad-\allowbreak isolate-\allowbreak WGS & \href{https://doi.org/10.15252/emmm.201910264}{doi:\allowbreak 10.\allowbreak 15252/\allowbreak emmm.\allowbreak 2019102\allowbreak 64} \\
\addlinespace[5pt]
126 & \textbf{Poulsen 2019 (nine-strain Tn-seq core essentiality)}\,\textsuperscript{\textbf{M}} & \org{P. aeruginosa} & Essentiality & 321-\allowbreak gene core set across nine strains & five media & 321 & 321 & S & essentiality call, three calling methods & PA14+\allowbreak transposon & \href{https://www.biorxiv.org/content/10.1101/396689v2.full.pdf}{biorxiv.\allowbreak org/\allowbreak content\allowbreak /\allowbreak 10.\allowbreak 1101/\allowbreak 396689v\allowbreak 2.\allowbreak full.\allowbreak pdf} \\
\addlinespace[5pt]
127 & \textbf{Carrera 2014 EcoMAC and EcoPhe compendium}\,\textsuperscript{\textbf{M}} & \org{E. coli} & Transcriptome & 31 strains including cataloged knockouts, with 141 perturbed genes across the compendium & more than 15 media, as the source states it & $2.2\times 10^{3}$ & $9.2\times 10^{6}$ & I & transcript abundance, with growth for a subset & ecoli-\allowbreak KO & \href{https://doi.org/10.15252/msb.20145108}{doi:\allowbreak 10.\allowbreak 15252/\allowbreak msb.\allowbreak 2014510\allowbreak 8} \\
\addlinespace[5pt]
128 & \textbf{PAO1 arrayed transposon libraries (Jacobs 2003, Held 2012)}\,\textsuperscript{\textbf{M}} & \org{P. aeruginosa} & Modality / backbone & 42,240 archived mutants, reduced to a sequence-\allowbreak verified two-\allowbreak allele set of 9,437 & none; a genotype resource & $9.4\times 10^{3}$ & $9.4\times 10^{3}$ & V & insertion position per mutant, sequence-\allowbreak verified & PAO1+\allowbreak transposon & \href{https://doi.org/10.1128/JB.01479-12}{doi:\allowbreak 10.\allowbreak 1128/\allowbreak JB.\allowbreak 01479-\allowbreak 12} \\
\addlinespace[5pt]
129 & \textbf{KT2440 pooled transposon library (Putida\_ML5 metadata)} & \org{P. putida} & Modality / backbone & 4,778 non-\allowbreak essential genes, about 100,000 unique insertion mutants in one pool & LB plus kanamycin, time-\allowbreak zero cultures & $4.8\times 10^{3}$ & $4.8\times 10^{3}$ & S & insertion position and barcode per mutant & KT2440+\allowbreak transposon & \href{https://nrel.gov/docs/fy22osti/83183.pdf}{nrel.\allowbreak gov/\allowbreak docs/\allowbreak fy22ost\allowbreak i/\allowbreak 83183.\allowbreak pdf} \\
\addlinespace[5pt]
130 & \textbf{Gao 2020 (KT2440 targeted SRM proteomics)} & \org{P. putida} & Proteome & KT2440 wild type & 8 growth conditions, 30 biological samples & 132 & $1.1\times 10^{3}$ & I & absolute protein quantification, with units & KT2440-\allowbreak reference & \href{https://doi.org/10.3389/fbioe.2020.603488}{doi:\allowbreak 10.\allowbreak 3389/\allowbreak fbioe.\allowbreak 2020.\allowbreak 603488} \\
\addlinespace[5pt]
131 & \textbf{BACTOME (101 clinical genomes with eight phenotypes)}\,\textsuperscript{\textbf{R}} & \org{P. aeruginosa} & Natural variation & 99 clinical isolates plus the two reference strains; 18,319 pangenome genes; 96 transcriptomes & LB at OD 2.0 & 101 & 808 & I & Galleria virulence, biofilm, colony morphology, resistance to 5 drugs & pseudomonad-\allowbreak isolate-\allowbreak WGS & \href{https://doi.org/10.1093/nar/gky895}{doi:\allowbreak 10.\allowbreak 1093/\allowbreak nar/\allowbreak gky895} \\
\addlinespace[5pt]
132 & \textbf{Kohlstedt 2019 (13C metabolic flux analysis, two hosts)}\,\textsuperscript{\textbf{M}} & \org{P. putida} & Metabolome / flux & KT2440 and PAO1, both wild type & glucose & 534 & 534 & S & intracellular flux distribution from mass isotopomers & KT2440-\allowbreak reference & \href{https://doi.org/10.1016/j.ymben.2019.01.008}{doi:\allowbreak 10.\allowbreak 1016/\allowbreak j.\allowbreak ymben.\allowbreak 2019.\allowbreak 01.\allowbreak 008} \\
\addlinespace[5pt]
133 & \textbf{Turner 2015 (PA14 and PAO1 essentiality across four arms)} & \org{P. aeruginosa} & Essentiality & 352 general plus 199 condition-\allowbreak specific essential genes; 336 of 5,606 PAO1 genes & PA14 in brain-\allowbreak heart-\allowbreak infusion agar and sputum; PAO1 in the agar and in MOPS-\allowbreak succinate & 352 & 352 & S & essentiality call, general and condition-\allowbreak specific & PA14+\allowbreak transposon & \href{https://pmc.ncbi.nlm.nih.gov/articles/PMC4386324}{PMC43863\allowbreak 24} \\
\addlinespace[5pt]
134 & \textbf{Liberati 2006 (PA14 non-redundant transposon library and essentiality)}\,\textsuperscript{\textbf{M}} & \org{P. aeruginosa} & Essentiality & arrayed non-\allowbreak redundant mutant set; 335 genes called essential & LB & 335 & 335 & V & essentiality call & PA14+\allowbreak transposon & \href{https://doi.org/10.1073/pnas.0511100103}{doi:\allowbreak 10.\allowbreak 1073/\allowbreak pnas.\allowbreak 0511100\allowbreak 103} \\
\addlinespace[5pt]
135 & \textbf{GOLD\_115 and GOLD\_84 essential-gene sets}\,\textsuperscript{\textbf{M}} & \org{P. aeruginosa} & Essentiality & 115 PA14-\allowbreak in-\allowbreak LB plus 84 cross-\allowbreak strain gene calls & GOLD\_115 is PA14 in LB; GOLD\_84 is cross-\allowbreak strain and cross-\allowbreak medium & 199 & 199 & V & essentiality call, scoped by strain, medium and calling method & PA14+\allowbreak transposon & \href{https://doi.org/10.1371/journal.pcbi.1013945}{doi:\allowbreak 10.\allowbreak 1371/\allowbreak journal\allowbreak .\allowbreak pcbi.\allowbreak 1013945} \\
\addlinespace[5pt]
136 & \textbf{Werner 2023 (beta-ketoadipate from lignin)}\,\textsuperscript{\textbf{M}} & \org{P. putida} & Production campaign & 10 named engineered genotypes plus wild type, with named deletions and chromosomal insertions & shake flask, bioreactor on 3 feedstocks, two fed-\allowbreak batch modes, pretreatment liquor, corn stover & 30$\ddagger$ & 90 & V & titer, rate and molar yield & engineered-\allowbreak chassis & \href{https://doi.org/10.1126/sciadv.adj0053}{doi:\allowbreak 10.\allowbreak 1126/\allowbreak sciadv.\allowbreak adj0053} \\
\addlinespace[5pt]
137 & \textbf{Wilkes 2026 (muconate multi-omics with fluxomics)}\,\textsuperscript{\textbf{M}} & \org{P. putida} & Metabolome / flux & KT2440 wild type and the muconate producer CJ781 & glucose and 4-\allowbreak hydroxybenzoate & 4$\dagger$ & 12 & V & proteome, metabolome and 13C flux per strain and condition & engineered-\allowbreak chassis & \href{https://doi.org/10.1016/j.ymben.2026.102573}{doi:\allowbreak 10.\allowbreak 1016/\allowbreak j.\allowbreak ymben.\allowbreak 2026.\allowbreak 102573} \\
\addlinespace[5pt]
138 & \textbf{GSE213139 (low-dose ionizing radiation)} & \org{P. putida} & Transcriptome & not reported & radiation doses, count not reported & 30 & $1.7\times 10^{5}$ & V & RNA-\allowbreak seq expression & KT2440-\allowbreak reference & \href{https://www.ncbi.nlm.nih.gov/geo/query/acc.cgi?acc=GSE213139}{ncbi.\allowbreak nlm.\allowbreak nih.\allowbreak gov/\allowbreak geo/\allowbreak query/\allowbreak acc.\allowbreak cgi?\allowbreak acc=\allowbreak GSE2131\allowbreak 39} \\
\addlinespace[5pt]
139 & \textbf{GSE131229 (non-antibiotic pharmaceutical stress)} & \org{P. putida} & Transcriptome & not reported & pharmaceutical exposures, count not reported & 30 & $1.7\times 10^{5}$ & V & RNA-\allowbreak seq expression & KT2440-\allowbreak reference & \href{https://www.ncbi.nlm.nih.gov/geo/query/acc.cgi?acc=GSE131229}{ncbi.\allowbreak nlm.\allowbreak nih.\allowbreak gov/\allowbreak geo/\allowbreak query/\allowbreak acc.\allowbreak cgi?\allowbreak acc=\allowbreak GSE1312\allowbreak 29} \\
\addlinespace[5pt]
140 & \textbf{GSE248909 (micro and nanoplastic stress)} & \org{P. putida} & Transcriptome & not reported & plastic particle exposures, count not reported & 28 & $1.6\times 10^{5}$ & V & RNA-\allowbreak seq expression & KT2440-\allowbreak reference & \href{https://www.ncbi.nlm.nih.gov/geo/query/acc.cgi?acc=GSE248909}{ncbi.\allowbreak nlm.\allowbreak nih.\allowbreak gov/\allowbreak geo/\allowbreak query/\allowbreak acc.\allowbreak cgi?\allowbreak acc=\allowbreak GSE2489\allowbreak 09} \\
\addlinespace[5pt]
141 & \textbf{GSE193493 (putidaPRECISE321 second deposit)} & \org{P. putida} & Transcriptome & KT2440 and engineered strains; genotype count not reported & not reported & 22 & $1.2\times 10^{5}$ & V & RNA-\allowbreak seq expression & engineered-\allowbreak chassis & \href{https://www.ncbi.nlm.nih.gov/geo/query/acc.cgi?acc=GSE193493}{ncbi.\allowbreak nlm.\allowbreak nih.\allowbreak gov/\allowbreak geo/\allowbreak query/\allowbreak acc.\allowbreak cgi?\allowbreak acc=\allowbreak GSE1934\allowbreak 93} \\
\addlinespace[5pt]
142 & \textbf{GSE221044 (kaiC deletion)} & \org{P. putida} & Transcriptome & wild type and a kaiC deletion & not reported & 22 & $1.2\times 10^{5}$ & V & RNA-\allowbreak seq expression & KT2440-\allowbreak KO & \href{https://www.ncbi.nlm.nih.gov/geo/query/acc.cgi?acc=GSE221044}{ncbi.\allowbreak nlm.\allowbreak nih.\allowbreak gov/\allowbreak geo/\allowbreak query/\allowbreak acc.\allowbreak cgi?\allowbreak acc=\allowbreak GSE2210\allowbreak 44} \\
\addlinespace[5pt]
143 & \textbf{GSE144045 (solvent-tolerance evolution in P. putida S12)} & \org{P. putida} & Transcriptome & evolved S12 clones; count not reported & solvent challenge, count not reported & 18 & $1.0\times 10^{5}$ & V & RNA-\allowbreak seq expression & pseudomonad-\allowbreak reference & \href{https://www.ncbi.nlm.nih.gov/geo/query/acc.cgi?acc=GSE144045}{ncbi.\allowbreak nlm.\allowbreak nih.\allowbreak gov/\allowbreak geo/\allowbreak query/\allowbreak acc.\allowbreak cgi?\allowbreak acc=\allowbreak GSE1440\allowbreak 45} \\
\addlinespace[5pt]
144 & \textbf{Jessop-Fabre 2019 (Crabtree-negative hydroxy-acid strain transcriptomes)} & \org{S. cerevisiae} & Transcriptome & 5 hydroxy-\allowbreak acid-\allowbreak producing strains in a TAM pdc-\allowbreak null background & one growth condition, three replicates & 15$\dagger$ & $9.9\times 10^{4}$ & S & transcript abundance & S288C+\allowbreak designed-\allowbreak edit & \href{https://doi.org/10.1002/biot.201900013}{doi:\allowbreak 10.\allowbreak 1002/\allowbreak biot.\allowbreak 2019000\allowbreak 13} \\
\addlinespace[5pt]
145 & \textbf{GSE275525 (HexR regulatory landscape)} & \org{P. putida} & Transcriptome & not reported & not reported & 16 & $8.9\times 10^{4}$ & V & RNA-\allowbreak seq expression & KT2440-\allowbreak KO & \href{https://www.ncbi.nlm.nih.gov/geo/query/acc.cgi?acc=GSE275525}{ncbi.\allowbreak nlm.\allowbreak nih.\allowbreak gov/\allowbreak geo/\allowbreak query/\allowbreak acc.\allowbreak cgi?\allowbreak acc=\allowbreak GSE2755\allowbreak 25} \\
\addlinespace[5pt]
146 & \textbf{GSE309641 (fluoride adaptive evolution)} & \org{P. putida} & Transcriptome & evolved clones; count not reported & fluoride, count not reported & 16 & $8.9\times 10^{4}$ & V & RNA-\allowbreak seq expression & evolved-\allowbreak WGS & \href{https://www.ncbi.nlm.nih.gov/geo/query/acc.cgi?acc=GSE309641}{ncbi.\allowbreak nlm.\allowbreak nih.\allowbreak gov/\allowbreak geo/\allowbreak query/\allowbreak acc.\allowbreak cgi?\allowbreak acc=\allowbreak GSE3096\allowbreak 41} \\
\addlinespace[5pt]
147 & \textbf{GSE155767 (evolved KT2440 on non-native sugars)} & \org{P. putida} & Transcriptome & evolved clones; count not reported & non-\allowbreak native sugars, count not reported & 14 & $7.8\times 10^{4}$ & V & RNA-\allowbreak seq expression & evolved-\allowbreak WGS & \href{https://www.ncbi.nlm.nih.gov/geo/query/acc.cgi?acc=GSE155767}{ncbi.\allowbreak nlm.\allowbreak nih.\allowbreak gov/\allowbreak geo/\allowbreak query/\allowbreak acc.\allowbreak cgi?\allowbreak acc=\allowbreak GSE1557\allowbreak 67} \\
\addlinespace[5pt]
148 & \textbf{GSE147420 (ammonium against nitrate)} & \org{P. putida} & Transcriptome & not reported & ammonium and nitrate & 14 & $7.8\times 10^{4}$ & V & RNA-\allowbreak seq expression & KT2440-\allowbreak reference & \href{https://www.ncbi.nlm.nih.gov/geo/query/acc.cgi?acc=GSE147420}{ncbi.\allowbreak nlm.\allowbreak nih.\allowbreak gov/\allowbreak geo/\allowbreak query/\allowbreak acc.\allowbreak cgi?\allowbreak acc=\allowbreak GSE1474\allowbreak 20} \\
\addlinespace[5pt]
149 & \textbf{GSE229763 (one-carbon metabolism)} & \org{P. putida} & Transcriptome & not reported & one-\allowbreak carbon substrates, count not reported & 12 & $6.7\times 10^{4}$ & V & RNA-\allowbreak seq expression & engineered-\allowbreak chassis & \href{https://www.ncbi.nlm.nih.gov/geo/query/acc.cgi?acc=GSE229763}{ncbi.\allowbreak nlm.\allowbreak nih.\allowbreak gov/\allowbreak geo/\allowbreak query/\allowbreak acc.\allowbreak cgi?\allowbreak acc=\allowbreak GSE2297\allowbreak 63} \\
\addlinespace[5pt]
150 & \textbf{KT2440 aromatics RB-TnSeq (2019 preprint)} & \org{P. putida} & Transposon fitness & not reported & coumarate plus seven related aromatics & 8 & $4.5\times 10^{4}$ & S & RB-\allowbreak TnSeq gene fitness & KT2440+\allowbreak transposon & \href{https://www.biorxiv.org/content/10.1101/753889v1.full-text}{biorxiv.\allowbreak org/\allowbreak content\allowbreak /\allowbreak 10.\allowbreak 1101/\allowbreak 753889v\allowbreak 1.\allowbreak full-\allowbreak text} \\
\addlinespace[5pt]
151 & \textbf{PMkbase v1.0 (phenotype-microarray database)}\,\textsuperscript{\textbf{M}} & \org{P. putida} & Natural variation & strain list not confirmed; whether any pseudomonad is included is unconfirmed & carbon substrates, count not reported per strain & $4.2\times 10^{4}$ & $4.2\times 10^{4}$ & S & phenotype-\allowbreak microarray carbon utilization & pseudomonad-\allowbreak reference & \href{https://www.ncbi.nlm.nih.gov/pmc/articles/PMC12584640/}{PMC12584\allowbreak 640} \\
\addlinespace[5pt]
152 & \textbf{IPCD (International P. aeruginosa Consortium Database)}\,\textsuperscript{\textbf{R}} & \org{P. aeruginosa} & Natural variation & 3,038 isolates with draft genomes for all of them & phenotype fields not enumerated & $3.0\times 10^{3}$ & $3.0\times 10^{3}$ & V & linked phenotype, genotype and clinical data; fields not enumerated & pseudomonad-\allowbreak isolate-\allowbreak WGS & \href{https://doi.org/10.3389/fmicb.2015.01036}{doi:\allowbreak 10.\allowbreak 3389/\allowbreak fmicb.\allowbreak 2015.\allowbreak 01036} \\
\addlinespace[5pt]
153 & \textbf{Pirnay 2009 (328-isolate population survey)}\,\textsuperscript{\textbf{B}} & \org{P. aeruginosa} & Natural variation & 328 unrelated isolates from 69 localities in 30 countries & 8 parameters per isolate, type not resolved & 328 & $2.6\times 10^{3}$ & V & 8 parameters per isolate; whether they are measurements or typing markers is unresolved & pseudomonad-\allowbreak isolate-\allowbreak WGS & \href{https://doi.org/10.1371/journal.pone.0007740}{doi:\allowbreak 10.\allowbreak 1371/\allowbreak journal\allowbreak .\allowbreak pone.\allowbreak 0007740} \\
\addlinespace[5pt]
154 & \textbf{Kim 2013 (472 knockouts, five fermentation products)}\,\textsuperscript{\textbf{M}} & \org{E. coli} & Fitness / chemical genomics & 472 single-\allowbreak gene knockouts from a cataloged collection & anaerobic fermentation & 472$\dagger$ & $2.4\times 10^{3}$ & S & secreted lactate, succinate, acetate, formate and ethanol & ecoli-\allowbreak KO & \href{https://doi.org/10.1016/j.ymben.2013.03.004}{doi:\allowbreak 10.\allowbreak 1016/\allowbreak j.\allowbreak ymben.\allowbreak 2013.\allowbreak 03.\allowbreak 004} \\
\addlinespace[5pt]
155 & \textbf{n-butanol KT2440 proteome (2015)} & \org{P. putida} & Proteome & KT2440 wild type & n-\allowbreak butanol challenge & $1.5\times 10^{3}$ & $1.5\times 10^{3}$ & S & relative protein abundance & KT2440-\allowbreak reference & \href{https://biochemie.uni-hohenheim.de/en/organization/publication/analysis-of-the-molecular-response-of-pseudomonas-putida-kt2440-to-the-next-generation-biofuel-n-butanol}{biochemi\allowbreak e.\allowbreak uni-\allowbreak hohenhe\allowbreak im.\allowbreak de/\allowbreak en/\allowbreak organiz\allowbreak ation/\allowbreak publica\allowbreak tion/\allowbreak analysi\allowbreak s-\allowbreak of-\allowbreak the-\allowbreak molecul\allowbreak ar-\allowbreak respons\allowbreak e-\allowbreak of-\allowbreak pseudom\allowbreak onas-\allowbreak putida-\allowbreak kt2440-\allowbreak to-\allowbreak the-\allowbreak next-\allowbreak generat\allowbreak ion-\allowbreak biofuel\allowbreak -\allowbreak n-\allowbreak butanol} \\
\addlinespace[5pt]
156 & \textbf{Li 2025 (xylose deregulation panel with 3-HP readout)} & \org{S. cerevisiae} & Production campaign & 177 labeled strains off the IMX581 CEN.PK113-\allowbreak 5D lineage, each a stated module of central-\allowbreak carbon deregulations & xylose and glucose & $1.2\times 10^{3}$$\dagger$ & $1.2\times 10^{3}$ & I & 3-\allowbreak HP titer with xylose and glucose consumption & S288C+\allowbreak designed-\allowbreak edit & \href{https://doi.org/10.1038/s41467-025-59966-x}{doi:\allowbreak 10.\allowbreak 1038/\allowbreak s41467-\allowbreak 025-\allowbreak 59966-\allowbreak x} \\
\addlinespace[5pt]
157 & \textbf{Lazar 2014 (antibiotic resistance evolution with b-numbered mutations)} & \org{E. coli} & Tolerance / robustness & 63 sequenced evolved lines, from 120 phenotyped populations & 12 antibiotics & 63 & 756 & V & cross-\allowbreak resistance and collateral sensitivity & evolved-\allowbreak WGS & \href{https://doi.org/10.1038/ncomms5352}{doi:\allowbreak 10.\allowbreak 1038/\allowbreak ncomms5\allowbreak 352} \\
\addlinespace[5pt]
158 & \textbf{Skurnik 2013 (PA14 essentiality)}\,\textsuperscript{\textbf{M}} & \org{P. aeruginosa} & Essentiality & 634 genes called essential & LB & 634 & 634 & V & essentiality call & PA14+\allowbreak transposon & \href{https://doi.org/10.1371/journal.pcbi.1013945}{doi:\allowbreak 10.\allowbreak 1371/\allowbreak journal\allowbreak .\allowbreak pcbi.\allowbreak 1013945} \\
\addlinespace[5pt]
159 & \textbf{Valle 2015 (150 knockouts, hydrogen and ethanol on glycerol)} & \org{E. coli} & Fitness / chemical genomics & 150 single-\allowbreak gene knockouts drawn from both the Keio and the Yale CGSC collections, each in triplicate & glycerol-\allowbreak based medium & 150$\dagger$ & 450 & V & hydrogen, ethanol and glycerol consumption & ecoli-\allowbreak KO & \href{https://doi.org/10.1186/s12934-015-0285-6}{doi:\allowbreak 10.\allowbreak 1186/\allowbreak s12934-\allowbreak 015-\allowbreak 0285-\allowbreak 6} \\
\addlinespace[5pt]
160 & \textbf{Lee 2015 (PAO1 essentiality across three media)}\,\textsuperscript{\textbf{M}} & \org{P. aeruginosa} & Essentiality & 352 genes called essential & three media & 352 & 352 & V & essentiality call & PAO1+\allowbreak transposon & \href{https://doi.org/10.1371/journal.pcbi.1013945}{doi:\allowbreak 10.\allowbreak 1371/\allowbreak journal\allowbreak .\allowbreak pcbi.\allowbreak 1013945} \\
\addlinespace[5pt]
161 & \textbf{P. putida S12 phenotype microarray (190 carbon sources)} & \org{P. putida} & Natural variation & S12, not the reference strain & 190 carbon sources & 190 & 190 & S & phenotype-\allowbreak microarray carbon utilization & pseudomonad-\allowbreak reference & \href{https://doi.org/10.1186/s12864-023-09940-y}{doi:\allowbreak 10.\allowbreak 1186/\allowbreak s12864-\allowbreak 023-\allowbreak 09940-\allowbreak y} \\
\addlinespace[5pt]
162 & \textbf{Qin 2024 (CO2 fixation route to 3-HP)} & \org{S. cerevisiae} & Production campaign & 71 strains, QLW1 through QLW71 with no gaps, off CEN.PK113-\allowbreak 11C, each with a cumulative genotype string & glucose with carbon-\allowbreak fixation modules & 71 & 142 & V & titer and carbon yield & S288C+\allowbreak designed-\allowbreak edit & \href{https://doi.org/10.1038/s41467-024-45557-9}{doi:\allowbreak 10.\allowbreak 1038/\allowbreak s41467-\allowbreak 024-\allowbreak 45557-\allowbreak 9} \\
\addlinespace[5pt]
163 & \textbf{Chandler 2023 (128 functionally phenotyped isolates)}\,\textsuperscript{\textbf{M}} & \org{P. aeruginosa} & Natural variation & 128 isolates sequenced with functional phenotyping & not reported & 128 & 128 & S & functional phenotypes, not enumerated & pseudomonad-\allowbreak isolate-\allowbreak WGS & \href{https://doi.org/10.1128/spectrum.01556-23}{doi:\allowbreak 10.\allowbreak 1128/\allowbreak spectru\allowbreak m.\allowbreak 01556-\allowbreak 23} \\
\addlinespace[5pt]
164 & \textbf{Breidenstein 2008 (ciprofloxacin resistome)} & \org{P. aeruginosa} & Tolerance / robustness & 114 mutants with altered susceptibility, 35 more and 79 less susceptible & ciprofloxacin & 114 & 114 & V & altered antimicrobial susceptibility & PA14+\allowbreak transposon & \href{https://doi.org/10.1128/AAC.00222-08}{doi:\allowbreak 10.\allowbreak 1128/\allowbreak AAC.\allowbreak 00222-\allowbreak 08} \\
\addlinespace[5pt]
165 & \textbf{Remaining P. putida GEO expression series}\,\textsuperscript{\textbf{A}} & \org{P. putida} & Transcriptome & varies per series & varies per series & 78$\ddagger$ & 78 & V & RNA-\allowbreak seq or array expression & KT2440-\allowbreak reference & \href{https://www.ncbi.nlm.nih.gov/gds/?term=\%22Pseudomonas+putida\%22\%5BOrganism\%5D}{ncbi.\allowbreak nlm.\allowbreak nih.\allowbreak gov/\allowbreak gds/\allowbreak ?\allowbreak term=\allowbreak \%\allowbreak 22Pseud\allowbreak omonas+\allowbreak putida\%\allowbreak 22\%\allowbreak 5BOrgan\allowbreak ism\%\allowbreak 5D} \\
\addlinespace[5pt]
166 & \textbf{KT2440 mini-Tn5 knockout collection (2010)} & \org{P. putida} & Essentiality & 24 mutants in genes proposed conditionally essential plus 24 newly implicated genes, against 68 predicted essential & minimal glucose medium & 68 & 68 & S & conditional essentiality call & KT2440+\allowbreak transposon & \href{https://pmc.ncbi.nlm.nih.gov/articles/PMC5814926}{PMC58149\allowbreak 26} \\
\addlinespace[5pt]
167 & \textbf{Xu 2023 (heterologous carbon-fixation route, 3-HP and mesaconate)} & \org{S. cerevisiae} & Production campaign & 59 tabulated strains off CEN.PK2-\allowbreak 1D, with the explicit genotype of each in Supplementary Table S2 & glucose & 42 & 42 & I & 3-\allowbreak HP titer, with mesaconate titer on a subset & S288C+\allowbreak designed-\allowbreak edit & \href{https://doi.org/10.1186/s40643-023-00652-5}{doi:\allowbreak 10.\allowbreak 1186/\allowbreak s40643-\allowbreak 023-\allowbreak 00652-\allowbreak 5} \\
\addlinespace[5pt]
168 & \textbf{Nunez-Garcia 2024 (39 isolates over ten years)}\,\textsuperscript{\textbf{M}} & \org{P. aeruginosa} & Natural variation & 39 isolates collected over ten years & growth-\allowbreak rate assay & 39 & 39 & S & growth rate & pseudomonad-\allowbreak isolate-\allowbreak WGS & \href{https://pmc.ncbi.nlm.nih.gov/articles/PMC11619361/}{PMC11619\allowbreak 361} \\
\addlinespace[5pt]
169 & \textbf{Borodina 2015 (beta-alanine route strain series)}\,\textsuperscript{\textbf{R}} & \org{S. cerevisiae} & Production campaign & 35 strains carrying the beta-\allowbreak alanine route, with at least 38 once Supplementary Figure S3 is counted & glucose & 35 & 35 & I & titer & S288C+\allowbreak designed-\allowbreak edit & \href{https://doi.org/10.1016/j.ymben.2014.10.003}{doi:\allowbreak 10.\allowbreak 1016/\allowbreak j.\allowbreak ymben.\allowbreak 2014.\allowbreak 10.\allowbreak 003} \\
\addlinespace[5pt]
170 & \textbf{Remaining P. putida PRIDE proteomics projects}\,\textsuperscript{\textbf{A}} & \org{P. putida} & Proteome & varies per project; three projects hold several organisms and must be split & carbon sources, fluoride, phage infection, pathway engineering & 26$\ddagger$ & 26 & S & relative protein abundance & KT2440-\allowbreak reference & \href{https://www.ebi.ac.uk/pride/ws/archive/v3/search/projects?keyword=Pseudomonas\%20putida}{ebi.\allowbreak ac.\allowbreak uk/\allowbreak pride/\allowbreak ws/\allowbreak archive\allowbreak /\allowbreak v3/\allowbreak search/\allowbreak project\allowbreak s?\allowbreak keyword\allowbreak =\allowbreak Pseudom\allowbreak onas\%\allowbreak 20putid\allowbreak a} \\
\addlinespace[5pt]
171 & \textbf{Zhang 2023 (mitochondrial 3-HP compartment panel)} & \org{S. cerevisiae} & Production campaign & 38 strains off the CEN.PK113-\allowbreak 5D lineage, differing by which pathway steps are targeted to the mitochondrion & glucose & 22 & 22 & I & titer & S288C+\allowbreak designed-\allowbreak edit & \href{https://doi.org/10.1186/s13068-023-02309-z}{doi:\allowbreak 10.\allowbreak 1186/\allowbreak s13068-\allowbreak 023-\allowbreak 02309-\allowbreak z} \\
\addlinespace[5pt]
172 & \textbf{Chen 2018 (Acc1 phosphorylation-site mutant panel)} & \org{S. cerevisiae} & Deep mutational scan & 13 ACC1 phosphorylation-\allowbreak site substitutions off CEN.PK102-\allowbreak 5B & glucose with a malonyl-\allowbreak CoA biosensor readout & 13 & 13 & S & biosensor output per substitution & S288C+\allowbreak designed-\allowbreak edit & \href{https://doi.org/10.3389/fmicb.2018.00047}{doi:\allowbreak 10.\allowbreak 3389/\allowbreak fmicb.\allowbreak 2018.\allowbreak 00047} \\
\addlinespace[5pt]
173 & \textbf{Qiu 2022 (biosensor-coupled in vivo mutagenesis, malonyl-CoA)} & \org{S. cerevisiae} & Production campaign & 12 mutants recovered from biosensor-\allowbreak coupled in vivo mutagenesis & glucose & 12 & 12 & S & product titer with paired transcriptome and genome sequence & evolved-\allowbreak WGS & \href{https://doi.org/10.1128/msystems.01366-21}{doi:\allowbreak 10.\allowbreak 1128/\allowbreak msystem\allowbreak s.\allowbreak 01366-\allowbreak 21} \\
\addlinespace[5pt]
174 & \textbf{Kildegaard 2015 (3-HP from glucose and xylose)}\,\textsuperscript{\textbf{M}} & \org{S. cerevisiae} & Production campaign & 3 strains carrying the beta-\allowbreak alanine route with xylose utilization added & glucose and xylose & 6$\dagger$ & 12 & S & titer and yield & S288C+\allowbreak designed-\allowbreak edit & \href{https://doi.org/10.1016/j.meteno.2015.10.001}{doi:\allowbreak 10.\allowbreak 1016/\allowbreak j.\allowbreak meteno.\allowbreak 2015.\allowbreak 10.\allowbreak 001} \\
\addlinespace[5pt]
175 & \textbf{Shi 2014 (Snf1-dependent Acc1 regulation abolished)} & \org{S. cerevisiae} & Production campaign & 4 strains carrying Acc1 phosphosite substitutions off CEN.PK113-\allowbreak 11C & glucose & 4 & 4 & I & titer & S288C+\allowbreak designed-\allowbreak edit & \href{https://doi.org/10.1128/mbio.01130-14}{doi:\allowbreak 10.\allowbreak 1128/\allowbreak mbio.\allowbreak 01130-\allowbreak 14} \\
\addlinespace[5pt]
176 & \textbf{Kohlstedt 2022 (genome-reduced MA-10 muconate for PET)}\,\textsuperscript{\textbf{M}} & \org{P. putida} & Production campaign & genome-\allowbreak reduced MA-\allowbreak 10 lineage; distinct-\allowbreak genotype count unconfirmed & catechol plus glucose & 3$\ddagger$ & 3 & S & titer and productivity & engineered-\allowbreak chassis & \href{https://doi.org/10.1016/j.ymben.2022.05.001}{doi:\allowbreak 10.\allowbreak 1016/\allowbreak j.\allowbreak ymben.\allowbreak 2022.\allowbreak 05.\allowbreak 001} \\
\addlinespace[5pt]
177 & \textbf{Kohlstedt 2018 (MA-6 muconate to nylon)}\,\textsuperscript{\textbf{M}} & \org{P. putida} & Production campaign & MA-\allowbreak 6 lineage; distinct-\allowbreak genotype count unconfirmed & catechol, fed-\allowbreak batch & 3$\ddagger$ & 3 & S & titer and productivity & engineered-\allowbreak chassis & \href{https://doi.org/10.1016/j.ymben.2018.03.003}{doi:\allowbreak 10.\allowbreak 1016/\allowbreak j.\allowbreak ymben.\allowbreak 2018.\allowbreak 03.\allowbreak 003} \\
\addlinespace[5pt]
178 & \textbf{PAO1 ppk1 and ppk2 Biolog panels} & \org{P. aeruginosa} & Natural variation & ppk1 and ppk2 mutants plus the parent & Biolog PM1 and PM2 carbon plus PM3 nitrogen & 3 & 3 & S & Biolog respiration per well & PAO1-\allowbreak reference & \href{https://doi.org/10.1186/s40659-015-0012-0}{doi:\allowbreak 10.\allowbreak 1186/\allowbreak s40659-\allowbreak 015-\allowbreak 0012-\allowbreak 0} \\
\addlinespace[5pt]
179 & \textbf{Kildegaard 2014 (3-HP tolerance evolution, resequenced clones)}\,\textsuperscript{\textbf{M}} & \org{S. cerevisiae} & Tolerance / robustness & 3 resequenced clones from 30 isolated across 3 independent lineages, with 8 mutated genes listed & 3-\allowbreak HP in the medium & 3 & 3 & I & 3-\allowbreak HP tolerance & evolved-\allowbreak WGS & \href{https://doi.org/10.1016/j.ymben.2014.09.004}{doi:\allowbreak 10.\allowbreak 1016/\allowbreak j.\allowbreak ymben.\allowbreak 2014.\allowbreak 09.\allowbreak 004} \\
\addlinespace[5pt]
180 & \textbf{Vardon 2015 (adipic acid from lignin)}\,\textsuperscript{\textbf{M}} & \org{P. putida} & Production campaign & not reported & a model lignin monomer & 2$\ddagger$ & 2 & S & titer over time & engineered-\allowbreak chassis & \href{https://doi.org/10.1039/C4EE03230F}{doi:\allowbreak 10.\allowbreak 1039/\allowbreak C4EE032\allowbreak 30F} \\
\addlinespace[5pt]
181 & \textbf{PXD074559 (JAG1 proteome, acetate against glucose)} & \org{P. putida} & Proteome & JAG1, a KT2440-\allowbreak background strain & acetate and glucose & 2 & 2 & S & relative protein abundance & engineered-\allowbreak chassis & \href{https://www.ebi.ac.uk/pride/archive/projects/PXD074559}{ebi.\allowbreak ac.\allowbreak uk/\allowbreak pride/\allowbreak archive\allowbreak /\allowbreak project\allowbreak s/\allowbreak PXD0745\allowbreak 59} \\
\addlinespace[5pt]
182 & \textbf{Batianis 2020 (expanded CRISPRi toolbox)} & \org{P. putida} & CRISPR library screen & knockdowns of pyrF and ftsZ & not reported & 2 & 2 & S & growth rate and cell morphology under knockdown & KT2440+\allowbreak guide & \href{https://www.biorxiv.org/content/10.1101/2020.03.04.975888v1}{biorxiv.\allowbreak org/\allowbreak content\allowbreak /\allowbreak 10.\allowbreak 1101/\allowbreak 2020.\allowbreak 03.\allowbreak 04.\allowbreak 975888v\allowbreak 1} \\
\addlinespace[5pt]
183 & \textbf{KT2440 lignin-stream tolerance RB-TnSeq (2023)} & \org{P. putida} & Transposon fitness & not reported & lignin-\allowbreak related compounds, count not reported & --$\ddagger$ & -- & S & RB-\allowbreak TnSeq fitness under stress & KT2440+\allowbreak transposon & \href{https://research-hub.nrel.gov/en/publications/rb-tnseq-identifies-genetic-targets-for-improved-tolerance-of-seu/}{research\allowbreak -\allowbreak hub.\allowbreak nrel.\allowbreak gov/\allowbreak en/\allowbreak publica\allowbreak tions/\allowbreak rb-\allowbreak tnseq-\allowbreak identif\allowbreak ies-\allowbreak genetic\allowbreak -\allowbreak targets\allowbreak -\allowbreak for-\allowbreak improve\allowbreak d-\allowbreak toleran\allowbreak ce-\allowbreak of-\allowbreak seu} \\
\addlinespace[5pt]
184 & \textbf{KT2440 metal-excess Tn-seq (2024)} & \org{P. putida} & Transposon fitness & not reported & metal excess, count not reported & --$\ddagger$ & -- & S & conditional essentiality under metal excess & KT2440+\allowbreak transposon & \href{https://doi.org/10.1101/2024.04.12.589247}{doi:\allowbreak 10.\allowbreak 1101/\allowbreak 2024.\allowbreak 04.\allowbreak 12.\allowbreak 589247} \\
\addlinespace[5pt]
185 & \textbf{Woodworth 2021 (syringyl lignin to pyrone dicarboxylic acid)}\,\textsuperscript{\textbf{M}} & \org{P. putida} & Production campaign & not reported & a lignin-\allowbreak monomer mixture & 1$\ddagger$ & 1 & S & molar conversion to product & engineered-\allowbreak chassis & \href{https://www.sciencedirect.com/science/article/abs/pii/S1096717621000276}{scienced\allowbreak irect.\allowbreak com/\allowbreak science\allowbreak /\allowbreak article\allowbreak /\allowbreak abs/\allowbreak pii/\allowbreak S109671\allowbreak 76210002\allowbreak 76} \\
\addlinespace[5pt]
186 & \textbf{PXD067010 (lignocellulosic sugars to an aviation-fuel precursor)} & \org{P. putida} & Proteome & not reported & lignocellulosic sugar feeds, count not reported & --$\ddagger$ & -- & S & relative protein abundance & engineered-\allowbreak chassis & \href{https://www.ebi.ac.uk/pride/archive/projects/PXD067010}{ebi.\allowbreak ac.\allowbreak uk/\allowbreak pride/\allowbreak archive\allowbreak /\allowbreak project\allowbreak s/\allowbreak PXD0670\allowbreak 10} \\
\addlinespace[5pt]
187 & \textbf{PXD054609 (isoprenol tolerance, evolution plus omics)} & \org{P. putida} & Proteome & evolved clones; count not reported & isoprenol challenge, count not reported & --$\ddagger$ & -- & S & relative protein abundance & evolved-\allowbreak WGS & \href{https://www.ebi.ac.uk/pride/archive/projects/PXD054609}{ebi.\allowbreak ac.\allowbreak uk/\allowbreak pride/\allowbreak archive\allowbreak /\allowbreak project\allowbreak s/\allowbreak PXD0546\allowbreak 09} \\
\addlinespace[5pt]
188 & \textbf{PXD045993 (polyhydroxyalkanoate biology in KT2440)} & \org{P. putida} & Proteome & not reported & not reported & --$\ddagger$ & -- & S & relative protein abundance & KT2440-\allowbreak reference & \href{https://doi.org/10.1016/j.mcpro.2024.100765}{doi:\allowbreak 10.\allowbreak 1016/\allowbreak j.\allowbreak mcpro.\allowbreak 2024.\allowbreak 100765} \\
\addlinespace[5pt]
189 & \textbf{QconCAT absolute quantification workflow (2023)} & \org{P. putida} & Proteome & KT2440 & not reported & --$\ddagger$ & -- & S & absolute protein quantification & KT2440-\allowbreak reference & \href{https://pmc.ncbi.nlm.nih.gov/articles/PMC9881405}{PMC98814\allowbreak 05} \\
\addlinespace[5pt]
190 & \textbf{MTBLS12468 (GAPDH homologs on glucose against ferulate)} & \org{P. putida} & Metabolome / flux & not reported & glucose and ferulate & --$\ddagger$ & -- & S & metabolite intensities & KT2440-\allowbreak reference & \href{https://www.ebi.ac.uk/metabolights/MTBLS12468}{ebi.\allowbreak ac.\allowbreak uk/\allowbreak metabol\allowbreak ights/\allowbreak MTBLS12\allowbreak 468} \\
\addlinespace[5pt]
191 & \textbf{MTBLS1715 (13C metabolomics under iron limitation)} & \org{P. putida} & Metabolome / flux & several pseudomonads including KT2440 and S12 & iron-\allowbreak replete and iron-\allowbreak limited & --$\ddagger$ & -- & S & labeled metabolite intensities & KT2440-\allowbreak reference & \href{https://www.ebi.ac.uk/metabolights/MTBLS1715}{ebi.\allowbreak ac.\allowbreak uk/\allowbreak metabol\allowbreak ights/\allowbreak MTBLS17\allowbreak 15} \\
\addlinespace[5pt]
192 & \textbf{MTBLS319 and MTBLS320 (DOT-T1E solvent and drug exposure)} & \org{P. putida} & Metabolome / flux & DOT-\allowbreak T1E, not the reference strain & toluene and propranolol & --$\ddagger$ & -- & S & metabolite intensities & pseudomonad-\allowbreak reference & \href{https://www.ebi.ac.uk/metabolights/MTBLS319}{ebi.\allowbreak ac.\allowbreak uk/\allowbreak metabol\allowbreak ights/\allowbreak MTBLS31\allowbreak 9} \\
\addlinespace[5pt]
193 & \textbf{Sigma-factor regulon reconstruction (2026)} & \org{P. putida} & Transcriptome & three sigma-\allowbreak factor backgrounds & not reported & --$\ddagger$ & -- & S & binding occupancy plus expression & KT2440-\allowbreak KO & \href{https://doi.org/10.1186/s13036-026-00769-z}{doi:\allowbreak 10.\allowbreak 1186/\allowbreak s13036-\allowbreak 026-\allowbreak 00769-\allowbreak z} \\
\addlinespace[5pt]
194 & \textbf{Fenster 2022 (KT2440 CRISPRi toolset with a product readout)} & \org{P. putida} & CRISPR library screen & guide constructs against named targets; count not reported & p-\allowbreak coumarate & --$\ddagger$ & -- & S & beta-\allowbreak ketoadipate production plus single-\allowbreak cell repression & engineered-\allowbreak chassis+\allowbreak guide & \href{https://doi.org/10.1016/j.mec.2022.e00204}{doi:\allowbreak 10.\allowbreak 1016/\allowbreak j.\allowbreak mec.\allowbreak 2022.\allowbreak e00204} \\
\addlinespace[5pt]
195 & \textbf{Rhizosphere CRISPRi (2024)} & \org{P. putida} & CRISPR library screen & essential-\allowbreak gene knockdowns; count not reported & plant rhizosphere & --$\ddagger$ & -- & S & cells per root & KT2440+\allowbreak guide & \href{https://pubmed.ncbi.nlm.nih.gov/39163848/}{PMID 391\allowbreak 63848} \\
\addlinespace[5pt]
196 & \textbf{Pediatric cystic fibrosis susceptibility panel (2025)}\,\textsuperscript{\textbf{M}} & \org{P. aeruginosa} & Tolerance / robustness & isolate count unconfirmed & EUCAST panel, drug list not read & --$\ddagger$ & -- & S & susceptibility per a named standard & pseudomonad-\allowbreak isolate-\allowbreak WGS & \href{https://doi.org/10.3390/antibiotics14030302}{doi:\allowbreak 10.\allowbreak 3390/\allowbreak antibio\allowbreak tics1403\allowbreak 0302} \\
\addlinespace[5pt]
197 & \textbf{PA14 non-redundant phage-resistance screen (2024)}\,\textsuperscript{\textbf{B}}\,\textsuperscript{\textbf{M}} & \org{P. aeruginosa} & Tolerance / robustness & arrayed PA14 mutants; count unconfirmed & phage challenge, count unconfirmed & --$\ddagger$ & -- & S & phage resistance per mutant & PA14+\allowbreak transposon & \href{https://www.ncbi.nlm.nih.gov/pmc/articles/PMC11241198/}{PMC11241\allowbreak 198} \\
\addlinespace[5pt]
198 & \textbf{OSTI barcoded P. fluorescens transposon libraries}\,\textsuperscript{\textbf{M}} & \org{P. fluorescens} & Transposon fitness & over 100,000 independent insertions; strain not stated & root colonization & --$\ddagger$ & -- & S & colonization fitness & pseudomonad-\allowbreak reference & \href{https://www.osti.gov/biblio/1236050}{osti.\allowbreak gov/\allowbreak biblio/\allowbreak 1236050} \\
\addlinespace[5pt]
199 & \textbf{Raman 2014 (multiplexed mutagenesis, naringenin and glucaric acid)} & \org{E. coli} & Production campaign & clones selected from a library the record gives as 10,000 members; how many carry a released phenotype was not established & one selection condition per product & --$\ddagger$ & -- & S & naringenin and glucaric acid production & evolved-\allowbreak WGS & \href{https://doi.org/10.1073/pnas.1409523111}{doi:\allowbreak 10.\allowbreak 1073/\allowbreak pnas.\allowbreak 1409523\allowbreak 111} \\
\addlinespace[5pt]
200 & \textbf{Tong 2021 (oxaloacetate route to 3-HP)}\,\textsuperscript{\textbf{M}} & \org{S. cerevisiae} & Production campaign & a strain series decarboxylating oxaloacetate through KdcA; counts not reachable & glucose & --$\ddagger$ & -- & S & titer and yield & S288C+\allowbreak designed-\allowbreak edit & \href{https://doi.org/10.1039/d0gc04431h}{doi:\allowbreak 10.\allowbreak 1039/\allowbreak d0gc044\allowbreak 31h} \\
\addlinespace[5pt]
\midrule \multicolumn{12}{@{}l}{\textbf{Long-run cut. Rows above are the 200 this wave adds; rows below are the ranked reserve.}}\\ \midrule
201 & \textbf{Liu 2026 (transcription-factor rewiring for 3-HP)}\,\textsuperscript{\textbf{M}} & \org{S. cerevisiae} & Production campaign & a transcription-\allowbreak factor mutagenesis series; counts not reachable behind the paywall & glucose & --$\ddagger$ & -- & S & titer & S288C+\allowbreak designed-\allowbreak edit & \href{https://doi.org/10.1016/j.ymben.2026.102527}{doi:\allowbreak 10.\allowbreak 1016/\allowbreak j.\allowbreak ymben.\allowbreak 2026.\allowbreak 102527} \\
\addlinespace[5pt]
\end{longtable}
\endgroup
\end{landscape}


\clearpage
%% GENERATED FILE -- do not hand-edit.
%% SOURCE: experiments/database/scripts/build_hosts_candidate_datasets_table.py
\begin{landscape}
\begingroup
\footnotesize
\setlength{\tabcolsep}{4pt}
\renewcommand{\arraystretch}{1.15}
\begin{longtable}{@{}r@{\hspace{4pt}} L{78mm} L{72mm} L{76mm}@{}}
\caption[]{Sources for the whole pool of Table~\ref{tab:candidates}, in the same order.
\emph{Why it is ordered here} is the band reason for a host-bootstrap, regulatory-DNA or
cross-host-fitness row, and the tier and measurement rule for a scale row. \emph{Where
the per-record data lives} is the deposit, with its retrieval class: a plain entry is
fetched programmatically, \emph{browser} needs a human because the host returns 403 or
renders in JavaScript, and \emph{ask} means no public copy exists and the request is in
Table~\ref{tab:asks}. Every link is clickable.}
\label{tab:sources-all}\\
\toprule
\textbf{\#} & \textbf{Dataset, citation and link} & \textbf{Why it is ordered here} & \textbf{Where the per-record data lives} \\
\midrule
\endfirsthead
\multicolumn{4}{@{}l}{\footnotesize\emph{Table~\ref{tab:sources-all}, continued}}\\
\toprule
\textbf{\#} & \textbf{Dataset, citation and link} & \textbf{Why it is ordered here} & \textbf{Where the per-record data lives} \\
\midrule
\endhead
\bottomrule
\endfoot
1 & \textbf{de Boer 2020 (GPRA sort-seq, 100M random promoters)} \org{S. cerevisiae}\newline de Boer CG, Vaishnav ED, Sadeh R, Abeyta EL, Friedman N, Regev A. Nat Biotechnol 2020;38:56-65.\newline \href{https://doi.org/10.1038/s41587-019-0315-8}{doi:\allowbreak 10.\allowbreak 1038/\allowbreak s41587-\allowbreak 019-\allowbreak 0315-\allowbreak 8} & Regulatory DNA band: one row carrying more measurements than every other row in this wave combined, on a schema that already serves its child, so the admission is a download decision rather than a design one. & GEO GSE1\allowbreak 04878 an\allowbreak d GSE104\allowbreak 903; Add\allowbreak gene 127\allowbreak 546; whi\allowbreak ch acces\allowbreak sion hol\allowbreak ds which\allowbreak  library\allowbreak  is unre\allowbreak solved a\allowbreak nd one p\allowbreak age chec\allowbreak k settle\allowbreak s it \\
\addlinespace[5pt]
2 & \textbf{Vaishnav 2022 deferred training libraries} \org{S. cerevisiae}\newline Vaishnav ED, de Boer CG, Molinet J, Yassour M, Fan L, Adiconis X, Thompson DA, Levin JZ, Cubillos FA, Regev A. Nature 2022;603:455-463.\newline \href{https://doi.org/10.1038/s41586-022-04506-6}{doi:\allowbreak 10.\allowbreak 1038/\allowbreak s41586-\allowbreak 022-\allowbreak 04506-\allowbreak 6} & Regulatory DNA band, and the single lowest-effort large win in the wave: every piece of machinery it needs is already written and serving its subset. & Zenodo 4\allowbreak 436477 (\allowbreak complex\_\allowbreak media\_\allowbreak trainin\allowbreak g\_\allowbreak data\_\allowbreak Glu.\allowbreak txt 3.\allowbreak 59 GB, \allowbreak defined\_\allowbreak media\_\allowbreak trainin\allowbreak g\_\allowbreak data\_\allowbreak SC\_\allowbreak Ura.\allowbreak txt 2.\allowbreak 47 GB);\allowbreak  GEO GSE\allowbreak 163045 a\allowbreak nd GSE16\allowbreak 3866 \\
\addlinespace[5pt]
3 & \textbf{DREAM 2022 random promoter challenge (Rafi 2024)} \org{S. cerevisiae}\newline Rafi AM, et al. A community effort to optimize sequence-based deep learning models of gene regulation. Nat Biotechnol 2024.\newline \href{https://www.ncbi.nlm.nih.gov/geo/query/acc.cgi?acc=GSE254493}{ncbi.\allowbreak nlm.\allowbreak nih.\allowbreak gov/\allowbreak geo/\allowbreak query/\allowbreak acc.\allowbreak cgi?\allowbreak acc=\allowbreak GSE2544\allowbreak 93} & Regulatory DNA band: the benchmark slice, which is what makes a model comparison possible at all, and it carries its own eight-category test design. & GEO GSE2\allowbreak 54493; Z\allowbreak enodo 10\allowbreak 633252 (\allowbreak v2, abou\allowbreak t 4.\allowbreak 0 GB), \allowbreak supersed\allowbreak ing Zeno\allowbreak do 73953\allowbreak 97 \\
\addlinespace[5pt]
4 & \textbf{Kotopka and Smolke 2020 (artificial promoter libraries)} \org{S. cerevisiae}\newline Kotopka BJ, Smolke CD. Model-driven generation of artificial yeast promoters. Nat Commun 2020;11:2113.\newline \href{https://doi.org/10.1038/s41467-020-15977-4}{doi:\allowbreak 10.\allowbreak 1038/\allowbreak s41467-\allowbreak 020-\allowbreak 15977-\allowbreak 4} & Regulatory DNA band: the only million-scale promoter library outside the GPRA lineage, which is what makes it a check on it rather than more of it. & Zenodo 3\allowbreak 735426 (\allowbreak Promoter\allowbreak  Design \allowbreak Data\_\allowbreak 2020-\allowbreak 03.\allowbreak zip, ab\allowbreak out 1.\allowbreak 1 GB) \\
\addlinespace[5pt]
5 & \textbf{Zhu 2024 (Pol II MASTER, TSS initiation efficiency)} \org{S. cerevisiae}\newline Zhu C, Vvedenskaya IO, Sze SH, Nickels BE, Kaplan CD. Nat Struct Mol Biol 2024;31:190-202.\newline \href{https://doi.org/10.1038/s41594-023-01171-9}{doi:\allowbreak 10.\allowbreak 1038/\allowbreak s41594-\allowbreak 023-\allowbreak 01171-\allowbreak 9} & Regulatory DNA band: a different readout from reporter level, so it extends the axis rather than deepening it, and at a million records it is the third-largest row in the wave. & GEO GSE1\allowbreak 85290 (p\allowbreak rocessed\allowbreak ); BioPr\allowbreak oject PR\allowbreak JNA76662\allowbreak 4 (raw) \\
\addlinespace[5pt]
6 & \textbf{Cuperus 2017 (random 5' UTR library)} \org{S. cerevisiae}\newline Cuperus JT, Groves B, Kuchina A, Rosenberg AB, Jojic N, Fields S, Seelig G. Genome Res 2017;27:2015-2028.\newline \href{https://genome.cshlp.org/content/27/12/2015}{genome.\allowbreak cshlp.\allowbreak org/\allowbreak content\allowbreak /\allowbreak 27/\allowbreak 12/\allowbreak 2015} & Regulatory DNA band: the only large 5' UTR library, and the only one in the band whose readout is growth rather than fluorescence. & GEO GSE1\allowbreak 04252 \\
\addlinespace[5pt]
7 & \textbf{Gui 2021 (position effects with transcriptomes)} \org{S. cerevisiae}\newline Gui, et al. Transcriptome analysis in yeast reveals the externality of position effects. Mol Biol Evol 2021.\newline \href{https://www.biorxiv.org/content/10.1101/2020.04.02.021162}{biorxiv.\allowbreak org/\allowbreak content\allowbreak /\allowbreak 10.\allowbreak 1101/\allowbreak 2020.\allowbreak 04.\allowbreak 02.\allowbreak 021162} & Regulatory DNA band: the sequence variable is purely the integration coordinate, and this is the one row in that tier whose phenotype is a vector. & no acces\allowbreak sion fou\allowbreak nd; need\allowbreak s a brow\allowbreak ser chec\allowbreak k of the\allowbreak  publish\allowbreak ed suppl\allowbreak ement\newline {\scriptsize (browser)} \\
\addlinespace[5pt]
8 & \textbf{Keren 2016 (promoter panel x named gene, pooled fitness)} \org{S. cerevisiae}\newline Keren L, et al. Massively parallel interrogation of the effects of gene expression levels on fitness. Cell 2016;166:1282-1294.\newline \href{https://www.ncbi.nlm.nih.gov/geo/query/acc.cgi?acc=GSE83936}{ncbi.\allowbreak nlm.\allowbreak nih.\allowbreak gov/\allowbreak geo/\allowbreak query/\allowbreak acc.\allowbreak cgi?\allowbreak acc=\allowbreak GSE8393\allowbreak 6} & Regulatory DNA band: the only large tabulated panel that links expression level to fitness gene by gene, and the only one whose perturbation is a swap at a native locus. & GEO GSE8\allowbreak 3936; Ar\allowbreak rayExpre\allowbreak ss E-\allowbreak GEOD-\allowbreak 83936; \allowbreak BioProje\allowbreak ct PRJNA\allowbreak 329129 \\
\addlinespace[5pt]
9 & \textbf{Shalem 2015 (3' end FACS-seq library)} \org{S. cerevisiae}\newline Shalem O, et al. 3' end determinants. PLoS Genet 2015;11:e1005147.\newline \href{https://doi.org/10.1371/journal.pgen.1005147}{doi:\allowbreak 10.\allowbreak 1371/\allowbreak journal\allowbreak .\allowbreak pgen.\allowbreak 1005147} & Regulatory DNA band: the terminator end of the same axis, with the sequence and the value in one released file. & PLoS Gen\allowbreak et S3 Ta\allowbreak ble (XLS\allowbreak X), no a\allowbreak ccession\allowbreak ; needs \allowbreak a supple\allowbreak ment fet\allowbreak ch and a\allowbreak  sha256 \allowbreak pin\newline {\scriptsize (browser)} \\
\addlinespace[5pt]
10 & \textbf{Renganaath 2020 (native promoter variant MPRA)} \org{S. cerevisiae}\newline Renganaath K, Cheung R, Day L, Kosuri S, Kruglyak L, Albert FW. eLife 2020;9:e62669.\newline \href{https://doi.org/10.7554/eLife.62669}{doi:\allowbreak 10.\allowbreak 7554/\allowbreak eLife.\allowbreak 62669} & Regulatory DNA band: real alleles rather than sampled sequence, measured on the same readout, which is the bridge between the synthetic and the natural axes. & GEO GSE1\allowbreak 55944 (r\allowbreak aw plus \allowbreak barcode \allowbreak assignme\allowbreak nts); th\allowbreak e MPRA d\allowbreak esign an\allowbreak d proces\allowbreak sed tabl\allowbreak es are e\allowbreak Life Sup\allowbreak plementa\allowbreak ry Files \\
\addlinespace[5pt]
11 & \textbf{Arg-CCU tRNA deep mutational scan} \org{S. cerevisiae}\newline Deep mutational scan of a single-copy Arg-CCU tRNA over 14 substitutions seen across post-whole-genome-duplication Saccharomyces. PMC6193533.\newline \href{https://pmc.ncbi.nlm.nih.gov/articles/PMC6193533}{PMC61935\allowbreak 33} & Regulatory DNA band: point substitutions in a non-coding gene, which is the one DMS shape the protein aggregates cannot hold. & journal \allowbreak suppleme\allowbreak nt, no a\allowbreak ccession\allowbreak  located\allowbreak ; needs \allowbreak a supple\allowbreak ment fet\allowbreak ch\newline {\scriptsize (browser)} \\
\addlinespace[5pt]
12 & \textbf{Keren 2013 (native promoter panel across conditions)} \org{S. cerevisiae}\newline Keren L, et al. Mol Syst Biol 2013.\newline \href{https://spell.yeastgenome.org/}{spell.\allowbreak yeastge\allowbreak nome.\allowbreak org} & Regulatory DNA band: the only row in it with a real condition axis on native sequence. & not loca\allowbreak ted; rea\allowbreak ched onl\allowbreak y throug\allowbreak h the Zr\allowbreak imec and\allowbreak  Karollu\allowbreak s valida\allowbreak tion set\allowbreak s\newline {\scriptsize (browser)} \\
\addlinespace[5pt]
13 & \textbf{Lubliner 2015 (core promoter variant panel)} \org{S. cerevisiae}\newline Lubliner S, Keren L, Segal E, et al. Core promoter sequence in yeast is a major determinant of expression level. Genome Res 2015;25:1008-1017.\newline \href{https://genome.cshlp.org/content/25/7/1008}{genome.\allowbreak cshlp.\allowbreak org/\allowbreak content\allowbreak /\allowbreak 25/\allowbreak 7/\allowbreak 1008} & Regulatory DNA band: the core promoter is a distinct regulatory element from the upstream activating region every other row varies. & Genome R\allowbreak es Suppl\allowbreak emental \allowbreak Table 1;\allowbreak  needs a\allowbreak  supplem\allowbreak ent fetc\allowbreak h\newline {\scriptsize (browser)} \\
\addlinespace[5pt]
14 & \textbf{Sharon 2012 (designed promoter library)} \org{S. cerevisiae}\newline Sharon E, Kalma Y, Sharp A, et al., Segal E. Nat Biotechnol 2012;30:521-530.\newline \href{https://doi.org/10.1038/nbt.2205}{doi:\allowbreak 10.\allowbreak 1038/\allowbreak nbt.\allowbreak 2205} & Regulatory DNA band: the controlled counterpart the random libraries lack, and the platform's own origin. & Nat Biot\allowbreak echnol j\allowbreak ournal s\allowbreak upplemen\allowbreak t, no ac\allowbreak cession;\allowbreak  needs a\allowbreak  supplem\allowbreak ent fetc\allowbreak h\newline {\scriptsize (browser)} \\
\addlinespace[5pt]
15 & \textbf{Yamanishi 2013 (genomic terminatome)} \org{S. cerevisiae}\newline Yamanishi, et al. Characterization of heterologous protein expression using over 3,000 yeast genomic terminators. SGD reference S000206762.\newline \href{https://www.yeastgenome.org/reference/S000206762}{yeastgen\allowbreak ome.\allowbreak org/\allowbreak referen\allowbreak ce/\allowbreak S000206\allowbreak 762} & Regulatory DNA band: the largest native-terminator panel, which is the one element of the axis no library in the band randomizes. & primary \allowbreak table no\allowbreak t locate\allowbreak d; needs\allowbreak  the pap\allowbreak er ident\allowbreak ified, t\allowbreak hen a su\allowbreak pplement\allowbreak  fetch\newline {\scriptsize (browser)} \\
\addlinespace[5pt]
16 & \textbf{Dvir 2013 (5' UTR positions -10 to -1)} \org{S. cerevisiae}\newline Dvir S, Velten L, Sharon E, Zeevi D, Carey LB, Weinberger A, Segal E. Proc Natl Acad Sci USA 2013;110:E2792-E2801.\newline \href{https://doi.org/10.1073/pnas.1222534110}{doi:\allowbreak 10.\allowbreak 1073/\allowbreak pnas.\allowbreak 1222534\allowbreak 110} & Regulatory DNA band: a controlled ten-base window, which is the designed counterpart to the half-million random 5' UTRs above it. & PNAS jou\allowbreak rnal sup\allowbreak plement,\allowbreak  no acce\allowbreak ssion; n\allowbreak eeds a s\allowbreak upplemen\allowbreak t fetch\newline {\scriptsize (browser)} \\
\addlinespace[5pt]
17 & \textbf{Wu 2017 (genome-wide position effects, 1,044 loci)} \org{S. cerevisiae}\newline Wu XL, et al. Genome-wide landscape of position effects on heterogeneous gene expression in Saccharomyces cerevisiae. Biotechnol Biofuels 2017.\newline \href{https://www.ncbi.nlm.nih.gov/pmc/articles/PMC5516366/}{PMC55163\allowbreak 66} & Regulatory DNA band by its sequence variable, the integration coordinate, even though the readout is one reporter value. & Table S2\allowbreak , stated\allowbreak  availab\allowbreak le from \allowbreak the corr\allowbreak espondin\allowbreak g author\allowbreak  on reas\allowbreak onable r\allowbreak equest\newline {\scriptsize (browser)} \\
\addlinespace[5pt]
18 & \textbf{ProteinGym yeast subset} \org{S. cerevisiae}\newline ProteinGym: deep mutational scanning scores from 217 assays over 186 proteins, distributed through ProteinGymR and ProteinGymPy.\newline \href{https://proteingym.org}{proteing\allowbreak ym.\allowbreak org} & Regulatory DNA band as the aggregate route into the mutational-scan tier, which is the sequence-to-function axis in coding space rather than regulatory space. & ProteinG\allowbreak ymR and \allowbreak ProteinG\allowbreak ymPy pac\allowbreak kage loa\allowbreak ders; ar\allowbreak chived o\allowbreak n Zenodo\allowbreak , record\allowbreak  number \allowbreak not capt\allowbreak ured \\
\addlinespace[5pt]
19 & \textbf{Zhang 2016 (genome-wide GFP position scan)} \org{S. cerevisiae}\newline Zhang, et al. Genome-wide GFP position-effect scan using the yeast knockout collection. Chin J Biotech 2016.\newline \href{https://wsws.ijournals.cn/cjbcn/article/pdf/gc16070901}{wsws.\allowbreak ijourna\allowbreak ls.\allowbreak cn/\allowbreak cjbcn/\allowbreak article\allowbreak /\allowbreak pdf/\allowbreak gc16070\allowbreak 901} & Regulatory DNA band by its sequence variable, and the only FACS-read position-effect scan. & no acces\allowbreak sion loc\allowbreak ated; Ch\allowbreak inese-\allowbreak languag\allowbreak e journa\allowbreak l PDF, n\allowbreak eeds a b\allowbreak rowser f\allowbreak etch\newline {\scriptsize (browser)} \\
\addlinespace[5pt]
20 & \textbf{Kita 2017 (context-stratified yeast eQTLs)} \org{S. cerevisiae}\newline Kita R, et al. Proc Natl Acad Sci USA 2017;114.\newline \href{https://www.pnas.org/doi/10.1073/pnas.1717421114}{pnas.\allowbreak org/\allowbreak doi/\allowbreak 10.\allowbreak 1073/\allowbreak pnas.\allowbreak 1717421\allowbreak 114} & Regulatory DNA band: variants stratified by the regulatory element they sit in, which is the only released set with that stratification. & PNAS sup\allowbreak plementa\allowbreak ry data \allowbreak file pna\allowbreak s.\allowbreak 1717421\allowbreak 114.\allowbreak sd01; n\allowbreak eeds a s\allowbreak upplemen\allowbreak t fetch\newline {\scriptsize (browser)} \\
\addlinespace[5pt]
21 & \textbf{Chen (earlier position-effect study, 482 loci)} \org{S. cerevisiae}\newline An earlier Chen study characterizing position effects at 482 loci, reported inside Wu 2017; primary citation not located.\newline \href{https://www.ncbi.nlm.nih.gov/pmc/articles/PMC5516366/}{PMC55163\allowbreak 66} & Regulatory DNA band by its sequence variable, the integration coordinate. & not loca\allowbreak ted; the\allowbreak  primary\allowbreak  paper m\allowbreak ust be i\allowbreak dentifie\allowbreak d first\newline {\scriptsize (browser)} \\
\addlinespace[5pt]
22 & \textbf{Wei 2017 (synthetic terminator panel)} \org{S. cerevisiae}\newline Wei, et al. Yeast synthetic terminators: fine regulation of strength through linker sequences. ChemBioChem 2017.\newline \href{https://doi.org/10.1002/cbic.201700516}{doi:\allowbreak 10.\allowbreak 1002/\allowbreak cbic.\allowbreak 2017005\allowbreak 16} & Regulatory DNA band: the terminator parts tier, whose values are relative to a standard part rather than absolute. & ChemBioC\allowbreak hem jour\allowbreak nal supp\allowbreak lement, \allowbreak no repos\allowbreak itory de\allowbreak posit fo\allowbreak und\newline {\scriptsize (browser)} \\
\addlinespace[5pt]
23 & \textbf{Lubliner 2013 (native core promoter maximal strengths)} \org{S. cerevisiae}\newline Lubliner S, Keren L, Segal E. Nucleic Acids Res 2013.\newline \href{https://academic.oup.com/nar}{academic\allowbreak .\allowbreak oup.\allowbreak com/\allowbreak nar} & Regulatory DNA band: the native counterpart of the core-promoter tier, and the only source for its reference values. & not loca\allowbreak ted\newline {\scriptsize (browser)} \\
\addlinespace[5pt]
24 & \textbf{Morse (short synthetic terminators)} \org{S. cerevisiae}\newline Morse NJ, et al. Short synthetic terminators, 35 to 70 bp; primary citation not located.\newline \href{https://scite.ai/authors/nicholas-j-morse-1g5Qrn}{scite.\allowbreak ai/\allowbreak authors\allowbreak /\allowbreak nichola\allowbreak s-\allowbreak j-\allowbreak morse-\allowbreak 1g5Qrn} & Regulatory DNA band: the short-terminator corner of the parts tier, held so the lead is not lost. & no repos\allowbreak itory de\allowbreak posit fo\allowbreak und; pri\allowbreak mary pap\allowbreak er not i\allowbreak dentifie\allowbreak d\newline {\scriptsize (browser)} \\
\addlinespace[5pt]
25 & \textbf{MaveDB yeast score sets} \org{S. cerevisiae}\newline Esposito D, et al. MaveDB: an open-source platform to distribute and interpret data from multiplexed assays of variant effect. Genome Biol 2019;20:223.\newline \href{https://doi.org/10.1186/s13059-019-1845-6}{doi:\allowbreak 10.\allowbreak 1186/\allowbreak s13059-\allowbreak 019-\allowbreak 1845-\allowbreak 6} & Regulatory DNA band as the second aggregate route into the mutational-scan tier, held because the two aggregates do not hold the same assays. & mavedb.\allowbreak org; th\allowbreak e docume\allowbreak nted sea\allowbreak rch API \allowbreak endpoint\allowbreak  returns\allowbreak  HTTP 40\allowbreak 4, so th\allowbreak e organi\allowbreak sm filte\allowbreak r needs \allowbreak a live q\allowbreak uery\newline {\scriptsize (browser)} \\
\addlinespace[5pt]
26 & \textbf{Fitness Browser RCH2 (Stutzerimonas stutzeri\_AE)} \org{S. stutzeri}\newline Price MN, Wetmore KM, Waters RJ, Callaghan M, Ray J, Liu H, Kuehl JV, Melnyk RA, Lamson JS, Suh Y, et al., Deutschbauer AM, Arkin AP. Nature 2018;557:503-509.\newline \href{https://doi.org/10.1038/s41586-018-0124-0}{doi:\allowbreak 10.\allowbreak 1038/\allowbreak s41586-\allowbreak 018-\allowbreak 0124-\allowbreak 0} & Cross-host fitness: nine new hosts arrive inside one artifact under one schema, one condition vocabulary and an identifier bridge the source supplies, so the band exists to keep them together rather than scattered down a measurement ranking. & feba.\allowbreak db at 1\allowbreak 0,433,70\allowbreak 8,032 by\allowbreak tes, Las\allowbreak t-\allowbreak Modifie\allowbreak d 2026-\allowbreak 10-\allowbreak 27, HTT\allowbreak P 200 to\allowbreak  a plain\allowbreak  curl; o\allowbreak rgId psR\allowbreak CH2; ref\allowbreak seq\_\allowbreak mapping\allowbreak .\allowbreak tsv bri\allowbreak dges loc\allowbreak usId to \allowbreak RefSeq \\
\addlinespace[5pt]
27 & \textbf{Fitness Browser FW300-N2E3 (Aquipseudomonas helvetica)} \org{P. fluorescens}\newline Price MN, Wetmore KM, Waters RJ, et al., Deutschbauer AM, Arkin AP. Nature 2018;557:503-509.\newline \href{https://doi.org/10.1038/s41586-018-0124-0}{doi:\allowbreak 10.\allowbreak 1038/\allowbreak s41586-\allowbreak 018-\allowbreak 0124-\allowbreak 0} & Cross-host fitness: part of the single-artifact block of nine hosts, whose shared schema is the reason to admit them together. & feba.\allowbreak db, org\allowbreak Id pseud\allowbreak o3\_\allowbreak N2E3; t\allowbreak he histo\allowbreak rical bi\allowbreak nomial i\allowbreak s P.\allowbreak  fluore\allowbreak scens an\allowbreak d the cu\allowbreak rrent on\allowbreak e is A.\allowbreak  helvet\allowbreak ica \\
\addlinespace[5pt]
28 & \textbf{Fitness Browser FW300-N2E2 (Aquipseudomonas fluorescens\_Q)} \org{P. fluorescens}\newline Price MN, Wetmore KM, Waters RJ, et al., Deutschbauer AM, Arkin AP. Nature 2018;557:503-509.\newline \href{https://doi.org/10.1038/s41586-018-0124-0}{doi:\allowbreak 10.\allowbreak 1038/\allowbreak s41586-\allowbreak 018-\allowbreak 0124-\allowbreak 0} & Cross-host fitness: one of four isolates of one species in the block, which is the within-species comparison the band makes possible. & feba.\allowbreak db, org\allowbreak Id pseud\allowbreak o6\_\allowbreak N2E2; 1\allowbreak 16 updat\allowbreak ed annot\allowbreak ations \\
\addlinespace[5pt]
29 & \textbf{Fitness Browser WCS417 (Aquipseudomonas simiae)} \org{P. simiae}\newline Price MN, Wetmore KM, Waters RJ, et al., Deutschbauer AM, Arkin AP. Nature 2018;557:503-509.\newline \href{https://doi.org/10.1038/s41586-018-0124-0}{doi:\allowbreak 10.\allowbreak 1038/\allowbreak s41586-\allowbreak 018-\allowbreak 0124-\allowbreak 0} & Cross-host fitness: the plant-associated member of the single-artifact block, with its widest carbon panel. & feba.\allowbreak db, org\allowbreak Id WCS41\allowbreak 7; histo\allowbreak rical bi\allowbreak nomial P\allowbreak .\allowbreak  simiae\allowbreak , curren\allowbreak t A.\allowbreak  simiae \\
\addlinespace[5pt]
30 & \textbf{Fitness Browser FW300-N2C3 (Aquipseudomonas fluorescens\_AA)} \org{P. fluorescens}\newline Price MN, Wetmore KM, Waters RJ, et al., Deutschbauer AM, Arkin AP. Nature 2018;557:503-509.\newline \href{https://doi.org/10.1038/s41586-018-0124-0}{doi:\allowbreak 10.\allowbreak 1038/\allowbreak s41586-\allowbreak 018-\allowbreak 0124-\allowbreak 0} & Cross-host fitness: part of the four-isolate within-species block, and the row that taught the most unannotated genes. & feba.\allowbreak db, org\allowbreak Id pseud\allowbreak o5\_\allowbreak N2C3\_\allowbreak 1 \\
\addlinespace[5pt]
31 & \textbf{Fitness Browser FW300-N1B4 (Aquipseudomonas sp.)} \org{P. fluorescens}\newline Price MN, Wetmore KM, Waters RJ, et al., Deutschbauer AM, Arkin AP. Nature 2018;557:503-509.\newline \href{https://doi.org/10.1038/s41586-018-0124-0}{doi:\allowbreak 10.\allowbreak 1038/\allowbreak s41586-\allowbreak 018-\allowbreak 0124-\allowbreak 0} & Cross-host fitness: completes the four-isolate within-species block, and names the case of a host with no species assignment. & feba.\allowbreak db, org\allowbreak Id pseud\allowbreak o1\_\allowbreak N1B4 \\
\addlinespace[5pt]
32 & \textbf{Fitness Browser B728a (Aquipseudomonas fragariae pv. syringae)} \org{P. syringae}\newline Price MN, Wetmore KM, Waters RJ, et al., Deutschbauer AM, Arkin AP. Nature 2018;557:503-509.\newline \href{https://doi.org/10.1038/s41586-018-0124-0}{doi:\allowbreak 10.\allowbreak 1038/\allowbreak s41586-\allowbreak 018-\allowbreak 0124-\allowbreak 0} & Cross-host fitness: the pathogen member of the block, and half of the one strain in it screened in two genetic backgrounds. & feba.\allowbreak db, org\allowbreak Id Syrin\allowbreak gaeB728a\allowbreak ; the GT\allowbreak DB renam\allowbreak e to A.\allowbreak  fragar\allowbreak iae pv.\allowbreak  syring\allowbreak ae means\allowbreak  both bi\allowbreak nomials \allowbreak are need\allowbreak ed as al\allowbreak iases \\
\addlinespace[5pt]
33 & \textbf{Fitness Browser B728a mexB deletion background} \org{P. syringae}\newline Price MN, Wetmore KM, Waters RJ, et al., Deutschbauer AM, Arkin AP. Nature 2018;557:503-509.\newline \href{https://doi.org/10.1038/s41586-018-0124-0}{doi:\allowbreak 10.\allowbreak 1038/\allowbreak s41586-\allowbreak 018-\allowbreak 0124-\allowbreak 0} & Cross-host fitness: the background-crossed half of the B728a pair, which is the only interaction design in the block. & feba.\allowbreak db, org\allowbreak Id Syrin\allowbreak gaeB728a\_\allowbreak mexBdel\allowbreak ta \\
\addlinespace[5pt]
34 & \textbf{Fitness Browser PA14 (Pseudomonas aeruginosa)} \org{P. aeruginosa}\newline Price MN, Wetmore KM, Waters RJ, et al., Deutschbauer AM, Arkin AP. Nature 2018;557:503-509.\newline \href{https://doi.org/10.1038/s41586-018-0124-0}{doi:\allowbreak 10.\allowbreak 1038/\allowbreak s41586-\allowbreak 018-\allowbreak 0124-\allowbreak 0} & Cross-host fitness: the same artifact and schema, kept visible because its size is the finding about this species rather than an omission. & feba.\allowbreak db, org\allowbreak Id PA14;\allowbreak  the bro\allowbreak wser doe\allowbreak s not re\allowbreak name thi\allowbreak s organi\allowbreak sm \\
\addlinespace[5pt]
35 & \textbf{Dubinkina 2024 (12-strain low-pH transcriptomic atlas)} \org{P. kudriavzevii}\newline Dubinkina V, Bhogale S, Hsieh PH, Dibaeinia P, Nambiar S, Maslov S, Yoshikuni Y, Sinha S. Microbiol Spectr 2024.\newline \href{https://doi.org/10.1128/spectrum.02536-23}{doi:\allowbreak 10.\allowbreak 1128/\allowbreak spectru\allowbreak m.\allowbreak 02536-\allowbreak 23} & Host bootstrap: the only multi-strain expression matrix this host has, and the only deposit that already carries processed values rather than reads. & Illinois\allowbreak  Data Ba\allowbreak nk DOI 1\allowbreak 0.\allowbreak 13012/\allowbreak B2IDB-\allowbreak 1336177\_\allowbreak V1, hol\allowbreak ding str\allowbreak ain regu\allowbreak lators, \allowbreak strains \allowbreak and gene\allowbreak  sets, a\allowbreak nd gene \allowbreak expressi\allowbreak on count\allowbreak s; the p\allowbreak age retu\allowbreak rns 403 \allowbreak to every\allowbreak  automat\allowbreak ed clien\allowbreak t\newline {\scriptsize (browser)} \\
\addlinespace[5pt]
36 & \textbf{MAESTRO 2025 (170-isolate WGS, 57-stressor growth panel)} \org{P. kudriavzevii}\newline Hsieh PH, Sasaki Y, Yang CI, Wu ZY, Steindorff AS, Haridas S, Ke J, Fatma Z, Zhao Z, Opulente DA, Deng S, Hittinger CT, Grigoriev IV, Dien B, Zhao H, Li YP, Yoshikuni Y. bioRxiv 2025;2025.08.29.673034.\newline \href{https://doi.org/10.1101/2025.08.29.673034}{doi:\allowbreak 10.\allowbreak 1101/\allowbreak 2025.\allowbreak 08.\allowbreak 29.\allowbreak 673034} & Host bootstrap: this host has no genome-wide screen of any kind, so its largest resource is a 161-isolate panel that the measurement rule ranks below every pseudomonad screen. The band is what keeps it in the first tranche. & no acces\allowbreak sion ret\allowbreak rieved:\allowbreak  bioRxi\allowbreak v return\allowbreak s 403 an\allowbreak d Europe\allowbreak  PMC ful\allowbreak lTextXML\allowbreak  for PPR\allowbreak 1078985 \allowbreak returns \allowbreak 500; per\allowbreak -\allowbreak isolate\allowbreak  IO\_\allowbreak NNN Bio\allowbreak Projects\allowbreak  are inf\allowbreak erred, n\allowbreak ot confi\allowbreak rmed\newline {\scriptsize (ask)} \\
\addlinespace[5pt]
37 & \textbf{Fatma 2023 (landing-pad locus screen)} \org{P. kudriavzevii}\newline Fatma Z, Tan SI, Boob A, Zhao H. A landing pad system for multicopy gene integration in Issatchenkia orientalis. bioRxiv 2023;2023.05.21.541627.\newline \href{https://doi.org/10.1101/2023.05.21.541627}{doi:\allowbreak 10.\allowbreak 1101/\allowbreak 2023.\allowbreak 05.\allowbreak 21.\allowbreak 541627} & Host bootstrap: the host's only measured per-locus fitness, which no scale ranking would ever place in a first tranche at 188 records. & a 23 KB \allowbreak package \allowbreak at uofi.\allowbreak box.\allowbreak com (st\allowbreak rains, p\allowbreak lasmids,\allowbreak  guide R\allowbreak NAs, pri\allowbreak mers) pl\allowbreak us githu\allowbreak b.\allowbreak com/\allowbreak HuiminZ\allowbreak hao/\allowbreak Landing\allowbreak -\allowbreak pad-\allowbreak model; \allowbreak no DOI a\allowbreak nd no ar\allowbreak chival d\allowbreak eposit\newline {\scriptsize (ask)} \\
\addlinespace[5pt]
38 & \textbf{Douglass 2018 (32 sequenced isolates, 4-antifungal MIC)} \org{P. kudriavzevii}\newline Douglass AP, Offei B, Braun-Galleani S, Coughlan AY, Martos AAR, Ortiz-Merino RA, Byrne KP, Wolfe KH. PLoS Pathog 2018;14:e1007138.\newline \href{https://www.ncbi.nlm.nih.gov/pmc/articles/PMC6053246/}{PMC60532\allowbreak 46} & Host bootstrap: at 128 values it is the host's best public benchmark, which is the measure of how thin the host is and why a band rather than a score has to carry it. & SRA SRP1\allowbreak 39299 (3\allowbreak 2 strain\allowbreak s); GenB\allowbreak ank CP02\allowbreak 8773-\allowbreak CP02877\allowbreak 8 (CBS57\allowbreak 3 refere\allowbreak nce) and\allowbreak  CP02853\allowbreak 1-\allowbreak CP02853\allowbreak 5 (CBS51\allowbreak 47); MIC\allowbreak  table i\allowbreak n the op\allowbreak en-\allowbreak access \allowbreak article \allowbreak body \\
\addlinespace[5pt]
39 & \textbf{GSE282019 (Delhi NICU five-year outbreak)} \org{P. kudriavzevii}\newline Genomic analyses reveal high diversity and rapid evolution of Pichia kudriavzevii within a neonatal intensive care unit in Delhi, India. GEO series GSE282019, 2024; PMID 39853119.\newline \href{https://www.ncbi.nlm.nih.gov/geo/query/acc.cgi?acc=GSE282019}{ncbi.\allowbreak nlm.\allowbreak nih.\allowbreak gov/\allowbreak geo/\allowbreak query/\allowbreak acc.\allowbreak cgi?\allowbreak acc=\allowbreak GSE2820\allowbreak 19} & Host bootstrap: the host's only published drug-stratified expression series, and openly retrievable. & GEO GSE2\allowbreak 82019, p\allowbreak latform \allowbreak GPL34351\allowbreak ; progra\allowbreak mmatical\allowbreak ly retri\allowbreak evable \\
\addlinespace[5pt]
40 & \textbf{GSE202655 (L-malic acid, EcSthA transhydrogenase)} \org{P. kudriavzevii}\newline Metabolic engineering of the acid-tolerant yeast Pichia kudriavzevii for efficient L-malic acid production at low pH. GEO series GSE202655, 2023; PMID 36566973.\newline \href{https://www.ncbi.nlm.nih.gov/geo/query/acc.cgi?acc=GSE202655}{ncbi.\allowbreak nlm.\allowbreak nih.\allowbreak gov/\allowbreak geo/\allowbreak query/\allowbreak acc.\allowbreak cgi?\allowbreak acc=\allowbreak GSE2026\allowbreak 55} & Host bootstrap: an engineered-strain expression series in a non-SD108 background, which the host otherwise lacks. & GEO GSE2\allowbreak 02655, p\allowbreak latform \allowbreak GPL32243\allowbreak ; progra\allowbreak mmatical\allowbreak ly retri\allowbreak evable \\
\addlinespace[5pt]
41 & \textbf{GSE159457 (CY902 C4 dicarboxylic acid transporters)} \org{P. kudriavzevii}\newline Identification and characterization of C4 dicarboxylic acid transporters in Pichia kudriavzevii CY902. GEO series GSE159457, 2021; PMID 33629807.\newline \href{https://www.ncbi.nlm.nih.gov/geo/query/acc.cgi?acc=GSE159457}{ncbi.\allowbreak nlm.\allowbreak nih.\allowbreak gov/\allowbreak geo/\allowbreak query/\allowbreak acc.\allowbreak cgi?\allowbreak acc=\allowbreak GSE1594\allowbreak 57} & Host bootstrap: a carbon-source axis on expression in a third background, fully deposited. & GEO GSE1\allowbreak 59457, p\allowbreak latform \allowbreak GPL29244\allowbreak ; raw re\allowbreak ads at S\allowbreak RA SRP28\allowbreak 7255; pr\allowbreak ogrammat\allowbreak ically r\allowbreak etrievab\allowbreak le \\
\addlinespace[5pt]
42 & \textbf{GSE262836 (Delhi NICU, fluconazole-resistance expression)} \org{P. kudriavzevii}\newline Genomic characterization and transcriptomic profile of Candida krusei in a neonatal intensive care unit in Delhi, India. GEO series GSE262836, 2025.\newline \href{https://www.ncbi.nlm.nih.gov/geo/query/acc.cgi?acc=GSE262836}{ncbi.\allowbreak nlm.\allowbreak nih.\allowbreak gov/\allowbreak geo/\allowbreak query/\allowbreak acc.\allowbreak cgi?\allowbreak acc=\allowbreak GSE2628\allowbreak 36} & Host bootstrap: one of two rows in this host pairing a drug phenotype with expression, and openly retrievable. & GEO GSE2\allowbreak 62836, p\allowbreak latform \allowbreak GPL34351\allowbreak ; progra\allowbreak mmatical\allowbreak ly retri\allowbreak evable \\
\addlinespace[5pt]
43 & \textbf{Thompson 2019 (KT2440 lysine metabolism)} \org{P. putida}\newline Thompson MG, et al. Massively parallel fitness profiling reveals multiple novel enzymes in Pseudomonas putida lysine metabolism. mBio 2019.\newline \href{https://pmc.ncbi.nlm.nih.gov/articles/PMC6509195}{PMC65091\allowbreak 95} & Scale band: tier 3, ordered by measurements (22,256). & Table S1\allowbreak , a CSV \allowbreak suppleme\allowbreak ntary fi\allowbreak le; no r\allowbreak epositor\allowbreak y\newline {\scriptsize (browser)} \\
\addlinespace[5pt]
44 & \textbf{Isoprenol CRISPRi DBTL proteome ladder (7 PRIDE deposits)} \org{P. putida}\newline Automation of a CRISPRi platform for enhanced isoprenol production in Pseudomonas putida; PRIDE projects PXD063733, PXD063737, PXD063738, PXD063740, PXD063743, PXD063744, PXD063746.\newline \href{https://www.ebi.ac.uk/pride/archive/projects/PXD063733}{ebi.\allowbreak ac.\allowbreak uk/\allowbreak pride/\allowbreak archive\allowbreak /\allowbreak project\allowbreak s/\allowbreak PXD0637\allowbreak 33} & Scale band: tier 3, ordered by measurements (38,948). & PRIDE PX\allowbreak D063733,\allowbreak  PXD0637\allowbreak 37, PXD0\allowbreak 63738, P\allowbreak XD063740\allowbreak , PXD063\allowbreak 743, PXD\allowbreak 063744 a\allowbreak nd PXD06\allowbreak 3746, al\allowbreak l submit\allowbreak ted 2025\allowbreak -\allowbreak 05-\allowbreak 07 and \allowbreak 08 on an\allowbreak  Orbitra\allowbreak p Explor\allowbreak is 480 \\
\addlinespace[5pt]
45 & \textbf{Thompson 2020 (KT2440 fatty acids and alcohols)} \org{P. putida}\newline Thompson MG, Incha MR, Pearson AN, et al. Appl Environ Microbiol 2020;86:e01665-20; correction at 10.1128/AEM.00177-21.\newline \href{https://pubmed.ncbi.nlm.nih.gov/32826213}{PMID 328\allowbreak 26213} & Scale band: tier 3, ordered by measurements (127,972). & not stat\allowbreak ed; the \allowbreak values m\allowbreak ost like\allowbreak ly reach\allowbreak  feba.\allowbreak db \\
\addlinespace[5pt]
46 & \textbf{GSE198395 (putidaPRECISE321 carbon sources)} \org{P. putida}\newline GEO series GSE198395, KT2440 and engineered strains on different carbon sources, part of putidaPRECISE321, released 2022-04-28.\newline \href{https://www.ncbi.nlm.nih.gov/geo/query/acc.cgi?acc=GSE198395}{ncbi.\allowbreak nlm.\allowbreak nih.\allowbreak gov/\allowbreak geo/\allowbreak query/\allowbreak acc.\allowbreak cgi?\allowbreak acc=\allowbreak GSE1983\allowbreak 95} & Scale band: tier 3, ordered by measurements (450,684). & GEO GSE1\allowbreak 98395; p\allowbreak rogramma\allowbreak tically \allowbreak retrieva\allowbreak ble \\
\addlinespace[5pt]
47 & \textbf{Mueller 2025 (24-strain P. putida Biolog panel)} \org{P. putida}\newline Mueller J, Krishnan J, Wei Q, Hefner Y, Monk JM, Verkler H, Tibocha-Bonilla JD, et al. bioRxiv 2025;2025.11.17.688870.\newline \href{https://doi.org/10.1101/2025.11.17.688870}{doi:\allowbreak 10.\allowbreak 1101/\allowbreak 2025.\allowbreak 11.\allowbreak 17.\allowbreak 688870} & Scale band: tier 2, ordered by measurements (4,920). & unconfir\allowbreak med:\allowbreak  bioRxi\allowbreak v return\allowbreak ed HTTP \allowbreak 429 on f\allowbreak our atte\allowbreak mpts, so\allowbreak  no BioP\allowbreak roject, \allowbreak Zenodo r\allowbreak ecord or\allowbreak  machine\allowbreak -\allowbreak readabl\allowbreak e Table \allowbreak S2 is co\allowbreak nfirmed\newline {\scriptsize (ask)} \\
\addlinespace[5pt]
48 & \textbf{Schmidt 2022 (KT2440 nitrogen BarSeq)} \org{P. putida}\newline Schmidt M, Pearson AN, Incha MR, et al. Appl Environ Microbiol 2022;88:e02430-21.\newline \href{https://pmc.ncbi.nlm.nih.gov/articles/PMC9004399}{PMC90043\allowbreak 99} & Scale band: tier 2, ordered by measurements (395,044). & no acces\allowbreak sion sta\allowbreak ted; the\allowbreak  study l\allowbreak inks int\allowbreak o the Fi\allowbreak tness Br\allowbreak owser, s\allowbreak o its ex\allowbreak periment\allowbreak s are mo\allowbreak st likel\allowbreak y inside\allowbreak  feba.\allowbreak db \\
\addlinespace[5pt]
49 & \textbf{Borchert 2024 (KT2440 fitness compendium)} \org{P. putida}\newline Borchert AJ, Bleem AC, Lim HG, Rychel K, Dooley KD, Kellermyer ZA, Hodges TL, Palsson BO, Beckham GT. mSystems 2024;9:e00942-23.\newline \href{https://doi.org/10.1128/msystems.00942-23}{doi:\allowbreak 10.\allowbreak 1128/\allowbreak msystem\allowbreak s.\allowbreak 00942-\allowbreak 23} & Scale band: tier 1, ordered by measurements (1,571,024). & github.\allowbreak com/\allowbreak beckham\allowbreak -\allowbreak lab/\allowbreak fModule\allowbreak  (fitnes\allowbreak s values\allowbreak , statis\allowbreak tics, me\allowbreak tadata);\allowbreak  the non\allowbreak -\allowbreak browser\allowbreak  samples\allowbreak  are at \allowbreak SRA PRJN\allowbreak A809672,\allowbreak  PRJNA85\allowbreak 6070, PR\allowbreak JNA10112\allowbreak 87 \\
\addlinespace[5pt]
50 & \textbf{putidaPRECISE321 (Lim 2022)} \org{P. putida}\newline Lim HG, Rychel K, Sastry AV, Mueller J, Niu W, Feist AM, Palsson BO. Metab Eng 2022;72.\newline \href{https://doi.org/10.1016/j.ymben.2022.04.004}{doi:\allowbreak 10.\allowbreak 1016/\allowbreak j.\allowbreak ymben.\allowbreak 2022.\allowbreak 04.\allowbreak 004} & Scale band: tier 1, ordered by measurements (1,786,044). & github.\allowbreak com/\allowbreak SBRG/\allowbreak modulom\allowbreak e\_\allowbreak ppu (MI\allowbreak T); name\allowbreak d GEO de\allowbreak posits G\allowbreak SE198395\allowbreak  (81 sam\allowbreak ples) an\allowbreak d GSE193\allowbreak 493 (22)\allowbreak ; the fu\allowbreak ll per-\allowbreak sample \allowbreak GEO list\allowbreak  is in t\allowbreak he SI sp\allowbreak readshee\allowbreak t \\
\addlinespace[5pt]
51 & \textbf{Cao 2020 (pooled four-condition RNA-seq)} \org{P. kudriavzevii}\newline Cao M, et al., Zhao H. Metab Eng 2020;59:87-97, promoter-nomination RNA-seq.\newline \href{https://doi.org/10.1016/j.ymben.2020.01.005}{doi:\allowbreak 10.\allowbreak 1016/\allowbreak j.\allowbreak ymben.\allowbreak 2020.\allowbreak 01.\allowbreak 005} & Host bootstrap: the host's reference expression layer, and the one row that forces a decision about an environment that is a pool. & Illinois\allowbreak  Data Ba\allowbreak nk IDB-\allowbreak 8429399\allowbreak , Transc\allowbreak riptomic\allowbreak s tag; n\allowbreak o SRA or\allowbreak  GEO acc\allowbreak ession n\allowbreak amed; th\allowbreak e page r\allowbreak eturns 4\allowbreak 03\newline {\scriptsize (browser)} \\
\addlinespace[5pt]
52 & \textbf{PXD069498 (extracellular vesicle proteome)} \org{P. kudriavzevii}\newline PRIDE project PXD069498, submitter Chiara D'Ambrosio, April 2026.\newline \href{https://www.ebi.ac.uk/pride/archive/projects/PXD069498}{ebi.\allowbreak ac.\allowbreak uk/\allowbreak pride/\allowbreak archive\allowbreak /\allowbreak project\allowbreak s/\allowbreak PXD0694\allowbreak 98} & Host bootstrap: the host's only vesicle proteome, held for completeness at the bottom of the band. & PRIDE PX\allowbreak D069498;\allowbreak  retriev\allowbreak able thr\allowbreak ough the\allowbreak  PRIDE A\allowbreak PI \\
\addlinespace[5pt]
53 & \textbf{Douglass 2018 (per-strain SNP density and loss of heterozygosity)} \org{P. kudriavzevii}\newline Douglass AP, et al., Wolfe KH. PLoS Pathog 2018;14:e1007138, S2 Table.\newline \href{https://www.ncbi.nlm.nih.gov/pmc/articles/PMC6053246/}{PMC60532\allowbreak 46} & Host bootstrap: a genome-architecture layer for the host's only sequenced panel, which nothing else supplies. & S2 Table\allowbreak  of the \allowbreak open-\allowbreak access \allowbreak article \\
\addlinespace[5pt]
54 & \textbf{Tran 2023 (succinate end-to-end pipeline)} \org{P. kudriavzevii}\newline Tran VG, Mishra S, Bhagwat SS, Shafaei S, Shen Y, Allen JL, Crosly BA, Tan SI, Fatma Z, Rabinowitz JD, Guest JS, Singh V, Zhao H. Nat Commun 2023;14.\newline \href{https://pmc.ncbi.nlm.nih.gov/articles/PMC10547785/}{PMC10547\allowbreak 785} & Host bootstrap: the host's deepest named-genotype record, which is worth ingesting for the genotype side rather than for record count. & no seque\allowbreak nce or r\allowbreak aw-\allowbreak data ac\allowbreak cession;\allowbreak  all val\allowbreak ues in t\allowbreak he Suppl\allowbreak ementary\allowbreak  Informa\allowbreak tion and\allowbreak  the Sou\allowbreak rce Data\allowbreak  file, w\allowbreak ith plas\allowbreak mid sequ\allowbreak ences in\allowbreak  Supplem\allowbreak entary D\allowbreak ata 4\newline {\scriptsize (browser)} \\
\addlinespace[5pt]
55 & \textbf{Cao 2020 (constitutive promoter-GFP panel)} \org{P. kudriavzevii}\newline Cao M, Fatma Z, Song X, Hsieh PH, Tran VG, Lyon WL, Sayadi M, Shao Z, Yoshikuni Y, Zhao H. Metab Eng 2020;59:87-97.\newline \href{https://doi.org/10.1016/j.ymben.2020.01.005}{doi:\allowbreak 10.\allowbreak 1016/\allowbreak j.\allowbreak ymben.\allowbreak 2020.\allowbreak 01.\allowbreak 005} & Host bootstrap: the host's only parts panel, and the source of its gene-name to model-id map. & Illinois\allowbreak  Data Ba\allowbreak nk IDB-\allowbreak 8429399\allowbreak  (part s\allowbreak equences\allowbreak  plus a \allowbreak top-\allowbreak 10 prom\allowbreak oter lis\allowbreak t) and t\allowbreak he Elsev\allowbreak ier supp\allowbreak lement z\allowbreak ip; both\allowbreak  return \allowbreak 403 to a\allowbreak utomated\allowbreak  clients\newline {\scriptsize (browser)} \\
\addlinespace[5pt]
56 & \textbf{Xiao 2014 (SD108 reductive-TCA succinate)} \org{P. kudriavzevii}\newline Xiao H, et al., Zhao H. Microb Cell Fact 2014;13:121.\newline \href{https://pmc.ncbi.nlm.nih.gov/articles/PMC4244060}{PMC42440\allowbreak 60} & Host bootstrap: the strain-provenance root of every other engineering row in this host. & no acces\allowbreak sion nam\allowbreak ed; open\allowbreak -\allowbreak access \allowbreak article \\
\addlinespace[5pt]
57 & \textbf{Xylose utilization across four natural backgrounds} \org{P. kudriavzevii}\newline Cas9 metabolic-engineering work placing xylose reductase, xylitol dehydrogenase and xylulokinase from Scheffersomyces stipitis into four I. orientalis backgrounds; primary citation not located.\newline \href{https://core.ac.uk/works/286100561}{core.\allowbreak ac.\allowbreak uk/\allowbreak works/\allowbreak 2861005\allowbreak 61} & Host bootstrap: one construct in four backgrounds is the background-effect design, and it exists nowhere else in this wave. & not loca\allowbreak ted; onl\allowbreak y an agg\allowbreak regator \allowbreak record w\allowbreak as reach\allowbreak ed, so t\allowbreak he prima\allowbreak ry paper\allowbreak  must be\allowbreak  identif\allowbreak ied firs\allowbreak t\newline {\scriptsize (browser)} \\
\addlinespace[5pt]
58 & \textbf{Illinois Data Bank IDB-9063411 (high-yield 3-HP)} \org{P. kudriavzevii}\newline Data for high yield production of 3-hydroxypropionic acid using Issatchenkia orientalis. Illinois Data Bank, updated April 2026.\newline \href{https://databank.illinois.edu/datasets/IDB-9063411}{databank\allowbreak .\allowbreak illinoi\allowbreak s.\allowbreak edu/\allowbreak dataset\allowbreak s/\allowbreak IDB-\allowbreak 9063411} & Host bootstrap: a second product line on the host, with the highest yield it has recorded. & Illinois\allowbreak  Data Ba\allowbreak nk IDB-\allowbreak 9063411\allowbreak ; the pa\allowbreak ge retur\allowbreak ns 403\newline {\scriptsize (browser)} \\
\addlinespace[5pt]
59 & \textbf{Douglass 2018 (CBS573 RNA-seq)} \org{P. kudriavzevii}\newline Douglass AP, et al., Wolfe KH. PLoS Pathog 2018;14:e1007138; SRA SRP139056.\newline \href{https://www.ncbi.nlm.nih.gov/sra/?term=SRP139056}{ncbi.\allowbreak nlm.\allowbreak nih.\allowbreak gov/\allowbreak sra/\allowbreak ?\allowbreak term=\allowbreak SRP1390\allowbreak 56} & Host bootstrap: the only expression data on the host's clean reference, and the anchor for the namespace crosswalk. & SRA SRP1\allowbreak 39056 \\
\addlinespace[5pt]
60 & \textbf{MAESTRO 2025 (gene4260/gene4261 paralog-pair deletion)} \org{P. kudriavzevii}\newline Hsieh PH, et al., Yoshikuni Y. bioRxiv 2025;2025.08.29.673034, CRISPR validation arm.\newline \href{https://doi.org/10.1101/2025.08.29.673034}{doi:\allowbreak 10.\allowbreak 1101/\allowbreak 2025.\allowbreak 08.\allowbreak 29.\allowbreak 673034} & Host bootstrap: a measured gene perturbation in a host that has none at genome scale. & same unr\allowbreak esolved \allowbreak MAESTRO \allowbreak deposit;\allowbreak  needs t\allowbreak he brows\allowbreak er fetch\allowbreak  of the \allowbreak preprint\newline {\scriptsize (ask)} \\
\addlinespace[5pt]
61 & \textbf{Illinois Data Bank IDB-0538107 (succinate pipeline deposit)} \org{P. kudriavzevii}\newline Illinois Data Bank deposit, authors include Huimin Zhao; published March 2026.\newline \href{https://databank.illinois.edu/datasets/IDB-0538107}{databank\allowbreak .\allowbreak illinoi\allowbreak s.\allowbreak edu/\allowbreak dataset\allowbreak s/\allowbreak IDB-\allowbreak 0538107} & Host bootstrap: the host's deposited process data, which the journal pages do not carry. & Illinois\allowbreak  Data Ba\allowbreak nk IDB-\allowbreak 0538107\allowbreak ; the pa\allowbreak ge retur\allowbreak ns 403 t\allowbreak o automa\allowbreak ted clie\allowbreak nts\newline {\scriptsize (browser)} \\
\addlinespace[5pt]
62 & \textbf{CABBI inorganic carbon sources (low-pH succinate kinetics)} \org{P. kudriavzevii}\newline Data for influence of inorganic carbon sources on low-pH succinic acid production by Issatchenkia orientalis: process insights and kinetic analysis. CABBI dataset page.\newline \href{https://cabbi.bio/datasets/}{cabbi.\allowbreak bio/\allowbreak dataset\allowbreak s} & Host bootstrap: a process-condition axis on the host's main product. & CABBI da\allowbreak taset pa\allowbreak ge; the \allowbreak deposit \allowbreak target w\allowbreak as not r\allowbreak ead\newline {\scriptsize (browser)} \\
\addlinespace[5pt]
63 & \textbf{Xylose 3-HP from cellulosic hydrolysates (2026)} \org{P. kudriavzevii}\newline Xylose metabolic engineering of Issatchenkia orientalis for 3-hydroxypropionic acid production. Bioresour Technol 2026.\newline \href{https://doi.org/10.1016/j.biortech.2026.134783}{doi:\allowbreak 10.\allowbreak 1016/\allowbreak j.\allowbreak biortec\allowbreak h.\allowbreak 2026.\allowbreak 134783} & Host bootstrap: a feedstock axis on the host, which the glucose rows cannot supply. & companio\allowbreak n CABBI \allowbreak dataset \allowbreak page; de\allowbreak posit ta\allowbreak rget not\allowbreak  read; n\allowbreak o PMCID\newline {\scriptsize (browser)} \\
\addlinespace[5pt]
64 & \textbf{Jeong 2025 (acetate metabolism during xylose fermentation)} \org{P. kudriavzevii}\newline Jeong D, et al. Acetate metabolism during xylose fermentation enhances 3-HP production in acid-tolerant I. orientalis. bioRxiv 2025;2025.04.17.649204.\newline \href{https://doi.org/10.1101/2025.04.17.649204}{doi:\allowbreak 10.\allowbreak 1101/\allowbreak 2025.\allowbreak 04.\allowbreak 17.\allowbreak 649204} & Host bootstrap: a third carbon arm on the host's second product. & bioRxiv \allowbreak only; ar\allowbreak ticle pa\allowbreak ges retu\allowbreak rn 403 t\allowbreak o automa\allowbreak ted fetc\allowbreak h\newline {\scriptsize (browser)} \\
\addlinespace[5pt]
65 & \textbf{Glycolic acid overproduction (2026)} \org{P. kudriavzevii}\newline Xylose-assisted glycolic acid overproduction in engineered Issatchenkia orientalis. J Fungi 2026;12:751.\newline \href{https://doi.org/10.3390/jof12100751}{doi:\allowbreak 10.\allowbreak 3390/\allowbreak jof1210\allowbreak 0751} & Host bootstrap: a fourth product class on the host, and one of the few rows in it that is openly retrievable. & open-\allowbreak access \allowbreak article \allowbreak (PMC1370\allowbreak 5920); n\allowbreak o separa\allowbreak te depos\allowbreak it named \\
\addlinespace[5pt]
66 & \textbf{Wu 2023 (piggyBac citramalate copy-number and locus variants)} \org{P. kudriavzevii}\newline Wu, et al. Metab Eng Commun 2023;16:e00220.\newline \href{https://pmc.ncbi.nlm.nih.gov/articles/PMC9969067}{PMC99690\allowbreak 67} & Host bootstrap: a copy-number by locus dose axis, which is the only quantitative genotype axis the host offers. & Illinois\allowbreak  Data Ba\allowbreak nk IDB-\allowbreak 6723545\allowbreak  (publis\allowbreak hed Octo\allowbreak ber 2025\allowbreak , contac\allowbreak t Yoshik\allowbreak uni); th\allowbreak e page r\allowbreak eturns 4\allowbreak 03; the \allowbreak article \allowbreak is open \allowbreak access\newline {\scriptsize (browser)} \\
\addlinespace[5pt]
67 & \textbf{Mitochondrial decompartmentalization (2025, with GNPS metabolomics)} \org{P. kudriavzevii}\newline Decompartmentalization of the yeast mitochondrial metabolism to improve chemical production in Issatchenkia orientalis. Nat Commun 2025.\newline \href{https://doi.org/10.1038/s41467-025-62304-w}{doi:\allowbreak 10.\allowbreak 1038/\allowbreak s41467-\allowbreak 025-\allowbreak 62304-\allowbreak w} & Host bootstrap: the host's only metabolomics layer. & a GNPS d\allowbreak epositio\allowbreak n named \allowbreak on the C\allowbreak ABBI dat\allowbreak aset pag\allowbreak e but no\allowbreak t read; \allowbreak a furthe\allowbreak r Illino\allowbreak is Data \allowbreak Bank dep\allowbreak osit sur\allowbreak faces on\allowbreak ly throu\allowbreak gh OSTI \allowbreak as DOI 1\allowbreak 0.\allowbreak 13012/\allowbreak B2IDB-\allowbreak 9575028\_\allowbreak V1\newline {\scriptsize (browser)} \\
\addlinespace[5pt]
68 & \textbf{GSE39558 (ethanol-stress expression)} \org{P. kudriavzevii}\newline GEO series GSE39558, ethanol-stress RNA-seq; species assignment reported by a search index only.\newline \href{https://www.ncbi.nlm.nih.gov/geo/query/acc.cgi?acc=GSE39558}{ncbi.\allowbreak nlm.\allowbreak nih.\allowbreak gov/\allowbreak geo/\allowbreak query/\allowbreak acc.\allowbreak cgi?\allowbreak acc=\allowbreak GSE3955\allowbreak 8} & Host bootstrap: a stress axis the host does not otherwise have in expression, and the cheapest row in the wave to verify. & GEO GSE3\allowbreak 9558; on\allowbreak e call v\allowbreak erifies \allowbreak contents\allowbreak , sample\allowbreak  count a\allowbreak nd strai\allowbreak n \\
\addlinespace[5pt]
69 & \textbf{PXD041801 (flocculin expression and floating biofilms)} \org{P. kudriavzevii}\newline Raw proteomics for a 2023 PLoS Pathogens study of flocculin expression and floating biofilms; submitted via MassIVE.\newline \href{https://www.ebi.ac.uk/pride/archive/projects/PXD041801}{ebi.\allowbreak ac.\allowbreak uk/\allowbreak pride/\allowbreak archive\allowbreak /\allowbreak project\allowbreak s/\allowbreak PXD0418\allowbreak 01} & Host bootstrap: a biofilm axis the host has nowhere else. & Proteome\allowbreak Xchange \allowbreak PXD04180\allowbreak 1 via Ma\allowbreak ssIVE; r\allowbreak etrievab\allowbreak le \\
\addlinespace[5pt]
70 & \textbf{PXD039973 (succinylation and amphotericin B resistance)} \org{P. kudriavzevii}\newline PRIDE/iProX project PXD039973, protein succinylation linked to amphotericin B resistance.\newline \href{https://www.omicsdi.org/dataset/iprox/PXD039973}{omicsdi.\allowbreak org/\allowbreak dataset\allowbreak /\allowbreak iprox/\allowbreak PXD0399\allowbreak 73} & Host bootstrap: the host's only modification-site layer, and the only one tied to a drug phenotype. & PRIDE/\allowbreak iProX P\allowbreak XD039973\allowbreak ; iProX \allowbreak retrieva\allowbreak bility u\allowbreak nverifie\allowbreak d \\
\addlinespace[5pt]
71 & \textbf{PXD075355 (Pichia against Botrytis cinerea)} \org{P. kudriavzevii}\newline PRIDE project PXD075355, native Pichia against Botrytis cinerea, 2026; lists both P. kudriavzevii and P. kluyveri.\newline \href{https://www.ebi.ac.uk/pride/archive/projects/PXD075355}{ebi.\allowbreak ac.\allowbreak uk/\allowbreak pride/\allowbreak archive\allowbreak /\allowbreak project\allowbreak s/\allowbreak PXD0753\allowbreak 55} & Host bootstrap: the only biotic-interaction environment in the whole wave. & PRIDE PX\allowbreak D075355;\allowbreak  retriev\allowbreak able, bu\allowbreak t the de\allowbreak posit ho\allowbreak lds two \allowbreak species \\
\addlinespace[5pt]
72 & \textbf{Pk E1 toolkit expansion (mycoprotein, 2026)} \org{P. kudriavzevii}\newline Expanding the genetic toolkit of Pichia kudriavzevii E1 for enhanced mycoprotein production. Microb Biotechnol 2026;19.\newline \href{https://doi.org/10.1111/1751-7915.70447}{doi:\allowbreak 10.\allowbreak 1111/\allowbreak 1751-\allowbreak 7915.\allowbreak 70447} & Host bootstrap: a second parts panel in a second background, which is what tests whether the first one transfers. & open-\allowbreak access \allowbreak article \allowbreak (PMC1357\allowbreak 8925); n\allowbreak o separa\allowbreak te depos\allowbreak it named \\
\addlinespace[5pt]
73 & \textbf{China fermented-vegetable and clinical panel (2024)} \org{P. kudriavzevii}\newline Comparative genomics and susceptibility of fermented-vegetable and clinical Candida krusei isolates. Virulence 2024;15.\newline \href{https://www.ncbi.nlm.nih.gov/pmc/articles/PMC11487970/}{PMC11487\allowbreak 970} & Host bootstrap: the host's only second clinical panel, held because one supplement check decides whether it is a dataset. & open-\allowbreak access \allowbreak article;\allowbreak  per-\allowbreak isolate\allowbreak  MIC tab\allowbreak les and \allowbreak genome a\allowbreak ccession\allowbreak s not ve\allowbreak rified\newline {\scriptsize (browser)} \\
\addlinespace[5pt]
74 & \textbf{Indonesian clinical genomes (PRJNA819536)} \org{P. kudriavzevii}\newline Genome note on whole-genome sequencing of 16 clinical isolates from Jakarta hospitals. F1000Research 2022;11:593.\newline \href{https://www.ncbi.nlm.nih.gov/pmc/articles/PMC10320322/}{PMC10320\allowbreak 322} & Host bootstrap: it would add half again as many sequenced isolates to the host's only genotyped panel. & BioProje\allowbreak ct PRJNA\allowbreak 819536 f\allowbreak or the s\allowbreak equence;\allowbreak  the phe\allowbreak notype t\allowbreak able's l\allowbreak ocation \allowbreak is unkno\allowbreak wn and m\allowbreak ay be an\allowbreak  ask\newline {\scriptsize (browser)} \\
\addlinespace[5pt]
75 & \textbf{Unidentified Illinois Data Bank CRISPR/Cas9 gene-deletion deposit} \org{P. kudriavzevii}\newline An Illinois Data Bank CRISPR/Cas9 gene-deletion dataset surfacing only through an OSTI bibliographic record; deposit identifier not confirmed.\newline \href{https://www.osti.gov/biblio/3013685}{osti.\allowbreak gov/\allowbreak biblio/\allowbreak 3013685} & Host bootstrap: the one lead that could overturn the finding that this host has no multi-gene deletion resource. & an Illin\allowbreak ois Data\allowbreak  Bank de\allowbreak posit wh\allowbreak ose iden\allowbreak tifier i\allowbreak s unconf\allowbreak irmed; n\allowbreak eeds an \allowbreak OSTI and\allowbreak  a repos\allowbreak itory br\allowbreak owser vi\allowbreak sit\newline {\scriptsize (browser)} \\
\addlinespace[5pt]
76 & \textbf{Khaledi 2020 (414-isolate clinical transcriptome panel)} \org{P. aeruginosa}\newline Khaledi A, Weimann A, Schniederjans M, Asgari E, Kuo TH, Oliver A, Cabot G, Kola A, Gastmeier P, Hogardt M, Jonas D, Mofrad MRK, Bremges A, McHardy AC, Haussler S. EMBO Mol Med 2020;12:e10264.\newline \href{https://doi.org/10.15252/emmm.201910264}{doi:\allowbreak 10.\allowbreak 15252/\allowbreak emmm.\allowbreak 2019102\allowbreak 64} & Scale band: tier 1, ordered by measurements (2,442,600). & GEO GSE1\allowbreak 23544 (R\allowbreak NA-\allowbreak seq); B\allowbreak ioProjec\allowbreak t PRJNA5\allowbreak 26797 (D\allowbreak NA-\allowbreak seq); Z\allowbreak enodo 10\allowbreak .\allowbreak 5281/\allowbreak zenodo.\allowbreak 3464542\allowbreak  (classi\allowbreak fier inp\allowbreak uts); re\allowbreak ference \allowbreak NC\_\allowbreak 008463.\allowbreak 1 \\
\addlinespace[5pt]
77 & \textbf{Kaczmarczyk 2025 (PAO1 genome-wide CRISPRi)} \org{P. aeruginosa}\newline Kaczmarczyk A, Klotz A, Manfredi P, Jenal U. bioRxiv 2025;2025.08.12.669819.\newline \href{https://doi.org/10.1101/2025.08.12.669819}{doi:\allowbreak 10.\allowbreak 1101/\allowbreak 2025.\allowbreak 08.\allowbreak 12.\allowbreak 669819} & Scale band: tier 2, ordered by measurements (80,000). & Zenodo 1\allowbreak 0.\allowbreak 5281/\allowbreak zenodo.\allowbreak 1559174\allowbreak 2, .\allowbreak 1559223\allowbreak 8 and .\allowbreak 1559336\allowbreak 1, conte\allowbreak nts unla\allowbreak beled; n\allowbreak o SRA or\allowbreak  GEO acc\allowbreak ession\newline {\scriptsize (ask)} \\
\addlinespace[5pt]
78 & \textbf{Dotsch 2009 (arrayed PA14 library against 19 antibiotics)} \org{P. aeruginosa}\newline Dotsch A, Becker T, Pommerenke C, Janowska Z, Jansch L, Haussler S. Antimicrob Agents Chemother 2009;53:2522-2531.\newline \href{https://doi.org/10.1128/AAC.00035-09}{doi:\allowbreak 10.\allowbreak 1128/\allowbreak AAC.\allowbreak 00035-\allowbreak 09} & Scale band: tier 2, ordered by measurements (5,459). & value lo\allowbreak cation u\allowbreak nconfirm\allowbreak ed; the \allowbreak PA14 non\allowbreak -\allowbreak redunda\allowbreak nt libra\allowbreak ry backb\allowbreak one is a\allowbreak t the PA\allowbreak TIMDB po\allowbreak rtal, wh\allowbreak ose down\allowbreak loads pa\allowbreak ge was n\allowbreak ot read\newline {\scriptsize (ask)} \\
\addlinespace[5pt]
79 & \textbf{Bartell 2019 (longitudinal cystic fibrosis panel)} \org{P. aeruginosa}\newline Bartell JA, Sommer LM, Haagensen JAJ, Loch A, Espinosa R, Molin S, Johansen HK. Nat Commun 2019;10:629.\newline \href{https://doi.org/10.1038/s41467-019-08504-7}{doi:\allowbreak 10.\allowbreak 1038/\allowbreak s41467-\allowbreak 019-\allowbreak 08504-\allowbreak 7} & Scale band: tier 2, ordered by measurements (3,544). & Suppleme\allowbreak ntary Da\allowbreak ta 1 (xl\allowbreak sx) hold\allowbreak s the en\allowbreak tire dat\allowbreak aset and\allowbreak  the iso\allowbreak late acc\allowbreak essions;\allowbreak  genomes\allowbreak  at SRA \allowbreak study ER\allowbreak P004853\newline {\scriptsize (browser)} \\
\addlinespace[5pt]
80 & \textbf{Genome Medicine 2024 (1,877-strain genome and MIC set)} \org{P. aeruginosa}\newline Paired whole-genome and susceptibility data for 1,877 diverse strains plus a 102-strain validation set, with MICs reclassified per CLSI M100S-Ed32:2022. Genome Med 2024;16.\newline \href{https://doi.org/10.1186/s13073-024-01346-z}{doi:\allowbreak 10.\allowbreak 1186/\allowbreak s13073-\allowbreak 024-\allowbreak 01346-\allowbreak z} & Scale band: tier 2, ordered by measurements (1,979). & unconfir\allowbreak med; the\allowbreak  publish\allowbreak er redir\allowbreak ects to \allowbreak an authe\allowbreak nticatio\allowbreak n endpoi\allowbreak nt\newline {\scriptsize (browser)} \\
\addlinespace[5pt]
81 & \textbf{Penaranda 2025 (125-isolate diversity survey)} \org{P. aeruginosa}\newline Penaranda C, Brenner EP, Clatworthy AE, Cosimi LA, Ravi J, Hung DT. bioRxiv 2025;2025.09.22.677895.\newline \href{https://doi.org/10.1101/2025.09.22.677895}{doi:\allowbreak 10.\allowbreak 1101/\allowbreak 2025.\allowbreak 09.\allowbreak 22.\allowbreak 677895} & Scale band: tier 2, ordered by measurements (1,664). & Suppleme\allowbreak ntal Dat\allowbreak a 1 (xls\allowbreak x) holds\allowbreak  genome \allowbreak accessio\allowbreak ns, meta\allowbreak data and\allowbreak  phenoty\allowbreak pes; the\allowbreak  BioProj\allowbreak ect numb\allowbreak er is in\allowbreak side tha\allowbreak t file a\allowbreak nd uncon\allowbreak firmed\newline {\scriptsize (browser)} \\
\addlinespace[5pt]
82 & \textbf{Khaledi 2020 (1,656-record agar-dilution MIC panel)} \org{P. aeruginosa}\newline Khaledi A, et al., Haussler S. EMBO Mol Med 2020;12:e10264, susceptibility layer.\newline \href{https://doi.org/10.15252/emmm.201910264}{doi:\allowbreak 10.\allowbreak 15252/\allowbreak emmm.\allowbreak 2019102\allowbreak 64} & Scale band: tier 2, ordered by measurements (1,656). & phenotyp\allowbreak e table \allowbreak as Datas\allowbreak et EV1; \allowbreak Zenodo 1\allowbreak 0.\allowbreak 5281/\allowbreak zenodo.\allowbreak 3464542\allowbreak  for the\allowbreak  classif\allowbreak ier inpu\allowbreak ts \\
\addlinespace[5pt]
83 & \textbf{Poulsen 2019 (nine-strain Tn-seq core essentiality)} \org{P. aeruginosa}\newline Poulsen BE, et al. Defining the core essential genome of Pseudomonas aeruginosa. bioRxiv 2019;396689.\newline \href{https://www.biorxiv.org/content/10.1101/396689v2.full.pdf}{biorxiv.\allowbreak org/\allowbreak content\allowbreak /\allowbreak 10.\allowbreak 1101/\allowbreak 396689v\allowbreak 2.\allowbreak full.\allowbreak pdf} & Scale band: tier 2, ordered by measurements (321). & not give\allowbreak n by the\allowbreak  compili\allowbreak ng paper\allowbreak ; the pr\allowbreak eprint i\allowbreak s open\newline {\scriptsize (browser)} \\
\addlinespace[5pt]
84 & \textbf{PAO1 arrayed transposon libraries (Jacobs 2003, Held 2012)} \org{P. aeruginosa}\newline Jacobs MA, et al. Proc Natl Acad Sci USA 2003;100:14339-14344; Held K, Ramage E, Jacobs M, Gallagher L, Manoil C. J Bacteriol 2012;194.\newline \href{https://doi.org/10.1128/JB.01479-12}{doi:\allowbreak 10.\allowbreak 1128/\allowbreak JB.\allowbreak 01479-\allowbreak 12} & Scale band: tier 3, ordered by measurements (9,437). & insertio\allowbreak n-\allowbreak positio\allowbreak n metada\allowbreak ta from \allowbreak the libr\allowbreak ary inde\allowbreak x page; \allowbreak strains \allowbreak are orde\allowbreak r-\allowbreak only th\allowbreak rough a \allowbreak cost cen\allowbreak ter and \allowbreak distribu\allowbreak tion has\allowbreak  moved l\allowbreak abs\newline {\scriptsize (browser)} \\
\addlinespace[5pt]
85 & \textbf{KT2440 pooled transposon library (Putida\_ML5 metadata)} \org{P. putida}\newline Pooled RB-TnSeq library metadata for Pseudomonas putida KT2440, one mapped version named Putida\_ML5.\newline \href{https://nrel.gov/docs/fy22osti/83183.pdf}{nrel.\allowbreak gov/\allowbreak docs/\allowbreak fy22ost\allowbreak i/\allowbreak 83183.\allowbreak pdf} & Scale band: tier 3, ordered by measurements (4,778). & NREL met\allowbreak hods doc\allowbreak umentati\allowbreak on plus \allowbreak the mapp\allowbreak ing arch\allowbreak ives shi\allowbreak pped wit\allowbreak h the bi\allowbreak gfit rel\allowbreak ease \\
\addlinespace[5pt]
86 & \textbf{Gao 2020 (KT2440 targeted SRM proteomics)} \org{P. putida}\newline Gao Y, Fillmore TL, Munoz N, Bentley GJ, Johnson CW, et al. Front Bioeng Biotechnol 2020;8:603488.\newline \href{https://doi.org/10.3389/fbioe.2020.603488}{doi:\allowbreak 10.\allowbreak 3389/\allowbreak fbioe.\allowbreak 2020.\allowbreak 603488} & Scale band: tier 3, ordered by measurements (1,056). & Panorama\allowbreak  Public,\allowbreak  no PXD \allowbreak accessio\allowbreak n; monit\allowbreak ored pep\allowbreak tides in\allowbreak  Supplem\allowbreak entary T\allowbreak ables S1\allowbreak  and S2,\allowbreak  limits \allowbreak of quant\allowbreak ificatio\allowbreak n in S3 \\
\addlinespace[5pt]
87 & \textbf{BACTOME (101 clinical genomes with eight phenotypes)} \org{P. aeruginosa}\newline Hornischer K, et al. BACTOME: a reference database to explore the sequence- and gene expression-variation landscape of Pseudomonas aeruginosa clinical isolates. Nucleic Acids Res 2019;47:D716.\newline \href{https://doi.org/10.1093/nar/gky895}{doi:\allowbreak 10.\allowbreak 1093/\allowbreak nar/\allowbreak gky895} & Scale band: tier 3, ordered by measurements (808). & no bulk \allowbreak download\allowbreak  and no \allowbreak ENA, SRA\allowbreak  or GEO \allowbreak accessio\allowbreak n on the\allowbreak  landing\allowbreak  page; n\allowbreak ot confi\allowbreak rmed bul\allowbreak k-\allowbreak retriev\allowbreak able\newline {\scriptsize (ask)} \\
\addlinespace[5pt]
88 & \textbf{Kohlstedt 2019 (13C metabolic flux analysis, two hosts)} \org{P. putida}\newline Kohlstedt M, Wittmann C. Metab Eng 2019;54:35-53.\newline \href{https://doi.org/10.1016/j.ymben.2019.01.008}{doi:\allowbreak 10.\allowbreak 1016/\allowbreak j.\allowbreak ymben.\allowbreak 2019.\allowbreak 01.\allowbreak 008} & Scale band: tier 3, ordered by measurements (534). & unconfir\allowbreak med:\allowbreak  the pu\allowbreak blisher \allowbreak returned\allowbreak  403 and\allowbreak  the ins\allowbreak titution\allowbreak al PDF r\allowbreak eset the\allowbreak  connect\allowbreak ion\newline {\scriptsize (browser)} \\
\addlinespace[5pt]
89 & \textbf{Turner 2015 (PA14 and PAO1 essentiality across four arms)} \org{P. aeruginosa}\newline Turner KH, Wessel AK, Palmer GC, Murray JL, Whiteley M. Proc Natl Acad Sci USA 2015.\newline \href{https://pmc.ncbi.nlm.nih.gov/articles/PMC4386324}{PMC43863\allowbreak 24} & Scale band: tier 3, ordered by measurements (352). & SRA PRJN\allowbreak A265367 \\
\addlinespace[5pt]
90 & \textbf{Liberati 2006 (PA14 non-redundant transposon library and essentiality)} \org{P. aeruginosa}\newline Liberati NT, Urbach JM, Miyata S, Lee DG, Drenkard E, Wu G, Villanueva J, Wei T, Ausubel FM. Proc Natl Acad Sci USA 2006;103:2833-2838.\newline \href{https://doi.org/10.1073/pnas.0511100103}{doi:\allowbreak 10.\allowbreak 1073/\allowbreak pnas.\allowbreak 0511100\allowbreak 103} & Scale band: tier 3, ordered by measurements (335). & the PATI\allowbreak MDB port\allowbreak al, with\allowbreak  Search,\allowbreak  BLAST a\allowbreak nd Inser\allowbreak tion Sit\allowbreak e Map pa\allowbreak ges; whe\allowbreak ther per\allowbreak -\allowbreak mutant \allowbreak tables a\allowbreak re bulk-\allowbreak downloa\allowbreak dable is\allowbreak  unconfi\allowbreak rmed\newline {\scriptsize (browser)} \\
\addlinespace[5pt]
91 & \textbf{GOLD\_115 and GOLD\_84 essential-gene sets} \org{P. aeruginosa}\newline Van Maele C, Caboche S, Nicolau-Guillaumet N, Muggeo A, Guillard T. PLoS Comput Biol 2026;22:e1013945.\newline \href{https://doi.org/10.1371/journal.pcbi.1013945}{doi:\allowbreak 10.\allowbreak 1371/\allowbreak journal\allowbreak .\allowbreak pcbi.\allowbreak 1013945} & Scale band: tier 3, ordered by measurements (199). & S2 Table\allowbreak , Gold S\allowbreak tandard \allowbreak List; th\allowbreak e paper'\allowbreak s own re\allowbreak ads at S\allowbreak RA PRJNA\allowbreak 1240204;\allowbreak  referen\allowbreak ce NC\_\allowbreak 008463.\allowbreak 1\newline {\scriptsize (browser)} \\
\addlinespace[5pt]
92 & \textbf{Werner 2023 (beta-ketoadipate from lignin)} \org{P. putida}\newline Werner AZ, Cordell WT, Lahive CW, Klein BC, Singer CA, Tan ECD, et al., Beckham GT. Sci Adv 2023;9:eadj0053.\newline \href{https://doi.org/10.1126/sciadv.adj0053}{doi:\allowbreak 10.\allowbreak 1126/\allowbreak sciadv.\allowbreak adj0053} & Scale band: tier 3, ordered by measurements (90). & sciadv.\allowbreak adj0053\_\allowbreak data\_\allowbreak s1.\allowbreak zip plu\allowbreak s the SI\allowbreak  PDF; no\allowbreak  reposit\allowbreak ory; pla\allowbreak smids de\allowbreak ferred t\allowbreak o Tables\allowbreak  S1 to S\allowbreak 3\newline {\scriptsize (browser)} \\
\addlinespace[5pt]
93 & \textbf{Wilkes 2026 (muconate multi-omics with fluxomics)} \org{P. putida}\newline Wilkes RA, Suthers PF, Borchert AJ, Callaghan MM, Thusoo E, Giannone RJ, et al. Metab Eng 2026;99:102573.\newline \href{https://doi.org/10.1016/j.ymben.2026.102573}{doi:\allowbreak 10.\allowbreak 1016/\allowbreak j.\allowbreak ymben.\allowbreak 2026.\allowbreak 102573} & Scale band: tier 3, ordered by measurements (12). & none con\allowbreak firmed; \allowbreak the bibl\allowbreak iographi\allowbreak c record\allowbreak  reports\allowbreak  no data\allowbreak  and no \allowbreak suppleme\allowbreak ntary fi\allowbreak les\newline {\scriptsize (browser)} \\
\addlinespace[5pt]
94 & \textbf{GSE213139 (low-dose ionizing radiation)} \org{P. putida}\newline GEO series GSE213139, low-dose ionizing radiation, released 2025-03-31.\newline \href{https://www.ncbi.nlm.nih.gov/geo/query/acc.cgi?acc=GSE213139}{ncbi.\allowbreak nlm.\allowbreak nih.\allowbreak gov/\allowbreak geo/\allowbreak query/\allowbreak acc.\allowbreak cgi?\allowbreak acc=\allowbreak GSE2131\allowbreak 39} & Scale band: tier 4, ordered by measurements (166,920). & GEO GSE2\allowbreak 13139; p\allowbreak rogramma\allowbreak tically \allowbreak retrieva\allowbreak ble \\
\addlinespace[5pt]
95 & \textbf{GSE131229 (non-antibiotic pharmaceutical stress)} \org{P. putida}\newline GEO SuperSeries GSE131229, non-antibiotic pharmaceutical stress.\newline \href{https://www.ncbi.nlm.nih.gov/geo/query/acc.cgi?acc=GSE131229}{ncbi.\allowbreak nlm.\allowbreak nih.\allowbreak gov/\allowbreak geo/\allowbreak query/\allowbreak acc.\allowbreak cgi?\allowbreak acc=\allowbreak GSE1312\allowbreak 29} & Scale band: tier 4, ordered by measurements (166,920). & GEO GSE1\allowbreak 31229; p\allowbreak rogramma\allowbreak tically \allowbreak retrieva\allowbreak ble \\
\addlinespace[5pt]
96 & \textbf{GSE248909 (micro and nanoplastic stress)} \org{P. putida}\newline GEO series GSE248909, micro and nanoplastic stress.\newline \href{https://www.ncbi.nlm.nih.gov/geo/query/acc.cgi?acc=GSE248909}{ncbi.\allowbreak nlm.\allowbreak nih.\allowbreak gov/\allowbreak geo/\allowbreak query/\allowbreak acc.\allowbreak cgi?\allowbreak acc=\allowbreak GSE2489\allowbreak 09} & Scale band: tier 4, ordered by measurements (155,792). & GEO GSE2\allowbreak 48909; p\allowbreak rogramma\allowbreak tically \allowbreak retrieva\allowbreak ble \\
\addlinespace[5pt]
97 & \textbf{GSE193493 (putidaPRECISE321 second deposit)} \org{P. putida}\newline GEO series GSE193493, part of putidaPRECISE321, released 2022-04-28.\newline \href{https://www.ncbi.nlm.nih.gov/geo/query/acc.cgi?acc=GSE193493}{ncbi.\allowbreak nlm.\allowbreak nih.\allowbreak gov/\allowbreak geo/\allowbreak query/\allowbreak acc.\allowbreak cgi?\allowbreak acc=\allowbreak GSE1934\allowbreak 93} & Scale band: tier 4, ordered by measurements (122,408). & GEO GSE1\allowbreak 93493; p\allowbreak rogramma\allowbreak tically \allowbreak retrieva\allowbreak ble \\
\addlinespace[5pt]
98 & \textbf{GSE221044 (kaiC deletion)} \org{P. putida}\newline GEO series GSE221044, a kaiC deletion background.\newline \href{https://www.ncbi.nlm.nih.gov/geo/query/acc.cgi?acc=GSE221044}{ncbi.\allowbreak nlm.\allowbreak nih.\allowbreak gov/\allowbreak geo/\allowbreak query/\allowbreak acc.\allowbreak cgi?\allowbreak acc=\allowbreak GSE2210\allowbreak 44} & Scale band: tier 4, ordered by measurements (122,408). & GEO GSE2\allowbreak 21044; p\allowbreak rogramma\allowbreak tically \allowbreak retrieva\allowbreak ble \\
\addlinespace[5pt]
99 & \textbf{GSE144045 (solvent-tolerance evolution in P. putida S12)} \org{P. putida}\newline GEO series GSE144045, solvent-tolerance adaptive evolution in Pseudomonas putida S12.\newline \href{https://www.ncbi.nlm.nih.gov/geo/query/acc.cgi?acc=GSE144045}{ncbi.\allowbreak nlm.\allowbreak nih.\allowbreak gov/\allowbreak geo/\allowbreak query/\allowbreak acc.\allowbreak cgi?\allowbreak acc=\allowbreak GSE1440\allowbreak 45} & Scale band: tier 4, ordered by measurements (100,152). & GEO GSE1\allowbreak 44045; p\allowbreak rogramma\allowbreak tically \allowbreak retrieva\allowbreak ble \\
\addlinespace[5pt]
100 & \textbf{GSE275525 (HexR regulatory landscape)} \org{P. putida}\newline GEO series GSE275525, the HexR regulatory landscape, released 2026-08-20.\newline \href{https://www.ncbi.nlm.nih.gov/geo/query/acc.cgi?acc=GSE275525}{ncbi.\allowbreak nlm.\allowbreak nih.\allowbreak gov/\allowbreak geo/\allowbreak query/\allowbreak acc.\allowbreak cgi?\allowbreak acc=\allowbreak GSE2755\allowbreak 25} & Scale band: tier 4, ordered by measurements (89,024). & GEO GSE2\allowbreak 75525; p\allowbreak rogramma\allowbreak tically \allowbreak retrieva\allowbreak ble \\
\addlinespace[5pt]
101 & \textbf{GSE309641 (fluoride adaptive evolution)} \org{P. putida}\newline GEO series GSE309641, fluoride adaptive evolution, released 2026-02-26.\newline \href{https://www.ncbi.nlm.nih.gov/geo/query/acc.cgi?acc=GSE309641}{ncbi.\allowbreak nlm.\allowbreak nih.\allowbreak gov/\allowbreak geo/\allowbreak query/\allowbreak acc.\allowbreak cgi?\allowbreak acc=\allowbreak GSE3096\allowbreak 41} & Scale band: tier 4, ordered by measurements (89,024). & GEO GSE3\allowbreak 09641; p\allowbreak rogramma\allowbreak tically \allowbreak retrieva\allowbreak ble \\
\addlinespace[5pt]
102 & \textbf{GSE155767 (evolved KT2440 on non-native sugars)} \org{P. putida}\newline GEO series GSE155767, evolved KT2440 on non-native sugars.\newline \href{https://www.ncbi.nlm.nih.gov/geo/query/acc.cgi?acc=GSE155767}{ncbi.\allowbreak nlm.\allowbreak nih.\allowbreak gov/\allowbreak geo/\allowbreak query/\allowbreak acc.\allowbreak cgi?\allowbreak acc=\allowbreak GSE1557\allowbreak 67} & Scale band: tier 4, ordered by measurements (77,896). & GEO GSE1\allowbreak 55767; p\allowbreak rogramma\allowbreak tically \allowbreak retrieva\allowbreak ble \\
\addlinespace[5pt]
103 & \textbf{GSE147420 (ammonium against nitrate)} \org{P. putida}\newline GEO series GSE147420, ammonium against nitrate.\newline \href{https://www.ncbi.nlm.nih.gov/geo/query/acc.cgi?acc=GSE147420}{ncbi.\allowbreak nlm.\allowbreak nih.\allowbreak gov/\allowbreak geo/\allowbreak query/\allowbreak acc.\allowbreak cgi?\allowbreak acc=\allowbreak GSE1474\allowbreak 20} & Scale band: tier 4, ordered by measurements (77,896). & GEO GSE1\allowbreak 47420; p\allowbreak rogramma\allowbreak tically \allowbreak retrieva\allowbreak ble \\
\addlinespace[5pt]
104 & \textbf{GSE229763 (one-carbon metabolism)} \org{P. putida}\newline GEO series GSE229763, one-carbon metabolism.\newline \href{https://www.ncbi.nlm.nih.gov/geo/query/acc.cgi?acc=GSE229763}{ncbi.\allowbreak nlm.\allowbreak nih.\allowbreak gov/\allowbreak geo/\allowbreak query/\allowbreak acc.\allowbreak cgi?\allowbreak acc=\allowbreak GSE2297\allowbreak 63} & Scale band: tier 4, ordered by measurements (66,768). & GEO GSE2\allowbreak 29763; p\allowbreak rogramma\allowbreak tically \allowbreak retrieva\allowbreak ble \\
\addlinespace[5pt]
105 & \textbf{KT2440 aromatics RB-TnSeq (2019 preprint)} \org{P. putida}\newline RB-TnSeq of Pseudomonas putida KT2440 on coumarate and seven related aromatics. bioRxiv 2019;753889.\newline \href{https://www.biorxiv.org/content/10.1101/753889v1.full-text}{biorxiv.\allowbreak org/\allowbreak content\allowbreak /\allowbreak 10.\allowbreak 1101/\allowbreak 753889v\allowbreak 1.\allowbreak full-\allowbreak text} & Scale band: tier 4, ordered by measurements (44,512). & not repo\allowbreak rted \\
\addlinespace[5pt]
106 & \textbf{PMkbase v1.0 (phenotype-microarray database)} \org{P. putida}\newline PMkbase v1.0, a phenotype-microarray knowledge base made interoperable with sequence information, 2025.\newline \href{https://www.ncbi.nlm.nih.gov/pmc/articles/PMC12584640/}{PMC12584\allowbreak 640} & Scale band: tier 4, ordered by measurements (41,664). & unconfir\allowbreak med:\allowbreak  the ar\allowbreak ticle pa\allowbreak ge serve\allowbreak d a chal\allowbreak lenge, s\allowbreak o no ide\allowbreak ntifier,\allowbreak  strain \allowbreak list, re\allowbreak source a\allowbreak ddress o\allowbreak r downlo\allowbreak ad forma\allowbreak t was re\allowbreak ad\newline {\scriptsize (browser)} \\
\addlinespace[5pt]
107 & \textbf{IPCD (International P. aeruginosa Consortium Database)} \org{P. aeruginosa}\newline Freschi L, Jeukens J, Kukavica-Ibrulj I, Boyle B, Dupont MJ, Laroche J, et al., Levesque RC. Front Microbiol 2015;6:1036.\newline \href{https://doi.org/10.3389/fmicb.2015.01036}{doi:\allowbreak 10.\allowbreak 3389/\allowbreak fmicb.\allowbreak 2015.\allowbreak 01036} & Scale band: tier 4, ordered by measurements (3,038). & no adver\allowbreak tised bu\allowbreak lk downl\allowbreak oad; sig\allowbreak n-\allowbreak in requ\allowbreak ired; no\allowbreak t confir\allowbreak med publ\allowbreak icly bul\allowbreak k-\allowbreak downloa\allowbreak dable\newline {\scriptsize (ask)} \\
\addlinespace[5pt]
108 & \textbf{Pirnay 2009 (328-isolate population survey)} \org{P. aeruginosa}\newline Pirnay JP, Bilocq F, Pot B, Cornelis P, Zizi M, Van Eldere J, et al. PLoS ONE 2009;4:e7740.\newline \href{https://doi.org/10.1371/journal.pone.0007740}{doi:\allowbreak 10.\allowbreak 1371/\allowbreak journal\allowbreak .\allowbreak pone.\allowbreak 0007740} & Scale band: tier 4, ordered by measurements (2,624). & open-\allowbreak access \allowbreak article;\allowbreak  whether\allowbreak  the 8 p\allowbreak arameter\allowbreak s are ph\allowbreak enotypes\allowbreak  is unre\allowbreak solved \\
\addlinespace[5pt]
109 & \textbf{n-butanol KT2440 proteome (2015)} \org{P. putida}\newline Analysis of the molecular response of Pseudomonas putida KT2440 to the next-generation biofuel n-butanol. J Proteomics 2015.\newline \href{https://biochemie.uni-hohenheim.de/en/organization/publication/analysis-of-the-molecular-response-of-pseudomonas-putida-kt2440-to-the-next-generation-biofuel-n-butanol}{biochemi\allowbreak e.\allowbreak uni-\allowbreak hohenhe\allowbreak im.\allowbreak de/\allowbreak en/\allowbreak organiz\allowbreak ation/\allowbreak publica\allowbreak tion/\allowbreak analysi\allowbreak s-\allowbreak of-\allowbreak the-\allowbreak molecul\allowbreak ar-\allowbreak respons\allowbreak e-\allowbreak of-\allowbreak pseudom\allowbreak onas-\allowbreak putida-\allowbreak kt2440-\allowbreak to-\allowbreak the-\allowbreak next-\allowbreak generat\allowbreak ion-\allowbreak biofuel\allowbreak -\allowbreak n-\allowbreak butanol} & Scale band: tier 4, ordered by measurements (1,467). & not repo\allowbreak rted \\
\addlinespace[5pt]
110 & \textbf{Skurnik 2013 (PA14 essentiality)} \org{P. aeruginosa}\newline Skurnik D, et al. Transposon sequencing of Pseudomonas aeruginosa PA14. 2013; full citation not captured.\newline \href{https://doi.org/10.1371/journal.pcbi.1013945}{doi:\allowbreak 10.\allowbreak 1371/\allowbreak journal\allowbreak .\allowbreak pcbi.\allowbreak 1013945} & Scale band: tier 4, ordered by measurements (634). & not give\allowbreak n by the\allowbreak  compili\allowbreak ng paper\allowbreak ; the ac\allowbreak cession \allowbreak must be \allowbreak chased f\allowbreak rom the \allowbreak primary \allowbreak study\newline {\scriptsize (browser)} \\
\addlinespace[5pt]
111 & \textbf{Lee 2015 (PAO1 essentiality across three media)} \org{P. aeruginosa}\newline Lee SA, et al. Transposon sequencing of Pseudomonas aeruginosa PAO1 across three media. 2015; full citation not captured.\newline \href{https://doi.org/10.1371/journal.pcbi.1013945}{doi:\allowbreak 10.\allowbreak 1371/\allowbreak journal\allowbreak .\allowbreak pcbi.\allowbreak 1013945} & Scale band: tier 4, ordered by measurements (352). & not give\allowbreak n by the\allowbreak  compili\allowbreak ng paper\allowbreak ; the ac\allowbreak cession \allowbreak must be \allowbreak chased f\allowbreak rom the \allowbreak primary \allowbreak study\newline {\scriptsize (browser)} \\
\addlinespace[5pt]
112 & \textbf{P. putida S12 phenotype microarray (190 carbon sources)} \org{P. putida}\newline Genome-scale model validation against a 190-carbon-source phenotype microarray in Pseudomonas putida S12. BMC Genomics 2023;24.\newline \href{https://doi.org/10.1186/s12864-023-09940-y}{doi:\allowbreak 10.\allowbreak 1186/\allowbreak s12864-\allowbreak 023-\allowbreak 09940-\allowbreak y} & Scale band: tier 4, ordered by measurements (190). & open-\allowbreak access \allowbreak article;\allowbreak  deposit\allowbreak  and for\allowbreak mat not \allowbreak reported \\
\addlinespace[5pt]
113 & \textbf{Chandler 2023 (128 functionally phenotyped isolates)} \org{P. aeruginosa}\newline Chandler CE, et al. Microbiol Spectr 2023;11:e01556-23.\newline \href{https://doi.org/10.1128/spectrum.01556-23}{doi:\allowbreak 10.\allowbreak 1128/\allowbreak spectru\allowbreak m.\allowbreak 01556-\allowbreak 23} & Scale band: tier 4, ordered by measurements (128). & unconfir\allowbreak med\newline {\scriptsize (browser)} \\
\addlinespace[5pt]
114 & \textbf{Breidenstein 2008 (ciprofloxacin resistome)} \org{P. aeruginosa}\newline Breidenstein EBM, Khaira BK, Wiegand I, Overhage J, Hancock REW. Antimicrob Agents Chemother 2008;52:4486-4491.\newline \href{https://doi.org/10.1128/AAC.00222-08}{doi:\allowbreak 10.\allowbreak 1128/\allowbreak AAC.\allowbreak 00222-\allowbreak 08} & Scale band: tier 4, ordered by measurements (114). & per-\allowbreak mutant \allowbreak values i\allowbreak n Tables\allowbreak  1 and 2\allowbreak  of the \allowbreak article;\allowbreak  no repo\allowbreak sitory \\
\addlinespace[5pt]
115 & \textbf{Remaining P. putida GEO expression series} \org{P. putida}\newline 118 GEO series are returned for this organism; 40 identifiers were enumerated, 20 of them with 12 or more samples, and 78 series remain unexamined.\newline \href{https://www.ncbi.nlm.nih.gov/gds/?term=\%22Pseudomonas+putida\%22\%5BOrganism\%5D}{ncbi.\allowbreak nlm.\allowbreak nih.\allowbreak gov/\allowbreak gds/\allowbreak ?\allowbreak term=\allowbreak \%\allowbreak 22Pseud\allowbreak omonas+\allowbreak putida\%\allowbreak 22\%\allowbreak 5BOrgan\allowbreak ism\%\allowbreak 5D} & Scale band: tier 4, ordered by measurements (78). & GEO, by \allowbreak organism\allowbreak  query; \allowbreak 78 of 11\allowbreak 8 series\allowbreak  unexami\allowbreak ned \\
\addlinespace[5pt]
116 & \textbf{KT2440 mini-Tn5 knockout collection (2010)} \org{P. putida}\newline Genome-wide mini-Tn5 knockout collection of Pseudomonas putida KT2440 screened on minimal glucose medium; the cited record should be re-checked against the intended work.\newline \href{https://pmc.ncbi.nlm.nih.gov/articles/PMC5814926}{PMC58149\allowbreak 26} & Scale band: tier 4, ordered by measurements (68). & not repo\allowbreak rted; th\allowbreak e cited \allowbreak identifi\allowbreak er is fl\allowbreak agged fo\allowbreak r re-\allowbreak checkin\allowbreak g \\
\addlinespace[5pt]
117 & \textbf{Nunez-Garcia 2024 (39 isolates over ten years)} \org{P. aeruginosa}\newline Nunez-Garcia L, et al. Microbiol Spectr 2024;12.\newline \href{https://pmc.ncbi.nlm.nih.gov/articles/PMC11619361/}{PMC11619\allowbreak 361} & Scale band: tier 4, ordered by measurements (39). & unconfir\allowbreak med\newline {\scriptsize (browser)} \\
\addlinespace[5pt]
118 & \textbf{Remaining P. putida PRIDE proteomics projects} \org{P. putida}\newline About 35 PRIDE projects returned by a keyword search of the PRIDE v3 REST API; 26 were read and named, 7 remain unverified, and no second pass was made on the strain-name keywords.\newline \href{https://www.ebi.ac.uk/pride/ws/archive/v3/search/projects?keyword=Pseudomonas\%20putida}{ebi.\allowbreak ac.\allowbreak uk/\allowbreak pride/\allowbreak ws/\allowbreak archive\allowbreak /\allowbreak v3/\allowbreak search/\allowbreak project\allowbreak s?\allowbreak keyword\allowbreak =\allowbreak Pseudom\allowbreak onas\%\allowbreak 20putid\allowbreak a} & Scale band: tier 4, ordered by measurements (26). & PRIDE v3\allowbreak  REST AP\allowbreak I; seven\allowbreak  candida\allowbreak te ident\allowbreak ifiers a\allowbreak re unver\allowbreak ified an\allowbreak d must b\allowbreak e re-\allowbreak checked\allowbreak  before \allowbreak any inge\allowbreak st \\
\addlinespace[5pt]
119 & \textbf{Kohlstedt 2022 (genome-reduced MA-10 muconate for PET)} \org{P. putida}\newline Kohlstedt M, Weimer A, Weiland F, Stolzenberger J, Selzer M, Sanz M, Kramps L, Wittmann C. Metab Eng 2022;72:337-352.\newline \href{https://doi.org/10.1016/j.ymben.2022.05.001}{doi:\allowbreak 10.\allowbreak 1016/\allowbreak j.\allowbreak ymben.\allowbreak 2022.\allowbreak 05.\allowbreak 001} & Scale band: tier 4, ordered by measurements (3). & unconfir\allowbreak med; pub\allowbreak lisher r\allowbreak eturned \allowbreak 403\newline {\scriptsize (browser)} \\
\addlinespace[5pt]
120 & \textbf{Kohlstedt 2018 (MA-6 muconate to nylon)} \org{P. putida}\newline Kohlstedt M, Starck S, Barton N, Stolzenberger J, Selzer M, Mehlmann K, et al., Wittmann C. Metab Eng 2018;47:279-293.\newline \href{https://doi.org/10.1016/j.ymben.2018.03.003}{doi:\allowbreak 10.\allowbreak 1016/\allowbreak j.\allowbreak ymben.\allowbreak 2018.\allowbreak 03.\allowbreak 003} & Scale band: tier 4, ordered by measurements (3). & unconfir\allowbreak med; pub\allowbreak lisher r\allowbreak eturned \allowbreak 403\newline {\scriptsize (browser)} \\
\addlinespace[5pt]
121 & \textbf{PAO1 ppk1 and ppk2 Biolog panels} \org{P. aeruginosa}\newline Polyphosphate kinase mutants of Pseudomonas aeruginosa PAO1 on Biolog PM1 to PM3. Biol Res 2015;48.\newline \href{https://doi.org/10.1186/s40659-015-0012-0}{doi:\allowbreak 10.\allowbreak 1186/\allowbreak s40659-\allowbreak 015-\allowbreak 0012-\allowbreak 0} & Scale band: tier 4, ordered by measurements (3). & open-\allowbreak access \allowbreak article;\allowbreak  deposit\allowbreak  and for\allowbreak mat not \allowbreak reported \\
\addlinespace[5pt]
122 & \textbf{Vardon 2015 (adipic acid from lignin)} \org{P. putida}\newline Vardon DR, et al. Adipic acid production from lignin. Energy Environ Sci 2015.\newline \href{https://doi.org/10.1039/C4EE03230F}{doi:\allowbreak 10.\allowbreak 1039/\allowbreak C4EE032\allowbreak 30F} & Scale band: tier 4, ordered by measurements (2). & unconfir\allowbreak med; pub\allowbreak lisher r\allowbreak eturned \allowbreak 403\newline {\scriptsize (browser)} \\
\addlinespace[5pt]
123 & \textbf{PXD074559 (JAG1 proteome, acetate against glucose)} \org{P. putida}\newline PRIDE project PXD074559, JAG1 proteome on acetate against glucose in a KT2440 background, 2026-02-17.\newline \href{https://www.ebi.ac.uk/pride/archive/projects/PXD074559}{ebi.\allowbreak ac.\allowbreak uk/\allowbreak pride/\allowbreak archive\allowbreak /\allowbreak project\allowbreak s/\allowbreak PXD0745\allowbreak 59} & Scale band: tier 4, ordered by measurements (2). & PRIDE PX\allowbreak D074559;\allowbreak  retriev\allowbreak able thr\allowbreak ough the\allowbreak  PRIDE A\allowbreak PI \\
\addlinespace[5pt]
124 & \textbf{Batianis 2020 (expanded CRISPRi toolbox)} \org{P. putida}\newline Batianis C, et al. An expanded CRISPRi toolbox for tunable control of gene expression in Pseudomonas putida. bioRxiv 2020;2020.03.04.975888.\newline \href{https://www.biorxiv.org/content/10.1101/2020.03.04.975888v1}{biorxiv.\allowbreak org/\allowbreak content\allowbreak /\allowbreak 10.\allowbreak 1101/\allowbreak 2020.\allowbreak 03.\allowbreak 04.\allowbreak 975888v\allowbreak 1} & Scale band: tier 4, ordered by measurements (2). & preprint\allowbreak  only; n\allowbreak o deposi\allowbreak t named \\
\addlinespace[5pt]
125 & \textbf{KT2440 lignin-stream tolerance RB-TnSeq (2023)} \org{P. putida}\newline RB-TnSeq identifies genetic targets for improved tolerance of lignin-related streams. Metab Eng 2023.\newline \href{https://research-hub.nrel.gov/en/publications/rb-tnseq-identifies-genetic-targets-for-improved-tolerance-of-seu/}{research\allowbreak -\allowbreak hub.\allowbreak nrel.\allowbreak gov/\allowbreak en/\allowbreak publica\allowbreak tions/\allowbreak rb-\allowbreak tnseq-\allowbreak identif\allowbreak ies-\allowbreak genetic\allowbreak -\allowbreak targets\allowbreak -\allowbreak for-\allowbreak improve\allowbreak d-\allowbreak toleran\allowbreak ce-\allowbreak of-\allowbreak seu} & Scale band: tier 4. & not repo\allowbreak rted \\
\addlinespace[5pt]
126 & \textbf{KT2440 metal-excess Tn-seq (2024)} \org{P. putida}\newline Genes required for growth of Pseudomonas putida KT2440 under metal excess. bioRxiv 2024;2024.04.12.589247; journal version in Environ Microbiol 2025.\newline \href{https://doi.org/10.1101/2024.04.12.589247}{doi:\allowbreak 10.\allowbreak 1101/\allowbreak 2024.\allowbreak 04.\allowbreak 12.\allowbreak 589247} & Scale band: tier 4. & not repo\allowbreak rted \\
\addlinespace[5pt]
127 & \textbf{Woodworth 2021 (syringyl lignin to pyrone dicarboxylic acid)} \org{P. putida}\newline Woodworth BD, et al., Beckham GT. Metabolism of syringyl lignin-derived compounds in Pseudomonas putida enables convergent production of 2-pyrone-4,6-dicarboxylic acid. Metab Eng 2021.\newline \href{https://www.sciencedirect.com/science/article/abs/pii/S1096717621000276}{scienced\allowbreak irect.\allowbreak com/\allowbreak science\allowbreak /\allowbreak article\allowbreak /\allowbreak abs/\allowbreak pii/\allowbreak S109671\allowbreak 76210002\allowbreak 76} & Scale band: tier 4, ordered by measurements (1). & unconfir\allowbreak med; ide\allowbreak ntifier \allowbreak not conf\allowbreak irmed\newline {\scriptsize (browser)} \\
\addlinespace[5pt]
128 & \textbf{PXD067010 (lignocellulosic sugars to an aviation-fuel precursor)} \org{P. putida}\newline PRIDE project PXD067010, multi-layered metabolic remodeling for lignocellulosic sugars to an aviation-fuel precursor, 2025-08-06.\newline \href{https://www.ebi.ac.uk/pride/archive/projects/PXD067010}{ebi.\allowbreak ac.\allowbreak uk/\allowbreak pride/\allowbreak archive\allowbreak /\allowbreak project\allowbreak s/\allowbreak PXD0670\allowbreak 10} & Scale band: tier 4. & PRIDE PX\allowbreak D067010;\allowbreak  retriev\allowbreak able thr\allowbreak ough the\allowbreak  PRIDE A\allowbreak PI \\
\addlinespace[5pt]
129 & \textbf{PXD054609 (isoprenol tolerance, evolution plus omics)} \org{P. putida}\newline PRIDE project PXD054609, adaptive laboratory evolution plus omics for isoprenol tolerance engineering, 2024-08-05.\newline \href{https://www.ebi.ac.uk/pride/archive/projects/PXD054609}{ebi.\allowbreak ac.\allowbreak uk/\allowbreak pride/\allowbreak archive\allowbreak /\allowbreak project\allowbreak s/\allowbreak PXD0546\allowbreak 09} & Scale band: tier 4. & PRIDE PX\allowbreak D054609;\allowbreak  retriev\allowbreak able thr\allowbreak ough the\allowbreak  PRIDE A\allowbreak PI \\
\addlinespace[5pt]
130 & \textbf{PXD045993 (polyhydroxyalkanoate biology in KT2440)} \org{P. putida}\newline Kelly S, et al. Comprehensive polyhydroxyalkanoate biology in Pseudomonas putida KT2440. Mol Cell Proteomics 2024;23:100765; PRIDE PXD045993.\newline \href{https://doi.org/10.1016/j.mcpro.2024.100765}{doi:\allowbreak 10.\allowbreak 1016/\allowbreak j.\allowbreak mcpro.\allowbreak 2024.\allowbreak 100765} & Scale band: tier 4. & PRIDE PX\allowbreak D045993;\allowbreak  retriev\allowbreak able thr\allowbreak ough the\allowbreak  PRIDE A\allowbreak PI \\
\addlinespace[5pt]
131 & \textbf{QconCAT absolute quantification workflow (2023)} \org{P. putida}\newline QconCAT absolute-quantification workflow for Pseudomonas putida KT2440. STAR Protoc 2023.\newline \href{https://pmc.ncbi.nlm.nih.gov/articles/PMC9881405}{PMC98814\allowbreak 05} & Scale band: tier 4. & code at \allowbreak biosusta\allowbreak in/\allowbreak QconCAT\allowbreak quantSTA\allowbreak R; no PR\allowbreak IDE acce\allowbreak ssion fo\allowbreak und \\
\addlinespace[5pt]
132 & \textbf{MTBLS12468 (GAPDH homologs on glucose against ferulate)} \org{P. putida}\newline MetaboLights study MTBLS12468, GAPDH homologs switching between glucose and ferulate metabolism in KT2440, LC-MS negative mode, 2025-11-03.\newline \href{https://www.ebi.ac.uk/metabolights/MTBLS12468}{ebi.\allowbreak ac.\allowbreak uk/\allowbreak metabol\allowbreak ights/\allowbreak MTBLS12\allowbreak 468} & Scale band: tier 4. & MetaboLi\allowbreak ghts MTB\allowbreak LS12468;\allowbreak  unverif\allowbreak ied, rea\allowbreak ched thr\allowbreak ough an \allowbreak aggregat\allowbreak or \\
\addlinespace[5pt]
133 & \textbf{MTBLS1715 (13C metabolomics under iron limitation)} \org{P. putida}\newline MetaboLights study MTBLS1715, 13C metabolomics of iron-replete against iron-limited pseudomonads including KT2440 and S12, 2020-12-05.\newline \href{https://www.ebi.ac.uk/metabolights/MTBLS1715}{ebi.\allowbreak ac.\allowbreak uk/\allowbreak metabol\allowbreak ights/\allowbreak MTBLS17\allowbreak 15} & Scale band: tier 4. & MetaboLi\allowbreak ghts MTB\allowbreak LS1715; \allowbreak unverifi\allowbreak ed, reac\allowbreak hed thro\allowbreak ugh an a\allowbreak ggregato\allowbreak r \\
\addlinespace[5pt]
134 & \textbf{MTBLS319 and MTBLS320 (DOT-T1E solvent and drug exposure)} \org{P. putida}\newline MetaboLights studies MTBLS319 (toluene) and MTBLS320 (propranolol) in Pseudomonas putida DOT-T1E, 2016-08-23.\newline \href{https://www.ebi.ac.uk/metabolights/MTBLS319}{ebi.\allowbreak ac.\allowbreak uk/\allowbreak metabol\allowbreak ights/\allowbreak MTBLS31\allowbreak 9} & Scale band: tier 4. & MetaboLi\allowbreak ghts MTB\allowbreak LS319 an\allowbreak d MTBLS3\allowbreak 20; unve\allowbreak rified, \allowbreak reached \allowbreak through \allowbreak an aggre\allowbreak gator \\
\addlinespace[5pt]
135 & \textbf{Sigma-factor regulon reconstruction (2026)} \org{P. putida}\newline System-level reconstruction of RpoD, RpoS and RpoN regulatory networks in Pseudomonas putida KT2440. J Biol Eng 2026.\newline \href{https://doi.org/10.1186/s13036-026-00769-z}{doi:\allowbreak 10.\allowbreak 1186/\allowbreak s13036-\allowbreak 026-\allowbreak 00769-\allowbreak z} & Scale band: tier 4. & unconfir\allowbreak med; GSE\allowbreak 322955 i\allowbreak s a plau\allowbreak sible bu\allowbreak t unveri\allowbreak fied mat\allowbreak ch \\
\addlinespace[5pt]
136 & \textbf{Fenster 2022 (KT2440 CRISPRi toolset with a product readout)} \org{P. putida}\newline Fenster J, Werner A, Tay JW, et al. Metab Eng Commun 2022;15:e00204.\newline \href{https://doi.org/10.1016/j.mec.2022.e00204}{doi:\allowbreak 10.\allowbreak 1016/\allowbreak j.\allowbreak mec.\allowbreak 2022.\allowbreak e00204} & Scale band: tier 4. & open-\allowbreak access \allowbreak article;\allowbreak  no depo\allowbreak sit name\allowbreak d \\
\addlinespace[5pt]
137 & \textbf{Rhizosphere CRISPRi (2024)} \org{P. putida}\newline Optimized CRISPR interference system for investigating Pseudomonas alloputida genes involved in rhizosphere microbiome assembly. 2024; PMID 39163848.\newline \href{https://pubmed.ncbi.nlm.nih.gov/39163848/}{PMID 391\allowbreak 63848} & Scale band: tier 4. & not repo\allowbreak rted \\
\addlinespace[5pt]
138 & \textbf{Pediatric cystic fibrosis susceptibility panel (2025)} \org{P. aeruginosa}\newline Pediatric cystic fibrosis isolates with EUCAST susceptibility testing and whole-genome sequencing. Antibiotics 2025;14:302.\newline \href{https://doi.org/10.3390/antibiotics14030302}{doi:\allowbreak 10.\allowbreak 3390/\allowbreak antibio\allowbreak tics1403\allowbreak 0302} & Scale band: tier 4. & unconfir\allowbreak med\newline {\scriptsize (browser)} \\
\addlinespace[5pt]
139 & \textbf{PA14 non-redundant phage-resistance screen (2024)} \org{P. aeruginosa}\newline A phage-resistance screen of the PA14 non-redundant transposon library, 2024; authors and identifier unconfirmed.\newline \href{https://www.ncbi.nlm.nih.gov/pmc/articles/PMC11241198/}{PMC11241\allowbreak 198} & Scale band: tier 4. & unconfir\allowbreak med; the\allowbreak  article\allowbreak  page se\allowbreak rved a c\allowbreak hallenge\newline {\scriptsize (browser)} \\
\addlinespace[5pt]
140 & \textbf{OSTI barcoded P. fluorescens transposon libraries} \org{P. fluorescens}\newline A conference record describing barcoded transposon insertion libraries in Pseudomonas fluorescens with over 100,000 independent insertions, used to study root colonization; the strain is not stated.\newline \href{https://www.osti.gov/biblio/1236050}{osti.\allowbreak gov/\allowbreak biblio/\allowbreak 1236050} & Scale band: tier 4. & not repo\allowbreak rted; th\allowbreak e strain\allowbreak  is not \allowbreak stated\newline {\scriptsize (browser)} \\
\addlinespace[5pt]
\end{longtable}
\endgroup
\end{landscape}


\clearpage
%% GENERATED FILE -- do not hand-edit.
%% SOURCE: experiments/database/scripts/build_hosts_candidate_datasets_table.py
\begingroup
\footnotesize
\begin{longtable}{@{}r L{76mm} L{20mm} L{36mm} L{24mm}@{}}
\caption[]{Candidate-gate verdict per ranked row, rendered from the verdict store
(\texttt{database/candidates/}). \emph{no verdict} means the row has not been gated, and
every row reads that way here: the stored verdict's \emph{table} field admits only the
two existing pools, so a row of this table cannot yet carry one
(Table~\ref{tab:schema}).}
\label{tab:verdicts}\\
\toprule
Rank & Dataset & Host & Verdict & Decided \\
\midrule
\endfirsthead
\toprule
Rank & Dataset & Host & Verdict & Decided \\
\midrule
\endhead
\bottomrule
\endfoot
1 & de Boer 2020 (GPRA sort-seq, 100M random promoters) & \org{S. cerevisiae} & no verdict &  \\
2 & Vaishnav 2022 deferred training libraries & \org{S. cerevisiae} & no verdict &  \\
3 & DREAM 2022 random promoter challenge (Rafi 2024) & \org{S. cerevisiae} & no verdict &  \\
4 & Kotopka and Smolke 2020 (artificial promoter libraries) & \org{S. cerevisiae} & no verdict &  \\
5 & Zhu 2024 (Pol II MASTER, TSS initiation efficiency) & \org{S. cerevisiae} & no verdict &  \\
6 & Cuperus 2017 (random 5' UTR library) & \org{S. cerevisiae} & no verdict &  \\
7 & Gui 2021 (position effects with transcriptomes) & \org{S. cerevisiae} & no verdict &  \\
8 & Keren 2016 (promoter panel x named gene, pooled fitness) & \org{S. cerevisiae} & no verdict &  \\
9 & Shalem 2015 (3' end FACS-seq library) & \org{S. cerevisiae} & no verdict &  \\
10 & Renganaath 2020 (native promoter variant MPRA) & \org{S. cerevisiae} & no verdict &  \\
11 & Arg-CCU tRNA deep mutational scan & \org{S. cerevisiae} & no verdict &  \\
12 & Keren 2013 (native promoter panel across conditions) & \org{S. cerevisiae} & no verdict &  \\
13 & Lubliner 2015 (core promoter variant panel) & \org{S. cerevisiae} & no verdict &  \\
14 & Sharon 2012 (designed promoter library) & \org{S. cerevisiae} & no verdict &  \\
15 & Yamanishi 2013 (genomic terminatome) & \org{S. cerevisiae} & no verdict &  \\
16 & Dvir 2013 (5' UTR positions -10 to -1) & \org{S. cerevisiae} & no verdict &  \\
17 & Wu 2017 (genome-wide position effects, 1,044 loci) & \org{S. cerevisiae} & no verdict &  \\
18 & ProteinGym yeast subset & \org{S. cerevisiae} & no verdict &  \\
19 & Zhang 2016 (genome-wide GFP position scan) & \org{S. cerevisiae} & no verdict &  \\
20 & Kita 2017 (context-stratified yeast eQTLs) & \org{S. cerevisiae} & no verdict &  \\
21 & Chen (earlier position-effect study, 482 loci) & \org{S. cerevisiae} & no verdict &  \\
22 & Wei 2017 (synthetic terminator panel) & \org{S. cerevisiae} & no verdict &  \\
23 & Lubliner 2013 (native core promoter maximal strengths) & \org{S. cerevisiae} & no verdict &  \\
24 & Morse (short synthetic terminators) & \org{S. cerevisiae} & no verdict &  \\
25 & MaveDB yeast score sets & \org{S. cerevisiae} & no verdict &  \\
26 & Fitness Browser RCH2 (Stutzerimonas stutzeri\_AE) & \org{S. stutzeri} & no verdict &  \\
27 & Fitness Browser FW300-N2E3 (Aquipseudomonas helvetica) & \org{P. fluorescens} & no verdict &  \\
28 & Fitness Browser FW300-N2E2 (Aquipseudomonas fluorescens\_Q) & \org{P. fluorescens} & no verdict &  \\
29 & Fitness Browser WCS417 (Aquipseudomonas simiae) & \org{P. simiae} & no verdict &  \\
30 & Fitness Browser FW300-N2C3 (Aquipseudomonas fluorescens\_AA) & \org{P. fluorescens} & no verdict &  \\
31 & Fitness Browser FW300-N1B4 (Aquipseudomonas sp.) & \org{P. fluorescens} & no verdict &  \\
32 & Fitness Browser B728a (Aquipseudomonas fragariae pv. syringae) & \org{P. syringae} & no verdict &  \\
33 & Fitness Browser B728a mexB deletion background & \org{P. syringae} & no verdict &  \\
34 & Fitness Browser PA14 (Pseudomonas aeruginosa) & \org{P. aeruginosa} & no verdict &  \\
35 & Dubinkina 2024 (12-strain low-pH transcriptomic atlas) & \org{P. kudriavzevii} & no verdict &  \\
36 & MAESTRO 2025 (170-isolate WGS, 57-stressor growth panel) & \org{P. kudriavzevii} & no verdict &  \\
37 & Fatma 2023 (landing-pad locus screen) & \org{P. kudriavzevii} & no verdict &  \\
38 & Douglass 2018 (32 sequenced isolates, 4-antifungal MIC) & \org{P. kudriavzevii} & no verdict &  \\
39 & GSE282019 (Delhi NICU five-year outbreak) & \org{P. kudriavzevii} & no verdict &  \\
40 & GSE202655 (L-malic acid, EcSthA transhydrogenase) & \org{P. kudriavzevii} & no verdict &  \\
41 & GSE159457 (CY902 C4 dicarboxylic acid transporters) & \org{P. kudriavzevii} & no verdict &  \\
42 & GSE262836 (Delhi NICU, fluconazole-resistance expression) & \org{P. kudriavzevii} & no verdict &  \\
43 & Thompson 2019 (KT2440 lysine metabolism) & \org{P. putida} & no verdict &  \\
44 & Isoprenol CRISPRi DBTL proteome ladder (7 PRIDE deposits) & \org{P. putida} & no verdict &  \\
45 & Thompson 2020 (KT2440 fatty acids and alcohols) & \org{P. putida} & no verdict &  \\
46 & GSE198395 (putidaPRECISE321 carbon sources) & \org{P. putida} & no verdict &  \\
47 & Mueller 2025 (24-strain P. putida Biolog panel) & \org{P. putida} & no verdict &  \\
48 & Schmidt 2022 (KT2440 nitrogen BarSeq) & \org{P. putida} & no verdict &  \\
49 & Borchert 2024 (KT2440 fitness compendium) & \org{P. putida} & no verdict &  \\
50 & putidaPRECISE321 (Lim 2022) & \org{P. putida} & no verdict &  \\
51 & Cao 2020 (pooled four-condition RNA-seq) & \org{P. kudriavzevii} & no verdict &  \\
52 & PXD069498 (extracellular vesicle proteome) & \org{P. kudriavzevii} & no verdict &  \\
53 & Douglass 2018 (per-strain SNP density and loss of heterozygosity) & \org{P. kudriavzevii} & no verdict &  \\
54 & Tran 2023 (succinate end-to-end pipeline) & \org{P. kudriavzevii} & no verdict &  \\
55 & Cao 2020 (constitutive promoter-GFP panel) & \org{P. kudriavzevii} & no verdict &  \\
56 & Xiao 2014 (SD108 reductive-TCA succinate) & \org{P. kudriavzevii} & no verdict &  \\
57 & Xylose utilization across four natural backgrounds & \org{P. kudriavzevii} & no verdict &  \\
58 & Illinois Data Bank IDB-9063411 (high-yield 3-HP) & \org{P. kudriavzevii} & no verdict &  \\
59 & Douglass 2018 (CBS573 RNA-seq) & \org{P. kudriavzevii} & no verdict &  \\
60 & MAESTRO 2025 (gene4260/gene4261 paralog-pair deletion) & \org{P. kudriavzevii} & no verdict &  \\
61 & Illinois Data Bank IDB-0538107 (succinate pipeline deposit) & \org{P. kudriavzevii} & no verdict &  \\
62 & CABBI inorganic carbon sources (low-pH succinate kinetics) & \org{P. kudriavzevii} & no verdict &  \\
63 & Xylose 3-HP from cellulosic hydrolysates (2026) & \org{P. kudriavzevii} & no verdict &  \\
64 & Jeong 2025 (acetate metabolism during xylose fermentation) & \org{P. kudriavzevii} & no verdict &  \\
65 & Glycolic acid overproduction (2026) & \org{P. kudriavzevii} & no verdict &  \\
66 & Wu 2023 (piggyBac citramalate copy-number and locus variants) & \org{P. kudriavzevii} & no verdict &  \\
67 & Mitochondrial decompartmentalization (2025, with GNPS metabolomics) & \org{P. kudriavzevii} & no verdict &  \\
68 & GSE39558 (ethanol-stress expression) & \org{P. kudriavzevii} & no verdict &  \\
69 & PXD041801 (flocculin expression and floating biofilms) & \org{P. kudriavzevii} & no verdict &  \\
70 & PXD039973 (succinylation and amphotericin B resistance) & \org{P. kudriavzevii} & no verdict &  \\
71 & PXD075355 (Pichia against Botrytis cinerea) & \org{P. kudriavzevii} & no verdict &  \\
72 & Pk E1 toolkit expansion (mycoprotein, 2026) & \org{P. kudriavzevii} & no verdict &  \\
73 & China fermented-vegetable and clinical panel (2024) & \org{P. kudriavzevii} & no verdict &  \\
74 & Indonesian clinical genomes (PRJNA819536) & \org{P. kudriavzevii} & no verdict &  \\
75 & Unidentified Illinois Data Bank CRISPR/Cas9 gene-deletion deposit & \org{P. kudriavzevii} & no verdict &  \\
76 & Khaledi 2020 (414-isolate clinical transcriptome panel) & \org{P. aeruginosa} & no verdict &  \\
77 & Kaczmarczyk 2025 (PAO1 genome-wide CRISPRi) & \org{P. aeruginosa} & no verdict &  \\
78 & Dotsch 2009 (arrayed PA14 library against 19 antibiotics) & \org{P. aeruginosa} & no verdict &  \\
79 & Bartell 2019 (longitudinal cystic fibrosis panel) & \org{P. aeruginosa} & no verdict &  \\
80 & Genome Medicine 2024 (1,877-strain genome and MIC set) & \org{P. aeruginosa} & no verdict &  \\
81 & Penaranda 2025 (125-isolate diversity survey) & \org{P. aeruginosa} & no verdict &  \\
82 & Khaledi 2020 (1,656-record agar-dilution MIC panel) & \org{P. aeruginosa} & no verdict &  \\
83 & Poulsen 2019 (nine-strain Tn-seq core essentiality) & \org{P. aeruginosa} & no verdict &  \\
84 & PAO1 arrayed transposon libraries (Jacobs 2003, Held 2012) & \org{P. aeruginosa} & no verdict &  \\
85 & KT2440 pooled transposon library (Putida\_ML5 metadata) & \org{P. putida} & no verdict &  \\
86 & Gao 2020 (KT2440 targeted SRM proteomics) & \org{P. putida} & no verdict &  \\
87 & BACTOME (101 clinical genomes with eight phenotypes) & \org{P. aeruginosa} & no verdict &  \\
88 & Kohlstedt 2019 (13C metabolic flux analysis, two hosts) & \org{P. putida} & no verdict &  \\
89 & Turner 2015 (PA14 and PAO1 essentiality across four arms) & \org{P. aeruginosa} & no verdict &  \\
90 & Liberati 2006 (PA14 non-redundant transposon library and essentiality) & \org{P. aeruginosa} & no verdict &  \\
91 & GOLD\_115 and GOLD\_84 essential-gene sets & \org{P. aeruginosa} & no verdict &  \\
92 & Werner 2023 (beta-ketoadipate from lignin) & \org{P. putida} & no verdict &  \\
93 & Wilkes 2026 (muconate multi-omics with fluxomics) & \org{P. putida} & no verdict &  \\
94 & GSE213139 (low-dose ionizing radiation) & \org{P. putida} & no verdict &  \\
95 & GSE131229 (non-antibiotic pharmaceutical stress) & \org{P. putida} & no verdict &  \\
96 & GSE248909 (micro and nanoplastic stress) & \org{P. putida} & no verdict &  \\
97 & GSE193493 (putidaPRECISE321 second deposit) & \org{P. putida} & no verdict &  \\
98 & GSE221044 (kaiC deletion) & \org{P. putida} & no verdict &  \\
99 & GSE144045 (solvent-tolerance evolution in P. putida S12) & \org{P. putida} & no verdict &  \\
100 & GSE275525 (HexR regulatory landscape) & \org{P. putida} & no verdict &  \\
101 & GSE309641 (fluoride adaptive evolution) & \org{P. putida} & no verdict &  \\
102 & GSE155767 (evolved KT2440 on non-native sugars) & \org{P. putida} & no verdict &  \\
103 & GSE147420 (ammonium against nitrate) & \org{P. putida} & no verdict &  \\
104 & GSE229763 (one-carbon metabolism) & \org{P. putida} & no verdict &  \\
105 & KT2440 aromatics RB-TnSeq (2019 preprint) & \org{P. putida} & no verdict &  \\
106 & PMkbase v1.0 (phenotype-microarray database) & \org{P. putida} & no verdict &  \\
107 & IPCD (International P. aeruginosa Consortium Database) & \org{P. aeruginosa} & no verdict &  \\
108 & Pirnay 2009 (328-isolate population survey) & \org{P. aeruginosa} & no verdict &  \\
109 & n-butanol KT2440 proteome (2015) & \org{P. putida} & no verdict &  \\
110 & Skurnik 2013 (PA14 essentiality) & \org{P. aeruginosa} & no verdict &  \\
111 & Lee 2015 (PAO1 essentiality across three media) & \org{P. aeruginosa} & no verdict &  \\
112 & P. putida S12 phenotype microarray (190 carbon sources) & \org{P. putida} & no verdict &  \\
113 & Chandler 2023 (128 functionally phenotyped isolates) & \org{P. aeruginosa} & no verdict &  \\
114 & Breidenstein 2008 (ciprofloxacin resistome) & \org{P. aeruginosa} & no verdict &  \\
115 & Remaining P. putida GEO expression series & \org{P. putida} & no verdict &  \\
116 & KT2440 mini-Tn5 knockout collection (2010) & \org{P. putida} & no verdict &  \\
117 & Nunez-Garcia 2024 (39 isolates over ten years) & \org{P. aeruginosa} & no verdict &  \\
118 & Remaining P. putida PRIDE proteomics projects & \org{P. putida} & no verdict &  \\
119 & Kohlstedt 2022 (genome-reduced MA-10 muconate for PET) & \org{P. putida} & no verdict &  \\
120 & Kohlstedt 2018 (MA-6 muconate to nylon) & \org{P. putida} & no verdict &  \\
121 & PAO1 ppk1 and ppk2 Biolog panels & \org{P. aeruginosa} & no verdict &  \\
122 & Vardon 2015 (adipic acid from lignin) & \org{P. putida} & no verdict &  \\
123 & PXD074559 (JAG1 proteome, acetate against glucose) & \org{P. putida} & no verdict &  \\
124 & Batianis 2020 (expanded CRISPRi toolbox) & \org{P. putida} & no verdict &  \\
125 & KT2440 lignin-stream tolerance RB-TnSeq (2023) & \org{P. putida} & no verdict &  \\
126 & KT2440 metal-excess Tn-seq (2024) & \org{P. putida} & no verdict &  \\
127 & Woodworth 2021 (syringyl lignin to pyrone dicarboxylic acid) & \org{P. putida} & no verdict &  \\
128 & PXD067010 (lignocellulosic sugars to an aviation-fuel precursor) & \org{P. putida} & no verdict &  \\
129 & PXD054609 (isoprenol tolerance, evolution plus omics) & \org{P. putida} & no verdict &  \\
130 & PXD045993 (polyhydroxyalkanoate biology in KT2440) & \org{P. putida} & no verdict &  \\
131 & QconCAT absolute quantification workflow (2023) & \org{P. putida} & no verdict &  \\
132 & MTBLS12468 (GAPDH homologs on glucose against ferulate) & \org{P. putida} & no verdict &  \\
133 & MTBLS1715 (13C metabolomics under iron limitation) & \org{P. putida} & no verdict &  \\
134 & MTBLS319 and MTBLS320 (DOT-T1E solvent and drug exposure) & \org{P. putida} & no verdict &  \\
135 & Sigma-factor regulon reconstruction (2026) & \org{P. putida} & no verdict &  \\
136 & Fenster 2022 (KT2440 CRISPRi toolset with a product readout) & \org{P. putida} & no verdict &  \\
137 & Rhizosphere CRISPRi (2024) & \org{P. putida} & no verdict &  \\
138 & Pediatric cystic fibrosis susceptibility panel (2025) & \org{P. aeruginosa} & no verdict &  \\
139 & PA14 non-redundant phage-resistance screen (2024) & \org{P. aeruginosa} & no verdict &  \\
140 & OSTI barcoded P. fluorescens transposon libraries & \org{P. fluorescens} & no verdict &  \\
\end{longtable}
\endgroup



\end{document}
